\delta_j := \|P_{\mathrm{nor}}T_j\|,
\qquad
r_j := \|C_j\|.
\]
This channel ledger matrix $\Lambda_{\mathrm{ch}}$ is not introduced for numerical decoration. Its function is to give a
single interface where lemma statements, appendix tables, and future simulations
can all report the same structural status.

\paragraph{Proposition 1 (matrix admissibility criterion).}
Suppose there exist constants $\alpha_j,\beta_j\ge0$ such that
\[
\delta_j^2 \le \alpha_j \Gamma + \beta_j \|\nabla\Gamma\|^2
\qquad\text{for each channel }j,
\]
and assume $\sum_j \beta_j < \kappa^2$. Then the channel-specialized master flow is
closure-admissible in a neighbourhood of the closure sector.

\begin{proof}
The normal contributions are exactly the defect-producing pieces entering the
closure balance. By Cauchy--Schwarz,
\[
|\langle \nabla\Gamma, B_j\rangle|
\le \|\nabla\Gamma\|\,\delta_j.
\]
Applying Young's inequality channel by channel and summing yields a bound of the
same type as in the v2.3.0 admissibility criterion, now expressed entirely through
the ledger entries. The compensator term dominates once the cumulative normal
budget stays below the compensator threshold.
\end{proof}

\subsection*{4. Probe-channel specialization}
The same channel grammar must hold for the two-probe sector. Write the probe drift
in the form
\[
\mathcal T_{\mathrm{probe}} =
\mathcal T_{\mathrm{tri}}^{(1)} +
\mathcal T_{\mathrm{tri}}^{(2)} +
\mathcal T_{\perp}^{(Q)} +
\mathcal T_{\mathrm{match}} +
\mathcal T_{\mathrm{res}}^{\mathrm{probe}}.
\]
Here the two intrinsic channels belong to the two probe triolen, the term
$\mathcal T_{\perp}^{(Q)}$ encodes the orthogonal search geometry selected by $Q$,
and $\mathcal T_{\mathrm{match}}$ measures the approach or mismatch between the two
probe traces. The new point of v2.4.0 is that probe matching is promoted to
a channel in its own right rather than being hidden implicitly inside a heuristic
search narrative.

\paragraph{Definition 2 (probe match defect).}
Let
\[
\Delta_{\mathrm{match}}(\Psi_1,\Psi_2)
:= \|\Pi(\Psi_1)-\Pi(\Psi_2)\|^2,
\]
where $\Pi$ is a chosen readout or projection interface. Then the match channel is
called admissible if along the compensated probe flow,
\[
\frac{d}{d\tau}\Delta_{\mathrm{match}}
\le -\eta\,\Delta_{\mathrm{match}} + c\,\Gamma_{\mathrm{probe}}
\]
for some $\eta>0$ and finite $c$.
This means that probe convergence is controlled up to the ambient closure defect.

\paragraph{Corollary 2 (probe convergence under closure control).}
If the probe closure functional remains in a forward-invariant small sublevel tube
and the match channel is admissible, then probe readout mismatch decays up to an
error proportional to the residual closure budget. Hence two probes can be treated
as a mathematically disciplined comparison device rather than only a heuristic
search metaphor.

\subsection*{5. Consequence for appendix tools and future derivations}
The specialization introduced in v2.4.0 changes the status of the appendix
again. From now on, every major table, tool, and shortcut should be readable as a
compressed report of one of the ledger entries or of one of the probe-channel
budgets. This gives a concrete criterion for deciding whether a future tool is
structural or merely expository: a genuine structural tool must reduce, estimate,
or transport one of the channel quantities.

The practical consequence is decisive. Future expansions no longer need to say only
that a section ``supports closure'' or ``helps probe reduction.'' They can specify
which channel is being sharpened, which defect budget is reduced, and which ledger
entry is stabilized. This is exactly the sort of disciplined mathematical
thickening that prepares the transition from structure-first architecture to a more
fully unfolded proof.


\section*{v2.5.0 Channel Budget Transport and Ledger Dynamics Core}
\addcontentsline{toc}{section}{v2.5.0 Channel Budget Transport and Ledger Dynamics Core}

\begin{quote}
\textbf{Normalization note (v2.17.1 local review).}
This block is retained as a \emph{historical transport and ledger precursor}. Its proof-bearing content should be read through the normalized grouped architecture and the later unfolding phases, rather than as an independent foundational core.
\end{quote}

The specialization step of v2.4.0 assigned every contribution to a disciplined
channel and recorded its tangential, normal, and compensator-facing sizes in the
ledger matrix. The next structural hardening step is to turn this static picture
into an evolution rule. In v2.5.0 the ledger is therefore promoted from a
reporting interface to a transport object. One no longer asks only whether a
channel is admissible at a fixed stage; one asks how admissibility is propagated,
recovered, or degraded along the proof flow.

\subsection*{1. Dynamic ledger variables}
For each channel $j\in\{\mathrm{tri},\perp,\mathrm{meas},\mathrm{res}\}$ define the
budget variables
\[
B_j(\tau):=\ell_j(\tau)^2+\delta_j(\tau)^2+r_j(\tau)^2,
\qquad
N_j(\tau):=\delta_j(\tau)^2,
\qquad
R_j(\tau):=r_j(\tau)^2.
\]
Here $B_j$ is the total structural budget carried by the channel, $N_j$ is the
normal leakage budget, and $R_j$ is the amount already routed through the
compensator-facing slot. The point is conceptual as much as technical: future
lemmas may now refer to budgets rather than to isolated terms, which makes it
possible to compare sections of the document without re-expanding the full
microscopic decomposition each time.

Let
\[
\mathfrak B(\tau):=\sum_j B_j(\tau),
\qquad
\mathfrak N(\tau):=\sum_j N_j(\tau),
\qquad
\mathfrak R(\tau):=\sum_j R_j(\tau).
\]
The quantity $\mathfrak B$ is the total visible channel load, while
$\mathfrak N$ and $\mathfrak R$ separate the defect-producing part from the
compensator-processed part.

\subsection*{2. Budget transport law}
The structural thesis of v2.5.0 is that channel evolution should be read as
transport with controlled exchange, not as unrelated term growth. This is encoded
in the schematic law
\[
\frac{d}{d\tau}B_j
= -\mu_j B_j + \sum_{k\neq j}\sigma_{kj}B_k + \rho_j\Gamma + \xi_j,
\]
where $\mu_j\ge0$ is the internal damping of channel $j$, the coefficients
$\sigma_{kj}\ge0$ measure admissible transfer from channel $k$ into channel $j$,
$\rho_j\Gamma$ records closure-driven forcing, and $\xi_j$ is a structured source
term that must be tagged by origin class. In particular, $\xi_j$ is not allowed to
hide unexplained growth. It must come from one of the already sanctioned origins:
projection conversion, probe matching, approximation residue, or phenomenological
placeholder.

\paragraph{Definition 3 (transport-admissible ledger).}
A channel ledger is called transport-admissible on an interval $I$ if there exist
nonnegative coefficients $(\mu_j,\sigma_{kj},\rho_j)$ and tagged sources $\xi_j$
such that the above transport law holds on $I$, the cumulative transfer matrix
$\Sigma=(\sigma_{kj})$ has finite operator norm, and every source term satisfies
\[
|\xi_j| \le c_j^{(0)}\Gamma + c_j^{(1)}\|\nabla\Gamma\|^2 + c_j^{(2)}\Delta_{\mathrm{match}}
+ c_j^{(3)}\mathrm{Res}_{\mathrm{tag}},
\]
where $\mathrm{Res}_{\mathrm{tag}}$ denotes a tagged residual quantity already
registered in the appendix grammar.

This definition codifies the project preference for auditable derivations: channel
budgets may move and mix, but they may not mutate mysteriously.

\subsection*{3. Global dissipative estimate}
The transport law becomes useful only once it closes against the compensated
closure sector.

\paragraph{Proposition 1 (global budget inequality).}
Assume the ledger is transport-admissible and suppose the transfer coefficients
satisfy
\[
\mu_* := \min_j \mu_j > \|\Sigma\|.
\]
Then there exist constants $A,B,C\ge0$ such that along the compensated flow,
\[
\frac{d}{d\tau}\mathfrak B
\le -A\mathfrak B + B\Gamma + C\Delta_{\mathrm{match}} + C\mathrm{Res}_{\mathrm{tag}}.
\]
In particular, if $\Gamma$, $\Delta_{\mathrm{match}}$, and the tagged residual budget
remain inside a forward-invariant small tube, then the total visible channel load
stays uniformly bounded and is asymptotically damped up to that tube.

\begin{proof}
Summing the channel transport laws gives
\[
\frac{d}{d\tau}\mathfrak B
= -\sum_j \mu_j B_j + \sum_j\sum_{k\neq j}\sigma_{kj}B_k
+ \Bigl(\sum_j \rho_j\Bigr)\Gamma + \sum_j \xi_j.
\]
The transfer contribution is bounded above by $\|\Sigma\|\mathfrak B$ after
collecting coefficients. The strict damping gap $\mu_*>\|\Sigma\|$ yields a net
negative coefficient in front of $\mathfrak B$. The source estimate from
transport-admissibility then produces the stated inequality after regrouping the
closure, match, and tagged residual terms.
\end{proof}

The structural meaning is important. Closure-admissibility is no longer a purely
local statement about one decomposition stage. It becomes a persistent budget
statement that survives transfer between channels.

\subsection*{4. Leakage absorption and residual discipline}
The next step is to distinguish visible channel mass from true defect leakage.
Because only the normal parts threaten closure, one needs a mechanism that proves
that leakage either decays or is absorbed into compensator-facing storage.

\paragraph{Definition 4 (absorption-efficient channel).}
A channel $j$ is called absorption-efficient if there exist constants
$\theta_j>0$ and $d_j\ge0$ such that
\[
\frac{d}{d\tau}N_j
\le -\theta_j N_j + d_j R_j + d_j\Gamma + d_j\Delta_{\mathrm{match}}.
\]
In words: normal leakage may temporarily feed on unresolved load, but it is
penalized unless supported by compensator storage or ambient closure stress.

\paragraph{Corollary 1 (aggregate leakage absorption).}
If every channel is absorption-efficient and if, in addition,
\[
\sum_j d_j R_j \le \lambda\,\mathfrak R
\qquad\text{with}\qquad
\lambda<\theta_*:=\min_j\theta_j,
\]
then there exist constants $E,F\ge0$ such that
\[
\frac{d}{d\tau}\mathfrak N
\le -E\mathfrak N + F\Gamma + F\Delta_{\mathrm{match}} + F\mathfrak R.
\]
Hence the total normal leakage is controlled by closure stress, probe mismatch,
and compensator-facing storage, rather than by an uncontrolled accumulation of
termwise errors.

This is exactly the desired middle statement between raw microlocal decomposition
and later field-equation language: the dangerous part of the proof can be
compressed into a leakage budget with explicit sinks and sources.

\subsection*{5. Probe-budget coupling}
The two-probe sector can now be written in the same grammar. Let
\[
\mathfrak B_{\mathrm{probe}}:=
B_{\mathrm{tri}}^{(1)}+B_{\mathrm{tri}}^{(2)}+B_{\perp}^{(Q)}+B_{\mathrm{match}}+B_{\mathrm{res}}^{\mathrm{probe}}.
\]
The new content of v\docversion\ is that probe comparison is no longer controlled
only through the mismatch scalar $\Delta_{\mathrm{match}}$, but through a full budget
that separates intrinsic probe cost, orthogonal search cost, match stabilization,
and residual search overhead.

\paragraph{Proposition 2 (probe budget stability).}
Assume the intrinsic probe channels and the orthogonal search channel are
transport-admissible, and assume the match channel satisfies
\[
\frac{d}{d\tau}B_{\mathrm{match}}
\le -\eta B_{\mathrm{match}} + c_1\Gamma_{\mathrm{probe}} + c_2\mathfrak N_{\mathrm{probe}}.
\]
Then the full probe budget obeys an estimate of the form
\[
\frac{d}{d\tau}\mathfrak B_{\mathrm{probe}}
\le -a\mathfrak B_{\mathrm{probe}} + b\Gamma_{\mathrm{probe}} + b\mathfrak N_{\mathrm{probe}} + b\mathrm{Res}^{\mathrm{probe}}_{\mathrm{tag}},
\]
for suitable $a,b>0$ whenever the transfer matrix in the probe sector has a
strict damping gap.

The interpretation is that two probes can consume structural budget while moving
through the search geometry, but this consumption is still accounted for by the
same ledger grammar as the main derivation. Nothing essentially new is hidden in
the probe picture; it is a structured subledger.

\subsection*{6. Consequence for appendix tools and proof memory}
The appendix can now be sharpened one step further. A legitimate tool is no longer
merely something that estimates a channel entry. In v\docversion\ a tool should,
whenever possible, declare one of four roles:
\begin{enumerate}[label=(\alph*)]
  \item budget estimator,
  \item transfer bound,
  \item leakage absorber, or
  \item source tagger.
\end{enumerate}
This classification is mathematically useful because it says where the tool enters
an actual differential inequality rather than only where it appears heuristically.

The same holds for proof memory. Once a lemma has been translated into budget form,
its content can be stored compactly as a transport or absorption rule without
repeating the entire derivation chain. This is the precise sense in which the
structure-first approach compresses thousands of calculations without erasing them:
it stores their net structural effect as a stable ledger law.

\subsection*{7. v\docversion\ readiness criterion}
The present section should be regarded as successful if the following can now be
done systematically:
\begin{itemize}
  \item every major future term can be assigned a channel budget rather than only a raw symbol;
  \item every appendix tool can be classified by its role in a budget inequality;
  \item probe arguments can be reported through a subledger instead of loose search language;
  \item residual terms can be traced by tagged sources rather than by informal comments; and
  \item later field equations can read these budget laws as compressed pre-equation structure.
\end{itemize}


This is the mathematical gain of v\docversion. The document now contains not only a
closure grammar and a channel grammar, but also a transport grammar that explains
how structural control is preserved as the proof unfolds.

\section*{v2.6.0 Lyapunov Ledger Closure and Forward-Invariant Regimes}
\addcontentsline{toc}{section}{v2.6.0 Lyapunov Ledger Closure and Forward-Invariant Regimes}

\begin{quote}
\textbf{Normalization note (v2.17.1 local review).}
This block is retained as a \emph{historical closure-stability and Lyapunov precursor}. Its role is supporting and developmental; the normalized proof-bearing closure chain is carried by the unfolding phases.
\end{quote}

The next logical step is to turn the budget grammar into a genuine closure device.
As long as one only writes transport inequalities, the proof still depends on a
stagewise reading. For the structure-first programme this is not yet strong enough.
One wants a single scalar functional that certifies that the admissible sector is
preserved under the whole transport chain. This is the role of the Lyapunov-ledger
closure introduced in v\docversion.

\subsection*{1. Weighted structural energy}
Choose positive weights
\[
\omega_j>0,
\qquad
\omega^{N}_j>0,
\qquad
\omega_R>0,
\qquad
\omega_\Gamma>0,
\]
and define the weighted structural energy
\[
\mathcal L
:=\sum_j \omega_j B_j + \sum_j \omega^{N}_j N_j + \omega_R\mathfrak R + \omega_\Gamma\Gamma.
\]
The aim is to choose the weights so that cross-channel transfer and residual spill
can be absorbed into one dissipative inequality. In this language, the ledger is no
longer just a bookkeeping device; it becomes the carrier of a monotonicity law.

\paragraph{Definition 1 (Lyapunov-admissible weight system).}
A family of weights is called Lyapunov-admissible if, after inserting the transport,
absorption, and source inequalities from the previous section, one obtains constants
$\kappa>0$ and $C_{\mathrm{tag}},C_{\mathrm{match}}\ge0$ such that
\[
\frac{d}{d\tau}\mathcal L
\le -\kappa \mathcal L
+ C_{\mathrm{tag}}\,\mathrm{Res}_{\mathrm{tag}}
+ C_{\mathrm{match}}\,\Delta_{\mathrm{match}}.
\]
This means that every admissible leakage path is paid for either by global decay or
by an explicitly tagged source of unresolved structure.

\subsection*{2. Weighted closure theorem}
The abstract transport statements from v2.5.0 are now compressed into a single
persistence theorem.

\paragraph{Theorem 1 (weighted ledger closure).}
Assume that the channel transport matrix has a strict damping gap, that every active
normal channel is absorption-efficient, and that residual production is source-tagged.
Then there exists a Lyapunov-admissible weight system for which
\[
\frac{d}{d\tau}\mathcal L
\le -\frac{\kappa}{2}\mathcal L
+ C_{\mathrm{tag}}\,\mathrm{Res}_{\mathrm{tag}}
+ C_{\mathrm{match}}\,\Delta_{\mathrm{match}}.
\]
In particular, if the tagged residual and probe mismatch stay integrably controlled,
then the whole derivation remains in a forward-invariant admissible regime.

\begin{proof}
Choose the weights recursively along the dominance order induced by the transfer
matrix, starting from the most dissipative channels and moving toward the weaker
ones. Because the damping gap is strict, the transfer contributions can be absorbed
by enlarging downstream weights finitely many times. The absorption-efficient normal
channels contribute additional negative terms for the dangerous normal mass. Residual
production enters only through tagged sources and therefore cannot create an
unidentified positive contribution. After summing the weighted inequalities and
absorbing the remaining off-diagonal terms, one obtains the stated estimate.
\end{proof}

This is the first point at which the document obtains a real forward-invariance
statement rather than only a local transport picture. The admissible proof sector is
now a dynamically protected cone.

\subsection*{3. Forward-invariant regime and restart stability}
One practical advantage of the Lyapunov-ledger closure is that it explains why the
proof can be resumed after interruptions without re-deriving the entire local chain.
If at some stage $\tau_0$ the state satisfies
\[
\mathcal L(\tau_0)\le M,
\qquad
\int_{\tau_0}^{\tau_1}\mathrm{Res}_{\mathrm{tag}}\,d\tau<\infty,
\qquad
\int_{\tau_0}^{\tau_1}\Delta_{\mathrm{match}}\,d\tau<\infty,
\]
then Gronwall's inequality implies a bound of the form
\[
\mathcal L(\tau)
\le e^{-\kappa(\tau-\tau_0)/2}M
+ \int_{\tau_0}^{\tau}e^{-\kappa(\tau-s)/2}
\bigl(C_{\mathrm{tag}}\,\mathrm{Res}_{\mathrm{tag}}(s)+C_{\mathrm{match}}\,\Delta_{\mathrm{match}}(s)\bigr)\,ds.
\]
Hence restarting the argument from an intermediate checkpoint is legitimate as long
as the checkpoint data are reported in ledger form. This gives a rigorous version of
proof memory: the document needs to preserve the weighted state, not the full local
micro-history.

\subsection*{4. Probe inheritance of invariance}
The two-probe picture inherits the same closure principle. Define the weighted probe
functional
\[
\mathcal L_{\mathrm{probe}}
:= \alpha_1 B^{(1)}_{\mathrm{tri}} + \alpha_2 B^{(2)}_{\mathrm{tri}}
+ \alpha_Q B^{(Q)}_{\perp} + \alpha_M B_{\mathrm{match}}
+ \alpha_R B^{\mathrm{probe}}_{\mathrm{res}}.
\]
If the probe transfer system has a strict damping gap and the mismatch source remains
ledger-tagged, then the same weighting argument yields
\[
\frac{d}{d\tau}\mathcal L_{\mathrm{probe}}
\le -\kappa_{\mathrm{probe}}\mathcal L_{\mathrm{probe}}
+ C^{\mathrm{probe}}_{\mathrm{tag}}\,\mathrm{Res}^{\mathrm{probe}}_{\mathrm{tag}}
+ C^{\mathrm{probe}}_{\mathrm{match}}\,\Delta_{\mathrm{match}}.
\]
Thus the probe sector is not merely compatible with the main ledger; it inherits the
same forward-invariant discipline. The search geometry therefore remains structurally
confined rather than opening an uncontrolled side channel.

\subsection*{5. Structural consequence for later field equations}
The future field-equation layer can now be read more sharply. A field equation will
not merely translate channel budgets into continuous variables; it must preserve the
forward-invariant regime induced by $\mathcal L$. In other words, admissible field
variables should be chosen so that the induced energy identity mirrors the weighted
ledger closure. This provides a direct criterion for accepting or rejecting later
field-equation ansaetze.

\subsection*{6. Consequence for appendix tools}
The tool appendix can be refined once more. Besides the roles introduced in v2.5.0,
a strong tool may now be marked as one of the following:
\begin{enumerate}[label=(\alph*)]
  \item weight constructor,
  \item off-diagonal absorber,
  \item invariant-region verifier, or
  \item restart certificate.
\end{enumerate}
This is useful because it identifies precisely which lemmas contribute to the global
Lyapunov closure and which only support local transport estimates.

\subsection*{7. v\docversion\ readiness criterion}
The present step should be regarded as successful if the following can now be done:
\begin{itemize}
  \item the entire admissible sector can be encoded in one weighted ledger functional;
  \item restart points can be justified by checkpoint data in ledger form;
  \item probe arguments inherit the same invariant-region control as the main proof;
  \item later field equations can be tested against a precise closure criterion; and
  \item appendix tools can be sorted by whether they contribute to local transport or global invariance.
\end{itemize}

This is the gain of v\docversion. The companion document now carries not only staged
closure, channel specialization, and budget transport, but also a genuine
forward-invariant structural regime that stabilizes the proof under continuation,
probe comparison, and later equation-level readout.


\section*{v2.7.0 Field-Energy Realization and Semigroup Readout Core}
\addcontentsline{toc}{section}{v2.7.0 Field-Energy Realization and Semigroup Readout Core}

\begin{quote}
\textbf{Normalization note (v2.17.1 local review).}
This block is retained as a \emph{historical realization and semigroup-readout precursor}. It supports the later field/probe realization layer and should not be read as a second independent primary core.
\end{quote}

The next structural step after forward-invariant ledger closure is to show that the
weighted ledger is not only a bookkeeping device but already the shadow of a genuine
field-energy balance. Version~\docversion\ therefore adds a realization layer: the
channel budgets are interpreted as coarse observables of an underlying evolution
system, and admissibility is reformulated as dissipative semigroup compatibility.
This still does not claim the final field equations, but it fixes the exact kind of
continuous dynamics those equations must generate if they are to remain faithful to
the structural proof architecture.

\subsection*{1. Abstract state space and energy readout}
Let
\[
X := X_{\mathrm{tri}}\oplus X_{\perp}\oplus X_{\mathrm{meas}}\oplus X_{\mathrm{res}}
\]
be an abstract Banach or Hilbert state space carrying the channel-resolved structural
variables. We write a state as
\[
U=(u_{\mathrm{tri}},u_{\perp},u_{\mathrm{meas}},u_{\mathrm{res}}).
\]
The ledger variables are no longer treated as primitive numbers but as readout
functionals
\[
B_i(U)\ge 0,
\qquad i\in\{\mathrm{tri},\perp,\mathrm{meas},\mathrm{res}\},
\]
with the weighted energy
\[
\mathcal E(U):=\alpha_{\mathrm{tri}}B_{\mathrm{tri}}(U)
+\alpha_{\perp}B_{\perp}(U)
+\alpha_{\mathrm{meas}}B_{\mathrm{meas}}(U)
+\alpha_{\mathrm{res}}B_{\mathrm{res}}(U).
\]
The content of this step is that $\mathcal E$ should reproduce the Lyapunov ledger
$\mathcal L$ up to uniform equivalence on the admissible sector. Thus the ledger is
interpreted as a compressed energy observable of the underlying state rather than an
external annotation.

\subsection*{2. Dissipative generator Ansatz}
Assume that the channel evolution is generated by an abstract equation
\[
\partial_\tau U = \mathcal A U + \mathcal N(U) + \mathcal S_{\mathrm{tag}}(\tau),
\]
where $\mathcal A$ is the linear transport-dissipation generator, $\mathcal N$ is a
structured nonlinear coupling compatible with HC confinement, and
$\mathcal S_{\mathrm{tag}}$ collects the tagged residual sources. The admissibility
requirement from the earlier sections now becomes the statement that
$\mathcal A+\mathcal N$ is energy-dissipative modulo tagged production and match
error, namely
\[
\frac{d}{d\tau}\mathcal E(U(\tau))
\le -c_0\mathcal E(U(\tau))
+ C_{\mathrm{tag}}\,\mathrm{Res}_{\mathrm{tag}}(\tau)
+ C_{\mathrm{match}}\,\Delta_{\mathrm{match}}(\tau)
\]
for every admissible trajectory. This is the semigroup form of the weighted ledger
closure derived in v2.6.0.

\subsection*{3. Realization principle}
The structural point can now be isolated as a single principle.

\paragraph{Principle 1 (field-energy realization).}
An eventual field equation is acceptable only if its induced state evolution admits a
channel readout $U\mapsto (B_i(U))_i$ and a weighted energy $\mathcal E$ such that:
\begin{enumerate}[label=(\alph*)]
  \item $\mathcal E$ is uniformly equivalent to the admissible ledger functional $\mathcal L$;
  \item the induced evolution preserves the forward-invariant admissible sector;
  \item all positive production terms are either dissipatively absorbed or explicitly tagged; and
  \item probe mismatch enters only through a controlled source term $\Delta_{\mathrm{match}}$.
\end{enumerate}
Any ansatz failing one of these four points is structurally too weak, even if it
looks formally plausible at equation level.

\subsection*{4. Semigroup consequence and proof memory}
If the above dissipative estimate holds, then the evolution operator $S(\tau)$
generated by the admissible dynamics satisfies an estimate of the form
\[
\mathcal E(S(\tau)U_0)
\le e^{-c_0\tau}\mathcal E(U_0)
+\int_0^\tau e^{-c_0(\tau-s)}
\bigl(C_{\mathrm{tag}}\,\mathrm{Res}_{\mathrm{tag}}(s)
+ C_{\mathrm{match}}\,\Delta_{\mathrm{match}}(s)\bigr)\,ds.
\]
Hence proof memory can now be reformulated in semigroup language: restarting from a
checkpoint is valid because the future evolution depends continuously on the current
energy state and the tagged source history, not on the unrecoverable microscopic path.
This upgrades the earlier restart argument into an abstract continuous-time stability
statement.

\subsection*{5. Probe-sector realization}
The same realization principle applies to the two-probe geometry. If
\[
U_{\mathrm{probe}}=(U^{(1)},U^{(2)},U^{(Q)},U^{(M)},U^{(R)})
\]
denotes the probe state and $\mathcal E_{\mathrm{probe}}$ the corresponding weighted
energy, then admissibility requires
\[
\frac{d}{d\tau}\mathcal E_{\mathrm{probe}}
\le -c_{\mathrm{probe}}\mathcal E_{\mathrm{probe}}
+ C^{\mathrm{probe}}_{\mathrm{tag}}\,\mathrm{Res}^{\mathrm{probe}}_{\mathrm{tag}}
+ C^{\mathrm{probe}}_{\mathrm{match}}\,\Delta_{\mathrm{match}}.
\]
Thus the probe mechanism is now seen as a dissipative comparison flow living in the
same semigroup class as the main derivation. This makes the search geometry more than
a metaphor: it becomes an admissible evolution subsystem with inherited energy
control.

\subsection*{6. Consequence for future field equations and appendix tools}
The field-equation programme is sharpened once again. A candidate PDE or abstract
operator equation must now be checked in the following order:
\begin{enumerate}[label=(\arabic*)]
  \item identify the state space and channel decomposition;
  \item construct the ledger readout functionals $B_i(U)$;
  \item verify equivalence between readout energy $\mathcal E$ and structural ledger $\mathcal L$;
  \item prove dissipativity modulo tagged sources; and
  \item verify probe inheritance and restart compatibility.
\end{enumerate}
Accordingly, appendix tools can receive one further classification layer. A strong
tool may now be marked as a state-space constructor, energy comparator, semigroup
stability lemma, or source-tagging certificate. This helps separate tools that only
support symbolic rearrangement from those that certify genuine dynamic admissibility.

\subsection*{7. v\docversion\ readiness criterion}
The present step should be regarded as successful if the following can now be done:
\begin{itemize}
  \item the ledger is interpretable as a uniformly equivalent field-energy observable;
  \item admissibility is reformulated as dissipative semigroup compatibility;
  \item proof memory is supported by a continuous-time stability estimate;
  \item the two-probe sector is embedded into the same realization framework; and
  \item future field equations can be accepted or rejected by an explicit realization checklist.
\end{itemize}

This is the gain of v\docversion. The companion monolith now carries not only
closure, channel specialization, transport dynamics, and forward invariance, but also
a first abstract field-energy realization principle. The later field-equation layer is
therefore no longer free-form: it has to generate the same dissipative proof memory,
probe inheritance, and tagged-source discipline already fixed at the structural level.


\section*{v2.8.0 Spectral Admissibility and Resolvent Gatekeeping Core}
\addcontentsline{toc}{section}{v2.8.0 Spectral Admissibility and Resolvent Gatekeeping Core}

\begin{quote}
\textbf{Normalization note (v2.17.1 local review).}
This block is retained as a \emph{historical spectral-admissibility precursor}. Its gatekeeping role is preserved, but the normalized admissibility logic is now carried by the semantic/tensor and grouped-layer unfolding phases.
\end{quote}

The semigroup realization from v2.7.0 still leaves one decisive structural question
open: which generators are admissible before one even starts evolving the system? In
this version the answer is sharpened by adding a spectral gatekeeping layer. The
state evolution is not accepted merely because it can be written as an abstract flow;
it must also possess a resolvent and spectrum compatible with the ledger discipline,
HC confinement, tagged-source accounting, and probe inheritance.

\subsection*{1. Spectral admissibility principle}
Let $\mathfrak X$ be the admissible state space from the previous realization step and
let $\mathcal G$ denote the effective generator of the dynamics, including its
controlled tagged-source coupling. We call $\mathcal G$ \emph{spectrally admissible}
if there exist $\omega_*>0$ and $C_*>0$ such that for every $\lambda$ with
$\Re \lambda\ge -\omega_*$ the resolvent exists on the admissible sector and obeys
\[
\| (\lambda-\mathcal G)^{-1} \|_{\mathfrak X\to\mathfrak X}
\le \frac{C_*}{1+|\lambda|},
\]
after projection to the HC-confined channel space. This estimate expresses that the
generator does not hide uncontrolled growth modes behind formal notation. The ledger
control from earlier versions is therefore promoted into a spectral entrance test.

\paragraph{Principle 2 (spectral gatekeeping).}
A candidate field equation or abstract operator model is structurally admissible only
if its effective generator is spectrally admissible in the above sense and the
corresponding resolvent bounds are compatible with the tagged-source decomposition.
The proof programme may therefore reject an operator before any detailed PDE analysis
if the spectrum already violates ledger closure.

\subsection*{2. Resolvent--ledger compatibility}
The weighted energy $\mathcal E$ from v2.7.0 should now be stable not only along the
flow but already at the resolvent level. Concretely, admissibility requires a bound
of the schematic form
\[
\mathcal E\bigl((\lambda-\mathcal G)^{-1}F\bigr)
\le C_{\mathrm{res}}
\Bigl( \mathcal E(F)
+ \mathrm{Tag}(F)
+ \mathrm{Match}(F)\Bigr)
\]
for every admissible forcing term $F$ and every $\Re\lambda\ge -\omega_*$. Here
$\mathrm{Tag}(F)$ denotes the explicitly source-tagged part and
$\mathrm{Match}(F)$ the mismatch component inherited from probe comparison. This is
the spectral analogue of the semigroup energy estimate: the resolvent may amplify
forcing, but only in a way that remains ledger-readable.

\subsection*{3. HC confinement as spectral projector}
The HC mechanism is now reinterpreted one step more sharply. Earlier versions treated
it as the structural compensator that confines evolution to the septinion closure.
At spectral level this becomes the statement that there exists a projector
$\Pi_{\mathrm{HC}}$ with
\[
(\lambda-\mathcal G)^{-1}
= \Pi_{\mathrm{HC}}(\lambda-\mathcal G)^{-1}\Pi_{\mathrm{HC}}
+ \mathcal R_{\mathrm{tag}}(\lambda),
\]
where $\mathcal R_{\mathrm{tag}}(\lambda)$ is either negligible in the admissible
regime or explicitly source-tagged. Hence uncontrolled extra channels are not merely
suppressed dynamically; they are prevented from entering the admissible resolvent
class except through a tagged remainder. This gives a clean spectral formulation of
HC confinement.

\subsection*{4. Spectral gap and restart robustness}
If the admissible generator satisfies the above bounds together with a positive gap
between the dominant admissible spectrum and the forbidden sector, then restart
stability acquires a stronger meaning. A checkpoint is not only dynamically safe but
spectrally robust: small reconstruction errors remain trapped inside the same
resolvent class and do not move the system into a hidden unstable branch. In this
sense proof memory is upgraded from semigroup continuity to spectral persistence.

\subsection*{5. Probe-sector comparison spectrum}
The two-probe subsystem inherits the same criterion. Let
$\mathcal G_{\mathrm{probe}}$ be the effective probe generator acting on the probe
state space. Then admissibility requires a comparison estimate
\[
\| (\lambda-\mathcal G_{\mathrm{probe}})^{-1} \| 
\le C_{\mathrm{probe,res}}(1+|\lambda|)^{-1},
\qquad \Re\lambda\ge -\omega_{\mathrm{probe}},
\]
with constants controlled by the same ledger hierarchy and mismatch tags. The search
geometry therefore becomes spectrally testable: a probe construction is not accepted
merely because it heuristically mirrors the main evolution, but because its spectrum
remains subordinate to the same admissible comparison regime.

\subsection*{6. Consequence for future field equations and appendix tools}
The field-equation admission checklist is sharpened again. A candidate model should
now be tested in the following order:
\begin{enumerate}[label=(\arabic*)]
  \item identify state space and admissible channel projector;
  \item construct the energy and ledger readout;
  \item verify semigroup dissipativity modulo tagged sources;
  \item prove spectral admissibility and resolvent compatibility; and
  \item verify probe-spectrum inheritance and restart robustness.
\end{enumerate}
This also refines the appendix tools. Strong tools may now additionally be marked as
resolvent estimator, spectral projector certificate, gap witness, or restart
persistence lemma. These labels distinguish genuinely admissibility-producing tools
from merely formal symbolic aids.

\subsection*{7. v\docversion\ readiness criterion}
The present step should be regarded as successful if the following can now be done:
\begin{itemize}
  \item admissible generators can be filtered by a spectral entrance test;
  \item resolvent bounds can be read in the same ledger language as the semigroup flow;
  \item HC confinement is expressible as a spectral projector statement;
  \item restart stability is strengthened to spectral persistence; and
  \item the probe sector inherits a controlled comparison spectrum.
\end{itemize}

This is the gain of v\docversion. The companion monolith now carries not only
closure, channel specialization, transport dynamics, forward invariance, and
field-energy realization, but also a first spectral admissibility layer. Future
field equations are therefore constrained twice: they must generate the correct
ledger-dissipative semigroup and they must belong to the correct resolvent class
before they are even admitted into the structural programme.


\section*{v2.9.0 Spectral Readout Filters and Trace-Compatible Mode Selection Core}
\addcontentsline{toc}{section}{v2.9.0 Spectral Readout Filters and Trace-Compatible Mode Selection Core}

\begin{quote}
\textbf{Normalization note (v2.17.1 local review).}
This block is retained as a \emph{historical spectral-readout and trace-filter precursor}. It contributes proof memory and specialized readout discipline, but should be read through the later normalized readout and interface phases.
\end{quote}

The spectral admissibility layer from v2.8.0 decides which generators may enter the
programme at all. The next structural question is subtler: once a generator is
admissible, which parts of its spectrum are actually readable in a way compatible
with the ledger, HC confinement, probe comparison, and later trace-style summary
principles? In this version the monolith therefore adds a spectral readout layer.
Not every admissible mode is automatically promoted to a structural observable. A
mode must also survive a filter test that distinguishes stable, tagged, and
non-readable contributions.

\subsection*{1. Readout filter principle}
Let $\mathcal G$ be spectrally admissible on the HC-confined state space
$\mathfrak X$, and let $\Sigma_{\mathrm{adm}}(\mathcal G)$ denote the admissible part
of the spectrum. We introduce a family of bounded readout filters
$\{\mathcal F_\alpha\}_{\alpha\in A}$ acting on the admissible sector, where each
$\mathcal F_\alpha$ is designed to isolate one structural mode class: triadic,
orthogonal, measurement, residual, or probe-induced. A mode contribution is called
\emph{trace-compatible} if there exists at least one filter $\mathcal F_\alpha$ such
that
\[
\mathcal F_\alpha \Pi_{\mathrm{HC}} u = \Pi_{\mathrm{HC}} u + r_\alpha,
\qquad
\|r_\alpha\|_{\mathfrak X} \le \varepsilon_\alpha \|u\|_{\mathfrak X},
\]
with either small remainder $r_\alpha$ or explicitly tagged remainder. Thus the
readout is not merely spectral isolation but structural readability modulo tagged
error.

\paragraph{Principle 3 (filtered observability).}
A spectral component may be used in the proof architecture only if it is both
spectrally admissible and trace-compatible under an approved readout filter family.
This prevents uncontrolled spectral clutter from being mistaken for genuine
structural information.

\subsection*{2. Trace-compatible decomposition}
The admissible spectrum should now admit a decomposition of the schematic form
\[
\Sigma_{\mathrm{adm}}(\mathcal G)
= \Sigma_{\mathrm{read}} \cup \Sigma_{\mathrm{tag}} \cup \Sigma_{\mathrm{mute}},
\]
where readable modes lie in $\Sigma_{\mathrm{read}}$, tagged but non-primary modes lie
in $\Sigma_{\mathrm{tag}}$, and the remainder $\Sigma_{\mathrm{mute}}$ is admissible but
not allowed to enter the principal structural trace. The point is not that mute
modes vanish, but that they cannot serve as primary carriers of theorem-level
readout. This sharpens the difference between existence and interpretability.

\subsection*{3. HC-filter compatibility}
The HC projector from v2.8.0 is now required to commute with readout at least up to
controlled tagged remainder. Concretely, one seeks estimates of the form
\[
\mathcal F_\alpha \Pi_{\mathrm{HC}}
= \Pi_{\mathrm{HC}} \mathcal F_\alpha + \mathcal T_\alpha,
\]
where $\mathcal T_\alpha$ is either small in the admissible norm or belongs to a
source-tagged correction class. This means HC confinement is preserved under mode
selection: filtering is not allowed to re-open forbidden channels except through a
ledger-visible correction.

\subsection*{4. Mode budgets and readable trace mass}
The budget language from v2.5.0 now acquires a spectral refinement. For every
approved filter $\mathcal F_\alpha$ define the readable trace mass
\[
\mathfrak M_\alpha(u) := \mathcal E(\mathcal F_\alpha u)
\]
with $\mathcal E$ the field-energy functional from v2.7.0. Structural readability
requires a balance inequality of the schematic form
\[
\sum_{\alpha\in A_{\mathrm{read}}} \mathfrak M_\alpha(u)
\le C_{\mathrm{tr}}
\Bigl( \mathcal E(u) + \mathrm{Tag}(u) + \mathrm{Leak}(u) \Bigr).
\]
Hence readable spectral mass cannot exceed the ledger budget except through already
visible tagged or leakage terms. This gives a quantitative bridge between spectral
filtering and earlier channel transport estimates.

\subsection*{5. Probe-sector mode inheritance}
For the two-probe subsystem the same distinction becomes a comparison principle.
Let $\mathcal G_{\mathrm{probe}}$ be the effective probe generator and
$\mathcal F^{\mathrm{probe}}_\beta$ its admissible filter family. Then every readable
probe mode should either descend from a readable ambient mode or be marked as a
probe-only tagged correction. In particular, one asks for comparison bounds of the
form
\[
\mathfrak M^{\mathrm{probe}}_\beta(v)
\le C_{\mathrm{cmp}}
\bigl( \mathfrak M_\alpha(u) + \mathrm{Mismatch}(u,v) \bigr),
\]
for matched ambient/probe states $(u,v)$. This prevents the probe sector from
manufacturing fake observables that have no ambient structural parent.

\subsection*{6. Consequence for future field equations and appendix tools}
The admission checklist for candidate field equations is now extended once more:
\begin{enumerate}[label=(\arabic*)]
  \item verify semigroup dissipativity and ledger closure;
  \item verify spectral admissibility and resolvent compatibility;
  \item construct an HC-compatible readout filter family;
  \item separate readable, tagged, and mute spectral sectors; and
  \item prove probe inheritance of readable modes.
\end{enumerate}
The appendix tool taxonomy is correspondingly refined. Besides resolvent and gap
certificates, a strong tool may now also be labelled as readout filter constructor,
trace-mass comparator, mute-sector exclusion lemma, or probe-inheritance witness.
These labels indicate that the tool produces not only existence information but also
structural legibility.

\subsection*{7. v\docversion\ readiness criterion}
The present step should be regarded as successful if the following can now be done:
\begin{itemize}
  \item admissible spectra can be split into readable, tagged, and mute parts;
  \item HC confinement survives mode filtering modulo controlled tagged remainder;
  \item readable trace mass is bounded in the same ledger language as channel budgets;
  \item probe readout inherits ambient readable modes rather than inventing new ones; and
  \item future field equations can be screened not only for existence, but for structural legibility.
\end{itemize}

This is the gain of v\docversion. The companion monolith now carries not only
closure, channel specialization, transport dynamics, forward invariance,
field-energy realization, semigroup readout, and spectral admissibility, but also a
first filtered-observability layer. Future field equations are therefore constrained
three times: they must generate the correct ledger-dissipative semigroup, they must
belong to the correct resolvent class, and their admissible modes must pass a
trace-compatible readout filter before they can count as structural content.


\section*{v2.10.0 Trace Reconstruction and Observable Extraction Core}
\addcontentsline{toc}{section}{v2.10.0 Trace Reconstruction and Observable Extraction Core}

\begin{quote}
\textbf{Normalization note (v2.17.1 local review).}
This block is retained as a \emph{historical reconstruction and observable-extraction precursor}. Its reconstruction language remains useful, but the normalized proof-bearing architecture now lies in the unfolding phases and the grouped transport logic.
\end{quote}

The filtered-observability layer from v2.9.0 separates admissible modes into
readable, tagged, and mute sectors. The next structural task is no longer only to
screen admissible modes, but to reconstruct theorem-level observables from the
readable part without double counting, without reopening forbidden channels, and
without losing the ledger discipline enforced by HC confinement. The present step
therefore adds a trace-reconstruction layer. It answers the question: once a mode
has passed admissibility and readability tests, how is it promoted to an actual
observable in the companion proof architecture?

\subsection*{1. Reconstruction operator for readable modes}
Let $\{\mathcal F_\alpha\}_{\alpha\in A_{\mathrm{read}}}$ be the approved readable
filter family from v2.9.0. We now introduce a reconstruction operator
\[
\mathcal R: \bigoplus_{\alpha\in A_{\mathrm{read}}} \mathfrak X_\alpha \to \mathfrak O,
\]
where $\mathfrak X_\alpha = \mathcal F_\alpha \mathfrak X$ denotes the readable mode
sector and $\mathfrak O$ is the structural observable space. The role of
$\mathcal R$ is not arbitrary synthesis: it is constrained reconstruction. For a
state $u\in\mathfrak X$ one seeks a decomposition
\[
\mathrm{Obs}(u) = \mathcal R\bigl((\mathcal F_\alpha u)_{\alpha\in A_{\mathrm{read}}}\bigr)
\]
such that every observable component remains HC-confined, ledger-bounded, and
source-tagged whenever the reconstruction is not exact. Thus observable extraction
is treated as a disciplined readout map rather than an unrestricted projection.

\paragraph{Principle 4 (reconstructive observability).}
A readable mode contributes to theorem-level content only after passing through a
reconstruction operator whose output is compatible with HC confinement, channel
budgets, and source tags. Readability alone is therefore necessary but not yet
sufficient for observable status.

\subsection*{2. Observable extraction at theorem level}
For each admissible theorem-context $\Theta$ let $\mathfrak O_\Theta$ denote the class
of observables allowed in that context. One then requires extraction estimates of
schematic form
\[
\|\mathcal R_\Theta((\mathcal F_\alpha u)_\alpha)\|_{\mathfrak O_\Theta}
\le C_\Theta
\Bigl(
\sum_{\alpha\in A_{\mathrm{read}}} \mathfrak M_\alpha(u)
+ \mathrm{Tag}(u)+\mathrm{Leak}(u)
\Bigr),
\]
with $\mathfrak M_\alpha$ the readable trace masses from v2.9.0. This means an
observable is only legitimate when its theorem-level norm is controlled by the same
ledger quantities already used to regulate channel transport and spectral readout.
The companion proof therefore acquires an explicit extraction stage between readable
modes and theorem statements.

\subsection*{3. Overlap ledger and anti-double-counting rule}
Once multiple filters contribute to one observable quantity, overlap must be made
visible. Introduce an overlap ledger
\[
\mathrm{OL}(u)= \sum_{\alpha\neq\beta}
\omega_{\alpha\beta}\,\langle \mathcal F_\alpha u,\mathcal F_\beta u\rangle_{\mathfrak X},
\]
where the coefficients $\omega_{\alpha\beta}$ encode admissible correlation or
subtraction rules. The purpose of $\mathrm{OL}(u)$ is not to forbid overlap but to
prevent silent duplication of structural mass. Observable reconstruction is required
to satisfy a compensated estimate of the form
\[
\|\mathrm{Obs}(u)\|_{\mathfrak O}
\le C_{\mathrm{obs}}
\Bigl(
\sum_{\alpha\in A_{\mathrm{read}}} \mathfrak M_\alpha(u)
- \mathrm{OL}(u)
+ \mathrm{Tag}(u)+\mathrm{Leak}(u)
\Bigr),
\]
whenever the overlap ledger is well defined. Thus one can read several filtered
channels at once without pretending they were disjoint when they are not.

\subsection*{4. HC-stable observable closure}
The reconstructed observable sector must itself remain HC-stable. Concretely, one
requires that
\[
\Pi_{\mathrm{HC}}\,\mathrm{Obs}(u)=\mathrm{Obs}(u)+\tau(u),
\qquad
\|\tau(u)\|_{\mathfrak O}\le C_{\mathrm{HC}}
\bigl(\mathrm{Tag}(u)+\mathrm{Leak}(u)\bigr),
\]
so that reconstruction does not reactivate forbidden directions except through
ledger-visible correction terms. In this sense HC confinement becomes not only a
state-space condition and not only a filter-compatibility condition, but now also an
observable-closure condition. Only outputs that survive this test may be used as
stable structural readings.

\subsection*{5. Probe-sector reconstruction inheritance}
For the two-probe subsystem the same reconstruction principle must descend to the
probe readout. Let $\mathcal R^{\mathrm{probe}}$ be the probe reconstruction operator.
Then every probe observable should either be inherited from an ambient observable
through the descent chain or be explicitly tagged as probe-local correction. One
asks schematically for bounds of the form
\[
\|\mathrm{Obs}^{\mathrm{probe}}(v)\|_{\mathfrak O_{\mathrm{probe}}}
\le C_{\mathrm{pr}}
\bigl(
\|\mathrm{Obs}(u)\|_{\mathfrak O}
+\mathrm{Mismatch}(u,v)
\bigr),
\]
for matched ambient/probe states $(u,v)$. This ensures the probe sector remains an
instrument of structural descent rather than an independent source of untracked
observables.

\subsection*{6. Consequence for field equations and appendix tools}
The admission checklist for future candidate field equations is therefore extended
once more:
\begin{enumerate}[label=(\arabic*)]
  \item verify semigroup dissipativity and ledger closure;
  \item verify spectral admissibility and resolvent compatibility;
  \item verify readable mode separation under HC-compatible filters;
  \item construct a reconstruction operator for theorem-level observables; and
  \item bound overlap, leakage, and probe descent in a visible ledger form.
\end{enumerate}
The appendix tool taxonomy is refined accordingly. Besides resolvent estimators,
readout filters, and gap witnesses, a strong tool may now also be labelled as
reconstruction operator, observable extractor, overlap-ledger bound, HC-stability
witness, or probe-descent certificate. These labels mark tools that do not merely
prove admissibility, but actually promote admissible readable content into stable
observables.

\subsection*{7. v\docversion\ readiness criterion}
The present step should be regarded as successful if the following can now be done:
\begin{itemize}
  \item readable spectral content can be reconstructed into theorem-level observables;
  \item overlaps between readable channels are accounted for instead of silently duplicated;
  \item observable extraction remains HC-stable modulo visible tagged/leakage correction;
  \item probe observables descend from ambient observables rather than appearing autonomously; and
  \item future field equations can be screened for observable reconstruction, not merely for readable spectra.
\end{itemize}

This is the gain of v\docversion. The companion monolith now carries not only
closure, channel specialization, transport dynamics, forward invariance,
field-energy realization, semigroup readout, spectral admissibility, and filtered
observability, but also a first disciplined reconstruction layer. Future field
equations are therefore constrained four times: they must generate the correct
ledger-dissipative semigroup, belong to the correct resolvent class, pass a
trace-compatible readout filter, and finally admit HC-stable observable
reconstruction before they can count as theorem-level structural content.


\section*{v2.11.0 Multi-Readout Consistency and Fusion Core}
\addcontentsline{toc}{section}{v2.11.0 Multi-Readout Consistency and Fusion Core}

\begin{quote}
\textbf{Normalization note (v2.17.1 local review).}
This block is retained as a \emph{historical multi-readout and fusion precursor}. It remains relevant as supporting proof memory, but no longer functions as an independent primary structural core in the normalized architecture.
\end{quote}

The reconstruction layer from v2.10.0 promotes readable modes into theorem-level
observables. The next structural question is what happens when several
reconstructions are all individually admissible but are intended to describe one
and the same structural content. The task of v\docversion\ is therefore not merely
to extract observables, but to fuse them without creating contradictory readouts,
without hiding overlap behind notation, and without allowing probe-local views to
break the ambient closure architecture. This introduces a multi-readout
consistency layer.

\subsection*{1. Readout families and consistency targets}
Let $\{\mathcal R_j\}_{j\in J}$ be a family of admissible reconstruction maps,
possibly associated with different filter systems, theorem contexts, or probe
resolutions. For a given state $u$ one obtains a readout family
\[
\mathrm{Obs}_j(u)=\mathcal R_j\bigl((\mathcal F^{(j)}_\alpha u)_\alpha\bigr).
\]
The purpose of the new layer is to decide when the collection
$\{\mathrm{Obs}_j(u)\}_{j\in J}$ should be regarded as several compatible views of one
structural object rather than several unrelated observables. Compatibility is
measured against the same closure, budget, and tagging rules that already govern
the single-readout setting.

\paragraph{Principle 5 (multi-readout consistency).}
A family of theorem-level observables may be fused only when their mutual
discrepancies are either controlled by a visible ledger term or assigned to
explicit source tags. Silent disagreement is not an allowed mode of structural
multiplicity.

\subsection*{2. Consistency ledger and fusion defect}
Introduce a fusion-defect functional
\[
\mathfrak D_{\mathrm{fus}}(u)
:= \sum_{j<k}\nu_{jk}\,\|\mathrm{Obs}_j(u)-\Gamma_{jk}\mathrm{Obs}_k(u)\|_{\mathfrak O_*},
\]
where $\Gamma_{jk}$ transports the $k$-th observable into the comparison frame of the
$j$-th observable and the coefficients $\nu_{jk}\ge 0$ encode structural relevance.
The associated consistency ledger is then
\[
\mathrm{CL}(u):=\mathfrak D_{\mathrm{fus}}(u)+\mathrm{Tag}(u)+\mathrm{Leak}(u).
\]
A family is called fusion-admissible when $\mathrm{CL}(u)$ stays below the threshold
prescribed by the theorem context. Thus disagreement is not denied; it is made
visible, budgeted, and therefore manageable.

\subsection*{3. Fused observable and anti-inflation rule}
Whenever the family is fusion-admissible, define a fused observable
\[
\mathrm{Obs}_{\mathrm{fus}}(u)=\sum_{j\in J}\lambda_j\,\Gamma_j\mathrm{Obs}_j(u)-\Xi(u),
\]
where the coefficients $\lambda_j$ are admissible weights, $\Gamma_j$ transports to a
common observable frame, and $\Xi(u)$ is an anti-inflation compensator correcting
for visible overlap and recorded fusion defect. The fused observable is required
to satisfy an estimate of the form
\[
\|\mathrm{Obs}_{\mathrm{fus}}(u)\|_{\mathfrak O_*}
\le C_{\mathrm{fus}}
\bigl(
\sum_{j\in J}\|\mathrm{Obs}_j(u)\|_{\mathfrak O_j}
-\mathrm{OL}_{\mathrm{fus}}(u)
+\mathrm{CL}(u)
\bigr),
\]
so that fusion cannot enlarge theorem-level content by counting the same readable
structure several times through nominally different channels.

\subsection*{4. HC-stable fusion and theorem transport}
The fused observable must still respect HC confinement. One therefore asks for
\[
\Pi_{\mathrm{HC}}\,\mathrm{Obs}_{\mathrm{fus}}(u)
=\mathrm{Obs}_{\mathrm{fus}}(u)+\rho(u),
\qquad
\|\rho(u)\|_{\mathfrak O_*}
\le C_{\rho}\,\mathrm{CL}(u).
\]
This means that multi-readout fusion is accepted only if the residual mismatch is
fully visible in the same ledger that already recorded comparison defects. In
addition, theorem transport must preserve fusion status: if a fused observable is
moved into a lemma, table, or appendix tool, the same fusion certificate must
remain available after transfer.

\subsection*{5. Probe/ambient coherence under fusion}
In the two-probe sector one may have ambient and probe reconstructions of the same
structural feature. The new requirement is that the probe fusion defect be bounded
by ambient fusion defect plus descent mismatch:
\[
\mathfrak D^{\mathrm{probe}}_{\mathrm{fus}}(v)
\le C_{\mathrm{coh}}
\bigl(
\mathfrak D_{\mathrm{fus}}(u)+\mathrm{Mismatch}(u,v)
\bigr).
\]
Hence probe readouts do not become a second autonomous truth layer; they inherit
coherence from the ambient structure, up to visible descent error.

\subsection*{6. Consequence for appendix tools}
The appendix tool taxonomy is refined again. A strong appendix item may now be
labelled as a fusion certificate, consistency-ledger estimate, anti-inflation
compensator, theorem-transport witness, or probe-coherence bound. This fits the
user's structural preference that tools be collected as reusable proof interfaces
rather than scattered informally through the running text.

\subsection*{7. v\docversion\ readiness criterion}
The present step should be regarded as successful if the following can now be
done:
\begin{itemize}
  \item several admissible reconstructions can be compared in one explicit ledger;
  \item disagreements between readouts are tagged or budgeted instead of hidden;
  \item fused observables remain protected against overlap inflation;
  \item probe and ambient readouts satisfy a visible coherence estimate; and
  \item appendix transfer preserves not only observables, but also their fusion certificates.
\end{itemize}

This is the gain of v\docversion. The companion monolith now carries not only
closure, channel specialization, transport dynamics, forward invariance,
field-energy realization, semigroup readout, spectral admissibility, filtered
observability, and trace reconstruction, but also a first disciplined fusion
layer for theorem-level observables. Future field equations are therefore
constrained five times: they must generate the correct ledger-dissipative
semigroup, belong to the correct resolvent class, pass a trace-compatible readout
filter, admit HC-stable observable reconstruction, and finally support
multi-readout fusion without hidden inconsistency.


% ------------------------------------------------------------------
% Appendix: Structural Tools and Regime Tests (v2.27.4 integration)
% ------------------------------------------------------------------

\section{Structural Tools and Regime Tests}
\label{app:structural-tools-regime-tests}
\addcontentsline{toc}{section}{Appendix: Structural Tools and Regime Tests}

This appendix collects the main structural tools extracted from the
proof architecture. Their purpose is not to replace the main
derivations, but to normalize frequently recurring arguments into
reusable tests and reduction principles.

\subsection*{Sector and closure tools}

\paragraph{Tool (Closure-transfer principle).}
Let $X$ be any operator on the septonic probe space. Its admissible
reduction is
\[
  X_{\mathrm{adm}} := \hat{P}_{C_7} X \hat{P}_{C_7}.
\]
For admissible states $\Psi \in \operatorname{Ran}(\hat{P}_{C_7})$, the
physically relevant action of $X$ is determined by $X_{\mathrm{adm}}$,
while the complementary blocks measure off-sector leakage.

\paragraph{Tool (Equivalent closure criteria).}
The following are equivalent
(cf.\ Lemma~\texttt{lem:closure-forcing-lemma} in the active-core PDF):
\begin{enumerate}[label=(\roman*)]
  \item $\mathfrak{R}_C = 0$;
  \item $[\mathcal{M}_{16}, \hat{P}_{C_7}] = 0$;
  \item $\mathcal{M}_{16}$ preserves the admissible sector and its complement.
\end{enumerate}
Hence exact closure is equivalent to block preservation, i.e. block-diagonalization of $\mathcal{M}_{16}$ with respect to
$\mathcal{H}_{\mathrm{sept}} = \hat{P}_{C_7}\mathcal{H}_{\mathrm{sept}}
\oplus (\mathbf{1} - \hat{P}_{C_7})\mathcal{H}_{\mathrm{sept}}$.

\paragraph{Tool (Closure-preservation test).}
For any new structural term $X$, first test whether
$[X, \hat{P}_{C_7}] = 0$.
If yes, then $X$ respects the admissible stabilized sector directly.
If not, then $X$ either produces leakage or must be controlled by a
compensator-mediated defect term
(cf.\ Hypothesis~\texttt{hyp:frame-dependent-defect-readout} in the active-core PDF).

\paragraph{Tool (Sectorial normal form).}
Every operator $X$ decomposes as
\[
  X = \hat{P}_{C_7}X\hat{P}_{C_7}
    + (\mathbf{1}-\hat{P}_{C_7})X\hat{P}_{C_7}
    + \hat{P}_{C_7}X(\mathbf{1}-\hat{P}_{C_7})
    + (\mathbf{1}-\hat{P}_{C_7})X(\mathbf{1}-\hat{P}_{C_7}),
\]
isolating: (i) admissible core block, (ii) incoming leakage,
(iii) outgoing leakage, (iv) complement-sector block.

\subsection*{Master-equation tools}

\paragraph{Tool (Core-defect normal form).}
The projected structural field operator
$\mathfrak{F}_{\mathrm{struct}}
= \hat{P}_{C_7}(i\Gamma^\mu\mathcal{D}_\mu - \Omega)\hat{P}_{C_7}$
splits into
\[
  \mathfrak{F}_{\mathrm{struct}}
  = \mathbb{D}_{\mathrm{core}} + \mathbb{D}_{\mathrm{defect}},
\]
where $\mathbb{D}_{\mathrm{core}}$ contains autonomous admissible
transport and resonance, and $\mathbb{D}_{\mathrm{defect}}$ collects
closure mismatch, compensator activity, and off-sector suppression.

\paragraph{Tool (Exact-closure normal form).}
If $\mathfrak{R}_C = 0$, $[\mathcal{M}_{16},\hat{P}_{C_7}] = 0$, and
$\mathcal{H}_C = 0$ for the minimal compensator ansatz, then
$\mathfrak{F}_{\mathrm{struct}} = \mathbb{D}_{\mathrm{core}}$ and
$\mathbb{D}_{\mathrm{defect}} = 0$.
The visible compensator field vanishes in this regime; the kernel
$C = \Pi^{(\mathcal{M})}_{\mathrm{spec}}$
(Theorem~\texttt{thm:hc-spectral-projector} in the active-core PDF) remains active
(Remark~\texttt{rem:hc-shadow-vs-core} in the active-core PDF).

\paragraph{Tool (Mass-residue test).}
A proposed mass contribution is classified according to whether it
arises from: (i) a closed-core scalar reading of the projected resonance
operator; (ii) a weak closure-defect correction; (iii) a genuine
shifted-clock or projector residue
$\mathfrak{m}_{\mathrm{comm}}(\psi)
= \frac{\hbar}{c}|\langle\psi|[\hat{\mathcal{R}}_\beta,\hat{P}_{C_7}]|\psi\rangle|$.

\subsection*{Projection and regime tools}

\paragraph{Tool (Regime-strength test).}
A reduction of the structural master equation is: \emph{strong} if
exact closure plus direct slot restriction suffices; \emph{intermediate}
if projected algebra closure is additionally needed; \emph{effective} if
bilinear reconstruction, coarse-graining, or emergent parameter
generation are also required.
These levels correspond to the three steps of
Proposition~\texttt{prop:structural-lift-principle} in the active-core PDF.

\paragraph{Tool (Dirac gate).}
A Dirac-type reduction is available if: (i) exact closure holds;
(ii) the projected connection freezes locally to a trivial sector;
(iii) the projected resonance operator reduces to a scalar inertial
reading.

\paragraph{Tool (Lie-closure gate).}
A Yang--Mills-type reduction is available only if the projected
connection algebra closes on a finite-dimensional Lie algebra
$\mathfrak{g}$. Without that closure, the gauge interpretation is
structurally premature.

\paragraph{Tool (Einstein pipeline gate).}
An Einstein-type reduction requires all of: (i) exact closure on the
admissible sector; (ii) bilinear one-forms $E^a{}_\mu$ from projected
probe dynamics; (iii) maximal rank of the effective coframe; (iv)
metric-compatible geometric connection; (v) coarse-grained local
effective action with Einstein--Hilbert leading term. If any step is
missing, the Einstein block remains effective rather than
first-principles
(cf.\ Proposition~\texttt{prop:expected-probe-lensing-behaviour} in the active-core PDF).

\subsection*{Equivalence and anti-redundancy tools}

\paragraph{Tool (Stone--von Neumann compression rule).}
If two sectors are irreducible unitary realizations of the same
Weyl-type algebra on separable Hilbert spaces, differing only by
representation change, symmetry implementation, or frame reading, then
one representative suffices
(Theorem~\texttt{thm:unitary-equiv} in the active-core PDF).

\paragraph{Tool (Stone--von Neumann scope test).}
Do \emph{not} apply Stone--von Neumann as a shortcut if the passage
between two sectors involves: (i) controlled projection discarding
slots; (ii) coarse-graining over internal modes; (iii) bilinear
geometric reconstruction; (iv) emergent effective constants; (v) loss
of irreducibility. Stone--von Neumann compresses representation
multiplicity, not regime multiplicity.

\paragraph{Tool (Two-layer equivalence test).}
For candidate sectors $S_1, S_2$: (1)~\emph{Representation layer} ---
test whether they are unitary realizations of the same irreducible
Weyl-type kernel; if yes, SvN-compressible. (2)~\emph{Regime layer} ---
test whether one sector is obtained from the other by projection,
coarse-graining, geometric readout, or effective reconstruction; if yes,
additional structural work is required beyond Stone--von Neumann
(Definition~\texttt{def:readout-relative-lemma} in the active-core PDF).

\subsection*{Geometry-readout tools}

\paragraph{Tool (Coframe extraction gate).}
Given a closed triadic probe field $\Psi_\triangle$, define
$E^a{}_\mu := \operatorname{Re}(\bar\Psi_\triangle\,\Gamma^a\,
\mathcal{D}_\mu\Psi_\triangle)$.
This defines an effective coframe only if the resulting $4\times 4$
matrix has maximal rank.

\paragraph{Tool (Closure-to-geometry cleanliness test).}
If exact closure holds, probe bilinears are computed entirely inside
the invariant $C_7$-sector, so reconstructed metric data is not
contaminated by leakage terms. This is the minimal cleanliness
criterion for any geometric readout.

\paragraph{Tool (Geometry-reconstruction warning).}
Closure by itself does not yet produce Einstein geometry. It provides:
(i) an autonomous admissible sector; (ii) a well-defined projected
connection; (iii) a stable basis for bilinear readout. It does
\emph{not} yet provide automatically: (i) nondegeneracy of the
coframe; (ii) derivation of the effective constants; (iii) explicit
identification of the collective closure fields; (iv) justified
coarse-graining to a local geometric action. The bidirectional probe
lensing hypothesis (Section~\texttt{sec:bidirectional-probe-lensing} in the active-core PDF)
proposes a two-step route that may reduce this gap.

% Appendix mode is already active above; the following block continues the appendix material without re-entering appendix mode.

\section*{Appendix B. Symbol Normalization for the Late Readout Layers}
\addcontentsline{toc}{section}{Appendix B. Symbol Normalization for the Late Readout Layers}

This appendix fixes the notation used in the late Companion expansion layers so that
historical version labels and structural symbols remain auditable across the v2.4.0--v2.11.0
chain.

\paragraph{Historical section labels.}
The blocks entitled v2.4.0 through v2.11.0 are to be read as historical expansion
layers. Their section headings record the version in which the corresponding layer was
introduced, not merely the current document version.

\paragraph{Ledger matrix.}
The channel ledger matrix is denoted by
\[
\Lambda_{\mathrm{ch}}(\Psi)
=
\begin{pmatrix}
\ell_{\mathrm{tri}} & \delta_{\mathrm{tri}} & r_{\mathrm{tri}}\\
\ell_{\perp} & \delta_{\perp} & r_{\perp}\\
\ell_{\mathrm{meas}} & \delta_{\mathrm{meas}} & r_{\mathrm{meas}}\\
\ell_{\mathrm{res}} & \delta_{\mathrm{res}} & r_{\mathrm{res}}
\end{pmatrix}.
\]
It is a bookkeeping object for channel-resolved tangential, normal, and compensator-facing
sizes.

\paragraph{Lyapunov ledger functional.}
The weighted scalar functional introduced in v2.6.0 keeps the notation
\(\mathcal L\). It is a Lyapunov-type control quantity and is not to be confused with the
matrix-valued ledger \(\Lambda_{\mathrm{ch}}\).

\paragraph{Budget totals.}
The aggregated scalar budgets remain
\(\mathfrak B\) for total visible load, \(\mathfrak N\) for normal leakage, and
\(\mathfrak R\) for compensator-facing storage.

\paragraph{Readout and fusion objects.}
The readable-mode trace mass, reconstruction operators, observable extraction maps,
and fusion-defect quantities introduced in v2.9.0--v2.11.0 are interpreted as
structure-preserving descendants of the ledger/budget layer and therefore inherit the
normalization fixed here.



\subsection*{Session-budget rule}
A further practical restriction is useful at the present stage: each local mathematics-first pass should also observe a bounded session budget. The purpose of this rule is not to slow work down, but to prevent a single pass from silently expanding from a targeted closure attempt into a diffuse rewrite wave. In concrete terms, a local subversion should normally do only one of the following as its primary task: (i) stabilize definitions for one A-priority block, (ii) discharge one dependency layer for that block, (iii) perform one focused reference/notation repair required by that block, or (iv) close one clearly delimited appendix-tool support step. If more than one such task would be required, the safer action is to end the current pass, classify the result, and continue by explicit next local subversion.

The session-budget rule therefore complements the anti-parallelization and re-entry rules. A pass is considered well-scoped if its main closure claim can be summarized in one sentence and audited against one active block. A pass should be considered over budget if it begins to change multiple theorem-status boundaries, broad notation regimes, or several major sections at once. In that case, the correct response is not to continue accumulating edits, but to stop, audit, and reopen the next step as a fresh local subversion. This rule is intentionally conservative: its goal is to make every local version legible as a finite mathematical move rather than as an opaque mass of partially mixed intentions.



\subsection*{Archived final release-decision template for v2.27.0}
At this stage the document no longer needs another diffuse local steering note, but a compact decision template that makes explicit under which reading the next public release is justified. The present template distinguishes three cases. In the first, the document is strong enough that the next action is \emph{release now}: this applies if the A-priority blocks are structurally frozen, the gate conditions remain green, and the remaining work consists only of optional proof-density enrichment. In the second, the document should \emph{hold locally for proof-density enrichment}: this applies if the structural spine is already in place, but one deliberately wants a denser local unfolding before exposing the next public release. In the third, the document must \emph{hold locally for unresolved blocker}: this applies only if a genuine release-critical dependency, hypothesis-control issue, or reference-instability problem is still active.

Under the current reading, the first A-block no longer counts as such a blocker. Its local dependency spine has been extended through admissible transport, matrix-sector realization, and admissible trace/branch persistence, and it has already been reconciled with the mathematics-first gate logic. Hence the residual caution is no longer structural incompleteness, but only whether one prefers to accept the present medium proof density or to invest one further local pass in additional unfolding. Accordingly, the document is now close to the decision boundary where a direct step to v2.27.0 becomes reasonable.

\paragraph{Decision states.}
\begin{itemize}[leftmargin=2em]
  \item \textbf{Release now:} the first A-block is structurally release-supporting, no primary blocker remains, and only optional densification is deferred.
  \item \textbf{Hold locally for proof-density enrichment:} no structural blocker remains, but one more local pass is preferred for denser exposition.
  \item \textbf{Hold locally for unresolved blocker:} a genuine reference, dependency, or hypothesis-control failure is still present.
\end{itemize}

\paragraph{Current reading.}
At the present local stage, the third state no longer appears to apply to the first A-block. The second and first states are now the relevant ones: either one accepts the current medium-density closure as sufficient for the next public release, or one performs at most one more local enrichment pass before the public step.

\paragraph{Release consequence.}
The practical consequence is that any further local subversion should no longer open a new mathematical front. It should either confirm release-readiness or add one last density-improving clarification and then stop.




\subsection*{2J. Open-questions closure pass and chain check}
The present stage is also strong enough to classify several previously open questions more explicitly. The purpose of this pass is not to pretend that every weak-density zone has disappeared, but to distinguish between (i) questions that can already be closed reliably at the current structural level, (ii) questions that remain structurally intact but only under an explicit assumption or still-thin proof layer, and (iii) questions that should remain openly deferred. This distinction is especially important in a structure-first setting, because a locally thin argument is not automatically broken: some passages remain admissible as conditional or contradiction-style steps, provided that their status is made explicit and no illegitimate theorem-promotion occurs.

\paragraph{Closure classification.}
At the current stage the following questions may be treated as \emph{closable now}: the dependency spine from ground condition to HC and confinement; the compatibility of closure with projector stabilization; the admissible transport continuation on the stabilized sector; the admissible-sector action of the Millennium matrix; and the persistence of admissible trace/branch continuation once the stabilized sector has been fixed. These are not yet maximally unfolded everywhere, but they are no longer merely heuristic. By contrast, the following questions are presently better read as \emph{intact but proof-density limited}: the strongest available mathematical reading of the Millennium matrix as converter/intertwiner; the full density of trace/branch continuation lemmas; and the extent to which some still-compact structural passages should already be promoted beyond proposition level. Finally, the following remain \emph{intentionally open}: any full universality claim for the computational strand, any finished proof-kernel or compiler reading, and any interpretation that would make the mathematics depend on the computational teaser architecture.

\paragraph{Chain-check statuses.}
To keep the proof chain honest without overcalling local weakness as failure, the document now adopts four chain-check statuses for future mathematical passes:
\begin{itemize}[leftmargin=2em]
  \item \textbf{intact:} the dependency chain carries without visible status or continuity defect;
  \item \textbf{intact, but proof-density weak:} the chain carries, but one or more transitions still need denser local unfolding;
  \item \textbf{intact under explicit assumption:} the argument is acceptable only because an assumption, contradiction setup, or provisional hypothesis is openly declared;
  \item \textbf{blocker:} a genuine dependency break, status violation, or contradiction prevents closure.
\end{itemize}

\paragraph{Current reading.}
Under this classification, the first major dependency spine is currently best read as \emph{intact, but proof-density weak} rather than as blocked. No primary structural contradiction is visible in the path from ground condition through HC, closure, admissible transport, Millennium-matrix realization, and admissible trace/branch persistence. The remaining caution is therefore mainly one of density and explicitness, not one of broken continuity.

\paragraph{Local consequence.}
The practical consequence is conservative but positive. Several earlier open questions may already be marked as \emph{reliably closable at the present stage}, provided that the document does not silently promote weak-density passages to stronger theorem status than the current support allows. The next mathematics-first passes should therefore use the chain-check statuses as a standing audit device: strengthen density where useful, keep assumption-based passages openly marked, and treat only genuine continuity failure as a blocker.




\subsection*{2K. Tool extraction and leverage-point scan}
The next mathematics-first pass should not treat every locally thin argument as something that must immediately be expanded line by line. In a structure-first document, some still-wobbly passages already contain reusable mathematical moves that can be extracted earlier as \emph{tools}, provided that their status is made explicit. The present scan therefore distinguishes between (i) tools that are already operationally useful now, even if their surrounding proof density is still moderate, (ii) tools that remain usable only under an explicit structural assumption, and (iii) a small number of leverage points whose clarification would simplify several later proof passages at once.

\paragraph{Provisional tool extraction.}
At the current stage the following reusable tools may already be extracted from the first dependency spine. First, a \emph{closure-transfer tool}: once the admissible sector has been fixed, closure may be transported across projector stabilization without reopening the sector boundary. Second, a \emph{projected transport tool}: admissible transport may be read modulo the stabilized projector rather than by re-proving closure from scratch after each transport step. Third, a \emph{matrix-sector tool}: the Millennium matrix may already be used as an admissible-sector-preserving operator on the stabilized transported sector. Fourth, a \emph{trace/branch persistence tool}: admissible continuation in the trace/branch layer may be used without reintroducing an external admissibility source. None of these tools should yet be oversold as maximal-density theorems everywhere; however, they are already strong enough to simplify downstream arguments by replacing repeated local closure checks with a smaller number of reusable invariance moves.

\paragraph{Status discipline for extracted tools.}
To keep the extraction mathematically honest, each such tool should be marked in one of three ways:
\begin{itemize}[leftmargin=2em]
  \item \textbf{operational tool:} reusable now, with the present support level;
  \item \textbf{assumption-backed tool:} reusable only under an explicitly declared hypothesis or contradiction setup;
  \item \textbf{blocked tool candidate:} not yet safe to reuse because a real dependency gap remains.
\end{itemize}
Under the current chain check, the closure-transfer, projected transport, matrix-sector, and trace/branch persistence tools are best read as \emph{operational tools}, though some still sit in the ``proof-density weak'' zone rather than the fully unfolded zone.

\paragraph{Leverage-point scan.}
The present scan suggests that the most useful immediate leverage point is not a new broad theorem family, but a single stronger organizing hypothesis: namely that the admissible stabilized sector behaves as the true fixed working domain for closure, transport, matrix action, and admissible continuation. Even if one initially formulates this only as a hypothesis-backed simplifier, it has unusually high leverage: if strengthened, it compresses several later arguments by replacing repeated local admissibility restoration with one global working-domain principle. In other words, many later proofs become shorter once they are allowed to start with ``assume the admissible stabilized sector has already been fixed'' and then work internally on that domain.

\paragraph{Use of the computational layer as a secondary aid.}
The computational chapter may already help here, but only in a subordinate role. MML/MMA/MCP should not drive the mathematics; however, the language layer can already serve as a \emph{notation and workflow aid} for tool extraction. In particular, compact MML-style descriptions can help identify whether a candidate move is really one reusable operation---for example \texttt{project -> transport -> trace}---or whether it still hides an unresolved mathematical gap. Thus the computational layer may already be used as a diagnostic simplifier for tool extraction, while the mathematical status of the extracted tool remains decided purely on the mathematics-first side.

\paragraph{Current reading and local consequence.}
The current reading is therefore positive. The first dependency spine is no longer merely a place where one asks whether the chain survives; it is already rich enough to yield a first family of reusable structural tools. At the same time, the highest-leverage next target is also visible: stabilize the admissible working-domain principle as far as possible, because this one point appears to simplify closure, transport, matrix action, and trace/branch continuation simultaneously. Any next local pass should therefore prefer a leverage-point clarification over broad diffuse proof growth.



\section*{Appendix C. Julia Boundary, Master Residual $\eta$, and Reversal-Probe Consolidation}
\addcontentsline{toc}{section}{Appendix C. Julia Boundary, Master Residual $\eta$, and Reversal-Probe Consolidation}

This appendix records the condensed 3.3 main line. The local residual sector is organized around the single scalar
\[
  \eta_f:=\frac13-\beta_f^2.
\]
The point of the compression is that the previously separated Hopf-gap, Julia-boundary, isotropic defect, and reversal language are different visible faces of the same transported branch quantity.

\begin{proposition}[Master residual transport and isotropic readout]
\label{lem:defect-transport}
\label{prop:master-defect-scalar}
Let $f$ range over admissible readouts of the same closure-stable kernel, and assume that on the admissible branch the visible compensated defect is isotropic,
\[
  D^{(f)}_{\mu\nu}=\eta_f\,g^{(f)}_{\mu\nu}.
\]
Assume moreover that admissible transport preserves the underlying defect class. Define
\[
  \eta_f:=\frac13-\beta_f^2.
\]
Then the following are equivalent encodings of the same local residual sector:
\begin{enumerate}[label=(\roman*)]
  \item the additive $\beta^2$-gap $\eta_f$;
  \item the continuum-limit deformation
  \[
    \delta_{\mathrm{CL},f}:=\frac{\beta_f}{1/\sqrt3}=\sqrt3\,\beta_f,
    \qquad
    \eta_f=\frac{1-\delta_{\mathrm{CL},f}^2}{3};
  \]
  \item the cosmological residuals
  \[
    \eta_f=\frac13-\Omega_M^{(f)}=\Omega_\Lambda^{(f)}-\frac23
  \]
  whenever $\Omega_M^{(f)}=\beta_f^2$ and $\Omega_\Lambda^{(f)}=1-\beta_f^2$;
  \item the isotropic defect tensor $D^{(f)}_{\mu\nu}=\eta_f g^{(f)}_{\mu\nu}$;
  \item the transported quadratic invariant
  \[
    \mathfrak I_f=\eta_f^2.
  \]
\end{enumerate}
In particular, every admissible readout carries the same local defect class through the single scalar $\eta_f$.
\end{proposition}

\begin{proof}
The displayed identities are algebraic. Once the visible defect is isotropic, every normalized quadratic contraction is determined by its scalar coefficient, hence by $\eta_f$. Transport preservation identifies the same defect class in different readouts, so the admissible residual sector is represented everywhere by the same one-scalar quantity and its powers.
\end{proof}

\begin{remark}[Symmetric value and the $10^{-3}$ scale]
For the symmetric matter/tachyon readout with $\beta_f=\beta_{\mathrm{sym}}\approx0.5493$ one gets
\[
  \eta_{\mathrm{sym}}=\frac13-\beta_{\mathrm{sym}}^2\approx 3.1603\times10^{-2},
  \qquad
  \eta_{\mathrm{sym}}^2\approx 9.987\times10^{-4}.
\]
So a number near $10^{-3}$ is most naturally read not as a primitive constant, but as the quadratic visible residual of the transported master defect scalar.
\end{remark}

\begin{proposition}[Julia boundary and exact-core principle]
\label{prop:julia-master-residual}
Assume that admissible branching does not introduce a new primitive local scalar beyond the continuum--discrete deformation of the same transport sector. Then the exact compensated core has two equivalent local descriptions:
\[
  \text{intrinsic core: }\zminus=0,
  \qquad
  \text{normalized visible core: }\eta_f=0.
\]
Moreover every admissible leading-order branch/Julia defect functional satisfies
\[
  \mathfrak J_f=c_f\,\eta_f+O(\eta_f^2)
\]
with a nonzero normalization constant $c_f$. Thus the observed Julia-boundary sensitivity is controlled, to leading order, by the same scalar $\eta_f$ that governs the isotropic defect readout.
\end{proposition}

\begin{proof}
The exact compensated core is representation-stable, but it is seen in a phase-shifted readout frame. Intrinsically the core is the null-kernel relation $\zminus=0$; after Hopf normalization of the visible branch it is the vanishing of the transported residual coefficient $\eta_f$. If admissible branching opens no new primitive scalar channel, then the first visible boundary defect must be a function of the already available scalar $\eta_f$, and its first nontrivial term is therefore linear in $\eta_f$.
\end{proof}

\begin{theorem}[Local $\eta$-completeness and restoring normal form]
\label{thm:eta-completeness-restoring}
Assume a branch-local admissible transport chart in which visible observables depend only on the local admissible branch parameter and no admissible transverse scalar channel reopens. Then:
\begin{enumerate}[label=(\roman*)]
  \item every branch-local scalar observable without derivatives is a function of $\eta$;
  \item every first-jet scalar observable is a function of $\eta$ as well, because
  \[
    \nabla_\mu \eta = A(\eta)\,U_\mu
  \]
  for some branch tangent $U$ and scalar amplitude $A$;
  \item every second-jet branch-local scalar observable is a function of $\eta$ and the longitudinal second derivative
  \[
    Y:=U^\mu U^\nu \nabla_\mu\nabla_\nu\eta.
  \]
\end{enumerate}
If, in addition, the branch is locally $\eta$-complete, then there exists an autonomous reduced equation
\[
  \frac{d^2\eta}{ds^2}=B(\eta)
\]
with
\[
  B(0)=0.
\]
If the branch realizes a closure-restoring return law, then near the exact core one has
\[
  \eta''+\omega^2\eta=O(\eta^2)
\]
with $\omega^2>0$ on the nondegenerate branch. Hence the informal ``Umkehrwalze'' is mathematically realized as a local restoring dynamics for the transported master residual.
\end{theorem}

\begin{proof}
The zeroth-order claim follows from isotropy of the visible defect sector: once $D_{\mu\nu}=\eta g_{\mu\nu}$, all algebraic scalar contractions are polynomials in $\eta$. Branch-local visibility forces the first derivative of any such scalar to be tangent to the same admissible branch direction, so the first jet remains one-dimensional. At second order the only new natural scalar is the longitudinal Hessian component $Y$, and local $\eta$-completeness identifies it as a function of the already visible branch value. The exact core is transport-preserved, so it must be an equilibrium of the reduced equation; closure-restoring return then fixes the linearized sign and yields the stated normal form.
\end{proof}

\begin{definition}[Hybrid reversal probe]
Define the condensed reversal probe by
\[
  \mathfrak S_{\mathrm{rev}}:=(\mathbb P_{\mathrm B},\mathcal L_{\mathrm L}),
\]
where $\mathbb P_{\mathrm B}$ is a hard admissibility/closure projector (the ``Bombe'' side) and $\mathcal L_{\mathrm L}$ is a fine residual functional on the $\eta$-branch (the ``Lorenz'' side).
\end{definition}

\begin{proposition}[Reversal-probe dichotomy]
Let two admissible probe trajectories on the same local residual branch have the same visible master residual value. Then locally only two outcomes remain:
\[
  \Delta\theta\equiv 0 \pmod{2\pi}
  \qquad\text{or}\qquad
  \Delta\theta\equiv \pi \pmod{2\pi}.
\]
Thus the two probes either converge to the same visible branch representative, or they lie on the same transported structure with nontrivial $\mathbb Z_2$ phase class and close only on the spinorial/M\"obius double cover.
\end{proposition}

\begin{proof}
The hard projector removes inadmissible branch continuations, while the fine functional identifies the remaining residual class through the same local variable $\eta$. Equality of the visible residual therefore leaves no generic room for an independent third local phase sector. What remains is either direct coincidence in the visible branch parameterization or the unique nontrivial $\mathbb Z_2$ lift compatible with spinorial sign reversal.
\end{proof}

\subsection*{C.1 First 3.4 extension: closure potential and global branch law}

\begin{definition}[Closure potential and branch energy]
Assume a connected admissible branch interval $I\ni 0$ on which the reduced residual equation is exact, i.e.
\[
  \eta'' = B(\eta) = -\mathcal U'(\eta)
\]
for some $\mathcal U\in C^2(I)$. We call $\mathcal U$ a closure potential for the branch and define the associated branch energy by
\[
  \mathcal E_\eta := \frac12 (\eta')^2 + \mathcal U(\eta).
\]
\end{definition}

\begin{theorem}[Global exact branch law under closure exactness]
\label{thm:global-exact-branch-law}
Assume the hypotheses of Theorem~\ref{thm:eta-completeness-restoring}, and assume in addition branch-local closure exactness in the above sense. Then:
\begin{enumerate}[label=(\roman*)]
  \item $\mathcal E_\eta$ is conserved along every admissible branch trajectory;
  \item the exact compensated core $\eta=0$ is an equilibrium with $\mathcal U'(0)=0$;
  \item if the branch is closure-restoring and nondegenerate, then $\mathcal U''(0)=\omega^2>0$, hence
  \[
    \mathcal U(\eta)=\frac12 \omega^2 \eta^2 + O(\eta^3);
  \]
  \item every admissible turning point of the visible residual branch satisfies
  \[
    \mathcal U(\eta)=\mathcal E_\eta.
  \]
\end{enumerate}
\end{theorem}

\begin{proof}
Differentiating $\mathcal E_\eta$ along a trajectory gives
\[
  \frac{d}{ds}\mathcal E_\eta = \eta'\eta''+\mathcal U'(\eta)\eta' = \eta'(\eta''+\mathcal U'(\eta))=0.
\]
The equilibrium statement is equivalent to $B(0)=0$, already forced by core preservation. The nondegenerate restoring sign from the local normal form gives $B(\eta)=-\omega^2\eta + O(\eta^2)$, hence $\mathcal U'(\eta)=\omega^2\eta + O(\eta^2)$ and therefore the stated Taylor expansion. Turning points are exactly the points where $\eta'=0$, so the conserved-energy identity reduces to $\mathcal U(\eta)=\mathcal E_\eta$.
\end{proof}

\begin{proposition}[Energy-form reversal-probe criterion]
\label{prop:energy-reversal-probe}
Assume the hypotheses of Theorem~\ref{thm:global-exact-branch-law}. Let two admissible probe trajectories be selected by the same hard projector $\mathbb P_{\mathrm B}$ and lie on the same connected regular level set $\mathcal E_\eta = E$. Assume moreover that this energy component carries no monodromy higher than the spinorial $\mathbb Z_2$-lift. Then exactly two global outcomes remain:
\[
\begin{aligned}
&\text{direct convergence on the visible branch}\\[0.2em]
&\qquad \text{or} \qquad
\text{closure only on the spinorial/M\"obius double cover}
\end{aligned}
\]
\end{proposition}

\begin{proof}
The hard projector excludes inadmissible components, while the fine residual functional identifies the same visible residual class through the common potential-energy level. On a connected regular admissible energy component, the remaining ambiguity is topological rather than scalar. By hypothesis no higher monodromy is available, so the only surviving alternatives are trivial monodromy and the unique nontrivial $\mathbb Z_2$ lift.
\end{proof}


\subsection*{C.2 Second 3.4 extension: reversible single-well class}

\begin{theorem}[Parity and single-well restriction of the closure potential]
\label{thm:appendix-single-well-341}
Assume the hypotheses of Theorem~\ref{thm:global-exact-branch-law}. Assume moreover that the admissible residual branch is invariant under the involution $R(\eta)=-\eta$, that the reduced branch law commutes with $R$, and that $\eta=0$ is the unique admissible critical core on the connected interval $I$ with no further admissible critical point of $\mathcal U$ on $I\setminus\{0\}$. Normalize $\mathcal U(0)=0$. Then:
\begin{enumerate}[label=(\roman*)]
  \item $B(-\eta)=-B(\eta)$ and $\mathcal U(-\eta)=\mathcal U(\eta)$ on $I$;
  \item near the exact core,
  \[
    \mathcal U(\eta)=\frac12\omega^2\eta^2+a_4\eta^4+O(\eta^6),
    \qquad
    B(\eta)=-\omega^2\eta-4a_4\eta^3+O(\eta^5);
  \]
  \item $\mathcal U'(\eta)>0$ for every $\eta>0$ in $I$ and $\mathcal U'(\eta)<0$ for every $\eta<0$ in $I$; in particular $\mathcal U(\eta)>0$ for all $\eta\neq0$.
\end{enumerate}
Consequently $\mathcal U$ belongs to the symmetric single-well class on $I$.
\end{theorem}

\begin{proof}
If $\eta(s)$ solves $\eta''=B(\eta)$, branch reversibility and commutation with $R$ imply that $-\eta(s)$ is again an admissible solution of the same reduced law. Hence
\[
  -\eta'' = B(-\eta),
\]
while applying $R$ to the original equation gives $-\eta''=-B(\eta)$, so $B(-\eta)=-B(\eta)$. Since $\mathcal U'(\eta)=-B(\eta)$, the derivative $\mathcal U'$ is odd. With the normalization $\mathcal U(0)=0$, integration yields $\mathcal U(-\eta)=\mathcal U(\eta)$.

Theorem~\ref{thm:global-exact-branch-law} already gives $\mathcal U''(0)=\omega^2>0$. Evenness forces the Taylor series of $\mathcal U$ at the exact core to contain only even powers, so
\[
  \mathcal U(\eta)=\frac12\omega^2\eta^2+a_4\eta^4+O(\eta^6),
\]
and differentiation gives the stated odd expansion of $B$.

Because $\mathcal U''(0)=\omega^2>0$, one has $\mathcal U'(\eta)>0$ for sufficiently small $\eta>0$. If $\mathcal U'$ vanished again at some positive point, continuity would force an additional admissible critical point of $\mathcal U$ on $I\cap(0,\infty)$, excluded by hypothesis. Hence $\mathcal U'(\eta)>0$ on the whole positive half-branch. The negative half-branch follows by oddness of $\mathcal U'$. Therefore $\eta=0$ is the unique well bottom and $\mathcal U(\eta)>0$ for all $\eta\neq0$.
\end{proof}

\begin{proposition}[Symmetric turning-pair criterion]
\label{prop:symmetric-turning-pair-341}
Under the hypotheses of Theorem~\ref{thm:appendix-single-well-341}, let $E>0$ be a regular branch-energy value. Then the turning-point equation
\[
  \mathcal U(\eta)=E
\]
has on $I$ either no solution or exactly the symmetric pair $\eta=\pm\eta_E$ with $\eta_E>0$. In particular every regular admissible energy shell carries at most the scalar ambiguity $\eta\leftrightarrow-\eta$.
\end{proposition}

\begin{proof}
By the preceding theorem, $\mathcal U$ is even and strictly increasing on each half-branch as a function of $|\eta|$. Therefore the equation $\mathcal U(\eta)=E$ can have at most one positive and at most one negative solution, and evenness forces them to occur as a symmetric pair whenever they exist.
\end{proof}

\subsection*{C.3 Third 3.4 extension: canonical quartic family}
\addcontentsline{toc}{subsection}{C.3 Third 3.4 extension: canonical quartic family}

\begin{proposition}[Quartic reduction under first nonlinear visibility]
\label{prop:appendix-342a}
Assume the hypotheses of Theorem~\ref{thm:mainline-342}. Then on a symmetric admissible neighborhood $J$ one has
\[
  \mathcal U(\eta)=\frac12\omega^2\eta^2+\frac{\lambda}{4}\eta^4+R_6(\eta),
  \qquad R_6(\eta)=O(\eta^6),
\]
for some $\lambda\ge 0$.
\end{proposition}

\begin{proof}
By Theorem~\ref{thm:mainline-341}, the potential is even, $C^6$, and has a strict minimum at $0$, so its Taylor expansion on $J$ has the form
\[
  \mathcal U(\eta)=\frac12\omega^2\eta^2+a_4\eta^4+a_6\eta^6+o(\eta^6).
\]
The hypotheses exclude any additional visible scalar or independent local scale at first nonlinear order, so the quartic coefficient is the only admissible datum at that order. Writing $\lambda:=4a_4$ and absorbing the higher terms into $R_6$ gives the claimed form. Since $0$ remains the unique admissible local minimum on $J$, shrinking $J$ if necessary excludes the sign choice that would open a new nearby admissible critical pair; therefore $\lambda\ge 0$ on the retained neighborhood.
\end{proof}

\begin{proposition}[Turning-value parameterization]
\label{prop:appendix-342b}
In the minimal truncation
\[
  \mathcal U_{\min}(\eta)=\frac12\omega^2\eta^2+\frac{\lambda}{4}\eta^4,
  \qquad \lambda\ge 0,
\]
regular energy shells are parameterized by the unique positive turning value $A(E)$ satisfying
\[
  E=\frac12\omega^2A(E)^2+\frac{\lambda}{4}A(E)^4.
\]
If $\lambda>0$, then
\[
  A(E)^2=\frac{-\omega^2+\sqrt{\omega^4+4\lambda E}}{\lambda},
\]
while for $\lambda=0$ one has $A(E)=\sqrt{2E}/\omega$.
\end{proposition}

\begin{proof}
On a regular shell, turning points are exactly the zeros of $\eta'$, so they satisfy $E=\mathcal U_{\min}(\eta)$. By even single-well monotonicity there is exactly one positive solution $A(E)$ and one negative solution $-A(E)$. Solving the quadratic equation in $x=A(E)^2$ gives the explicit formula when $\lambda>0$, and the harmonic formula when $\lambda=0$.
\end{proof}

\begin{proposition}[Appendix C.4: quartic closure coefficient as visible datum]
\label{prop:app-c-34-4}
Under the hypotheses of Theorem~\ref{thm:mainline-343}, the visible quartic branch coefficient $\lambda$ is not an independent nonlinear parameter. It is exactly the transported quartic closure coefficient $\mu c_4$, and every regular probe shell may therefore be parameterized either by its conserved energy $E$ or by its unique positive turning value $A(E)$ without introducing an additional local scalar.
\end{proposition}

\begin{proof}
Theorem~\ref{thm:mainline-343} gives $\lambda=\mu c_4$ with $\mu>0$, so the quartic correction is fixed by the same closure geometry that already determines the quadratic restoring core. Proposition~\ref{prop:appendix-342b} identifies the regular shell by either $E$ or $A(E)$, and Theorem~\ref{thm:mainline-342} excludes any further visible local scalar at quartic order. Hence the quartic visible datum is exactly the transported closure coefficient and does not enlarge the local scalar algebra.
\end{proof}


\begin{proposition}[Appendix C.5: quartic timing drift as probe readout]
\label{prop:app-c-34-5}
Under the hypotheses of Theorem~\ref{thm:mainline-344}, the first nonlinear timing drift of a regular periodic shell is already fixed by the transported quartic closure coefficient. More precisely,
\[
  T(E)=\frac{2\pi}{\omega}\Bigl(1-\frac{3c_4}{4\mu c_2^2}E+O(E^2)\Bigr),
\]
so the branch is locally isochronous at first nonlinear order if and only if $c_4=0$, while $c_4>0$ forces a strictly decreasing period for all sufficiently small regular energies.
\end{proposition}

\begin{proof}
Theorem~\ref{thm:mainline-344} gives the small-energy period expansion together with the closure rewrite $\omega^2=\mu c_2$ and $\lambda=\mu c_4$. Since $\mu>0$ and $c_2>0$, the first nonlinear period correction vanishes exactly when $c_4=0$ and has negative sign whenever $c_4>0$. Thus the timing drift is a direct probe readout of the transported quartic closure coefficient and does not introduce a second nonlinear scalar.
\end{proof}


\begin{proposition}[Appendix C.6: higher even hierarchy and delayed timing defect]
\label{prop:app-c-34-6}
Under the hypotheses of Theorem~\ref{thm:mainline-345}, the first nonzero higher nonlinear timing defect is governed by the first nonvanishing higher even closure coefficient. More precisely, if
\[
  c_4=c_6=\cdots=c_{2m-2}=0,
  \qquad c_{2m}\neq 0,
  \qquad m\ge 3,
\]
then
\[
  T(E)=\frac{2\pi}{\omega}\Bigl(1-\kappa_m(\omega)\,\mu c_{2m}\,E^{m-1}+O(E^m)\Bigr),
  \qquad \kappa_m(\omega)>0.
\]
Hence the branch is isochronous through all lower nonlinear orders, and the first visible higher timing drift appears exactly at order $E^{m-1}$ with sign determined by $c_{2m}$.
\end{proposition}

\begin{proof}
Theorem~\ref{thm:mainline-345} identifies the first visible higher nonlinear coefficient as $\alpha_{2m}=\mu c_{2m}$ and gives the corresponding small-energy period law. Because the intermediate even coefficients vanish, no lower-order contribution to the period expansion is available. Since $\kappa_m(\omega)>0$ and $\mu>0$, the sign of the first nonzero higher timing defect is exactly the sign of $c_{2m}$.
\end{proof}


\begin{proposition}[Appendix C.7: analytic generator form and coefficient transport]
\label{prop:app-c-34-7}
Under the hypotheses of Theorem~\ref{thm:mainline-346}, the local visible branch dynamics is canonically encoded by a single even generator $\Phi$. More precisely,
\[
  \mathcal U(\eta)=\Phi(\eta^2),
  \qquad
  B(\eta)=-2\eta\,\Phi'(\eta^2),
\]
with
\[
  \Phi(s)=\frac12\omega^2 s + \frac{\lambda}{4}s^2 + \sum_{m\ge 3}\frac{\alpha_{2m}}{2m}s^m,
\]
and coefficient transport
\[
  \omega^2=\mu c_2,
  \qquad
  \lambda=\mu c_4,
  \qquad
  \alpha_{2m}=\mu c_{2m}\quad (m\ge 3).
\]
Hence the entire local nonlinear shell hierarchy is fixed by the single generator germ $\Phi$, and every delayed timing defect is a readout of its coefficients rather than a new independent scalar.
\end{proposition}

\begin{proof}
Theorem~\ref{thm:mainline-346} gives the analytic factorization $\mathcal U(\eta)=\Phi(\eta^2)$ together with the coefficient identifications. Expanding $\Phi$ near $s=0$ yields the displayed formula, while differentiating gives the canonical branch law $B(\eta)=-2\eta\Phi'(\eta^2)$. The higher-order timing and turning-value readouts were already tied to the visible coefficients in Theorems~\ref{thm:mainline-344} and \ref{thm:mainline-345}, so the generator form introduces no additional scalar data.
\end{proof}


\begin{proposition}[Appendix C.8: generator-shell classifier]
\label{prop:app-c-34-8}
Under the hypotheses of Theorem~\ref{thm:mainline-347}, every sufficiently small regular admissible shell is completely determined by its generator radius $s(E)=\Phi^{-1}(E)$, and the visible shell ambiguity is exactly the transported sign pair
\[
  \eta=\pm\sqrt{s(E)}.
\]
Hence the hybrid reversal probe has on such a shell only two local transport classes: sign-preserving transport (direct visible convergence) and sign-reversing transport (closure only after the unique spinorial/M\"obius lift).
\end{proposition}

\begin{proof}
Strict monotonicity of $\Phi$ on the small-shell interval gives the inverse shell map $s(E)=\Phi^{-1}(E)$. Since the visible potential is $\mathcal U(\eta)=\Phi(\eta^2)$, the shell equation $\mathcal U(\eta)=E$ is equivalent to $\eta^2=s(E)$, so the shell consists exactly of the two visible representatives $\pm\sqrt{s(E)}$. Any admissible transport restricted to this shell acts on that pair by either preserving sign or reversing it. By the standing no-higher-monodromy hypothesis, these two sign actions are the only local shell transport classes.
\end{proof}


\begin{proposition}[Appendix C.9: shell parity and diagnostic M\"obius-reflected render]
\label{prop:app-c-34-9}
Under the hypotheses of Theorem~\ref{thm:mainline-348}, the local probe outcome on every sufficiently small regular shell is completely determined by the shell parity $\pi(E)\in\{+1,-1\}$. If an admissible shell cycle decomposes into regular segments and transverse reflector crossings, then
\[
  \pi(E)=(-1)^{N_E},
\]
where $N_E$ is the number of transverse reflector crossings in the cycle. Accordingly, even crossing parity yields direct visible convergence, while odd crossing parity yields closure only after the spinorial/M\"obius lift.
\end{proposition}

\begin{proof}
Theorem~\ref{thm:mainline-348} already shows that the shell transport acts on the visible turning pair by a unique sign parity $\pi(E)$. Under the stated decomposition, regular segments preserve sign and every transverse reflector crossing contributes the involution $\eta\mapsto-\eta$. Multiplying these contributions gives $\pi(E)=(-1)^{N_E}$, from which the direct-convergence versus lifted-closure criterion follows immediately.
\end{proof}

\begin{figure}[H]
\centering
\IfFileExists{julia_probe_residual_scales_v3841.pdf}{\includegraphics[width=0.96\textwidth]{julia_probe_residual_scales_v3841.pdf}}{\safegraphic[width=0.96\textwidth]{julia_probe_residual_scales_v3841.png}}
\caption{\textbf{Diagnostic M\"obius-reflected quadratic probe render at the current residual scales.} Comparison of the iterated map $z_{n+1}=1/\overline{(z_n^2+c_\gamma)}$ for the baseline $\gamma=0$ and for the three currently relevant OP~1 residual scales $\Delta_\beta\approx 7.50\times 10^{-5}$, $\Delta_\theta\approx 1.79\times 10^{-4}$, and $\Delta\mu\approx 3.38\times 10^{-4}$. For the present diagnostic render, the parameter family is instantiated by $c_\gamma=-1+i\gamma$, with the unit circle drawn as the M\"obius horizon. The global morphology remains stable across all four renders, while fine boundary filaments and local transition zones respond smoothly under the residual-scale shifts. The figure is included as a diagnostic visualization of structural robustness versus boundary-level response; the earlier marker value $0.0009$ is retained elsewhere only as a historical stress-test scale, not as the current main comparison value.}
\label{fig:app-c-julia-diagnostic}
\end{figure}

Figure~\ref{fig:app-c-julia-diagnostic} is included as a diagnostic probe image rather than as a first-principles derivation. Its role is to visualize a separation between global morphological stability and boundary-level sensitivity: the large-scale visible body remains stable, while the finest branch and transition regions react under the current OP~1 residual scales. The asset is now generated reproducibly by the local Python renderer \texttt{generate\_julia\_probe\_residual\_scales.py}, which reads the displayed values of $\Delta_\beta$, $\Delta_\theta$, and $\Delta\mu$ from the TeX source and produces both a PDF and a PNG version; the document includes the PDF render when available so that the figure can be rebuilt directly from the project material.

\begin{figure}[H]
\centering
\begin{minipage}{0.475\textwidth}
\centering
\safegraphic[width=\linewidth]{gamma0_julia_probe_v3841.png}
\\[-0.5ex]
{\footnotesize $\gamma=0$}
\end{minipage}\hfill
\begin{minipage}{0.47\textwidth}
\centering
\safegraphic[width=0.96\linewidth]{delta_beta_julia_probe_v3841.png}
\\[-0.5ex]
{\footnotesize $\gamma=\Delta_\beta$}
\end{minipage}

\medskip

\begin{minipage}{0.47\textwidth}
\centering
\safegraphic[width=0.95\linewidth]{delta_theta_julia_probe_v3841.png}
\\[-0.5ex]
{\footnotesize $\gamma=\Delta_\theta$}
\end{minipage}\hfill
\begin{minipage}{0.47\textwidth}
\centering
\safegraphic[width=0.95\linewidth]{delta_mu_julia_probe_v3841.png}
\\[-0.5ex]
{\footnotesize $\gamma=\Delta\mu$}
\end{minipage}
\caption{\textbf{Individual diagnostic M\"obius-reflected probe panels.} The four residual-scale renders from Figure~\ref{fig:app-c-julia-diagnostic} shown again as separate panels, so that the boundary-level response at $\gamma=0$, $\Delta_\beta$, $\Delta_\theta$, and $\Delta\mu$ can be inspected without the surrounding comparison frame.}
\label{fig:app-c-julia-panels}
\end{figure}


\subsection*{C.9 Tenth 3.4 extension: tunnel observables and quantitative readout interface}
\addcontentsline{toc}{subsection}{C.9 Tenth 3.4 extension: tunnel observables and quantitative readout interface}

\begin{proposition}[Appendix support for Proposition~\ref{prop:mainline-3410}]
\label{prop:appendix-3410}
Under the hypotheses of Proposition~\ref{prop:mainline-3410}, any admissible quantitative comparison between the diagnostic Julia/probe render and an external readout image may be organized through the tunnel-observable package
\[
  \mathfrak O_{\mathrm{diag}}(\gamma)
  =\bigl(w_{\mathrm{tun}}(\gamma;\tau),\ I_{\mathrm{edge}}(\gamma),\ \Theta_{\mathrm{open}}(\gamma),\ P_{\mathbb Z_2}(E)\bigr),
\]
with no additional local scalar required beyond the generator-shell radius and parity sector. In particular, if two rendering protocols produce the same $s(E)$ and the same parity law on the same thresholded shell family, then any discrepancy among the first three observables is attributable to rendering normalization rather than to a new local structural channel.
\end{proposition}

\begin{proof}
This is a direct restatement of Proposition~\ref{prop:mainline-3410} in appendix language. The tunnel, edge, and opening observables are post-processing functionals of the same thresholded shell geometry, and the parity response is already intrinsic to the shell classifier. Hence matching generator/parity data leaves only normalization choices, not an extra structural scalar.
\end{proof}


\subsection*{C.10 Eleventh 3.4 extension: inverse readout and shell-consistency defect}
\addcontentsline{toc}{subsection}{C.10 Eleventh 3.4 extension: inverse readout and shell-consistency defect}

\begin{theorem}[Appendix support for Theorem~\ref{thm:mainline-3411}]
\label{thm:appendix-3411}
Assume the calibrated tunnel-width law
\[
  w_{\mathrm{tun}}(\gamma;\tau)=W_\tau(s(E))
\]
with strictly monotone $W_\tau$ on a small-shell interval. Then the diagnostic shell data is locally equivalent to the pair
\[
  \bigl(w_{\mathrm{tun}}(\gamma;\tau),P_{\mathbb Z_2}(E)\bigr),
\]
with inverse reconstruction
\[
  s(E)=W_\tau^{-1}(w_{\mathrm{tun}}(\gamma;\tau)),
  \qquad
  \pi(E)=P_{\mathbb Z_2}(E).
\]
Any further shell metric extracted from the same thresholded render therefore factors through these two calibrated quantities.
\end{theorem}

\begin{proof}
Monotonicity of $W_\tau$ gives local invertibility. The shell-parity response already records the transported $\mathbb Z_2$ class. Since Proposition~\ref{prop:mainline-3410} shows that all remaining tunnel/edge/opening metrics are readout-functionals of the same shell geometry, they factor through the recovered pair.
\end{proof}


\subsection*{C.11 Twelfth 3.4 extension: inverse identifiability and stability support}
\addcontentsline{toc}{subsection}{C.11 Twelfth 3.4 extension: inverse identifiability and stability support}

\begin{theorem}[Appendix support for Theorem~\ref{thm:mainline-3412}]
\label{thm:appendix-3412}
Assume the calibrated inverse-readout hypotheses of Theorem~\ref{thm:mainline-3411} and fix a shell point $s_0$ with $W_\tau'(s_0)\neq 0$. Then on a sufficiently small shell neighborhood the calibration law is bi-Lipschitz,
\[
  m_0|s_1-s_2|
  \le |W_\tau(s_1)-W_\tau(s_2)|
  \le M_0|s_1-s_2|,
\]
for suitable constants $m_0,M_0>0$. Hence the calibrated tunnel width determines the shell radius uniquely on each fixed parity branch and does so with inverse condition number bounded by $1/m_0$.
\end{theorem}

\begin{proof}
Since $W_\tau'(s_0)\neq 0$, continuity of the derivative yields a neighborhood on which $W_\tau'$ stays bounded away from zero and from infinity. The mean-value theorem gives the displayed estimate, and the inverse condition-number statement is just the corresponding inverse Lipschitz bound.
\end{proof}


\subsection*{C.13 Fourteenth 3.4 extension: obstruction-shell transition support}
\addcontentsline{toc}{subsection}{C.13 Fourteenth 3.4 extension: obstruction-shell transition support}

\begin{theorem}[Appendix support for Theorem~\ref{thm:mainline-3414}]
\label{thm:appendix-3414}
Let $s_\ast$ be an isolated obstruction shell. If $s_\ast$ is a nondegenerate calibration critical shell, then the calibrated tunnel law has a local fold normal form with a two-sheet inverse exchange. If $s_\ast$ is instead a pure parity shell, then the shell radius continues uniquely and only the transported $\mathbb Z_2$ class jumps. In the mixed case the local transition is the direct product of these two involutions.
\end{theorem}

\begin{proof}
A nondegenerate critical shell of a one-dimensional smooth map is locally a fold, so the inverse consists of exactly two sheets exchanged by the local deck involution. A pure parity shell leaves the scalar calibration regular and therefore does not affect the shell coordinate, but it does flip the transported sign class by definition. Since one operation acts on the shell coordinate and the other on parity, the mixed case is their commuting composition.
\end{proof}

\paragraph{3.4.14 working interpretation.}
The obstruction shells are therefore organized rather than chaotic. Each isolated obstruction contributes only one of two primitive involutions---fold exchange or parity flip---and every local continuation across the inverse-readout atlas is built from these same $\mathbb Z_2$ pieces. This makes the spinorial/M\"obius branch no longer just a visual metaphor at the image layer, but the explicit parity part of the local obstruction monodromy.

\paragraph{3.4.12 working interpretation.}
The inverse-readout block now has a local stability test. Once the calibration slope on a shell interval is known to stay away from zero, tunnel width ceases to be a merely qualitative marker and becomes a controlled coordinate for the shell radius. The remaining local obstruction is discrete---the parity branch---together with the obvious continuous failure mode at critical points of the calibration law.

\subsection*{C.12 Thirteenth 3.4 extension: shell atlas and obstruction support}
\addcontentsline{toc}{subsection}{C.12 Thirteenth 3.4 extension: shell atlas and obstruction support}

\begin{theorem}[Appendix support for Theorem~\ref{thm:mainline-3413}]
\label{thm:appendix-3413}
Let $I$ be a shell interval for the calibrated tunnel-width law $W_\tau$, and let
\[
  \Sigma_{\mathrm{crit}}=\{W_\tau'=0\},
  \qquad
  \Sigma_{\mathrm{par}}=\{\text{parity switch shells}\}.
\]
Then every connected component of $I\setminus(\Sigma_{\mathrm{crit}}\cup\Sigma_{\mathrm{par}})$ is a parity-preserving inverse chart on which $W_\tau$ is one-to-one. On compact subintervals of such a chart, the inverse is bi-Lipschitz and therefore the shell-consistency defect is a well-defined internal diagnostic of model quality rather than of inverse failure.
\end{theorem}

\begin{proof}
Outside $\Sigma_{\mathrm{par}}$ the parity class is fixed. Outside $\Sigma_{\mathrm{crit}}$, the derivative of $W_\tau$ never vanishes and therefore keeps a constant sign on each connected component, so the calibration map is strictly monotone there. Compactness then upgrades the local inverse estimate to a uniform chartwise estimate. The shell-consistency defect is therefore an internal observable on such a chart, and only the two obstruction sets can destroy the chart itself.
\end{proof}

\paragraph{3.4.13 working interpretation.}
The inverse detector-style layer is now stratified. Stable readout happens chart by chart, with sharp obstruction shells where either the calibration slope collapses or the transported parity changes. This gives a clean separation between chart-internal mismatch and genuine breakdown of inverse control.

\paragraph{3.4.11 working interpretation.}
The diagnostic image can therefore be treated as a small inverse instrument. After threshold calibration, tunnel width and shell parity recover the local generator shell and predict the remaining edge/opening readouts. The natural residual mismatch is then measured by the shell-consistency defect $\Delta_{\mathrm{diag}}$, which quantifies how well the image closes back to the same shell/parity model.

\paragraph{3.4.10 working interpretation.}
The image layer is therefore now coupled back into the mathematics in a controlled way: the diagnostic render may be discussed through tunnel width, edge intensification, and opening angle without pretending that visual similarity alone proves a detector interpretation. This keeps the observational hypothesis testable while preserving the structure-first discipline.

\paragraph{3.4.8 working interpretation.}
The shell classifier can therefore be read as a parity law. Once the generator radius $s(E)$ is fixed, the only remaining local ambiguity is whether the admissible shell cycle returns with the same sign or with reversed sign. Under a shell decomposition into regular segments and transverse reflector crossings, this becomes an explicit even/odd criterion. The accompanying M\"obius-reflected diagnostic render is not used as a proof of the residual scalar; it is included only to make the stable-kernel / branch-frontier split visible at the level of the current probe geometry.


\paragraph{3.4.9 working hypothesis.}
The current diagnostic render also motivates a cautious readout hypothesis. Because the plot was not tuned to detector imagery, its visual similarity to some detector-like photon displays should not be dismissed as a mere decorative analogy, but it should likewise not be elevated to an identification claim. Within the standing phase-shift / frame-shift reading, the safer interpretation is that detector images may display readout-relative shadows of the same branch-sensitive residual geometry. This can be tested at the level of tunnel width, opening angle, boundary intensification, and shell-parity response, without assuming that the diagnostic probe already reproduces a full detector event.

\paragraph{3.4.7 working interpretation.}
The analytic generator now also controls the probe decision itself. On every sufficiently small regular shell, the energy first collapses to a unique generator radius $s(E)$, and the visible branch data then reduces to the transported sign pair $\pm\sqrt{s(E)}$. The direct-convergence versus spinorial/M\"obius-lift dichotomy is therefore the $\mathbb Z_2$ shell classifier of the same generator family rather than an additional independent mechanism.

\paragraph{3.4.6 working interpretation.}
The even hierarchy is therefore no longer just a ladder of separate coefficients. Once the transported closure germ is analytic, the visible branch law is packaged by a single local generator $\Phi(\eta^2)$. The quartic shell, the delayed higher timing drifts, and the turning-value corrections then become coordinated readouts of one canonical germ rather than separate nonlinear sectors.

\paragraph{3.4.5 working interpretation.}
The higher nonlinear hierarchy is therefore already strongly constrained: once the quartic channel disappears, the next visible correction cannot re-enter arbitrarily, but only through the first surviving even closure coefficient. Operationally, the probe remains isochronous through all lower orders and only begins to drift at the exact order at which the closure hierarchy itself reopens.

\paragraph{3.4.4 working interpretation.}
The quartic correction is therefore not only structurally present; it is already operationally visible. The first failure of exact isochrony on a regular periodic shell is controlled by the same closure coefficient $c_4$ that generates the quartic branch correction. The probe may therefore read the same admissible shell through energy, through its unique positive turning value, or through its first nonlinear period drift, without opening a new scalar channel.


\paragraph{3.4.3 working interpretation.}
The nonlinear visible coefficient is therefore no longer free: once the exact branch potential is the transported readout of the same closure geometry, the quartic branch coefficient $\lambda$ is simply the positive transport of the closure coefficient $c_4$. The regular shell can then be indexed either by its energy $E$ or by its unique positive turning value $A(E)$, while the nonlinear strength itself is already fixed by the closure side.

\paragraph{3.4.2 working interpretation.}
The first concrete nonlinear model is therefore already visible: once the admissible branch closes to a one-field exact potential and no additional scalar opens at first nonlinear order, the restoring branch law falls into the quartic reversible family. The global probe can then read the same admissible shell either as an energy level or as a unique positive turning value, while the residual ambiguity remains exhausted by the transported $\mathbb Z_2$ pair.

\paragraph{3.4 working interpretation.}
The first move toward 3.4 is therefore not a proliferation of new local lemmas, but a globalization principle: once the reduced branch equation is exact, the restoring branch law becomes a one-dimensional potential dynamics, and the reversal probe becomes an energy-component test for direct closure versus spinorial lift.

\paragraph{Condensed 3.3 status.}
The local 3.3 line is therefore fully captured by four linked claims: a single transported residual scalar $\eta$, a Julia-boundary readout of that same scalar, a restoring branch dynamics for the local reversal profile, and a hybrid reversal probe with the two-way outcome ``direct convergence or spinorial/M\"obius lift''. Relative to the longer 3.3.0 working appendix, this is a structure-preserving compression rather than a deletion: the intermediate defect, jet, and probe lemmas are absorbed into the four main statements above.

\section*{Appendix D. Historical Release Notes and Version Record}
\addcontentsline{toc}{section}{Appendix D. Historical Release Notes and Version Record}

This appendix gathers the immediate Companion version trail and historical release notes so that the mathematical main line can close before the documentation apparatus begins. The live bibliography itself is kept as the terminal back-matter block of the document.

\clearpage
\section*{Historical release notes and roadmap blocks}
\addcontentsline{toc}{section}{Historical release notes and roadmap blocks}
The following materials are retained deliberately as archival proof-memory and release genealogy.
They preserve earlier release-stage notes, roadmap controls, and checkpoint explanations, but they are no longer kept in the opening sequence of the public document.
Accordingly, the front of the release now moves directly from abstract and contents into the nomenclature, while the historical front-matter materials remain available here for reconstruction, audit, and dependency tracing.

\section*{Archived Reader's Guide and working roadmap from the public v3.13.0 release}
\addcontentsline{toc}{section}{Archived Reader's Guide and working roadmap from the public v3.13.0 release}

The following note is preserved from the public v3.13.0 release stage. It records the release-facing reader's guide and roadmap exactly for that checkpoint; in the present release continuation it is archival proof-memory rather than a current status block.
It presents the existing mathematical line in one monolithic text, preserves the stabilized diagnostic material already carried by the document, and keeps the operative rule:
\emph{do not shorten, do not overwrite, only sharpen, normalize, and preserve.}

\subsection*{Archived version identity for the public v3.13.0 checkpoint}
\begin{itemize}[leftmargin=2em]
  \item Historical version: 3.13.0
  \item Historical date: 2026-04-07
  \item Concept DOI: \href{https://doi.org/10.5281/zenodo.19210728}{10.5281/zenodo.19210728}
  \item Historical status: public Zenodo release v3.13.0
  \item Historical public release DOI (v3.13.0): \href{https://doi.org/10.5281/zenodo.19446872}{10.5281/zenodo.19446872}
  \item Historical version-specific DOI for v3.13.0: \href{https://doi.org/10.5281/zenodo.19446872}{10.5281/zenodo.19446872}
  \item Previous Zenodo releases: v3.3.5 (\href{https://doi.org/10.5281/zenodo.19396143}{10.5281/zenodo.19396143}), v3.4.0 (\href{https://doi.org/10.5281/zenodo.19410949}{10.5281/zenodo.19410949}), v3.7.1 (\href{https://doi.org/10.5281/zenodo.19412478}{10.5281/zenodo.19412478}), v3.9.0 (\href{https://doi.org/10.5281/zenodo.19423689}{10.5281/zenodo.19423689})
  \item Main-paper relation: Explicit mathematical companion to the current Sphaera master line~\cite{brendecke2026main}
\end{itemize}

\subsection*{Archived working note from v3.13.0}
This public release proceeds from the public v3.11.0 Companion release in one preserved monolithic text. Content-wise it now keeps together three central gains:
\begin{enumerate}[leftmargin=2em,label=(\roman*)]
  \item the stabilized residual-record interface and its associated structural tools,
  \item the reproducible M\"obius-shell diagnostic suite for the live residual scales, and
  \item the OP~5 packet, breathing-envelope, and frame-flicker visual audit line.
\end{enumerate}
Accordingly, that archived public release should be read as the stabilized public checkpoint of its then-current monolithic chain: it preserved the mathematical line of that stage while keeping the front cleaner and the live burdens more sharply routed.

\subsection*{Archived primary working goals recorded at v3.13.0}
\begin{enumerate}[leftmargin=2em]
  \item Present the current Companion mathematics in one public monolithic release.
  \item Keep the stabilized generator--atlas--holonomy core and the residual-record package $\mathfrak R(B)=\bigl(\Phi,\Delta_{\mathrm{diag}},\varepsilon_{\mathrm{par}},\chi_{\mathrm{par}},\widehat\chi_{\mathrm{par}}\bigr)$ intact as the faithful visible interface.
  \item Preserve the M\"obius-shell diagnostics as reproducible, active-PDF figures so that residual drift, continuation, sensitivity, and sector persistence remain directly auditable.
  \item Preserve the OP~5 packet-wave, breathing-envelope, and frame-flicker visual line in one release-safe appendix block.
  \item Preserve the monolithic document without shortening historical material.
\end{enumerate}

\subsection*{Archived release priorities recorded at v3.13.0 toward Zenodo deposition}
\begin{enumerate}[leftmargin=2em]
  \item Turn the current OP~5 packet/phase/epsilon analysis into one cleaner structural branch theorem with a small admissible branch load and a non-escape criterion.
  \item Keep OP~1 on the transported-branch reading and compress the remaining burden to the discrete correction channel without reopening settled continuum structure.
  \item Normalize the document language so that internal patch chronology remains archival only, while the main line reads as one coherent proof route.
  \item Collect the now-stabilized transport, readout, checkpoint, and epsilon-trap principles as explicit reusable tools for the next theorem blocks.
  \item Use the present v3.13.0 public release as the stabilized next checkpoint because the OP transport closure architecture is now integrated more consistently across the monolithic text rather than remaining split between new theorem blocks and unreconciled historical passages.
\end{enumerate}

\subsection*{Continuity of the v3.5.2 checkpoint in the present release}
The present file preserves the former \emph{v3.5.2 abstract-polish checkpoint} as compatible earlier material within the current release.
Concretely, the 3.5.2 achievements remain part of the present monolith:
\begin{enumerate}[leftmargin=2em,label=(\roman*)]
  \item the residual-record interface \(\mathfrak R(B)\) and the reconstruction normal form remain unchanged;
  \item the readout-relative lemma methodology, structural lift logic, and Theorem~7 / Hypotheses~7.4--7.6 remain preserved as dependencies;
  \item the former OP~2 operator burden is no longer carried as a separate primary front, but is preserved in reclassified form inside the later OP-first status reading;
  \item the abstract-polish framing of 3.5.2 is therefore not discarded, but sharpened and transported into the later residual-controlled and boundary-layer language.
\end{enumerate}
Accordingly, the development line may be read as:
\[
\begin{aligned}
&\text{v3.5.2 interface normalization}
\Longrightarrow
\text{v3.6 conservative frontend}
\Longrightarrow
\text{v3.7 OP-first stabilization}\\[0.2em]
&\Longrightarrow
\text{v3.8 residual / boundary-layer sharpening}.
\end{aligned}
\]


\subsection*{Archived current mathematical status at the public v3.13.0 release stage}
For the current release line, the reader should keep one status separation in mind before entering the preserved historical layers below:
\begin{itemize}[leftmargin=2em]
  \item The stabilized generator--atlas--holonomy sector from v3.4 is already closed at the structural level and has been compressed into the admissible residual record
  \[
  \mathfrak R(B)=\bigl(\Phi,\Delta_{\mathrm{diag}},\varepsilon_{\mathrm{par}},\chi_{\mathrm{par}},\widehat\chi_{\mathrm{par}}\bigr).
  \]
  \item The current v3.5.x step is therefore no longer a search for new visible primitives, but a normalization of comparison, transport, and interpretation through that record.
  \item Historical OP~2 should no longer be read as a separate live operator problem; its T1/T2/T3 blocks are closed, and the former T4 frame-offset part is absorbed into the residual quantitative completion of OP~1.
  \item Historical OP~3 is closed at the present structural level; it may still be enriched later, but it is no longer one of the primary live fronts.
  \item Historical OP~5 has been reduced to the photon-frame Step~(B) remainder only; Step~(A) is already closed in the present text.
  \item The genuinely live fronts are therefore the residual continuum-to-measurement correction for \(\bsym\), controlled emergence with explicit constants and bounds, the quantitative mass hierarchy, the three-generation splitting law, the compact Dixon-algebra reformulation, and the remaining Theta-sum cancellation step.
\end{itemize}

\paragraph{Release-stage caution.}
The present reading is strong enough for this release, but it is intentionally stated in stabilized form: the remaining live fronts are read as residual downstream tasks \emph{in the present stabilized reading}. This is stronger than a loose hypothesis, because it rests on the exact representability and conservative-frontend results of Sections~3.5--3.6; but it is still a mathematically revisable reading if a later residual OP were to exhibit a genuinely missing visible degree of freedom. No such counter-indication is currently visible in the Companion chain.

\paragraph{How to read the preserved historical roadmap blocks.}
The longer roadmap sections that follow from earlier cycles are retained deliberately as proof-memory and dependency control. They should now be read as \emph{historical support strata}: useful for reconstruction, closure auditing, and tool extraction, but no longer identical with the short list of primary live questions at the head of the present line.

\subsection*{What the 3.5.x line achieved, and what 3.6 now fixes first}
At this point the 3.5.x line should be read as a \emph{completed interface-normalization stage} rather than as an unfinished search for further visible structure.
Its concrete achievements are:
\begin{enumerate}[leftmargin=2em]
  \item the full 3.4 generator--atlas--holonomy package has been compressed into the admissible residual record
  \[
  \mathfrak R(B)=\bigl(\Phi,\Delta_{\mathrm{diag}},\varepsilon_{\mathrm{par}},\chi_{\mathrm{par}},\widehat\chi_{\mathrm{par}}\bigr);
  \]
  \item visible reconstruction has been normalized through the record and the chartwise shell data \((s,\pi)\);
  \item record-level comparison and transport have been fixed through admissible morphisms and residual-record isomorphisms;
  \item the faithful interpreter principle has made the stabilized visible sector exactly representable on record level;
  \item the historical OP-layer has been reclassified so that closed, reduced, and genuinely open fronts are no longer mixed.
\end{enumerate}
Accordingly, the next step should \emph{not} attempt to reopen the visible interface.
A mathematically clean 3.6 target is instead:
\begin{enumerate}[leftmargin=2em,label=(\roman*)]
  \item prove the conservative-frontend principle explicitly, i.e. every later language layer must factor through the residual-record interpreter and may not introduce any new visible primitive;
  \item sharpen the remaining live fronts only as downstream quantitative tasks: prime-lattice/continuum correction for \(\bsym\), explicit constants and bounds for controlled emergence, mass hierarchy and three-generation splitting, compact Dixon-algebra reformulation, and the unresolved Theta-sum cancellation block;
  \item prepare the eventual ZFM-/language-layer only after the conservative factorization theorem is in place.
\end{enumerate}
In short: \emph{3.5.x stabilized the interface; 3.6 begins the stabilization of the admissible languages above it.}

\subsection*{Reading discipline for the next stage}
Every new subsection introduced below is to be read under four status labels that must be kept distinct:
\begin{enumerate}[leftmargin=2em,label=(\alph*)]
  \item structurally forced,
  \item projectively read,
  \item interpretive overlay separated from theorem status,
  \item reserved for later proof expansion.
\end{enumerate}


\subsection*{Archived mathematics-first completion path realized in v2.27.0}

The present major release is to be read as a \emph{mathematics-first completion cycle}. The computational chapter opened in v2.24.0--v2.26.0 is to remain in place as a visible teaser and reserved architecture, but it is no longer the primary work driver for the immediate next version. For the next cycle, the operative rule is instead:
\begin{center}
\emph{close the mathematics first; let the proof engine wait for a sufficiently closed mathematical layer.}
\end{center}

Accordingly, the completion path realized in v2.27.0 is organized into five mathematical work packages:
\begin{enumerate}[leftmargin=2em,label=\textbf{WP\arabic*:}]
  \item \textbf{Core structural closure.} Complete and normalize the remaining structural chains around ground condition, HC emergence, closure admissibility, and the Millennium matrix as global HC realization.
  \item \textbf{Proof-chain hardening.} Close missing lemma-to-theorem dependencies, remove duplicated or merge-residual argument fragments, and restore explicit proof-memory where shortened passages still rely on hidden expansion.
  \item \textbf{Projector / transport mathematics.} Systematize projector logic, carrier-group-readout transport, and the semantic/tensorial kernel interface so that projection steps are mathematically explicit before they are computationally compressed.
  \item \textbf{Tools as mathematical appendices.} Gather structural tools, computational shortcuts, closure tricks, and trace/branch transfer rules as named mathematical instruments with scope conditions, rather than letting them remain scattered as implicit know-how.
  \item \textbf{Computational strand on hold-but-present.} Retain MML, MMA, MCP, certificates, and host export as a prepared future chapter, but do not let that layer outrun the mathematical substance it is supposed to compress.
\end{enumerate}

In practical terms, v2.27.0 therefore prioritizes mathematical completion over computational expansion. The computational chapter remains part of the Companion and should be preserved, cited, and gently normalized where needed, but only as a secondary architecture until the mathematical layer is sufficiently closed. The version-specific DOI for the present v2.27.0 release is \href{https://doi.org/10.5281/zenodo.19266666}{10.5281/zenodo.19266666}.


\subsection*{Operational mathematical closure agenda}

To make the mathematics-first rule operational rather than merely declarative, the next revision cycle is to be executed along the following closure agenda:
\begin{enumerate}[leftmargin=2em,label=\textbf{MA\arabic*:}]
  \item \textbf{Definition pass.} Re-read and normalize the basic definitions that the later chains depend on: collapse state, closure, admissibility, projector sector, compensator action, transport layer, carrier/group/readout distinction, and Millennium-matrix status.
  \item \textbf{Dependency pass.} Identify theorem blocks whose proofs still rely on compressed or implicit dependencies, then make the lemma chain explicit enough that each forward step can be reconstructed without external memory.
  \item \textbf{Transport pass.} Isolate the mathematics of projection, frame transport, Lorentz-mediated readout, and disjoint-sector return so that these are expressed first as mathematical maps and only later as computational metaphors.
  \item \textbf{Appendix-tools pass.} Move recurrent structural tricks into named mathematical tools with hypotheses and scope, so that later proof reuse becomes mathematically transparent rather than conversationally remembered.
  \item \textbf{Stability pass.} Check the normalized chapter sequence against merge-residual duplication, outdated internal cross-references, and proof steps that were historically preserved but not yet fully re-hardened.
\end{enumerate}

\subsection*{Priority clusters for v2.27.0}

The following clusters should receive mathematical priority before any further major proof-engine expansion:
\begin{enumerate}[leftmargin=2em,label=\textbf{PC\arabic*:}]
  \item \textbf{Ground condition $\Rightarrow$ HC $\Rightarrow$ confinement.} This chain should be rewritten as a clean mathematical dependency, with the Heisenberg compensator treated as emergent from the ground condition rather than as an independently floating postulate.
  \item \textbf{Projector and closure mathematics.} The admissible-sector logic should be made fully explicit, including which projectors are primitive, which are derived, and how closure is checked at the mathematical rather than merely intuitive level.
  \item \textbf{Millennium matrix as mathematical object.} The Millennium matrix should first be stabilized as a global mathematical realization of the HC/closure architecture before any stronger computational reading is allowed to dominate the exposition.
  \item \textbf{Trace / branch mathematics.} Probe traces, branch transfer, Hopf-transition logic, and return structure should be stated as mathematical transformations with hypotheses and admissibility conditions, not only as algorithmic motifs.
  \item \textbf{Proof-memory restoration.} Where earlier versions compressed long chains into structural shorthand, the next cycle should restore enough explicit mathematics that the shorthand remains auditable and publication-safe.
\end{enumerate}

\subsection*{What is deliberately deferred beyond v2.27.0}

The following lines are explicitly preserved but not treated as primary completion goals for the next release cycle:
\begin{enumerate}[leftmargin=2em,label=\textbf{D\arabic*:}]
  \item full MML language formalization,
  \item full MMA/MCP execution refinement,
  \item standalone proof-kernel extraction,
  \item stronger universality claims,
  \item operating-system style runtime architecture.
\end{enumerate}
These elements remain valuable, but in the current sequencing they are to follow mathematical closure rather than compete with it.

\subsection*{Mathematical work matrix for the next pass}
To make the mathematics-first agenda operational, the next pass is organized as a compact work matrix rather than as a loose wish list.
\begin{center}
{\scriptsize\setlength{\tabcolsep}{2pt}\begin{tabular}{@{}p{0.18\textwidth}p{0.24\textwidth}p{0.27\textwidth}p{0.06\textwidth}@{}}
\toprule
\textbf{Block} & \textbf{Target state} & \textbf{Main dependencies} & \textbf{Priority}\\
\midrule
Ground condition $\to$ HC $\to$ confinement & explicit statement chain, normalized notation, and no ambiguity about emergence vs. axiomatics & core structural assumptions and septinion confinement language & A\\
Projector and closure mathematics & sharpened lemmas, admissibility conditions, and stable projector notation & closure operator layer, readout discipline, sector decomposition & A\\
Millennium matrix as mathematical object & operator-level formulation, domain/codomain clarity, and clarified relation to HC & projector layer, transport layer, compensator interpretation & A\\
Trace / branch / Hopf transport mathematics & clean separation of trace, branching, return, and transport rules & probe framework, disjointness, Lorentz-mediated transport & B\\
Dependency chain hardening & theorem/lemma dependency map cleaned and checked for merge artifacts & current theorem chain, proposition and lemma references & A\\
Structural tools appendix & collect reusable mathematical shortcuts without shortening proofs & established proof moves, normalized notation & B\\
Proof-memory restoration / archival recovery & recover or reintegrate historically displaced proof material where needed & older Companion states, recovered large basis, appendix chronology & B\\
Computational teaser retention & keep chapter visible, but frozen except for compatibility edits & mathematical closure progress & C\\
\bottomrule
\end{tabular}}
\end{center}

A-priority blocks should drive the next version boundary. B-priority blocks should be developed in service of those A-blocks. C-priority blocks should be touched only when required for compatibility with the mathematical reorganization.

\subsection*{Dependency map of the mathematical core}
To prevent the next pass from drifting into local edits without structural order, the mathematical core is also organized as a dependency map. The intended reading is simple: if an upstream block remains unstable, all downstream blocks that consume it can only be provisionally edited.
\begin{center}
\begin{tabular}{p{0.28\textwidth}p{0.30\textwidth}p{0.30\textwidth}}
\toprule
\textbf{Core block} & \textbf{Direct prerequisite} & \textbf{Main downstream consumers}\\
\midrule
Ground condition and collapse picture & none; foundational starting point & HC emergence, confinement thesis, collapse-state language, admissibility reading\\
HC emergence and confinement chain & ground condition, collapse picture & projector logic, trace/branch admissibility, Millennium-matrix interpretation\\
Projector / closure mathematics & HC chain, admissibility normalization & branch filtering, Boolean readout, proof templates, certificate layer\\
Transport / Lorentz layer & collapse picture, projector sector discipline & return-map reading, branch transport, reflector-like structural reversals\\
Millennium matrix as mathematical object & HC chain, projector/closure layer, transport layer & computational teaser chapter, co-processor view, host-export discipline\\
Trace / branch mathematics & projector logic, transport layer, admissibility conventions & Hopf branching examples, trace reconciliation, branch templates\\
Tool and appendix lemmas & all stabilized upstream definitions & proof hardening, auditability, future proof-engine extraction\\
\bottomrule
\end{tabular}
\end{center}

In operational terms, this means that the next mathematical pass should proceed roughly in the order
\[
\begin{aligned}
&(\text{ground condition} \Rightarrow \text{HC} \Rightarrow \text{confinement}) \prec (\text{projector/closure})\\
&\prec (\text{transport}) \prec (\text{Millennium matrix, trace, branch})\\
&\prec (\text{appendix tools and supporting lemmas}).
\end{aligned}
\]
Any local repair that violates this order may still be useful as a note, but it should not be mistaken for structural closure.


\noindent\textbf{Definition-and-notation pass for the next mathematical stage.} Before the next major release, the Companion should also undergo a focused definition-and-notation pass. The purpose of this pass is not to shorten the document, but to reduce merge drift, ambiguous symbol reuse, and status confusion between theorem, lemma, tool, hypothesis, and teaser language.

\paragraph{Mathematically fixed for the current line.}
For the current Companion line, the following should be treated as fixed structural anchors unless explicitly revised in a later release: the collapse-state viewpoint as the local readout picture; the phase-shifted measurement-frame perspective; the ground-condition $\Rightarrow$ Heisenberg-compensator $\Rightarrow$ confinement chain; the projector/closure discipline for admissible sectors; the standing triolen architecture with one orthogonal component; and the imaginary convention already adopted for the relevant octonionic/triadic units.

\paragraph{Notation to normalize before v2.27.0.}
Notation should be normalized especially in the following zones: projector symbols and their admissibility role; closure notation versus stabilization notation; transport/Lorentz notation across layers; Millennium-matrix notation when used as a mathematical object rather than as a computational teaser; trace versus probe-trace notation; branch/Hopf notation; and the distinction between theorem-level mathematical statements and later computational analogies. In particular, the same symbol should not silently alternate between mathematical object, physical readout, and computational metaphor.

\paragraph{Items that must remain hypotheses.}
Some parts of the Companion should remain explicitly flagged as hypotheses, interpretive proposals, or deferred claims rather than being silently upgraded into theorems. This applies in particular to the Higgs/coupling interpretation, to stronger black-hole or event-horizon extrapolations, to stronger computational-universality language, and to any claim that would identify the computational chapter with a finished proof engine rather than a prepared architecture. The next mathematical pass should therefore favor closure and definition-hardening over conceptual inflation.

\paragraph{Operational rule for the next pass.}
For the next revision cycle, each touched block should be tagged mentally by the question: \emph{definition, theorem, lemma, tool, hypothesis, or teaser?} The next major mathematical gain is expected to come more from making these statuses explicit and stable across the Companion than from adding further new conceptual branches.


\subsection*{Target-state matrix for v2.27.0}

\small\begin{longtable}{p{0.22\textwidth}p{0.17\textwidth}p{0.24\textwidth}p{0.25\textwidth}}
\toprule
\textbf{Core block} & \textbf{Current status} & \textbf{Target state for v2.27.0} & \textbf{Minimal completion condition} \\
\midrule
\endfirsthead
\toprule
\textbf{Core block} & \textbf{Current status} & \textbf{Target state for v2.27.0} & \textbf{Minimal completion condition} \\
\midrule
\endhead
Ground condition / collapse picture & Structurally central, but still partly distributed across explanatory passages. & Gathered into a visibly controlling mathematical block. & The collapse-state picture is stated once in stable mathematical language and used consistently as the local readout frame. \\

HC emergence / confinement chain & Present as a guiding structural principle, but not yet fully normalized as a dependency spine. & Explicitly treated as a mathematical chain: ground condition $\Rightarrow$ HC $\Rightarrow$ confinement. & The chain appears in one normalized theorem/lemma sequence and is referenced consistently downstream. \\

Projector / closure mathematics & Strongly present, but notation and admissibility wording still need tightening. & Treated as the first closed admissibility layer of the framework. & Projector symbols, closure symbols, and admissibility language are normalized and no longer drift between sections. \\

Transport / Lorentz layer & Conceptually active, but still mixed between mathematical transport and physical interpretation. & Clarified as the transport layer relating frames without overloading its physical extrapolations. & A stable mathematical transport notation is fixed and separated from speculative physical remarks. \\

Millennium matrix as mathematical object & Partly present beside its computational reading. & Brought back under mathematical control before any stronger proof-engine use. & The matrix is stated first as a mathematical object/operator, with the computational reading clearly subordinate. \\

Trace / branch mathematics & Available in structural language, but still partly entangled with later computational analogies. & Recast as a mathematical trace/branch layer before proof-engine expansion resumes. & Trace, probe-trace, branch, and Hopf language are normalized and tied to mathematical dependence rather than metaphor. \\

Tool / appendix lemmas & Growing, but not yet fully separated from main-line dependencies. & Reorganized as supporting tools rather than semi-primary arguments. & Each tool is clearly marked as tool/appendix support and no longer silently carries theorem-level load without declaration. \\
\bottomrule
\end{longtable}

\paragraph{Use of the matrix.}
This target-state matrix is intended as an operational filter for the next revision cycle. A block should only be counted as \emph{mathematically advanced} if its current status has moved measurably toward the declared target state and if its minimal completion condition has been checked explicitly.

\subsection*{Proof-dependency checklist for v2.27.0}
To keep the next mathematical pass referentially stable, the following dependency checklist should be treated as a minimal proof-control layer. The point is not to produce exhaustive dependency graphs for every sentence, but to ensure that the main mathematical blocks are supported by a sufficiently explicit chain of definitions, lemmas, transport statements, and appendix tools.

\begin{itemize}
    \item \textbf{Ground condition / collapse picture.} Must rest on stable base definitions, a clear distinction between structural state and measurement readout, and an explicit statement of which consequences are structural and which remain interpretive.
    \item \textbf{HC emergence / confinement chain.} Must depend only on the ground condition and the confinement logic, not on later computational analogies. For v2.27.0 the chain ground condition $\Rightarrow$ HC $\Rightarrow$ confinement should be readable without hidden appeals to deferred sections.
    \item \textbf{Projector / closure mathematics.} Requires normalized projector notation, at least one closure lemma family, and stable cross-references into the appendix tools where routine projector identities are collected.
    \item \textbf{Transport / Lorentz layer.} Must state exactly which structural transformations are taken as transport rules, which are theorem-level, and which remain hypothesis-level interpretations. Dependencies on projector compatibility should be explicit.
    \item \textbf{Millennium matrix as mathematical object.} Should depend on the already stabilized projector/closure and transport layer, rather than carrying foundational burden by itself. For v2.27.0 the matrix should be mathematically positioned after the base closure apparatus, not before it.
    \item \textbf{Trace / branch mathematics.} Requires a small stable family of supporting lemmas on admissible continuation, branch exhaustion, and trace comparison. The next pass should avoid leaving branch claims supported only by prose.
    \item \textbf{Appendix tools and supporting lemmas.} These should be treated as genuine support infrastructure: enough to stabilize the main line, but not so invasive that the main thread becomes unreadable. The key requirement is reference stability and not merely accumulation.
\end{itemize}

For the purpose of the next release threshold, a proof block should count as \emph{referentially stable} only if its direct supporting definitions and lemmas can be named without ambiguity, found in the current Companion without version drift, and read in an order compatible with the mathematical dependency map already given above.


\subsection*{Reference-stability pass for v2.27.0}

For the mathematics-first pass toward v2.27.0, internal references are themselves treated as part of structural closure. The next public version should not only add arguments, but also stabilize the minimal web of labels, theorem references, and section anchors on which the A-priority mathematical blocks depend.

\begin{itemize}
  \item \textbf{Definition stability.} The core labels for the ground condition, Heisenberg-compensator emergence, confinement statement, projector/closure layer, transport layer, and Millennium-matrix mathematical section should remain fixed under the next hardening pass unless an explicit renumbering note is added.
  \item \textbf{Dependency stability.} Every A-priority theorem or proposition should point only to labels that are themselves already fixed at the level of definition, lemma, theorem, or supporting tool. Soft roadmap anchors should not be used as silent proof dependencies.
  \item \textbf{Appendix-tool stability.} If a proof uses an appendix tool, the tool must be referenced by a stable label and its status must remain visibly marked as tool, lemma, or hypothesis.
  \item \textbf{Hypothesis transparency.} Any block that still depends on a hypothesis label must keep that hypothesis status explicit in the surrounding prose and may not inherit theorem-strength merely by cross-reference.
  \item \textbf{Public-release threshold.} For v2.27.0, a block counts as reference-stable only if its internal labels, backward references, and appendix-tool dependencies survive a clean compilation pass without ad hoc renaming.
\end{itemize}

Operationally, the next mathematical stage should therefore include a short explicit reference-stability check before public release: (i) labels on A-priority definitions and theorems are frozen, (ii) proof dependencies resolve only to mathematically declared objects, (iii) appendix tools referenced in the main line remain stable and typed, and (iv) theorem-to-hypothesis drift is not introduced by accidental cross-linking.

\subsection*{Block-by-block closure checklist for the next pass}
For each A-priority block, the next pass should explicitly separate three statuses: what remains to be closed, what may count as done for v2.27.0, and what may remain open without blocking the release threshold.

\begin{itemize}
    \item \textbf{Ground condition / collapse picture.} Missing: final normalized wording and explicit distinction between structural and observational language. Done for v2.27.0: stable definition and one clear theorem-level consequence. May remain open: broader interpretive extensions.
    \item \textbf{HC emergence / confinement chain.} Missing: explicit derivational chain and at least one supporting lemma that fixes the confinement reading. Done for v2.27.0: readable chain ground condition $\Rightarrow$ HC $\Rightarrow$ confinement. May remain open: stronger physical extrapolations.
    \item \textbf{Projector / closure mathematics.} Missing: normalized notation and a compact closure lemma family. Done for v2.27.0: a stable projector/closure core with appendix support. May remain open: larger auxiliary projector calculus.
    \item \textbf{Transport / Lorentz layer.} Missing: clean split between structural transport statements and hypothesis-level interpretive claims. Done for v2.27.0: one stable transport layer attached to the projector/closure block. May remain open: more speculative extensions.
    \item \textbf{Millennium matrix as mathematical object.} Missing: definitive positioning relative to HC, projector structure, and transport. Done for v2.27.0: mathematically explicit role after closure/transport stabilization. May remain open: the computational reading beyond teaser level.
    \item \textbf{Trace / branch mathematics.} Missing: compact lemma support for admissible continuation, exhaustion, and trace comparison. Done for v2.27.0: one stable branch/trace family sufficient for the main line. May remain open: extended proof-program analogies.
    \item \textbf{Appendix tools / supporting lemmas.} Missing: sharper separation between indispensable support and purely archival material. Done for v2.27.0: tool appendix that stabilizes references without overgrowing the main line. May remain open: secondary tool families.
\end{itemize}

\subsection*{Explicit subversion discipline}

From this point onward, intermediate work states are to be recorded explicitly as local subversions whenever substantial planning, normalization, recovery, or integration work has been performed without yet justifying a new Zenodo release. The intended pattern is:
\begin{center}
\texttt{public release v2.27.0 -> local subversion v2.27.1 -> further local subversions as needed -> next public milestone}
\end{center}
This rule is not cosmetic. It is meant to prevent silent rollback, silent overwrite, and accidental compression of the project memory inside a single file name.


\subsection*{Concrete chapter scaffold to be unfolded in the next revision cycle}

\paragraph{Part I --- Structural Core}
\begin{enumerate}[leftmargin=2em,label=\arabic*.]
  \item \textbf{Introduction}
    \begin{enumerate}[label*=\arabic*., leftmargin=2.5em]
      \item Motivation
      \item Scope of this version
      \item Relation to the main paper
      \item Structural strategy: structure first, closure before projection
    \end{enumerate}
  \item \textbf{Structural Preliminaries}
    \begin{enumerate}[label*=\arabic*., leftmargin=2.5em]
      \item Basic structural units
      \item Orthogonality constraint
      \item Complex and imaginary structural setting
      \item Frame shift of measurement
      \item Structural closure intuition
    \end{enumerate}
  \item \textbf{Ground Condition and Emergent Closure}
    \begin{enumerate}[label*=\arabic*., leftmargin=2.5em]
      \item Ground condition
      \item From ground condition to HC
      \item HC as compensating constraint
      \item From HC to closure
      \item Protected invariants
      \item Proposition block: Ground Condition $\Rightarrow$ HC $\Rightarrow$ Closure
    \end{enumerate}
  \item \textbf{Septinion / Seption Structural Embedding}
    \begin{enumerate}[label*=\arabic*., leftmargin=2.5em]
      \item From triadic units to higher structural coupling
      \item Octonionic and higher assembly logic
      \item Septinion / seption confinement
      \item Structural consequences
    \end{enumerate}
  \item \textbf{Structural Master Equation --- First Layer}
    \begin{enumerate}[label*=\arabic*., leftmargin=2.5em]
      \item Why a master equation emerges
      \item Pure structural form
      \item Term-by-term structural origin
      \item Structural vs.\ projective terms
      \item First consistency remarks
    \end{enumerate}
\end{enumerate}

\paragraph{Part II --- Probe Sector}
\begin{enumerate}[leftmargin=2em,label=\arabic*.]
  \setcounter{enumi}{5}
  \item \textbf{Probe Geometry}
    \begin{enumerate}[label*=\arabic*., leftmargin=2.5em]
      \item Two probes moving toward each other
      \item Each probe as a triolic structure
      \item Orthogonal third group
      \item Definition of the search space
    \end{enumerate}
  \item \textbf{Probe Dynamics and Field Structure}
    \begin{enumerate}[label*=\arabic*., leftmargin=2.5em]
      \item Structural interaction regime
      \item Closure in the probe sector
      \item Probe sector equation
      \item Search, trace, and structural read-off
    \end{enumerate}
  \item \textbf{Probe Tools and Guided Search}
    \begin{enumerate}[label*=\arabic*., leftmargin=2.5em]
      \item Frequency and pattern detection
      \item Routing and candidate compression
      \item Transition guidance
      \item Convergence tools
      \item Optional statistical layer
      \item Interpretation as structural toolkit
    \end{enumerate}
\end{enumerate}

\paragraph{Part III --- Technical Normalization}
\begin{enumerate}[leftmargin=2em,label=\arabic*.]
  \setcounter{enumi}{8}
  \item \textbf{Structural Tools and Computational Shortcuts}
    \begin{enumerate}[label*=\arabic*., leftmargin=2.5em]
      \item Why tool normalization is needed
      \item Closure-tools
      \item Frame-shift-tools
      \item Projection-tools
      \item Probe-tools
      \item Equivalence-tools
      \item Read-off / trace-tools
      \item Standard form for every tool
    \end{enumerate}
\end{enumerate}

\paragraph{Part IV --- Interpretive Bridge}
\begin{enumerate}[leftmargin=2em,label=\arabic*.]
  \setcounter{enumi}{9}
  \item \textbf{Projection Logic and Interpretive Discipline}
    \begin{enumerate}[label*=\arabic*., leftmargin=2.5em]
      \item From structure to projection
      \item What is forced, what is read, what is hypothesized
      \item Dirac / QFT / Einstein-type readings
      \item Higgs as coupling instance or mean coupling point
    \end{enumerate}
\end{enumerate}

\paragraph{Part V --- Consolidation}
\begin{enumerate}[leftmargin=2em,label=\arabic*.]
  \setcounter{enumi}{10}
  \item \textbf{Consistency Layer}
    \begin{enumerate}[label*=\arabic*., leftmargin=2.5em]
      \item Internal structural consistency
      \item Probe/main compatibility
      \item Closure robustness
      \item Remaining soft spots
    \end{enumerate}
  \item \textbf{Outlook to v1.2.0}
    \begin{enumerate}[label*=\arabic*., leftmargin=2.5em]
      \item What comes next
      \item Planned lemma expansion
      \item Publication path
    \end{enumerate}
\end{enumerate}

\paragraph{Appendix scaffold}
\begin{enumerate}[leftmargin=2em,label=\Alph*.]
  \item Symbol table
  \item Structural tools catalog
  \item Technical lemmas
  \item Computational shortcuts
  \item Projection tables
  \item Probe sector schematics
\end{enumerate}

\subsection*{Minimum formal architecture reserved for v\docversion}
The next expansion pass should lock at least the following formal blocks:
\begin{itemize}[leftmargin=2em]
  \item Proposition: Ground Condition implies emergence of HC.
  \item Proposition: HC implies structural closure.
  \item Proposition: Closure confines the admissible dynamics to the septinion/seption regime.
  \item Proposition: The probe sector inherits the same closure discipline under orthogonal coupling.
  \item Proposition: The first-layer master equation is structurally forced by the closure architecture.
  \item Auxiliary lemmata: orthogonality, frame shift, confinement, probe compatibility, read-off consistency.
\end{itemize}

\subsection*{Practical build order}
For the actual writing sequence of the next cycle:
\begin{enumerate}[leftmargin=2em]
  \item Ground Condition / HC / Closure
  \item Probe sector
  \item Master equation (first layer)
  \item Tool normalization
  \item Interpretive bridge
  \item Consistency layer
  \item Appendix fill-in
\end{enumerate}


% ==================================================================
% v2.27.5 PROOF-CHAIN ROADMAP
% ==================================================================

\clearpage
\section*{v2.27.5 Proof-Chain Roadmap: Identified Gaps\\Closure Targets}
\label{sec:v2275-roadmap}
\addcontentsline{toc}{section}{v2.27.5 Proof-Chain Roadmap: Identified Gaps and Closure Targets}

\noindent\fbox{\begin{minipage}{0.94\textwidth}
\small\textbf{Status of this section.}
Version~2.27.4 completed a full structural scan of the companion.
This section records the identified proof-chain gaps in order of priority,
and defines the closure targets for the next local subversion sequence.
It is a working roadmap, not a release statement.
The DOI and version references inside this preserved v2.27.x roadmap belong to the original historical checkpoint and are retained only as archival proof-memory. They are not current release-status statements for the present Companion line.
\end{minipage}}

\bigskip

\subsection*{Proof-chain inventory (v2.27.4 scan)}

The structural scan of v2.27.4 identified the following main proof chain:

{\small
\[
\begin{aligned}
\text{Axiom}
&\Rightarrow \underbrace{\text{Semantic Kernel}}_{\text{def}}
\Rightarrow \underbrace{\text{HC Forcing}}_{\text{thm}}
\Rightarrow \underbrace{\text{Closure Stable}}_{\text{lem}}\\[2pt]
&\Rightarrow \underbrace{\text{Diagonal Traversal}}_{\text{lem}}
\Rightarrow \underbrace{C = \Pi^{(\mathcal{M})}_{\mathrm{spec}}}_{\text{thm}}
\Rightarrow \underbrace{\text{SM Realization}}_{\text{cor}}
\end{aligned}
\]
}

The following table rates each link on four axes:
\emph{proof type} (Full / Short / Sketch / Structural),
\emph{explicit refs} (references cited in proof body),
\emph{dependency gap} (missing intermediate step),
and \emph{priority} for the next closure pass.

\medskip
{\footnotesize
\setlength{\LTleft}{0pt}
\setlength{\LTright}{0pt}
\setlength{\tabcolsep}{4pt}
\renewcommand{\arraystretch}{1.18}
\begin{longtable}{>{\raggedright\arraybackslash}p{0.26\textwidth}>{\raggedright\arraybackslash}p{0.12\textwidth}>{\centering\arraybackslash}p{0.07\textwidth}>{\raggedright\arraybackslash}p{0.39\textwidth}>{\centering\arraybackslash}p{0.06\textwidth}}
\toprule
\textbf{Chain link} & \textbf{Type} & \textbf{Refs} & \textbf{Gap / note} & \textbf{Prio} \\
\midrule
\endfirsthead
\toprule
\textbf{Chain link} & \textbf{Type} & \textbf{Refs} & \textbf{Gap / note} & \textbf{Prio} \\
\midrule
\endhead
\midrule
\multicolumn{5}{r}{\emph{Continued on next page}}\\
\endfoot
\bottomrule
\endlastfoot
\shortstack[l]{\texttt{def:semantic-}\\\texttt{structural-kernel}} & Short (9L) & 0 & Definitions; no gap. & -- \\
\shortstack[l]{\texttt{prop:semantic-}\\\texttt{forcing-chain}} & Short & 0 & Structural; no proof. & B \\
\shortstack[l]{\texttt{thm:hc-}\\\texttt{forcing-theorem}} & Short (14L) & 2 & Cites \texttt{lem:re-identification-lemma}; sound. & B \\
\shortstack[l]{\texttt{lem:closure-}\\\texttt{forcing-lemma}} & Short (8L) & 0 & \textbf{No refs; pure prose; not auditable.} & \textbf{A} \\
\shortstack[l]{\texttt{prop:tensorial-}\\\texttt{intertwining-m}} & Short (8L) & 0 & \textbf{Necessity argument only; effectively a hypothesis.} & \textbf{A} \\
\shortstack[l]{\texttt{lem:diagonal-}\\\texttt{traversal-}\\\texttt{scalar-action}} & Full (24L) & 3 & Well-founded; Jordan-block argument solid. & -- \\
\shortstack[l]{\texttt{thm:hc-spectral-}\\\texttt{projector}} & Full (30L) & 5 & Depends on \texttt{prop:tensorial-intertwining-m} hypothesis. & B \\
\shortstack[l]{\texttt{cor:M-HC-}\\\texttt{fixpoint}} & Short & 3 & Follows cleanly from theorem. & -- \\
\shortstack[l]{\texttt{cor:SM-spectral-}\\\texttt{realization}} & Short & 4 & Depends on \texttt{thm:SM}; sound. & -- \\
\midrule
\shortstack[l]{\texttt{thm:koide-}\\\texttt{triolic} (OP~3)} & Short (7L) & 0 & \textbf{$\theta_K$ formula not derived; $K=2/3$ by C--S only.} & \textbf{A} \\
\shortstack[l]{\texttt{thm:bsym-hopf}\\(OP~1)} & Short (11L) & 0 & \textbf{5.1\% gap to $1/\sqrt{3}$ unaddressed.} & \textbf{A} \\
\shortstack[l]{\texttt{thm:unity}\\(central)} & Short (9L) & 5 & \texttt{prop:compensator} ref resolved; label confirmed. & B \\
\shortstack[l]{\texttt{prop:structural-}\\\texttt{lift-principle}} & Short (8L) & 1 & Circular: uses its own assumptions as proof. & B \\
\midrule
\shortstack[l]{$\mathrm{ClosureStable}$\\$\Leftrightarrow\,\mathrm{im}(C)$} & \emph{missing} & -- & \textbf{Bijection not stated; gap in thm:hc-spectral-projector Step~3.} & \textbf{A} \\
\end{longtable}
}

\subsection*{Priority A: Required for structural soundness}

\paragraph{A1. \texttt{lem:closure-forcing-lemma} hardening.}
The proof currently reads as structural prose with no citations. The required explicit chain is:
\[
\begin{aligned}
\mathrm{Triolic}(x) &\Rightarrow \texttt{lem:triolic-minimality-lemma},\\
\mathrm{Anchored}(x) &\Rightarrow \texttt{lem:orthogonal-anchor-lemma},\\
\mathrm{HCForced}(x) &\Rightarrow \texttt{lem:hc-admissibility-lemma}.
\end{aligned}
\]
Combined consequence: no leakage outside the common container
$\Rightarrow \mathrm{ClosureStable}(x)$. Each implication must cite its
lemma explicitly.

\paragraph{A2. \texttt{prop:tensorial-intertwining-m} status correction.}
$\mathcal{M}\mathcal{C} = \mathcal{C}\mathcal{M}$ is currently stated as a
proposition but proved only by necessity argument. It must either be:
(a) promoted to an axiom/hypothesis with explicit hypothesis-grade marker, or
(b) derived from \texttt{def:tensor-structural-kernel} by explicit index
calculation showing that the idempotent tensor $\mathcal{C}^\mu{}_\nu$ commutes
with any structure-preserving $\mathcal{M}^\mu{}_\nu$ satisfying the
admissibility conditions. Option (b) is preferred.

\paragraph{A3. $\mathrm{ClosureStable}(x) \Leftrightarrow x \in \mathrm{im}(C)$ as
explicit lemma.}
Step~3 of \texttt{thm:hc-spectral-projector} uses \texttt{lem:closure-forcing-lemma} to identify $\mathrm{im}(C)$ with $\mathcal S_0$, but the reverse implication is still missing. The required gap may be summarized as:
\begin{itemize}[leftmargin=1.8em]
\item $\mathrm{im}(C) \subseteq \mathcal S_0$ is already covered,
\item closure-stable states must still be shown to be compensator-normalized,
\item hence a dedicated closure-stable/im$(C)$ lemma is still needed.
\end{itemize}

\paragraph{A4. \texttt{thm:koide-triolic} (OP~3) numerical hardening.}
The existing proof derives $K = 2/3$ via Cauchy-Schwarz but does not
derive the angle formula
$\theta_K = \pi - \sqrt{2}\cdot\arcsin(\beta_{\mathrm{sym}})$
established in prior session work. This formula predicts the Koide angle
to within 0.115\% and must be incorporated as an explicit step.
The residual mass-ratio discrepancy ($\sim 7$--$16\%$, amplification
factor $\sim 60$ near the electron mass) must also be stated explicitly.

\paragraph{A5. \texttt{thm:bsym-hopf} (OP~1) gap statement.}
The Hopf-angle derivation gives $\beta_{\mathrm{sym}} \approx 1/\sqrt{3}$
(continuum limit) while the measured value is $\approx 0.5493$, a 5.1\% gap
in $\beta_{\mathrm{sym}}$ corresponding to a $\sim 61\%$ gap in $\gamma_L$
due to nonlinear amplification. The theorem currently does not state this gap
explicitly. A remark recording the gap and its structural interpretation
(continuum $S^3$ vs discrete prime lattice) must be added.

\subsection*{Priority B: Structural hardening (next passes)}

\paragraph{B1. \texttt{prop:semantic-forcing-chain} proof.}
Currently stated without proof body. The implication
$\mathrm{GroundLinked} \wedge \mathrm{FrameShifted} \Rightarrow \mathrm{HCForced}$
needs an explicit derivation from the semantic kernel axioms (S1)--(S6).

\paragraph{B2. \texttt{prop:structural-lift-principle} circularity fix.}
The proof assumes its conclusion. A non-circular argument requires explicit
reference to the admissible transport operations (for example, \texttt{lem:group-transport} and \texttt{lem:hc-admissibility-lemma}) as the
mechanism by which kernel-invariant content is preserved.

\paragraph{B3. \texttt{thm:hc-spectral-projector} conditional status.}
The theorem is sound \emph{under} the hypothesis that
$\mathcal{M}\mathcal{C} = \mathcal{C}\mathcal{M}$ (A2). Until A2 is resolved,
the theorem should carry an explicit conditional marker. A remark to this
effect should be inserted immediately after the theorem statement.

\subsection*{Closure sequence for next subversions}

\begin{enumerate}[label=\textbf{v2.27.\arabic*:}, leftmargin=4em, start=6]
  \item Close A1 (\texttt{lem:closure-forcing-lemma}), A3
    ($\mathrm{ClosureStable} \Leftrightarrow \mathrm{im}(C)$ lemma),
    and B3 (conditional marker on \texttt{thm:hc-spectral-projector}).
  \item Close A2 (\texttt{prop:tensorial-intertwining-m} index calculation
    or hypothesis-promotion) and B2 (structural lift circularity fix).
  \item Close A4 (\texttt{thm:koide-triolic} $\theta_K$ formula, OP~3)
    and A5 (\texttt{thm:bsym-hopf} gap remark, OP~1).
  \item Close B1 (\texttt{prop:semantic-forcing-chain} proof).
    Evaluate gate readiness for next public release.
\end{enumerate}

\subsection*{What is already sound (do not reopen)}

The following blocks passed the v2.27.4 scan with no structural issues and
should not be reopened unless a genuine contradiction is found:
\begin{itemize}[leftmargin=1.8em]
\item \texttt{lem:diagonal-traversal-scalar-action} (Jordan-block exclusion),
\item \texttt{thm:fractal} (fractal self-traversal),
\item \texttt{thm:symmetric-frame} (algebraically complete),
\item \texttt{thm:unitary-equiv} (Stone--von Neumann, cites literature),
\item \texttt{cor:M-HC-fixpoint}, \texttt{cor:SM-spectral-realization},
\item all of \texttt{sec:readout-relative-lemmas},
\item \texttt{thm:hc-readout-geometry} (conditional on A2).
\end{itemize}

% ==================================================================

\section*{v1.2.0 Consolidation Core}
\addcontentsline{toc}{section}{v1.2.0 Consolidation Core}

\paragraph{Normalization note.} This block is retained as a \emph{historical consolidation layer}. In the normalized grouped architecture of Version~\docversion, its active generative content is absorbed into the normative layer $\mathfrak G$; the present section therefore records the earlier consolidation surface and should no longer be read as an independent foundational core competing with the later normalized structural kernel.


This version closes the gap between the previously inserted roadmap and the
already existing structural sections of the monolith. Its purpose is to state,
in one contiguous place, the working discipline of the next release, the core
closure chain, the probe-sector architecture, the first-layer master-equation
logic, and the normalized tool catalogue. The guiding rule remains unchanged:
\emph{no shortening, no silent merging, no replacement of the leading line;
only sharpening, explicit classification, and append-only completion.}

\subsection*{1. Structural reading discipline}
Every major claim used below is assigned one of four statuses.
\begin{enumerate}[leftmargin=2em,label=(\alph*)]
  \item \textbf{Structurally forced:} the statement follows from the ground
  condition, the triolic organization, orthogonality, and closure.
  \item \textbf{Projectively read:} the statement is not primary, but appears
  after passage to a chosen observational or effective sector.
  \item \textbf{Heuristically interpreted:} the structural content is stable,
  while its physical naming is still a disciplined reading rather than a final
  derivation.
  \item \textbf{Open:} the location of the claim is fixed, but its expanded
  proof or explicit computation is postponed to the next lemma layer.
\end{enumerate}

This fourfold discipline is essential because the framework is intended to be a
\emph{structure-first} proof architecture. In particular, equations are not to
be introduced as autonomous starting axioms when they can instead be read as
forced compression formulas of a deeper closure mechanism.

\subsection*{2. Ground condition, HC, and closure}
The working core of v1.2.0 is the explicit chain
\[
\text{Ground Condition} \Longrightarrow \text{HC} \Longrightarrow \text{Closure}.
\]
Here the ground condition is the requirement that the underlying triolic body is
not free to unfold into arbitrary independent sectors. Rather, the threefold
organization is constrained so that one component remains orthogonal to the
other two, the admissible units are treated in the standing imaginary
convention, and every physical readout is already phase-shifted with respect to
the underlying structural frame.

Under this condition, the Heisenberg compensator is not an external correction
patched in after the fact. It is the minimal balancing mechanism that becomes
necessary once a common body must be represented consistently in dual or
phase-shifted ways without allowing the structure to drift apart. In the logic
of the present paper, HC therefore carries three simultaneous tasks: it
re-identifies dual or shifted representatives, it suppresses unstable escape
modes, and it transports the system toward a light/collapse-compatible output.

Closure is then the global consequence of this balancing action. Once the
compensator has become active, admissible dynamics are no longer free to leave
the confined regime; instead they remain trapped within the higher coupled
architecture represented in the paper by the septinion/seption perspective.
This is the sense in which HC is emergent, while closure is enforced.

\paragraph{Proposition A (ground condition induces HC).}
Given the triolic structural core, the orthogonality constraint, and the
phase-shifted measurement principle, a compensating identification mechanism is
forced whenever the same body is to admit more than one consistent appearance.
That mechanism is the Heisenberg compensator.

\paragraph{Proof sketch.}
Without compensation, dual or shifted presentations of the same body would
produce uncontrolled representation drift. Because the framework demands one
common object across all sectors, a balancing map is required. This balancing
map is precisely the HC in structural form.

\paragraph{Proposition B (HC induces closure).}
Whenever HC acts as the balancing constraint of the triolic system, the
admissible dynamics are confined to a closure-stable regime. In particular,
sectors may transform, couple, or be read out differently, but they cannot
escape the common structural container without violating the ground condition.

\paragraph{Proof sketch.}
HC identifies and neutralizes precisely those deviations that would split the
common body into incompatible local descriptions. The remaining degrees of
freedom therefore assemble into a closed class. This class is the closure
regime read throughout the monolith.

\paragraph{Protected invariants.}
At the level of v1.2.0, the most important protected invariants are:
(i) the persistence of one common underlying body, (ii) orthogonality of the
third triolic direction relative to the first two, (iii) frame-shift discipline
between structure and observer description, and (iv) confinement of admissible
dynamics to the closure-compatible sector. These invariants are not yet the
full spectrum-level output, but they are the stable backbone needed for all
later projections.

\subsection*{2A. Consolidated dependency spine for v2.27.0}
For the next mathematical pass, the previous informal chain
\[
	ext{Ground Condition} \Longrightarrow 	ext{HC} \Longrightarrow 	ext{Closure} \Longrightarrow 	ext{Confinement}
\]
should be treated as a single dependency spine rather than as four loosely
adjacent slogans. The point of the present local pass is not to over-claim a
finished theorem package, but to normalize what the monolith already uses as
its operative order of dependence.

\paragraph{Lemma C (dependency spine lemma).}
Assume the triolic core, the standing orthogonality condition, and the
phase-shifted readout discipline. Then every admissible multi-representative
reading of one and the same body requires a compensating identification map;
this yields HC. Once HC is active, admissible state evolution is restricted to
those sectors that remain compatible with the common body; this yields closure.
Whenever closure is read at the level of the higher coupled container, the
resulting admissible class is a confinement class.

\paragraph{Proof sketch.}
The first implication is the content of Proposition~A: without compensation,
multiple shifted or dual descriptions of one body drift apart. The second is
Proposition~B: once HC acts as balancing constraint, the admissible dynamics can
no longer leave the common structural container. The last implication is a
change of reading level rather than a new mechanism: closure in the higher
coupled setting is exactly the confinement statement used downstream in the
septinion/seption language.

\paragraph{Operational reading.}
The practical value of this dependency spine is that later sections may appeal
to one normalized arrow instead of reintroducing the same chain piecemeal. In
particular, the later projector, transport, and Millennium-matrix sections
should be read as refinements of this spine, not as independent foundations.
This is also the correct mathematics-first reading for that archived public release:
first the spine is stabilized, then the surrounding operator language is
hardened.

\paragraph{Local closure target.}
For the purposes of the v2.26.x local sequence, this block counts as improved
only if all later uses of HC emergence, closure, and confinement can be traced
back to this normalized dependency spine without silent appeals to the
computational teaser architecture.

\subsection*{2B. Projector/closure anchoring of the dependency spine}
The next hardening step is to make explicit that the closure arrow in the
normalized spine is not merely verbal. It is read downstream through projector
language. In other words, the chain
\[
\text{Ground Condition} \Longrightarrow \text{HC} \Longrightarrow \text{Closure}
\]
should already be understood as carrying an operator-level consequence: once HC
has been introduced as the compensating identification map, admissible readings
must be projectively stabilized inside the common structural sector.

\paragraph{Lemma D (projector anchoring lemma).}
Assume the dependency spine of Lemma~C. Then every admissible downstream use of
closure may be rewritten as a projector-compatible statement: there exists a
sector projector $P_{\mathrm{adm}}$ such that the HC-balanced reading is
admissible only through the image of $P_{\mathrm{adm}}$, while noncompatible
continuations lie outside this image.

\paragraph{Proof sketch.}
If closure is the restriction to the common structural container, then a
projector description is simply the operator-level reading of that restriction.
The role of $P_{\mathrm{adm}}$ is not to invent a new mechanism, but to express
in normalized language that the HC-balanced state is retained only in the
closure-compatible sector. Thus projector logic is not parallel to closure
logic; it is the stabilized formal reading of closure once the argument moves
into the operator layer.

\paragraph{Downstream consequence.}
This anchoring is what allows the later projector/closure and transport sections
to be treated as genuine mathematical refinements of the dependency spine.
Ground Condition $\Rightarrow$ HC $\Rightarrow$ Closure no longer terminates in
a slogan, but feeds directly into the admissible projector sector on which the
later transport and Millennium-matrix mathematics are allowed to act.

\paragraph{Local closure target.}
For the v2.26.x sequence, this block counts as improved only if later projector
uses can be traced back to the dependency-spine-plus-projector reading, rather
than being invoked as an independent auxiliary formalism.

\subsection*{2C. Transport/Lorentz continuation of the dependency spine}
The next hardening step is to show that the normalized spine does not stop at
closure and projector admissibility. Once the admissible sector has been fixed,
the later transport layer must be read as a constrained continuation of the same
chain. In particular, Lorentz-transport language is not allowed to appear as an
independent dynamical decoration; it must preserve the HC-balanced,
projector-stabilized structural identity already enforced upstream.

\paragraph{Lemma E (transport continuation lemma).}
Assume the dependency spine of Lemma~C together with the projector anchoring of
Lemma~D. Then every admissible transport between downstream readout frames must
preserve the HC-balanced closure sector. Equivalently, Lorentz-type transport is
admissible only as a continuation on the image of the stabilized projector, not
as a map into an unconstrained external sector.

\paragraph{Proof sketch.}
By Lemma~D, admissibility is already restricted to the projector-compatible
closure image. A transport map acting on downstream representatives is therefore
allowed only if it respects that image. Otherwise the map would reintroduce a
break between closure logic and readout logic, contradicting the role of HC as
the compensating identification mechanism. Hence transport does not weaken the
spine; it inherits and preserves it.

\paragraph{Downstream consequence.}
This continuation is what permits the later Lorentz layer to be read in a
mathematics-first way: not as a second origin of structure, but as the
frame-translation discipline internal to the already stabilized closure sector.
It also prepares the next pass in which the Millennium matrix is treated as a
mathematical object acting only on admissibly transported states.

\paragraph{Local closure target.}
For the v2.26.x sequence, this block counts as improved only if later transport
claims can be referred back to the dependency spine through the combined chain
Ground Condition $\Rightarrow$ HC $\Rightarrow$ Closure $\Rightarrow$
Projector Stabilization $\Rightarrow$ Admissible Transport.


\subsection*{2D. Millennium-matrix realization of the dependency spine}
The next closure step is to prevent the Millennium matrix from floating as a
separate computational or architectural metaphor. If the mathematics-first
sequence is to hold, then the matrix must be read downstream from the already
stabilized chain of Ground Condition, HC emergence, closure, projector
stabilization, and admissible transport. In that sense the Millennium matrix is
not an independent origin of admissibility; it is the global matrix-level
realization of the closure-preserving compensator architecture already fixed by
the earlier lemmas.

\paragraph{Lemma F (Millennium-matrix realization lemma).}
Assume the dependency spine of Lemma~C, the projector anchoring of Lemma~D, and
the transport continuation of Lemma~E. Then the Millennium matrix may be
introduced as a mathematically admissible operator only on the HC-balanced,
closure-compatible, projector-stabilized, and transport-preserving state space.
Equivalently, its admissible role is not to generate a new unconstrained sector,
but to realize at matrix level the already normalized confinement-preserving
structural action.

\paragraph{Proof sketch.}
By Lemma~C the structural dependence runs from ground condition to HC and from
HC to closure. By Lemma~D this closure is stabilized as an admissible projector
image, and by Lemma~E every downstream transport must preserve that image.
Hence a matrix-level realization can be admitted only if it acts on that same
stabilized transported image. If it acted outside this chain, the Millennium
matrix would become an additional primitive rather than a realization of the
existing architecture, breaking the mathematics-first normalization. Therefore
its mathematically acceptable status is that of a global operator reading of the
already enforced confinement structure.

\paragraph{Proposition I (admissible matrix-sector invariance).}
Let $\mathcal M$ denote the Millennium matrix restricted to the HC-balanced,
closure-compatible, projector-stabilized, and transport-preserving sector
identified by Lemmas~C--F. Then $\mathcal M$ preserves that admissible sector;
in particular,
\[
\mathcal M\,P_{\mathrm{adm}} = P_{\mathrm{adm}}\,\mathcal M
\qquad	ext{on the stabilized transported image.}
\]
Hence the Millennium-matrix realization is mathematically acceptable only as a
sector-preserving global operator, not as a generator of a fresh unconstrained
layer.

\paragraph{Proof sketch.}
By Lemma~F the admissible introduction of $\mathcal M$ is already downstream
from closure, projector stabilization, and admissible transport. Proposition~H
states that the transport layer commutes with the admissibility projector on the
stabilized sector. Therefore the only mathematically consistent matrix-level
realization is one that acts within that same sector and preserves its image.
If $\mathcal M$ failed to preserve $P_{\mathrm{adm}}$, the matrix would reopen a
non-admissible branch and would no longer realize the normalized confinement
architecture. The commutation relation above is thus the matrix-level form of
admissible sector preservation.

\paragraph{Mathematical significance.}
This proposition upgrades the matrix discussion from a merely allowed operator
reading to a first invariance statement. The Millennium matrix is not only
introduced downstream from the dependency spine; it is constrained to preserve
the same admissible image that closure, projector stabilization, and transport
have already fixed. This gives the matrix chapter a clearer mathematical status
inside the normalization program.

\paragraph{Downstream consequence.}
This is the step that allows the later Millennium-matrix chapter to be retained
without overtaking the mathematics. The matrix can now be cited as a downstream
operator realization of the dependency spine, rather than as a free-standing
meta-object. It also prepares the next A-block pass, where trace and branch
mathematics can be read as acting under the matrix-level realization of the same
closure discipline.

\paragraph{Local closure target.}
For the v2.26.x sequence, this block counts as improved only if later uses of
the Millennium matrix can be referred back to the combined chain
Ground Condition $\Rightarrow$ HC $\Rightarrow$ Closure $\Rightarrow$
Projector Stabilization $\Rightarrow$ Admissible Transport $\Rightarrow$
Millennium-Matrix Realization.


\subsection*{2E. Trace/branch continuation of the dependency spine}
The final step in this first A-block pass is to prevent trace and branch
mathematics from re-opening the stabilized chain as if branching introduced a
new primitive source of admissibility. In the mathematics-first sequence, trace
formation and branch resolution must be read downstream from the already fixed
matrix-level realization. Thus branching is not a separate beginning of
structure; it is the controlled readout of differentiated admissible paths
inside the same confinement-preserving spine.

\paragraph{Lemma G (trace/branch continuation lemma).}
Assume the dependency spine of Lemma~C, the projector anchoring of Lemma~D, the
transport continuation of Lemma~E, and the Millennium-matrix realization of
Lemma~F. Then trace formation and branch resolution are mathematically
admissible only as downstream differentiations of the already HC-balanced,
closure-compatible, projector-stabilized, transport-preserving, and
matrix-realized state space. Equivalently, admissible trace/branch mathematics
does not create a fresh unconstrained dynamics; it resolves path structure
inside the same confinement-preserving architecture.

\paragraph{Proof sketch.}
By Lemma~F the Millennium matrix is admitted only as an operator realization of
the already stabilized dependency spine. Therefore any trace extracted from the
resulting evolution must be a trace of that same admissible operator action,
not a record of an independent unnormalized sector. Likewise, branch resolution
can be accepted only if it separates compatible downstream continuations of the
same confined structure. If branching were allowed to act as a new primitive,
then the chain from ground condition through HC, closure, projector
stabilization, transport, and matrix realization would no longer control the
later readout layer. Hence trace and branch mathematics remain admissible only
as terminal differentiations internal to the same spine.

\paragraph{Downstream consequence.}
This closes the first mathematics-first pass across the initial A-block group.
The trace/branch layer can now be cited without drifting into an autonomous
search or computational metaphor: it is anchored in the same structural chain
that begins at the ground condition and ends in admissible differentiated
readout. This is the point at which the first local freeze judgement for the
A-block group becomes meaningful.

\paragraph{Local closure target.}
For the v2.26.x sequence, this block counts as improved only if later trace and
branch statements can be referred back to the combined chain
Ground Condition $\Rightarrow$ HC $\Rightarrow$ Closure $\Rightarrow$
Projector Stabilization $\Rightarrow$ Admissible Transport $\Rightarrow$
Millennium-Matrix Realization $\Rightarrow$ Trace/Branch Continuation.


\paragraph{Proposition J (admissible trace/branch persistence).}
Let $\Gamma_{\mathrm{adm}}$ denote an admissible trace family generated inside the stabilized admissible sector, and let $B_{\mathrm{adm}}$ denote a branch selector that only separates continuations already induced by the same admissible matrix-sector action. Then admissible trace continuation and admissible branch separation preserve the same normalized admissible sector; in particular they do not generate a new independent admissibility source outside the projector/transport/matrix spine.
\[
B_{\mathrm{adm}}\,\Gamma_{\mathrm{adm}}\subseteq P_{\mathrm{adm}}\mathcal H_{\mathrm{conf}}.
\]

\paragraph{Proof sketch.}
By Proposition~H, admissible transport acts inside the stabilized projector image, and by Proposition~I the Millennium-matrix realization preserves that same admissible sector. Hence every admissible trace produced downstream is already a differentiated record of motion internal to the projector-stabilized, transport-compatible, matrix-preserving sector. If a branch operation is admissible at all, it cannot introduce a fresh primitive sector; it can only separate compatible continuations already contained in that same image. Therefore both trace continuation and admissible branching persist inside the normalized admissible sector.

\paragraph{Mathematical significance.}
This upgrades the final step of the first A-block from a narrative endpoint to a persistence statement: the trace/branch layer is no longer merely attached downstream, but mathematically constrained not to reopen an external admissibility source. In that sense the first A-block now carries a full preservation chain from the ground condition through HC, closure, projector stabilization, transport, matrix realization, and finally trace/branch continuation.

\paragraph{Local closure target.}
For the v2.26.x sequence, this block counts as strengthened only if later trace/branch arguments can explicitly cite Proposition~J as the persistence rule preventing downstream drift outside the admissible spine.

\subsection*{2F. Local freeze check for the projector/transport core}
The first A-block is now strong enough that its middle core --- projector,
closure, and transport --- can be evaluated on its own rather than only as a
passing stage inside the longer dependency spine. This local freeze check asks
whether that middle layer has reached a condition in which later work may treat
it as stable for the next release cycle without pretending that every possible
auxiliary lemma has already been unfolded to maximal density.

\paragraph{Local freeze status.}
At the present stage, the projector/closure plus transport/Lorentz core is best
read as \emph{locally freeze-ready for v2.27.0}, provided that no later pass
re-opens the admissible image as an unconstrained external sector. The closure
logic now fixes the admissible image; the projector logic stabilizes that
image; and the transport continuation restricts downstream frame changes to
that already fixed sector. In that sense, the central middle layer of the first
A-block is no longer merely planned architecture: it has acquired a coherent
mathematical dependency order.

\paragraph{Residual blocker.}
What still limits the block is not the absence of an admissibility spine, but
rather the remaining density gap between the local lemmas and a fully expanded
proof presentation. In other words, the structure is already in place, but the
companion could still benefit from additional explicit micro-lemmas later if a
future release wants a denser audit trail.

\paragraph{Release consequence.}
For the purposes of the mathematics-first roadmap, the projector/transport core
may therefore count as \emph{closed for now at medium proof density}. This is
enough to support the next local passes and to contribute positively to the
v2.27.0 gate, but not yet enough to justify claiming that every supporting
auxiliary statement has already reached final archival density.

\paragraph{Local closure target.}
The block counts as satisfactorily frozen for the next public release if later
passes preserve the chain
Ground Condition $\Rightarrow$ HC $\Rightarrow$ Closure $\Rightarrow$
Projector Stabilization $\Rightarrow$ Admissible Transport
without reintroducing a competing admissibility source outside that normalized
middle layer.



\subsection*{2G. Consolidated freeze and gate pass for the first A-block}
The first A-block has now been hardened not only by the dependency-spine lemmas,
but also by the three preservation propositions on transport, matrix action, and
trace/branch continuation. It is therefore possible to perform a consolidated
freeze and gate pass that asks whether the block, taken as a whole, has reached
a level suitable for counting positively toward the v2.27.0 release threshold.

\paragraph{Consolidated freeze status.}
Taken together, Lemma~C through Lemma~G and Proposition~H through Proposition~J
justify reading the first A-block as \emph{structurally freeze-ready for
v2.27.0 at medium proof density}. The block now carries a continuous admissible
spine from ground condition through HC, closure, projector stabilization,
transport, Millennium-matrix realization, and admissible trace/branch
persistence. No later local pass should need to reopen the block at the level
of its basic dependency order unless a genuine contradiction is found.

\paragraph{Gate consequence.}
For the global v2.27.0 gate, this means that the first A-block can now count as
a \emph{positive readiness contributor}. It does not by itself open the public
release, because the document-wide gate still depends on the remaining A-blocks,
reference stability, and hypothesis-control discipline. But it does remove the
first major blocker from the mathematics-first release logic.

\paragraph{Residual blocker.}
The remaining limitation is not structural incompleteness, but proof density and
audit granularity. Later archival passes may still unfold more micro-lemmas,
but this is no longer the same as saying that the block is mathematically open.
The correct current reading is: \emph{closed in structure, not maximally dense in
presentation}.

\paragraph{Local closure target.}
For the v2.26.x sequence, this pass counts as successful only if later global
readiness summaries can cite the first A-block as freeze-ready without again
collapsing its seven-step admissible spine into vague architectural language.

\subsection*{3. Higher embedding and confinement logic}
The local triolic organization is not the terminal description. Its load-bearing
role is to provide the minimal unit from which higher coupling becomes possible.
The octonionic step is useful because it already carries nontrivial coupled
imaginary directions, but within the present project logic it is not yet the
final confinement container. The septinion/seption reading is introduced
precisely to express that the entire compensated dynamics lives inside a more
rigid higher assembly in which uncontrolled sector break-out is blocked.

This point should be read carefully: the higher embedding is not decorative
algebraic escalation. It is the formal language by which the paper states that
closure is a global architectural fact rather than a local numerical accident.
In this sense, the passage
\[
\text{local triolic unit} \to \text{compensated coupling} \to \text{higher confinement}
\]
is one of the main structural arrows of the version.

\subsection*{4. First-layer master-equation logic}
The monolith already contains the structural core section and the ToE-style
master-equation section. v1.2.0 now fixes how these are to be read.
The first-layer master equation is not yet the maximal explicit final law of
the framework. It is the minimal compressed form that must simultaneously carry
geometry, triolic sector dynamics, compensator action, closure residue, and
observer-frame phase shift.

Accordingly, the five structural blocks used later in the paper should be
classified as follows.
\begin{enumerate}[leftmargin=2em]
  \item The geometric block is \emph{projectively privileged}: it is forced as a
  large-scale closure reading, but not as the primitive starting point.
  \item The triolic block is \emph{structurally forced}: without it there is no
  load-bearing internal organization.
  \item The HC block is \emph{structurally forced}: without it, no stable
  identification of dual/shifted representatives exists.
  \item The closure block is \emph{structurally forced}: without it, the theory
  would have no admissibility criterion.
  \item The observer-frame block is \emph{projectively forced}: once the
  measurement frame is shifted, a readout term is unavoidable.
\end{enumerate}

The upshot is that the master equation should always be read from left to right
as a compression chain,
\[
\text{structure} \to \text{compensation} \to \text{closure} \to \text{projection}.
\]
Any later physical naming --- Einstein-like, Dirac-like, or QFT-like --- comes
only after this structural order has been respected.

\paragraph{Proposition C (first-layer master form is forced).}
Once triolic organization, HC, closure, and frame shift are all imposed, the
dynamics can no longer be represented by a purely geometric or purely sectorial
equation. A combined master form is forced at first layer.

\paragraph{Proof sketch.}
Each omitted block destroys one already fixed architectural demand: omitting
triolicity erases the internal organization; omitting HC erases compensation;
omitting closure erases admissibility; omitting the observer term erases the
readout discipline. Hence only a combined structural form remains compatible.

\subsection*{5. Probe sector as second access path}
The probe sector is not a side remark but the second major access path of
v1.2.0. Two probes are sent toward one another, and each probe itself is
to be treated as a triolic structure. The third group stands orthogonal to both
probes and thereby defines a constrained search region rather than an arbitrary
ambient space. This orthogonal third group is crucial because it prevents the
probe architecture from degenerating into a mere two-body scan.

The role of the probe picture is twofold. First, it gives a structural way to
read interaction and search without abandoning the closure discipline. Second,
it provides an explicit place where the framework's search tools can be
attached. In the present version, the probe geometry is therefore a controlled
laboratory for structural read-off.

\paragraph{Proposition D (probe sector inherits closure).}
If each probe carries triolic organization and the third group remains
orthogonal to both probes, then the same HC/closure discipline that governs the
main sector also governs the admissible probe dynamics.

\paragraph{Proof sketch.}
The probe architecture does not introduce a new ontology. It duplicates the
same structural ingredients in a constrained interaction geometry. Since the
closure mechanism depends on these ingredients rather than on the label
``background'' or ``probe'', the probe sector inherits the same admissibility
logic.

\paragraph{Search space.}
The search space is therefore not an unrestricted continuum of candidates. It
is the region left after imposing triolic self-organization of each probe,
orthogonality of the third group, compensator compatibility, and closure.
Everything outside that region is structurally irrelevant for the probe readout.

\subsection*{6. Normalized tool apparatus}
A central goal of v1.2.0 is to stop treating rechen-tricks as scattered
intuition and instead store them as a reusable apparatus. The following classes
are fixed for this version.

\begin{description}[leftmargin=2.2em, style=nextline]
  \item[Closure-tools.] Tools that test, restore, or compare closure-compatible
  configurations. Their purpose is to decide admissibility and leakage.
  \item[Frame-shift-tools.] Tools that translate between the underlying
  structure and the observer's projected readout.
  \item[Projection-tools.] Tools that extract effective sector equations or
  recognizable physical forms from the common structural body.
  \item[Probe-tools.] Tools that organize the constrained search in the probe
  sector.
  \item[Equivalence-tools.] Tools that identify when two differently written
  states belong to the same structural class.
  \item[Read-off / trace-tools.] Tools that recover results from the structure
  without expanding the entire calculation tree every time.
\end{description}

Every individual tool should be documented under the same six labels:
\emph{name, purpose, prerequisite, short formal scheme, typical use, boundary
of validity}. This normalization rule is itself one of the deliverables of
v1.2.0.

\paragraph{Canonical examples in the present project.}
Kasisky detection is stored as a frequency/pattern detector rather than as a
cryptographic side story. Hash and rainbow methods are stored as routing and
candidate-compression devices. Markov organization is stored as a transition
controller over admissible search paths. Newton is stored as a convergence
accelerator once a structural branch has already been selected. GARCH, where
used at all, remains optional and secondary.

\subsection*{7. Projection discipline and physics reading}
The structural reading of matter, gauge, mass, and Higgs later in the monolith
is to be understood under a strict hierarchy.
\begin{enumerate}[leftmargin=2em]
  \item Matter as closure-stable bound sector is a structural reading.
  \item Gauge symmetry as invariance of closure-compatible transport is a
  projective reading.
  \item Mass as closure/frame residue is a projective reading with strong
  structural motivation.
  \item Higgs as coupling instance or mean coupling point is currently a
  hypothesis block and must remain explicitly marked as such.
\end{enumerate}

This hierarchy matters because it prevents later criticism from collapsing the
whole framework into one category. The paper does not claim that every named
physical object is already fully derived at the same level. It claims instead
that the structural origin is common and that different physical sectors appear
at different distances from that origin.

\subsection*{8. Consistency status of v1.2.0}
The present version should therefore be regarded as \emph{finished in working
form} whenever the following points are simultaneously true:
\begin{enumerate}[leftmargin=2em]
  \item the DOI, version stamp, and roadmap are fixed;
  \item the chain Ground Condition $\Rightarrow$ HC $\Rightarrow$ Closure is
  explicitly stated;
  \item the probe sector is presented as a genuine structural sector;
  \item the master equation is classified as a first-layer forced compression;
  \item the tool catalogue is normalized;
  \item open points are named rather than hidden.
\end{enumerate}

These conditions are now met in the present monolith together with the later
sections \emph{Structural Core and Closure Principles}, \emph{Toward a
Structural ToE Master Equation}, and \emph{Structural Reading of Matter,
Gauge, Mass, and Higgs}. What remains for v1.2.0 is not the recovery of a
missing architecture, but the next proof-theoretic deepening: more explicit
lemma chains, stronger sector-by-sector projections, and a larger appendix of
normalized tools and shortcuts.

\section*{v1.4.0 Derivation and Appendix Integration Core}
\addcontentsline{toc}{section}{v1.4.0 Derivation and Appendix Integration Core}

\paragraph{Normalization note.} In Version~\docversion, this block is read as a \emph{mixed historical derivation layer}: its closure-ladder and micro-to-macro passages belong genealogically to $\mathfrak G$ and $\mathfrak P$, while its probe-reduction and appendix-coupling material remains operational support for the normalized architecture. It is therefore preserved for proof-memory and appendix linkage, but not as a second normative origin.


The specific task of v1.4.0 is to turn the earlier consolidation and
lemma-development layers into a more tightly chained proof surface. This means
that the argument is no longer presented only as a list of strategic claims.
Instead, the intermediate dependencies are written so that later appendix items
can be cited as reusable proof modules. The rule of this stage remains simple:
\emph{no shortening, no creative overwrite, only denser local chaining and
clearer main-text/appendix interlock.}

\subsection*{1. Closure ladder as a reusable derivation chain}
The closure architecture is now fixed as a ladder rather than a single jump.
The intended reading is
\[
\begin{aligned}
&\text{local triolic unit} \Longrightarrow \text{orthogonal stabilization}\\
&\Longrightarrow \text{compensated identification} \Longrightarrow \text{closure class}\\
&\Longrightarrow \text{effective projection}.
\end{aligned}
\]
The advantage of this ladder is methodological. It prevents the later sections
from collapsing several distinct proof moves into one rhetorical transition.
Each arrow can now be isolated, tested, and, where needed, exported into the
appendix as a technical lemma or shortcut.

\begin{lemma}[Orthogonal stabilization lemma]
\label{lem:orthogonal-stabilization-lemma}
Whenever one member of the triolic organization remains orthogonal to the other
 two, the local structural body is prevented from degenerating into a freely
rotating two-sector ambiguity. The third direction therefore acts as a
stabilizing anchor for all later compensator and closure statements.
\end{lemma}

\begin{proof}
Without the orthogonal anchor, the first two directions would admit a family of
relabelings that preserve local appearance while destroying global structural
identity. The orthogonal third component blocks this ambiguity by fixing a
non-coincident direction that all admissible couplings must respect.
\end{proof}

\begin{lemma}[Compensated identification lemma]
\label{lem:compensated-identification-lemma}
If the same structural body is observed through shifted or dual frames, then a
compensating map is forced in order to preserve common identity across those
frames. In the present framework this map is read as the Heisenberg
compensator.
\end{lemma}

\begin{proof}
Frame-shifted representatives are not independent ontological copies. They are
multiple presentations of one common body. A balancing operation is therefore
necessary to identify them as one object rather than a divergent family. That
operation is exactly what HC contributes at the structural level.
\end{proof}

\begin{lemma}[Closure-class lemma]
\label{lem:closure-class-lemma}
Once orthogonal stabilization and compensated identification are both active,
the admissible states assemble into a closure class: transition is allowed,
but leakage into an incompatible external regime is not.
\end{lemma}

\begin{proof}
The orthogonal anchor limits admissible couplings, while HC neutralizes those
frame-induced mismatches that would otherwise split the common body. The
surviving states therefore form a self-consistent class under the allowed
moves of the theory.
\end{proof}

\subsection*{2. Micro-to-macro chain tightened}
The earlier versions already stated that the macroscopic projection must be read
from the microscopic structural unit through HC and closure. In v1.4.0,
this chain is stated more sharply:
\[
\begin{aligned}
&\text{micro structural unit} \Longrightarrow \text{local compensated coupling}\\
&\Longrightarrow \text{sector stability} \Longrightarrow \text{macroscopic effective law}.
\end{aligned}
\]
No macroscopic term is treated as primitive if it can be located lower in this
chain. This applies in particular to geometric large-scale terms, matter/gauge
labels, and all observer-frame residues.

\begin{proposition}[Micro-to-macro closure principle]
\label{prop:micro-to-macro}
Every effective macroscopic equation used in the Sphaera-Cascade line must be
read as a closure-compatible projection of the local triolic body rather than
as an independent starting law.
\end{proposition}

\begin{proof}
The ground condition, orthogonality requirement, and HC mechanism already fix
what counts as an admissible local evolution. A macroscopic law that failed to
arise as a projection of that admissible class would describe a sector not
licensed by the common body. Hence effective macroscopic forms are subordinate
projections of the closure class.
\end{proof}

\paragraph{Interpretive consequence.}
This principle strengthens the earlier interpretive discipline. Einstein-like,
Dirac-like, and QFT-like forms remain useful projective readings, but their
status is now tethered to a single chain of legitimacy: they are admissible
only insofar as they compress the closure-compatible microstructure.

\subsection*{3. Probe-reduction and guided search principle}
The probe sector is now equipped with a more explicit reduction principle. Two
probes approaching one another do not explore an arbitrary ambient continuum;
they reduce the candidate space by repeated closure tests.

\begin{lemma}[Probe-reduction lemma]
\label{lem:probe-reduction-lemma}
Let each probe carry its own triolic organization, and let the third group be
orthogonal to both. Then the admissible interaction region is the subset of the
naive search space that survives orthogonality, compensator compatibility, and
closure tests simultaneously.
\end{lemma}

\begin{proof}
Each probe contributes an internal triolic admissibility condition. The shared
orthogonal third group cuts away probe pairings that cannot be embedded in one
common constrained geometry. HC removes residual representational drift, and
closure excludes the remaining leakage channels. What survives is exactly the
reduced search region relevant to structural read-off.
\end{proof}

\paragraph{Tool interlock.}
This is the point where the normalized tools become operational rather than
merely classificatory. Kasisky is attached to frequency-pattern selection.
Hash/rainbow devices compress routing across candidate branches. Markov control
organizes transitions inside the still-admissible region, and Newton
accelerates convergence once a branch has already been selected. The appendix
can therefore store these devices as modular search operators without severing
them from the main proof line.

\subsection*{4. Appendix coupling rules}
The appendix is no longer a passive storage area. In v1.4.0 it is given a
precise function: each appendix block must either (i) justify a compressed step
in the main text, (ii) codify a reusable tool, or (iii) collect a repeated
technical transformation that would otherwise bloat the main line.

The intended appendix map is therefore:
\begin{enumerate}[leftmargin=2em]
  \item \textbf{Symbol and sector tables} collect notation and prevent local
  restatement overhead.
  \item \textbf{Structural tool records} store normalized rechen-tricks in the
  six-field format already fixed earlier.
  \item \textbf{Technical lemmas} receive derivations whose content is
  load-bearing but too local or repetitive for the main line.
  \item \textbf{Projection tables and read-off schemes} record the disciplined
  passage from structural objects to effective sector forms.
\end{enumerate}

\begin{proposition}[Appendix transfer principle]
\label{prop:appendix-transfer}
A step may be moved from the main text to the appendix only if its logical role
is preserved through an explicit citation path. Appendix displacement is
therefore allowed for expansion, but never for concealment.
\end{proposition}

\begin{proof}
The project rule forbids silent loss of proof substance. Therefore, every step
outsourced to the appendix must remain recoverable as a named lemma, tool, or
transformation block. Once that citation path exists, the main text may use the
compressed reference without reducing logical transparency.
\end{proof}

\subsection*{5. v1.4.0 readiness criteria}
This release should be regarded as complete in working form when the following
conditions hold simultaneously:
\begin{enumerate}[leftmargin=2em]
  \item the version stamp and DOI are updated to v1.4.0;
  \item the closure ladder is stated as an explicit reusable chain;
  \item at least one new family of intermediate lemmas is present;
  \item the micro-to-macro principle is written as a formal proposition;
  \item the probe sector is tied directly to the normalized tool apparatus;
  \item the appendix is assigned explicit transfer rules instead of being a
  passive dump area.
\end{enumerate}

These conditions are now satisfied in the present monolith. Relative to v1.3.0,
v1.4.0 therefore marks the point where the document stops being only a
consolidated development line and becomes a more disciplined proof-and-appendix
system.


\section*{v1.5.0 Projection and Appendix Table Integration Core}
\addcontentsline{toc}{section}{v1.5.0 Projection and Appendix Table Integration Core}

\paragraph{Normalization note.} In Version~\docversion, this block is retained as a \emph{$\mathfrak P$-dominant transport layer}. It introduces no independent generative axioms; rather, it records how closure-licensed structural objects are transported into projection tables, appendix targets, and effective readout interfaces.


The task of v1.5.0 is to convert the earlier derivation-and-appendix
interlock into a more explicit transport layer between the main line and the
technical tables. Earlier versions had already introduced the closure ladder,
the micro-to-macro principle, and the appendix transfer rule. The present stage
adds a stricter rule: every repeated projection, reduction, or read-off step
should be capturable either as a named lemma or as a named table interface with
a declared upstream structural source and a declared downstream effective use.
The point is not stylistic expansion but recoverability. A compressed move must
remain reconstructible.

\subsection*{1. Projection tables as proof interfaces}
Projection tables are not merely descriptive summaries. In v1.5.0 they
receive a precise logical role: a projection table records how a structural
object, already licensed by the closure architecture, may be read as an
effective sector form without pretending that the effective form is primary.
The direction of legitimacy is therefore always
\[
\text{structural source} \Longrightarrow \text{closure-compatible projection}
\Longrightarrow \text{effective sector reading.}
\]
This prevents the later document from treating a Dirac-like, Einstein-like, or
QFT-like formula as an autonomous axiom. The table is only valid if each row
preserves that direction.

\begin{definition}[Projection-table interface]
A projection-table interface is a named record that assigns to a structural
source object its admissible effective readings together with the conditions
under which the passage is allowed, restricted, or still open.
\end{definition}

\begin{lemma}[Projection-interface lemma]
\label{lem:projection-interface-lemma}
Whenever a projection step occurs repeatedly in the main line, the step may be
compressed into a projection-table interface provided that the source object,
closure condition, and interpretation status are all explicitly recorded.
\end{lemma}

\begin{proof}
A repeated projection is proof-relevant only because it preserves a controlled
relation between source and effective reading. If that relation is named once in
a table together with its conditions, the main line may cite the named interface
without losing substance. Compression is legitimate precisely because the table
stores the missing local detail in recoverable form.
\end{proof}

\subsection*{2. Interpretation status discipline}
The earlier reading discipline distinguished between what is structurally
forced, projectively read, interpretively overlaid, and reserved for later proof expansion. In
v1.5.0 this becomes a transport rule for all summary tables and appendix
catalogues.

\begin{proposition}[Interpretation-status transport principle]
\label{prop:interpretation-status-transport}
Every entry transferred from the main line into a table, catalogue, or appendix
record must preserve its interpretation status. A structurally forced item may
not be downgraded into merely interpretive language, and an interpretive or open reading may not
be silently upgraded into a forced statement through tabulation.
\end{proposition}

\begin{proof}
Tables compress syntax, not logical rank. If the interpretation status were
allowed to drift during compression, then the appendix would alter rather than
preserve the proof surface. The transfer must therefore carry both content and
status labels together.
\end{proof}

\paragraph{Practical consequence.}
Projection tables must visibly mark at least three columns: structural source,
allowed effective reading, and status label. Optional columns may record frame
conditions, probe restrictions, or unresolved hypotheses, but the status label
is mandatory.

\subsection*{3. Appendix-citation lattice}
The appendix now receives a more explicit citation lattice. Instead of treating
the appendix as a linear storage area, v1.5.0 organizes it as a small
network of recoverable proof targets:
\begin{enumerate}[leftmargin=2em]
  \item \textbf{lemma targets} for local derivation moves;
  \item \textbf{tool targets} for normalized rechen-tricks;
  \item \textbf{table targets} for repeated projection or reduction patterns;
  \item \textbf{schematic targets} for probe geometry and search-space layouts.
\end{enumerate}
A compressed sentence in the main text is acceptable only when it points to one
of these four target types.

\begin{lemma}[Appendix-citation lattice lemma]
\label{lem:appendix-citation-lattice-lemma}
A compressed main-text step is proof-disciplined if and only if it can be
expanded along at least one explicit appendix citation path ending in a lemma,
tool, table, or schematic target.
\end{lemma}

\begin{proof}
The project rule forbids concealment by relocation. Therefore every compressed
step must have a recoverable destination. The four target types above exhaust
the relevant forms of local expansion used in the present monolith. If no such
path exists, the compression is not yet justified.
\end{proof}

\subsection*{4. Projection robustness across frames and probes}
A repeated source of ambiguity is the temptation to read the same effective term
as though it had identical status in all frames and in the probe sector. The
present release therefore states a robustness rule.

\begin{proposition}[Projection-robustness principle]
\label{prop:projection-robustness}
An effective reading is robust only if it survives both frame-shift transport and
probe-sector restriction without leaving the closure-compatible class.
\end{proposition}

\begin{proof}
A term read in one frame but destroyed by admissible frame transport cannot be a
stable structural projection. Likewise, a reading that collapses once the probe
sector imposes orthogonality and closure tests is not globally licensed. Robust
readings are exactly those that persist under both operations.
\end{proof}

\paragraph{Use in later versions.}
This principle is meant to control later expansions of matter, gauge, mass, and
Higgs readings. The question is no longer only whether a reading is suggestive,
but whether it survives transport across the disciplined structural operations
already fixed by the framework.

\subsection*{5. v1.5.0 readiness criteria}
This release should be regarded as complete in working form when the following
conditions hold simultaneously:
\begin{enumerate}[leftmargin=2em]
  \item the version stamp and DOI are updated to v1.5.0;
  \item projection tables are assigned an explicit proof-interface role;
  \item interpretation status is required to survive table and appendix transfer;
  \item the appendix-citation lattice is stated as a reusable compression rule;
  \item projection robustness is tied to both frame transport and probe tests;
  \item the new material is appended without shortening the inherited monolith.
\end{enumerate}

These conditions are now satisfied in the present monolith. Relative to v1.4.0,
v1.5.0 marks the point where appendix tables and projection summaries are
integrated into the proof discipline itself rather than serving only as support
material.

\clearpage
\clearpage


\clearpage
\section*{Computational completeness layer: QAM / MML / MMA}

The present block records the first explicit machine reading of the post-v3.31.0 line.  
In the current v3.34.0-r09 working candidate it should be read together with the local HC/OP interface registry of Definition~\texttt{def:r07-local-language-version-registry} in the active-core PDF, the compression map of Definition~\texttt{def:compression-map-from-inherited-text-to-mml-mma-operations} in the active-core PDF, and the reusable tool normal forms of Definition~\texttt{def:first-reusable-hcop-tool-normal-form} in the active-core PDF. That registry gives the active sublanguage versions actually used by the present OP-entry chain, the compression map gives the active audit reading for inherited HC/OP text, and the r09 tool layer gives the controlled invocation form:
\[
  \mathrm{QAM}_{\mathrm{HCOP}}^{\mathrm{loc}}
  \to
  \mathrm{MML}_{\mathrm{HCOP}}^{\mathrm{loc}}
  \to
  \mathrm{MMA}_{\mathrm{HCOP}}^{\mathrm{loc}}.
\]
The broader computational chapter remains a reserve architecture; the r07 interface is the narrow executable slice used for HC-routed OP-entry.

Its purpose is to separate three layers which were already implicit in the project architecture:
\begin{enumerate}[leftmargin=2em]
  \item QAM as the finite tagged object language;
  \item MML as the finite machine / instruction language;
  \item MMA as the matrix-realized execution layer.
\end{enumerate}

The block proceeds in three steps. First, it isolates finite machine completeness and finite universality (the Z3-style regime). Second, it exhibits the unbounded extension mechanism through extensible address towers and fresh-address allocation. Third, it formulates stage-stable universality and the resulting architectural Turing-style completeness reading.

This computational layer is to be read \emph{subordinate} to the earlier strength hierarchy
\[
\mathbf S_1\to \mathbf S_2\to \mathbf S_3(\mathcal R)\to \mathbf S_4(\mathcal R)\to \mathbf S_5
\]
from Sections~19U--19Z and the subsequent machine-status decoder. Accordingly, the present section does not by itself upgrade the unconditional public status of the line beyond \(\mathbf S_1\). Rather, it records an \emph{internal machine-reading package}: finite machine completeness and finite universality sharpen the internal route toward \(\mathbf S_2\), while explicit fresh-address extension and stage-stable universality formulate the internal architectural ingredients that would be needed later on any conditional route toward \(\mathbf S_4(\mathcal R)\) and \(\mathbf S_5\).

\begin{definition}[Finite machine store]\label{def:finite-machine-store-computational-layer}
A \emph{finite machine store} is a finite tagged family
\[
\mathfrak{S}_{\mathrm{fin}}
=
\{(a_i,x_i)\}_{i=1}^{N}
\]
where each \(a_i\) is an address tag and each \(x_i\) is a stored QAM-object.
\end{definition}

\begin{definition}[Finite machine state]\label{def:finite-machine-state-computational-layer}
A \emph{finite machine state} is a tuple
\[
\Sigma
=
\bigl(
\mathfrak{S}_{\mathrm{fin}},
p,
q,
r
\bigr),
\]
where:
\begin{enumerate}[leftmargin=2em]
    \item \(\mathfrak{S}_{\mathrm{fin}}\) is a finite machine store;
    \item \(p\) is the current program position;
    \item \(q\) is the current control state;
    \item \(r\) is the current route-faithful residual/control register.
\end{enumerate}
\end{definition}

\begin{definition}[QAM-object layer]\label{def:qam-object-layer-computational-layer}
The \emph{QAM-object layer} is the finite tagged object language in which the machine store contents \(x_i\) are represented.

A QAM-object may be:
\begin{enumerate}[leftmargin=2em]
    \item a scalar-like tagged object;
    \item a packet-like tagged object;
    \item a tuple/block object;
    \item a routed or normalized residual object;
    \item any finite object built from the current QAM grammar by the allowed formation rules.
\end{enumerate}
\end{definition}

\begin{remark}[QAM as object language]\label{rem:qam-as-object-language-computational-layer}
QAM is not yet used here as a new infinite set theory. At this stage it is used as the finite object language in which machine-addressable project data are encoded.
\end{remark}

\begin{definition}[MML instruction alphabet]\label{def:mml-instruction-alphabet-computational-layer}
The \emph{MML instruction alphabet} is a finite set
\[
\mathcal{I}_{\mathrm{MML}}
\]
of primitive machine instructions, divided into the following classes:
\begin{enumerate}[leftmargin=2em]
    \item \textbf{read/write instructions:}
    \[
    \mathrm{READ}(a),\qquad \mathrm{WRITE}(a,\xi);
    \]
    \item \textbf{routing/transfer instructions:}
    \[
    \mathrm{ROUTE}(\xi),\qquad \mathrm{PROJECT}(\xi),\qquad \mathrm{NORMALIZE}(\xi);
    \]
    \item \textbf{comparison/test instructions:}
    \[
    \mathrm{TEST}_{=},\qquad \mathrm{TEST}_{<},\qquad \mathrm{TEST}_{\mathrm{adm}};
    \]
    \item \textbf{branching/control instructions:}
    \[
    \mathrm{SELECT}(b_1,\dots,b_k),\qquad \mathrm{BRANCH}(j),\qquad \mathrm{RETURN};
    \]
    \item \textbf{cycle/termination instructions:}
    \[
    \mathrm{STEP},\qquad \mathrm{LOOP},\qquad \mathrm{HALT}.
    \]
\end{enumerate}
\end{definition}

\begin{remark}[MML as finite machine language]\label{rem:mml-as-finite-machine-language-computational-layer}
MML is the first finite machine language layer of the project. Its purpose is not yet unrestricted computation, but finite programmable symbolic control on route-faithful QAM-data.
\end{remark}

\begin{definition}[Finite MML program]\label{def:finite-mml-program-computational-layer}
A \emph{finite MML program} is a finite word
\[
P=(I_1,\dots,I_m),\qquad I_j\in\mathcal{I}_{\mathrm{MML}}.
\]
\end{definition}

\begin{definition}[MMA execution layer]\label{def:mma-execution-layer-computational-layer}
The \emph{MMA execution layer} is the matrix-realized execution semantics of MML on finite machine states.

Concretely, an MMA-step is a map
\[
\mathcal{E}_{\mathrm{MMA}}:(\Sigma,I)\longmapsto \Sigma^{+},
\]
sending a finite machine state \(\Sigma\) and one MML instruction \(I\) to the updated state \(\Sigma^{+}\), in a way compatible with the current Millennium-matrix / Heisenberg-compensator realization of the route-faithful control structure.
\end{definition}

\begin{remark}[Why MMA is the right execution layer]\label{rem:why-mma-is-the-right-execution-layer-computational-layer}
MML gives the symbolic instruction form. MMA gives the executable matrix-side realization. Thus QAM, MML, and MMA correspond respectively to object language, machine language, and execution dynamics.
\end{remark}

\begin{definition}[Finite branch emulation]\label{def:finite-branch-emulation-computational-layer}
A \emph{finite branch emulation rule} is any route-faithful realization of conditional continuation by means of:
\begin{enumerate}[leftmargin=2em]
    \item a test instruction;
    \item a finite selector family;
    \item a branch-resolution instruction;
    \item an updated program position and control state.
\end{enumerate}
\end{definition}

\begin{remark}[Why explicit jump is not required]\label{rem:why-explicit-jump-is-not-required-computational-layer}
For the finite completeness level pursued here, it is not necessary to require a primitive classical jump instruction. It is sufficient that the machine language can emulate conditional continuation by a finite branch-resolution mechanism.
\end{remark}

\begin{definition}[Finite machine completeness]\label{def:finite-machine-completeness-computational-layer}
The triple
\[
(\mathrm{QAM},\mathrm{MML},\mathrm{MMA})
\]
is called \emph{finitely machine-complete} if:
\begin{enumerate}[leftmargin=2em]
    \item every finite project-relevant symbolic state can be encoded as a finite QAM-store;
    \item every finite admissible symbolic transformation on such states can be represented by a finite MML program;
    \item every such finite MML program is executable in the MMA layer;
    \item conditional continuation is realizable by finite branch emulation;
    \item halting and finite return are representable on the same machine level.
\end{enumerate}
\end{definition}

\begin{definition}[Finite address space]\label{def:finite-address-space-computational-layer}
A \emph{finite address space} is a finite nonempty set
\[
\mathfrak{A}_{\mathrm{fin}}=\{a_1,\dots,a_N\}
\]
of address tags.
\end{definition}

\begin{definition}[Finite machine store as address map]\label{def:finite-machine-store-as-address-map-computational-layer}
A \emph{finite machine store} on the address space
\[
\mathfrak{A}_{\mathrm{fin}}
\]
is a map
\[
\mathsf{Mem}:\mathfrak{A}_{\mathrm{fin}}\longrightarrow \mathsf{QAM}_{\mathrm{fin}},
\]
where \(\mathsf{QAM}_{\mathrm{fin}}\) denotes the class of finite QAM-objects admissible at machine level.
\end{definition}

\begin{definition}[Store read operation]\label{def:store-read-operation-computational-layer}
For \(a\in\mathfrak{A}_{\mathrm{fin}}\), the \emph{store read operation} is
\[
\mathsf{READ}_{\mathsf{Mem}}(a):=\mathsf{Mem}(a).
\]
\end{definition}

\begin{definition}[Store write operation]\label{def:store-write-operation-computational-layer}
For \(a\in\mathfrak{A}_{\mathrm{fin}}\) and \(\xi\in\mathsf{QAM}_{\mathrm{fin}}\), the \emph{store write operation} is the updated store
\[
\mathsf{WRITE}_{\mathsf{Mem}}(a,\xi)=\mathsf{Mem}^{+},
\]
defined by
\[
\mathsf{Mem}^{+}(a)=\xi,
\qquad
\mathsf{Mem}^{+}(b)=\mathsf{Mem}(b)
\;\text{for all } b\neq a.
\]
\end{definition}

\begin{definition}[Address register]\label{def:address-register-computational-layer}
An \emph{address register} is a distinguished element \(\alpha\in\mathfrak{A}_{\mathrm{fin}}\) serving as the currently selected address for machine-level read/write access.
\end{definition}

\begin{definition}[Control register]\label{def:control-register-computational-layer}
A \emph{control register} is a finite tuple
\[
\mathcal{R}_{\mathrm{ctrl}}=(p,q,r,\alpha),
\]
where \(p\) is the current program position, \(q\) the current control state, \(r\) the current route-faithful residual/control datum, and \(\alpha\) the current address register.
\end{definition}

\begin{definition}[Finite machine configuration]\label{def:finite-machine-configuration-computational-layer}
A \emph{finite machine configuration} is a pair
\[
\mathcal{C}=(\mathsf{Mem},\mathcal{R}_{\mathrm{ctrl}})
\]
consisting of a finite machine store and a control register.
\end{definition}

\begin{definition}[Finite universality]\label{def:finite-universality-computational-layer}
The stack
\[
(\mathrm{QAM},\mathrm{MML},\mathrm{MMA})
\]
is called \emph{finitely universal} if every finite admissible symbolic task on finite machine configurations admits a finite MML realization executable in MMA.
\end{definition}

\begin{definition}[Extensible address tower]\label{def:extensible-address-tower-computational-layer}
An \emph{extensible address tower} is a nested family
\[
\mathfrak{A}_1\subset\mathfrak{A}_2\subset\mathfrak{A}_3\subset\cdots
\]
of finite address spaces such that each inclusion preserves the previously existing address tags and every finite set of address demands is realized at some stage.
\end{definition}

\begin{definition}[Store extension map]\label{def:store-extension-map-computational-layer}
Given two stages \(\mathfrak{A}_n\subset\mathfrak{A}_{n+1}\), a \emph{store extension map} is a map
\[
\mathrm{Ext}_{n}^{n+1}:\mathsf{Mem}_n\longmapsto \mathsf{Mem}_{n+1}
\]
such that old address contents are preserved and every new address is initialized by a fixed admissible default QAM-object.
\end{definition}

\begin{definition}[Fresh tag]\label{def:fresh-tag-computational-layer}
Let \(\mathfrak{A}_{\mathrm{fin}}\) be the currently active finite address/tag carrier set. A \emph{fresh tag} is a tag
\[
a_{\mathrm{new}}\notin\mathfrak{A}_{\mathrm{fin}}.
\]
\end{definition}

\begin{definition}[Fresh-address realization rule]\label{def:fresh-address-realization-rule-computational-layer}
A \emph{fresh-address realization rule} assigns to every fresh tag \(a_{\mathrm{new}}\) an extended finite address space
\[
\mathfrak{A}_{\mathrm{fin}}^{+}=\mathfrak{A}_{\mathrm{fin}}\cup\{a_{\mathrm{new}}\}.
\]
\end{definition}

\begin{definition}[Default initialization object]\label{def:default-initialization-object-computational-layer}
A \emph{default initialization object} is a distinguished finite QAM-object
\[
\mathbf{0}_{\mathrm{QAM}}\in\mathsf{QAM}_{\mathrm{fin}}
\]
used to initialize newly created addresses before any further write/update takes place.
\end{definition}

\begin{definition}[Fresh-address store extension]\label{def:fresh-address-store-extension-computational-layer}
Let
\[
\mathsf{Mem}:\mathfrak{A}_{\mathrm{fin}}\to\mathsf{QAM}_{\mathrm{fin}}
\]
be the current finite machine store, and let \(a_{\mathrm{new}}\notin\mathfrak{A}_{\mathrm{fin}}\) be a fresh tag.

The \emph{fresh-address store extension} is the extended store
\[
\mathsf{Mem}^{+}:\mathfrak{A}_{\mathrm{fin}}^{+}\to\mathsf{QAM}_{\mathrm{fin}}
\]
defined by
\[
\mathsf{Mem}^{+}(a)=\mathsf{Mem}(a)
\quad
\text{for all } a\in\mathfrak{A}_{\mathrm{fin}},
\]
and
\[
\mathsf{Mem}^{+}(a_{\mathrm{new}})=\mathbf{0}_{\mathrm{QAM}}.
\]
\end{definition}

\begin{definition}[Old-address semantic preservation]\label{def:old-address-semantic-preservation-computational-layer}
A fresh-address extension satisfies \emph{old-address semantic preservation} if every admissible instruction that only refers to addresses in the old domain has exactly the same effect on those old addresses before and after extension.
\end{definition}

\begin{definition}[Stage-stable universality]\label{def:stage-stable-universality-computational-layer}
The stack
\[
(\mathrm{QAM},\mathrm{MML},\mathrm{MMA})
\]
is called \emph{stage-stably universal} if every extension-compatible stage-stable task family admits a stable unbounded realization by one stage-stable finite program.
\end{definition}

\begin{remark}[Computational status reading]\label{rem:computational-status-reading-computational-layer}
At the present stage, the project computational architecture should be read in three levels:
\begin{enumerate}[leftmargin=2em]
    \item finite machine completeness;
    \item finite universality (Z3-style regime);
    \item stage-stable universality together with fresh-address extension, yielding an architectural Turing-style regime.
\end{enumerate}
\end{remark}

\begin{theorem}[First finite machine completeness reading theorem]\label{thm:first-finite-machine-completeness-computational-layer}
Assume:
\begin{enumerate}[leftmargin=2em]
    \item finite machine states are encoded in the QAM-object layer;
    \item finite admissible symbolic transformations are representable in MML;
    \item the MMA execution layer executes finite MML programs route-faithfully;
    \item finite branch emulation is available.
\end{enumerate}
Then the triple
\[
(\mathrm{QAM},\mathrm{MML},\mathrm{MMA})
\]
supports a finitely machine-complete internal reading.
\end{theorem}

\begin{proof}
Finite tagged QAM stores supply the data layer, finite MML programs supply the symbolic control layer, MMA supplies executable realization, and branch emulation supplies finite conditional continuation. Hence the machine architecture is complete at the finite programmable level.
\end{proof}

\begin{theorem}[Finite address/store layer theorem]\label{thm:finite-address-store-completeness-computational-layer}
Assume:
\begin{enumerate}[leftmargin=2em]
    \item the QAM-object layer encodes finite machine stores;
    \item the MML instruction alphabet contains finite read/write, test, branch, loop, and halt instructions;
    \item the MMA execution layer executes address-valid MML instructions on finite machine configurations;
    \item finite branch emulation is available.
\end{enumerate}
Then the project architecture supports an explicit complete finite address/store machine layer.
\end{theorem}

\begin{proof}
Finite stores are explicit address maps, read/write closure holds on the same finite address domain, finite control is preserved under address-valid execution, and finite conditional control is realizable without leaving the finite regime. Thus the address/store layer is complete.
\end{proof}

\begin{theorem}[First finite universality reading theorem]\label{thm:first-finite-universality-computational-layer}
Assume:
\begin{enumerate}[leftmargin=2em]
    \item the strengthened finite machine completeness theorem holds;
    \item every finite symbolic task decomposes into finitely many canonical finite programming primitives;
    \item finite branch emulation is available.
\end{enumerate}
Then the stack
\[
(\mathrm{QAM},\mathrm{MML},\mathrm{MMA})
\]
supports a finite-universality reading on the internal machine layer.
\end{theorem}

\begin{proof}
Every finite symbolic task can be decomposed into finite read/write, routing, projection, normalization, test, selection, branch, loop, and halt primitives; these are representable in MML and executable in MMA. Hence every finite admissible symbolic task admits a finite realization.
\end{proof}

\begin{theorem}[Fresh tag = fresh address reading theorem]\label{thm:fresh-tag-equals-fresh-address-computational-layer}
Assume:
\begin{enumerate}[leftmargin=2em]
    \item the architecture already contains an implicit extensible address-tower pattern;
    \item fresh tags are admissibly generable;
    \item fresh-address realization holds;
    \item default initialization is available;
    \item old-address semantic preservation holds.
\end{enumerate}
Then the current QAM/MML/MMA architecture supports an explicit extensible-address reading.
\end{theorem}

\begin{proof}
Fresh-address realization turns new tags into genuine new addresses, default initialization turns those addresses into usable machine cells, and old-address semantic preservation makes the extension conservative on the previous stage. Iterating this process makes the address tower explicit.
\end{proof}

\begin{theorem}[First stage-stable universality reading theorem]\label{thm:first-stage-stable-universality-computational-layer}
Assume:
\begin{enumerate}[leftmargin=2em]
    \item the first finite universality theorem holds;
    \item the fresh tag = fresh address theorem holds;
    \item uniform instruction persistence holds;
    \item uniform execution persistence holds;
    \item finite branch emulation persists along stage extension.
\end{enumerate}
Then the stack
\[
(\mathrm{QAM},\mathrm{MML},\mathrm{MMA})
\]
supports a stage-stably universal machine reading.
\end{theorem}

\begin{proof}
Finite universality gives programmable realization on each finite stage. Fresh-address extension makes store growth explicit and conservative. Uniform instruction and execution persistence keep the same finite program semantics valid on larger stages, and persistent branch emulation stabilizes conditional control. Therefore extension-compatible stage-stable task families admit stable unbounded realizations.
\end{proof}

\begin{theorem}[Conditional architectural Turing-style completeness reading theorem]\label{thm:architectural-turing-style-completeness-computational-layer}
Assume the first stage-stable universality reading theorem. Then the stack
\[
(\mathrm{QAM},\mathrm{MML},\mathrm{MMA})
\]
supports a conditional architectural Turing-style completeness reading along the explicit stage-stable route singled out by the present computational layer.
\end{theorem}

\begin{proof}
By stage-stable universality, the architecture supports programmable finite symbolic control, explicit unbounded store growth, stable finite program semantics across arbitrarily large stages, and persistent conditional control. These are exactly the architectural ingredients isolated for a Turing-style machine reading. Because the older Sections~19U--19Z still distinguish unconditional status \(\mathbf S_1\) from stronger conditional route-dependent upgrades toward \(\mathbf S_4(\mathcal R)\) and \(\mathbf S_5\), the present conclusion is recorded only as a conditional architectural reading rather than as a new unconditional global machine-status theorem.
\end{proof}

\begin{corollary}[Computational status embedding rule]\label{cor:computational-status-upgrade-rule-computational-layer}
The project computational status should now be embedded into the older strength hierarchy in the following disciplined way:
\begin{enumerate}[leftmargin=2em]
    \item unconditionally, the public machine-status reading of the line remains at \(\mathbf S_1\) unless the external route certificates of Sections~19U--19Z are invoked;
    \item internally, the present block sharpens the machine reading from a bare \(\mathbf S_1\) interpretation toward a finitely machine-complete / finitely universal (Z3-style) package, i.e. toward the primitive-core side of \(\mathbf S_2\);
    \item explicitly extensible fresh-address and stage-stable universality readings isolate the architectural ingredients that would later be needed on a conditional route toward \(\mathbf S_4(\mathcal R)\) and \(\mathbf S_5\), but they do not by themselves replace the earlier route-dependent certification ladder.
\end{enumerate}
\end{corollary}

\begin{proof}
Immediate from the preceding completeness theorems.
\end{proof}

\begin{remark}[No parallel machine-status hierarchy]\label{rem:no-parallel-machine-status-hierarchy-computational-layer}
This computational block should therefore not be read as introducing a second independent machine-status ladder beside Sections~19U--19Z. Its role is narrower and more local: it explicates the QAM/MML/MMA machine package that underlies the already existing hierarchy and sharpens the internal architectural route toward stronger conditional computational claims, while leaving the unconditional public status decoder unchanged.
\end{remark}

\subsection*{Version and DOI record}
The table below records the currently preserved Companion chain as a normalized release and working-state register. Zenodo releases are listed together with a publication-date column; earlier intermediate records and local working witnesses are retained so that the release chronology remains reconstructible even where the register groups recent working states before older archive rows.

The following chronology block is the active normalized release register for the present v\docversion\ release. The concept DOI remains the stable title-page pointer; version-specific DOI claims for public releases and pre-release freeze witnesses are recorded explicitly in this chronology.

{\small
\renewcommand{\arraystretch}{1.12}
\setlength{\LTleft}{0pt}
\setlength{\LTright}{0pt}
\begin{longtable}{|p{0.14\textwidth}|p{0.17\textwidth}|p{0.15\textwidth}|p{0.42\textwidth}|}
\hline
\textbf{Version / record} & \textbf{Status} & \textbf{Publication date} & \textbf{DOI / note} \\
\hline
\endfirsthead
\hline
\textbf{Version / record} & \textbf{Status} & \textbf{Publication date} & \textbf{DOI / note} \\
\hline
\endhead
Concept DOI & concept record & ongoing pointer & \href{https://doi.org/10.5281/zenodo.19210728}{10.5281/zenodo.19210728} \\
\hline
v3.28.0 & public Zenodo release & 2026-04-20 & \href{https://doi.org/10.5281/zenodo.19648743}{10.5281/zenodo.19648743} --- earlier public release; completes the deferred tool/QAM/decoder layer and fixes the route-control infrastructure that the present release inherits as its public baseline \\
\hline
v3.28.1-r39 & pre-release working continuation & 2026-04-20 & local working state above the public v3.28.0 release; assembled the certificate-aware extension/routing/probe layer, local retained-corridor closure of the live OP~1 and OP~5 Step~(B) kernels, QAM/decoder shadow reattachment, certificate-aware Monte-Carlo discipline, and restored full historical preservation; carried forward here as the final pre-release working witness for v3.29.0 \\
\hline
v3.29.0 & public Zenodo release & 2026-04-20 & \href{https://doi.org/10.5281/zenodo.19670829}{10.5281/zenodo.19670829} --- prior public release; promotes the certificate-aware extension/routing/probe layer, finite decision/action/word routing, matrix/crystal/field readout tags, admissibility and routing discipline, QAM/decoder shadow reattachment, certificate-aware Monte-Carlo discipline, and local retained-corridor closure of the live OP~1 and OP~5 Step~(B) kernels while preserving the full historical witness in Part~II \\
\hline
v3.29.1-r32 & release-freeze working continuation & 2026-04-21 & local freeze witness above the public v3.29.0 release; completes the conservative upgrade-plus-probe integration line through probe-compatible packets, work queues, executable conservative programs, compact readout/output/intake surfaces, OP-facing request/fulfillment/aggregation/dispatch/program/cycle/return logic, a completed second downstream loop, and an explicit release-freeze stabilization point; no version-specific DOI assigned \\
\hline
v3.30.0 & public Zenodo release & 2026-04-21 & \href{https://doi.org/10.5281/zenodo.19682951}{10.5281/zenodo.19682951} --- prior public release; promotes the conservative upgrade-plus-probe integration line above the public v3.29.0 release, including probe-compatible conservative upgrade packets, finite probe-ranked work queues, queue-extracted working programs, route-faithful lead tasks, executable conservative packets, compact conservative readout/output/intake surfaces, OP-facing request/fulfillment/aggregation/dispatch/program/cycle/return logic, a completed second downstream loop, and an explicit release-freeze stabilization point while preserving the full historical witness in Part~II \\
\hline
v3.31.0 & public Zenodo release & 2026-04-24 & \href{https://doi.org/10.5281/zenodo.19699827}{10.5281/zenodo.19699827} --- prior public release; freezes the groundwork-complete pre-OP frontier synthesis above the public v3.30.0 baseline and publishes the full monolith with the integrated coupled prime-lattice trigger correction and the QAM/MML/MMA machine-realization layer while preserving the full historical witness in Part~II \\
\hline
v3.32.0 & public Zenodo release & 2026-04-24 & \href{https://doi.org/10.5281/zenodo.19704489}{10.5281/zenodo.19704489} --- prior public release; current release final built above the public v3.31.0 baseline; integrates the first machine-fused coupled OP-entry transition line, the coupled prime-lattice machine realization block, and the strengthened finite-to-stage-stable universality reading while preserving the full historical witness in Part~II \\
v3.33.0 & public Zenodo release & 2026-04-24 & \href{https://doi.org/10.5281/zenodo.19716926}{10.5281/zenodo.19716926} --- prior public release; built above the public v3.32.0 baseline; carries the HC-captured local continuation line through lock-corridor capture, minimal HC-lock step, HC menu scan, plausible HC candidate, micro-lock update, and first local HC-capture chain while preserving the full historical witness in Part~II \\
\hline
v3.34.0 & public Zenodo release & 2026-04-25 & \href{https://doi.org/10.5281/zenodo.19746821}{10.5281/zenodo.19746821} --- prior public split-bundle release; promotes the locally audit-clean v3.34.0-r27 line into a reader-facing active core plus preserved historical witness bundle, adds the QAM/MML/MMA store-array layer, frontmatter reader architecture, and Sphaera array mirror/stripe/parity-residual infrastructure while preserving the non-closure status of OP~1/4 and OP~5(B) \\
\hline
v3.35.0 & public Zenodo release & 2026-04-25 & \href{https://doi.org/10.5281/zenodo.19750674}{10.5281/zenodo.19750674} --- current public pre-v4 QAM/tensor preparation release; adds the short/extended abstract architecture, Hilbert-hotel QAM reindexing plan for later v4, guarded mixed QAM readouts, optional binary mask registry, conservative two-sided structural tensor grammar, role-erasure and mixed-readout normal forms, Boolean selector grammar, and release-final DOI packaging while preserving the non-closure status of OP~1/4 and OP~5(B) \\
\hline
v3.34.0-r17 & local working candidate & 2026-04-24 & no Zenodo DOI assigned --- executes the first dispatched handoff-guard packet \(\mathsf W_{\mathrm{handoff}}^{\mathrm{guard}}\) after a stable lock/capture return, installs the handoff-guard defect register \(\Delta_{\mathrm{handoff}}^{\mathrm{guard}}\), and returns exactly one of handoff-ready source, bounded handoff-guard repair source, or smaller frontier source without manufacturing an OP closure certificate \\
v3.34.0-r18 & local working candidate & 2026-04-24 & no Zenodo DOI assigned --- executes the first dispatched OP-normalization/type packet \(\mathsf W_{\mathrm{norm}}^{\mathrm{type}}\) after a handoff-ready return, installs the normalization/type-coherence defect register \(\Delta_{\mathrm{norm}}^{\mathrm{type}}\), and returns exactly one of route/rank-ready typed OP input, bounded normalization/type repair source, or smaller frontier source without manufacturing an OP closure certificate \\
v3.34.0-r19 & local working candidate & 2026-04-24 & no Zenodo DOI assigned --- executes the first dispatched route/rank packet \(\mathsf W_{\mathrm{route}}^{\mathrm{rank}}\) after a normalization-ready return, installs the route/rank defect register \(\Delta_{\mathrm{route}}^{\mathrm{rank}}\), and returns exactly one of comparison-ready routed OP task, bounded route/rank repair source, or smaller frontier source without manufacturing an OP closure certificate \\
v3.34.0-r20 & local working candidate & 2026-04-24 & no Zenodo DOI assigned --- executes the first dispatched comparison/defect packet \(\mathsf W_{\mathrm{cmp}}^{\Delta}\) after a route/rank-ready return, installs the comparison/defect register \(\Delta_{\mathrm{cmp}}^{\Delta}\), and returns exactly one of guarded bind candidate, guarded split candidate, or smaller comparison frontier without manufacturing an OP closure certificate \\
v3.34.0-r21 & local working candidate & 2026-04-24 & no Zenodo DOI assigned --- executes the first dispatched resolve/return packet \(\mathsf W_{\mathrm{resolve}}^{\mathrm{next}}\) after a guarded comparison return, installs the resolve/return defect register \(\Delta_{\mathrm{resolve}}^{\mathrm{next}}\), and returns exactly one of guarded bind-work continuation, guarded split-work continuation, or smaller comparison-frontier descent without manufacturing an OP closure certificate \\
v3.34.0-r22 & local working candidate & 2026-04-24 & no Zenodo DOI assigned --- executes the first continuation-queue packet after resolve/return, installs the continuation-queue defect register \(\Delta_{\mathrm{queue}}^{\mathrm{cont}}\), and returns exactly one of a guarded common-origin bind queue, guarded ranked split queue, or protected frontier-descent queue without manufacturing an OP closure certificate \\
v3.34.0-r23 & local working candidate & 2026-04-24 & no Zenodo DOI assigned --- executes the first queue-head selection packet after continuation queue, installs the queue-head selection defect register \(\Delta_{\mathrm{head}}^{\mathrm{select}}\), and returns exactly one of a guarded common-origin-check head, guarded branch-rank-check head, or guarded frontier-descent-check head while retaining the finite tail and without manufacturing an OP closure certificate \\
v3.34.0-r24 & local working candidate & 2026-04-24 & no Zenodo DOI assigned --- executes the first selected queue-head execution packet, installs the selected queue-head execution defect register \(\Delta_{\mathrm{head}}^{\mathrm{exec}}\), and returns exactly one of an executed common-origin-check surface, executed branch-rank-check surface, or executed frontier-descent-check surface while retaining the finite queue tail and without manufacturing an OP closure certificate \\
v3.34.0-r25 & local working candidate & 2026-04-24 & no Zenodo DOI assigned --- executes the first executed-head tail-entry packet \(\mathsf W_{\mathrm{tail}}^{\mathrm{entry}}\), installs the tail-entry defect \(\Delta_{\mathrm{tail}}^{\mathrm{entry}}\), and returns exactly one of a guarded bind-tail, split-tail, or frontier-tail entry without manufacturing an OP closure certificate \\
v3.34.0-r26 & local audit-freeze candidate & 2026-04-24 & no Zenodo DOI assigned --- installs the first audit-freeze ledger \(F_{\mathrm{HCOP}}^{\mathrm{audit}}\) with audit-freeze defect \(\Delta_{\mathrm{audit}}^{\mathrm{freeze}}\), recording baseline, DOI state, local spine, tail-entry output, preservation obligations, reference/compilation audit data, terminal bibliography, historical witness, and the no-false-certificate marker \\
v3.34.0-r27 & local audit-freeze candidate & 2026-04-24 & no Zenodo DOI assigned --- installs the audit-freeze stabilization object \(F_{\mathrm{freeze}}^{\mathrm{audit}}\), permitting only audit, mechanical repair, metadata repair, release packaging, and Zenodo-description preparation before further theorem expansion \\
\hline
v3.34.0-r15 & local working candidate & 2026-04-24 & no Zenodo DOI assigned --- turns the r14 tail-normalized HC/OP proof skeleton into the first controlled skeleton-to-workstep dispatch object \(\mathsf D_{\mathrm{HCOP}}^{\mathrm{skel}}\), assigning finite admissible next work packets to lock/capture, handoff, normalization, routing, comparison, and resolution nodes while preserving dependency order, guards, witnesses, route tags, theorematic status, and the no-false-certificate rule \\
v3.34.0-r16 & local working candidate & 2026-04-24 & no Zenodo DOI assigned --- executes the first dispatched lock/capture stabilization packet \(\mathsf W_{\mathrm{lock}}^{\mathrm{stab}}\), installs the lock-stability defect register \(\Delta_{\mathrm{lock}}^{\mathrm{stab}}\), and returns exactly one of handoff-ready stable source, bounded local repair source, or smaller frontier source without manufacturing an OP closure certificate \\
\hline
v3.34.0-r14 & local working candidate & 2026-04-24 & no Zenodo DOI assigned --- extracts the first tail-normalized HC/OP proof skeleton \(\mathsf S_{\mathrm{HCOP}}^{\mathrm{tail}}\), organizing the active r01--r04 spine as LOCK/CAPTURE, HANDOFF, NORMALIZE, ROUTE, COMPARE, and RESOLVE source nodes plus later guarded citation tails while preserving dependency order and preventing false OP-closure certificates \\
\hline
v3.34.0-r13 & local working candidate & 2026-04-24 & no Zenodo DOI assigned --- turns the r12 ledger into the first guarded-reduction pilot normal form \(\mathrm{GRCite}_{\mathrm{HCOP}}(\ell;G,W,S)\), allowing only later repeated HC/OP tails to be rewritten as guarded tool citations while first occurrences, local guards, witnesses, certificate tests, proof-critical inequalities, theorematic status, and historical witness material remain explicit \\
\hline
v3.34.0-r12 & local working candidate & 2026-04-24 & no Zenodo DOI assigned --- applies the r11 guarded-reduction protocol for the first time by installing the finite ledger \(\mathcal L_{\mathrm{HCOP}}^{(1)}\) over the active HC/OP proof spine, allowing later repeated lock/capture, handoff, normalization, route/frontier, paired-comparison, probe-diagnostic, and archive-to-tool tails to be cited in guarded normal form while keeping first occurrences, local guards, certificate tests, proof-critical inequalities, theorematic status, and historical witness material explicit \\
\hline
v3.34.0-r11 & local working candidate & 2026-04-24 & no Zenodo DOI assigned --- keeps the r10 tool-contract application layer and adds the first tool-guarded reduction protocol \(\mathcal R_{\mathrm{HCOP}}^{(1)}\), defining reducible and non-reducible HC/OP occurrences, guarded citation normal form, and the partial simplification map \(\mathsf{Red}_{\mathrm{HCOP}}^{(1)}\) while preserving proof content, first definitions, local guards, certificate tests, theorematic status, and the full historical witness in Part~II \\
\hline
v3.34.0-r10 & local working candidate & 2026-04-24 & no Zenodo DOI assigned --- applies the r09 reusable tool contracts back to the active r01--r04 HC/OP proof spine as a conservative cross-reference layer \(\mathcal P_{\mathrm{HCOP}}^{\mathrm{tool}}\), allowing repeated route, frontier, and no-false-certificate safety prose to be replaced by typed tool-contract citations only where the original hypotheses and guards remain visible, while preserving proof content and the full historical witness in Part~II \\
\hline
v3.34.0-r09 & local working candidate & 2026-04-24 & no Zenodo DOI assigned --- keeps the r08 compression map unchanged and extracts the first reusable HC/OP tool normal forms \(\mathbb T_{\mathrm{HCOP}}^{(1)}\), giving typed contracts for lock/capture, handoff, normalization, routing, comparison, resolution, probe ranking, certificate tests, frontier descent, and archive-to-tool extraction, so later text can cite tool contracts instead of repeating safety prose while preserving proof content and the full historical witness in Part~II \\
\hline
v3.34.0-r08 & local working candidate & 2026-04-24 & no Zenodo DOI assigned --- keeps the r07 QAM/MML/MMA interface unchanged and adds \(\mathfrak C_{\mathrm{HCOP}}\), a conservative MML/MMA compression map for inherited HC/OP text, so repeated capture, handoff, normalization, routing, comparison, probe, certificate-test, route-preservation, and historical-tool paragraphs can be reread as typed operations without deleting proof content or manufacturing OP closure certificates, while preserving the full historical witness in Part~II \\
\hline
v3.34.0-r07 & local working candidate & 2026-04-24 & no Zenodo DOI assigned --- keeps the r04 paired OP-entry comparison execution and r06 bibliography hardening unchanged, assigns local interface versions \(\mathrm{QAM}_{\mathrm{HCOP}}^{\mathrm{loc}}\), \(\mathrm{MML}_{\mathrm{HCOP}}^{\mathrm{loc}}\), and \(\mathrm{MMA}_{\mathrm{HCOP}}^{\mathrm{loc}}\), and proves that the r01--r04 HC-routed OP-entry chain is executable as one finite typed MML/MMA program without manufacturing an OP closure certificate, while preserving the full historical witness in Part~II \\
v3.34.0-r06 & local working candidate & 2026-04-24 & no Zenodo DOI assigned --- keeps the r04/r05 mathematical and bibliography-placement state unchanged, expands the terminal bibliography after a whole-document literature scan, adds compact front/back literature anchor notes for Teichmuller/quasiconformal geometry, Riemann--GUE/RMT, Hopf/fibre-bundle topology, division algebras, gauge/QFT landmarks, CKM/PMNS mixing, Julia dynamics, set-theoretic foundations, and probe/cryptographic/computational diagnostics, while preserving the full historical witness in Part~II \\
v3.34.0-r05 & local working candidate & 2026-04-24 & no Zenodo DOI assigned --- keeps the r04 paired OP~1/4--OP~5(B) comparison execution layer unchanged and normalizes the bibliography/back-matter placement by moving the live bibliography to the terminal Final Bibliography block while preserving the full historical witness in Part~II \\
\hline
v3.34.0-r04 & local working candidate & 2026-04-24 & no Zenodo DOI assigned --- executes the paired OP~1/4--OP~5(B) comparison branch above the r03 task-routing layer; compares joint versus separated HC-compatible readouts and returns either a bound common-origin state, a split ranked single-branch state, or a smaller frontier-certified comparison state while preserving route tags and the full historical witness in Part~II \\
\hline
v3.34.0-r03 & local working candidate & 2026-04-24 & no Zenodo DOI assigned --- extracts the first finite OP-entry task-routing layer from the HC-derived normal form; maps the selected live residual branch into OP~1/4 continuum-limit work, OP~5(B) routed $\Theta$--$R$ balance work, anchor/admissibility refinement, or paired OP~1/4--OP~5(B) comparison while preserving route tags and the full historical witness in Part~II \\
\hline
v3.34.0-r02 & local working candidate & 2026-04-24 & no Zenodo DOI assigned --- extracts the first downstream-readable OP-entry normal form from the transport-stable HC-to-OP handoff packet; routes the terminal live residual triple into OP~1/4, OP~5(B), and anchor/admissibility request families while preserving route tags and the full historical witness in Part~II \\
\hline
v3.34.0-r01 & local working candidate & 2026-04-24 & no Zenodo DOI assigned --- begins the transport-stable HC-capture and OP-entry handoff line above the public v3.33.0 release; packages local HC capture into a route-faithful transport-stability test and an OP-facing handoff packet while preserving the full historical witness in Part~II \\
\hline
v3.25.0 & public Zenodo release & 2026-04-16 & \href{https://doi.org/10.5281/zenodo.19615952}{10.5281/zenodo.19615952} --- earlier public release; consolidates the optical-comparison and machine-theoretic overlay line, including compressed optical diagnostic control, bundled proof-use compression, conditional reference-core closure, and the capped machine-theoretic overlay line up to the present equivalence surface \\ 
\hline
v3.18.0 & public Zenodo release & 2026-04-13 & \href{https://doi.org/10.5281/zenodo.19548590}{10.5281/zenodo.19548590} --- earlier public release; carries the public v3.17.0 tensor/operator baseline forward into a first public OP~3/QAM closure package with shell certificates, shell-defect localization, shell-stable readout factorization, and a bundled OP~3$\to$OP~5/global interface while keeping stronger kernel-first activation internal \\
\hline
v3.17.0 & public Zenodo release & 2026-04-10 & \href{https://doi.org/10.5281/zenodo.19522352}{10.5281/zenodo.19522352} --- earlier public release; promotes the third public-shape QAM tensor/operator line above the public v3.16.0 release and keeps translation/round-trip statements in front-facing public form while leaving support-selection and wider compression internal \\
\hline
v3.16.0 & public Zenodo release & 2026-04-10 & \href{https://doi.org/10.5281/zenodo.19489324}{10.5281/zenodo.19489324} --- earlier public release; promotes the second public-shape QAM branch/shell/routed layer while leaving tensor/operator and support-selection expansion internal, and records a deferred kernel-first compression path with activation criteria, staged stress tests, and migration roadmap \\
\hline
v3.15.0 & public Zenodo release & 2026-04-08 & \href{https://doi.org/10.5281/zenodo.19463389}{10.5281/zenodo.19463389} --- earlier public release; freezes the first compressed public-shape QAM core on top of the public v3.14.0 branch/shell baseline \\
\hline
v3.14.0 & public Zenodo release & 2026-04-07 & \href{https://doi.org/10.5281/zenodo.19460589}{10.5281/zenodo.19460589} --- earlier public release; extends the public v3.13.0 line by freezing the first post-closure branch/shell hardening layer \\
\hline
v3.3.0 & local working state & 2026-04-03 & condensed Appendix-C integration of Julia-boundary, master-residual, Umkehrwalze, and reversal-probe sector \\
\hline
v3.3.1 & local working state & 2026-04-03 & clean condensed Appendix-C line; 16-page block reduced to a 3-page support appendix \\
\hline
v3.3.2 & local working state & 2026-04-03 & clean-mainline promotion: compact 3.3 synthesis theorem moved into main text; Appendix-C support retained \\
\hline
v3.3.3 & local final-polish state & 2026-04-03 & removed embedded editorial fragments in the new 3.3 synthesis line and normalized notation in the compressed summary \\
\hline
v3.3.4 & local editorial-clean state & 2026-04-03 & global editorial cleanup: removed remaining visible drafting fragments, normalized compressed remarks, and cleaned heuristic table notes \\
\hline
v3.3.5 & Zenodo release & 2026-04-03 & \href{https://doi.org/10.5281/zenodo.19396143}{10.5281/zenodo.19396143} --- strict wording cleanup: aligned version metadata, removed residual status-weakening phrasing in the 3.3 line, normalized the current-version record, and finalized the 3.3.5 release package \\
\hline
v3.4.0 & local working extension & 2026-04-03 & first 3.4 pass: added closure-potential globalization of the residual branch law, conserved branch energy, and an energy-form reading of the reversal-probe dichotomy \\
\hline
v3.4.1 & local working extension & 2026-04-03 & constrained the exact branch potential to the reversible single-well class and reduced the global scalar ambiguity to the transported $\mathbb Z_2$ pair \\
\hline
v3.4.2 & local working extension & 2026-04-03 & extracted the canonical quartic one-field family for the visible residual branch law and added turning-value parameterization of regular probe shells \\
\hline
v3.4.3 & local working extension & 2026-04-03 & tied the visible quartic branch coefficient to the transported quartic closure coefficient and replaced the technical turning-amplitude label by turning value \\
\hline
v3.4.4 & local working extension & 2026-04-03 & added the first nonlinear period-drift readout, showing that the timing defect on regular periodic shells is already fixed by the same quartic closure coefficient \\
\hline
v3.4.5 & local working extension & 2026-04-03 & extended the quartic readout to the full even hierarchy, proving that the first higher timing defect appears exactly at the order of the first surviving even closure coefficient \\
\hline
v3.4.6 & local working extension & 2026-04-03 & repackaged the transported even hierarchy into a single analytic generator $\Phi(\eta^2)$ for the visible branch potential and its nonlinear shell readouts \\
\hline
v3.4.7 & local working extension & 2026-04-03 & generator-shell probe classifier: small regular shells reduced to the generator radius $s(E)=\Phi^{-1}(E)$ and the transported sign pair $\pm\sqrt{s(E)}$ \\
\hline
v3.4.8 & local working extension & 2026-04-03 & shell-parity criterion: direct convergence versus spinorial/M\"obius lift determined by the even/odd reflector-crossing parity on sufficiently small regular shells; added first diagnostic Julia-type appendix render \\
\hline
v3.4.9 & local working extension & 2026-04-03 & added cautious frame-shifted detector-readout hypothesis for the diagnostic Julia/probe geometry, together with explicit candidate observables for later comparison work \\
\hline
v3.4.10 & local working extension & 2026-04-03 & converted the detector-readout hypothesis into a quantitative tunnel-observable interface: tunnel width, edge intensification, opening angle, and shell parity are now organized as a minimal diagnostic comparison package \\
\hline
v3.4.11 & local working extension & 2026-04-03 & added a calibrated inverse-readout theorem: tunnel width plus shell parity now recover the local shell radius, while edge intensification and opening angle become shell-consistency checks measured by $\Delta_{\mathrm{diag}}$ \\
\hline
v3.4.12 & local working extension & 2026-04-03 & added a local inverse identifiability/stability theorem: the calibrated tunnel-width law now carries a shell-wise stability margin and inverse condition number, so shell reconstruction is controlled whenever the calibration slope stays away from zero \\
\hline
v3.4.13 & local working extension & 2026-04-03 & organized the inverse readout globally into a shell atlas stratified by calibration-critical shells and parity-switch shells, so stable inversion now proceeds chart by chart away from the obstruction set \\
\hline
v3.4.14 & local working extension & 2026-04-03 & added an obstruction-shell transition law: isolated obstructions now carry only fold exchange, parity flip, or their composition, giving the inverse-readout layer a small $\mathbb Z_2$ transition algebra and local obstruction monodromy \\
\hline
v3.4.15 & local working extension & 2026-04-03 & organized the obstruction transitions into a global Klein-four transition cocycle on the shell atlas and showed that, after fold resolution, the remaining monodromy reduces to the parity local system underlying the spinorial/M\"obius lift \\
\hline
v3.4.16 & local working extension & 2026-04-03 & promoted the parity local system to a global double-cover criterion: direct visible convergence is globally available exactly when the parity class trivializes, while every nontrivial case becomes globally single-valued on the canonical parity double cover \\
\hline
v3.4.17 & local working extension & 2026-04-03 & recast the residual parity sector as a genuine cohomological obstruction class in $H^1(B,\mathbb Z_2)$, identifying direct visible convergence with the vanishing class and the canonical parity double cover with its associated classifying two-sheeted cover \\
\hline
v3.4.18 & local working extension & 2026-04-03 & promoted the parity obstruction to an operational global holonomy/readout character: admissible closed shell paths now evaluate the class $\varepsilon_{\mathrm{par}}$ directly, so direct visible return versus spinorial/M\"obius lift is read as a $\mathbb Z_2$ loop outcome factoring through mod-two homology \\
v3.4.19 & local working extension & 2026-04-03 & turned the parity holonomy character into an empirical loop-readout estimator: chartwise parity-jump counting across admissible obstruction shells now recovers the cohomological $\mathbb Z_2$ holonomy, giving a detector-side loop test for direct visible return versus spinorial/M\"obius lift \\
v3.4.20 & local closing state & 2026-04-03 & clean 3.4 closure: compressed the full extension line into a single generator-atlas-holonomy synthesis theorem, identified the minimal handoff package $(\Phi,\Delta_{\mathrm{diag}},\varepsilon_{\mathrm{par}},\chi_{\mathrm{par}},\widehat\chi_{\mathrm{par}})$, and marked the end of the 3.4 build-up before the next-stage simplification work \\
\hline
v3.5.0 & local working extension & 2026-04-03 & first compression-stage theorem: introduced the residual record $\mathfrak R(B)$ and proved the reconstruction normal form for the stabilized visible sector \\
\hline
v3.5.1 & local working extension & 2026-04-03 & record-level transport step: defined admissible residual-record morphisms and isomorphisms, and formalized the faithful translation principle on the stabilized visible sector \\
\hline
v3.5.2 & local abstract-polish state & 2026-04-03 & sharpened the abstract and the compression-stage framing: clarified the residual-record interface, the conservative-frontend requirement, and the split between structurally closed versus still-open global identification problems \\
\hline
v3.5.3 & local status-cleanup state & 2026-04-04 & reclassified the open-questions layer by actual status, absorbed OP~2(T4) into OP~1, marked OP~3 as no longer a primary open problem, and reduced OP~5 to Step~(B) only \\
\hline
v3.5.4 & local chronology-sync state & 2026-04-04 & synchronized the document head with the current 3.5.x line: added the missing v3.5.0--v3.5.4 chronology entries and replaced outdated roadmap goals by the residual-record/compression-stage agenda \\
\hline
v3.5.5 & local reader-guide alignment state & 2026-04-04 & aligned the Reader's Guide with the actual 3.5.x status: separated live fronts from preserved historical roadmap layers and made the OP~1/2/3/5 status split explicit at the document head \\
\hline
v3.5.6 & local main-text status-gloss state & 2026-04-04 & propagated the 3.5.x status split into the main compression-stage remarks and appendix support transitions: clarified that remaining fronts are downstream residual tasks and not invitations to reopen the visible interface \\
\hline
v3.5.7 & local stage-summary state & 2026-04-04 & added an explicit 3.5.x-versus-3.6 orientation block: recorded what the compression line has already achieved and fixed the next target as conservative language stabilization above the residual-record interface \\
\hline
v3.5.8 & local figure-integration state & 2026-04-04 & rebuilt the missing Julia/probe comparison render, integrated it into the text with a cautious diagnostic caption, and restored a clean PDF build with the image in place \\
\hline
v3.6.0 & local working extension & 2026-04-04 & introduced the conservative-frontend principle: any later language layer must factor through the residual-record interpreter and may not introduce new visible primitives \\
\hline
v3.6.1 & local frontend-hardening state & 2026-04-04 & hardened the conservative frontend line by adding closure under composition and a standalone primitive non-drift criterion; clarified that Section~3.6 fixes admissibility constraints for later language layers rather than introducing a second visible interface \\
\hline
v3.6.2 & local frontend-finalization state & 2026-04-04 & promoted the closure/non-drift hardening step into the main line and stabilized the 3.6 block as a completed admissibility framework for later language chains \\
\hline
v3.7.0 & local OP-first consolidation state & 2026-04-04 & formalized the current OP-first reading: the remaining live fronts are organized as residual downstream tasks over the fixed record interface, with direct language growth deliberately deferred \\
\hline
v3.7.1 & Zenodo release & 2026-04-04 & \href{https://doi.org/10.5281/zenodo.19412478}{10.5281/zenodo.19412478} --- public release of the OP-first 3.7 line; froze the present stabilized-reading claim and fixed 3.8 as an explicit OP-closure phase while keeping QAM deferred \\
\hline
v3.8.0 & local two-core-program state & 2026-04-04 & promoted 3.8 from a planning note to a working OP-closure stage: reduced the foundational historical burden to the OP~1 correction core and the OP~5 step~(B) Theta-balance core, with all other live fronts treated as secondary unless dependency is shown \\
\hline
v3.8.1 & local residual-emergence state & 2026-04-04 & recast the field-equation / controlled-emergence question into a residual-controlled physical interface form, concentrating the unresolved burden into the explicit remainder families \(\mathcal E_1\) and \(\mathcal E_5\) instead of introducing a new visible primitive basis \\
\hline
v3.8.2 & local OP~1 reduction state & 2026-04-04 & reduced the remaining OP~1 burden to a single scalar residual channel \(\mathcal R_1(\mu)\), making OP~1 a one-channel closure problem rather than a diffuse correction family \\
\hline
v3.8.3 & local OP~1 root-isolation state & 2026-04-04 & proved that the OP~1 scalar residual channel has a unique root in the physical window \(1.85\le \mu\le 1.90\), numerically \(\mu_\ast\approx 1.86366225\), thereby turning the remaining OP~1 burden from a broad numerical search into an exact structural-identification problem \\
v3.8.4 & local OP~1 frame-shift diagnostic state & 2026-04-04 & showed that the residual \(\gamma_L-\mu_\ast\) shrinks on the matter-frame offset scale, behaves like a local Newton/readout correction, and does not trigger a parity/horizon jump in the current Julia/probe family at the tested residual scales \\
v3.8.5 & local minimal neutrino-inflow test state & 2026-04-04 & converted the archived \(k(\nu)=4\) inflow hypothesis into a minimal OP~1 downstream correction ansatz \(\mathcal R_{1,\nu}(\mu;\eta_\nu)=\mathcal R_1(\mu)+\eta_\nu/\mu^4\), with calibrated \(\eta_\nu^\ast\approx -1.10\times10^{-3}\) and no reopening of the visible primitive basis \\
v3.8.6 & local semi-structural inflow estimate state & 2026-04-04 & showed that the calibrated coefficient satisfies \(\eta_\nu^\ast\approx -\varepsilon_\nu\gamma_L^4\beta_{\mathrm{sym}}^4\) with \(\varepsilon_\nu\approx 1.0012\times10^{-3}\), so the fitted inflow law reduces to a \(k=4\) matter-share factor times a per-mille fermionic leakage efficiency \\
v3.8.7 & local defect-squared inflow state & 2026-04-04 & identified the per-mille leakage efficiency with the already existing quadratic defect invariant \(\eta_{\mathrm{sym}}^2=(\frac13-\beta_{\mathrm{sym}}^2)^2\), yielding the sharper law \(\eta_\nu^\ast\approx -\eta_{\mathrm{sym}}^2\gamma_L^4\beta_{\mathrm{sym}}^4\) \\
v3.8.8 & local second-order fermionic leakage state & 2026-04-04 & used the \(J\)-antisymmetric fermionic sector \(\Hplusminus\) to justify why the first visible inflow scalar should be quadratic in the transported defect, leading to \(\eta_\nu\approx -\lambda_\nu\eta_{\mathrm{sym}}^2\gamma_L^4\beta_{\mathrm{sym}}^4\) with \(\lambda_\nu\approx 0.9990\) \\
v3.8.9 & local cosmological-offset link state & 2026-04-04 & showed that any cosmological contribution to the OP~1 inflow remainder is naturally mediated by the permanent frame-offset term \(\beta_{\mathrm{sym}}^2\) already used in the dark-offset theorem, not by a direct Hubble-rate coupling, yielding \(\eta_\nu^\ast\approx -(\frac13-\beta_{\mathrm{sym}}^2)^2\gamma_L^4\beta_{\mathrm{sym}}^4\) as the natural weak expansion reading \\
v3.8.10 & local one-background collapse state & 2026-04-04 & showed that the frame/readout, fermionic-leakage, and weak cosmological-offset readings of the smooth OP~1 remainder collapse numerically to the same background law \(-\eta_{\mathrm{sym}}^2\gamma_L^4\beta_{\mathrm{sym}}^4\), leaving the discrete prime-lattice mechanism as the only genuinely independent live burden \\
v3.8.11 & local readout-point cancellation state & 2026-04-04 & showed that the scalar prime/continuum residual \(\mathcal R_1(\gamma_L)\approx 9.1216\times10^{-5}\) and the collapsed smooth background term \(-\eta_{\mathrm{sym}}^2\beta_{\mathrm{sym}}^4\approx -9.1208\times10^{-5}\) cancel to within \(\Delta_{\mathrm{can}}\approx 8.68\times10^{-9}\), shifting the live OP~1 burden from finding a new correction term to explaining the near-identity itself \\
v3.8.12 & local background-corrected bridge state & 2026-04-04 & defined the bridge function \(\Phi_1(\mu)=\mathcal R_1(\mu)-\eta_{\mathrm{sym}}(\mu)^2\beta(\mu)^4\), proved that it has a unique root \(\gamma_{\mathrm{br}}\approx 1.86399996799643\) on the physical window, and showed that the Companion matter-frame value \(1.864\) already coincides with this bridge root to about \(3.2\times10^{-8}\) \\
v3.8.13 & local bridge-shape forcing state & 2026-04-04 & showed that the current Companion architecture forces any minimal admissible visible OP~1 bridge term to have the shape \(\lambda_{\mathrm{br}}\eta_{\mathrm{sym}}(\mu)^2\beta(\mu)^4\), leaving only the canonical normalization \(\lambda_{\mathrm{br}}\approx 1.00010\) as the live remaining freedom \\
v3.8.14 & local multiplicative bridge-renormalization state & 2026-04-04 & rewrote the forced bridge equation as the dressed readout law \(\beta(\mu)=\beta_{\mathrm{CL}}(\mu)\exp(\eta_{\mathrm{sym}}(\mu)^2\beta(\mu)^4)\), showing that the apparent non-canonicity of \(\lambda_{\mathrm{br}}\) at \(1.864\) is only the tiny \(8.68\times10^{-9}\) displacement from the exact bridge root rather than a missing normalization mechanism \\
v3.8.15 & local transport-balanced exponential state & 2026-04-04 & derived the exponential self-background dressing as the finite form of a repeated compensator-balanced transport law on the positive readout amplitude, reducing the remaining OP~1 burden to the path-constancy/normalization of the bridge source rather than to the appearance of the exponential itself \\
v3.8.16 & local bridge-contraction state & 2026-04-04 & recast the bridge law as the fixed-point equation \(T_\mu(\beta)=\beta\) and showed that the bridge map is strongly contractive at the distinguished root, with \(|\partial_b T_{\gamma_{\mathrm{br}}}(\beta_{\mathrm{br}})|\approx 0.00311\), thereby turning the source-freezing step into a local stability statement rather than a free assumption \\
v3.8.17 & local OP~1 closure state & 2026-04-04 & packaged the 3.8.2--3.8.16 line into a local closure theorem at the distinguished matter-frame point \(\gamma_{\mathrm{loc}}=\gamma_{\mathrm{br}}\approx 1.86399996799643\), relocating the remaining burden from local OP~1 readout to the question of global extension versus transition to OP~5 Step~(B) \\
v3.8.18 & final local OP~1 closure release state & 2026-04-04 & finalized the 3.8 release with version DOI \href{https://doi.org/10.5281/zenodo.19414221}{10.5281/zenodo.19414221}; the release message is local closure of OP~1 at the distinguished matter-frame point, with OP~5 Step~(B) promoted to the next primary live front \\
v3.8.19 & local OP~5 one-point reduction state & 2026-04-04 & used local OP~1 closure to fix the matter-frame and photon-frame background data, reducing historical OP~5 Step~(B) to the single scalar finite-cutoff identity \(\Omega_{5,\mathrm{loc}}=0\) at \(T_{\mathrm{loc}}=\gamma_{\mathrm{loc}}^{-3/2}\approx0.3929444410\) with target value \(\mathcal T(T_{\mathrm{loc}})\approx0.893255\) \\
v3.8.20 & local OP~5 exact-tail state & 2026-04-04 & rewrote historical OP~5 Step~(B) as the exact tail identity \(\mathcal E_N(T_{\mathrm{loc}})=0.893255-\mathcal T_N(T_{\mathrm{loc}})\) above an arbitrary truncation level \(N\); in particular, using the archived \(N=50\) diagnostic \(\mathcal T_{50}(T_{\mathrm{loc}})\approx-0.146\), the required tail contribution becomes \(\mathcal E_{50}(T_{\mathrm{loc}})\approx1.039255\) \\
v3.8.21 & local OP~5 counting-tail state & 2026-04-04 & rewrote the exact OP~5 tail as a Stieltjes integral against the photon-frame zero counting tail \(N_{\mathrm{ph}}^{>N}(u)\) and, after integration by parts, as the oscillatory counting-function integral with kernel \(K_{T_{\mathrm{loc}}}(u)=T_{\mathrm{loc}}\cos(T_{\mathrm{loc}}u)/u-\sin(T_{\mathrm{loc}}u)/u^2\), thereby isolating the exact analytic interaction that remains to be controlled \\
v3.8.22 & local OP~5 smooth-density split state & 2026-04-04 & decomposed the exact OP~5 tail into a smooth counting-density contribution \(\mathcal E_N^{\mathrm{sm}}(T_{\mathrm{loc}})\) plus a genuine fluctuation term \(\mathcal E_N^{\mathrm{fl}}(T_{\mathrm{loc}})\), turning Step~(B) into the precise question whether the distinguished-cutoff tail is density-dominated or fluctuation-dominated \\
v3.8.23 & local OP~5 leading-density test state & 2026-04-04 & inserted the leading transported photon-frame Riemann--von Mangoldt density, finding \(\mathcal E_{50,\mathrm{lead}}^{\mathrm{sm}}(T_{\mathrm{loc}})\approx 0.047867069\) versus the exact required tail \(\mathcal E_{50}(T_{\mathrm{loc}})\approx 1.039202153\); thus only about \(4.6\%\) of the tail is smooth and Step~(B) is decisively fluctuation-dominated \\
v3.8.24 & local OP~5 pure-fluctuation state & 2026-04-04 & subtracted the leading smooth counting density completely and rewrote Step~(B) as the pure fluctuation-tail identity \(-2\int_0^\infty S_{\mathrm{ph}}^{>N}(u)K_{T_{\mathrm{loc}}}(u)\,du=(0.893255-\mathcal T_N(T_{\mathrm{loc}}))-\mathcal E_{N,\mathrm{lead}}^{\mathrm{sm}}(T_{\mathrm{loc}})\); in particular, at \(N=50\) the fluctuation tail must contribute about \(0.991335084\) \\
v3.8.25 & local OP~5 pair-block state & 2026-04-04 & reorganized the pure fluctuation tail into adjacent photon-frame pair blocks \(P_{N,j}(T_{\mathrm{loc}})\), yielding the exact pair formula controlled by the pair means \(x_{N,j}\) and half-gaps \(d_{N,j}\); this turns the next OP~5 attack into a mesoscopic resonance test rather than another smooth-density refinement \\
v3.8.26 & local OP~5 quadratic-gap state & 2026-04-04 & showed that the exact adjacent pair-block profile is even in the half-gap \(d\), so the linear gap term vanishes identically and the first nontrivial local-gap correction is quadratic; consequently the first mesoscopic OP~5 layer is mean-phase dominated rather than gap-noise dominated \\
v3.8.27 & local OP~5 pair-mean quasi-lattice state & 2026-04-04 & showed that the adjacent pair means are well described by a slowly drifting density-controlled quasi-lattice, with observed step ratios near the prediction \(2/\rho_{\mathrm{ph}}(x_{N,j})\) having mean 0.997634; thus the first mesoscopic OP~5 burden is slow phase drift rather than linear gap disorder \\
v3.8.28 & local OP~5 four-block symmetry state & 2026-04-04 & grouped adjacent pair blocks into four-blocks and showed that the resulting super-gap dependence is even, so the linear pair-mean drift term vanishes identically; thus OP~5 now exhibits a two-stage mesoscopic suppression in which both the linear local-gap term and the linear pair-mean drift term are removed \\
v3.8.29 & local OP~5 super-mean quasi-lattice state & 2026-04-04 & showed that the four-block super-means themselves form a second-stage slowly drifting density-controlled quasi-lattice, with observed step ratios near the prediction \(4/\rho_{\mathrm{ph}}(X_{N,k})\) having mean 0.996050; thus the current OP~5 burden is best read on a hierarchy of quasi-lattices rather than as a raw oscillatory tail \\
v3.8.30 & local OP~5 quarter-turn super-mean state & 2026-04-04 & showed that the super-mean phases at the benchmark cutoff cycle through the four quarter-turn sectors in strict order for the first \(16\) super-means, while drifting on average by about \(-0.0652531\) relative to the exact \(\pi/2\) increment; thus the current OP~5 super-mean layer exhibits weak quarter-turn locking with slow slip \\
v3.8.31 & local OP~5 quarter-cycle packet state & 2026-04-04 & showed that one full quarter-turn cycle of the super-mean layer, i.e. a \(16\)-zero resonance packet, already exhibits strong cancellation: the first two packets leave only about \(2.21\%\) and \(3.99\%\) of their \(\ell^1\)-mass, so the mesoscopic OP~5 burden is best read in terms of packet drift and inter-packet residue rather than isolated block errors \\
v3.8.32 & local OP~5 packet-defect state & 2026-04-04 & derived the first-order packet-defect law, showing that once quarter-cycle cancellation is taken as the zeroth-order picture, the surviving OP~5 residue is carried only by slow packet amplitude drift and quarter-turn phase-slip, not by the ideal packet itself \\
v3.8.33 & local OP~5 packet-envelope state & 2026-04-04 & compressed the first-order packet-defect variables into a single scalar packet-envelope \(\mathfrak D_{N,m}\), proving that the packet residue is controlled by \(|R_{N,m}(T_{\mathrm{loc}})|\le \sqrt2\,\mathfrak D_{N,m}+O(\delta\varepsilon+\varepsilon^2)\); thus OP~5 is now a packet-envelope problem at first order \\
v3.8.34 & local OP~5 observable-envelope state & 2026-04-04 & introduced the observable lower packet-envelope proxy \(\underline{\mathfrak D}_{N,m}=|R_{N,m}(T_{\mathrm{loc}})|/\sqrt2\), giving benchmark values \(\underline{\mathfrak D}_{50,1}\approx 0.00273964\), \(\underline{\mathfrak D}_{50,2}\approx 0.00434351\), and \(\underline{\mathfrak D}_{50,3}\approx 0.00921439\); this moves Step~(B) onto a directly observable higher-scale scalar defect sequence \\
\hline
v3.8.50 & local merge-reconciliation checkpoint & 2026-04-04 & reconciled the former v3.5.2 abstract-polish checkpoint with the current v3.8.49 boundary-layer line; declared the 3.5.2 interface-normalization results fully absorbed into the active 3.8 working monolith and recorded the merge explicitly at the document head and in the main text \\
\hline
v3.8.51 & local reproducible-figure state & 2026-04-04 & regenerated the missing Julia/probe comparison figure from a local Python renderer, embedded the resulting PDF asset as the primary document figure with PNG fallback, and recorded the rebuild path directly in the appendix discussion \\
\hline
v3.8.52 & local current-scale figure state & 2026-04-04 & replaced the earlier standalone comparison marker by the current residual-scale diagnostic family $\gamma\in\{0,\Delta_\beta,\Delta_\theta,\Delta\mu\}$ and updated the appendix figure/caption accordingly \\
\hline
v3.8.53 & local caption/spacing fix state & 2026-04-04 & repaired caption spacing for the current residual-scale probe figure and stabilized the appendix-page layout for the diagnostic render \\
\hline
v3.8.55 & local deep-boundary diagnostic extension & 2026-04-04 & extended the Appendix-C diagnostics with a focused $\Delta\mu$ deep-zoom comparison, localized difference ladders, and edge-skeleton drift overlays to test whether the strongest current residual remains boundary-local under further magnification \\
v3.8.56 & local continuation diagnostic extension & 2026-04-04 & added a six-step continuation test from $0$ to $\Delta\mu$, including a boundary-belt strip, cumulative motion envelope, and horizon-profile continuation to check whether the active shell deforms continuously across the whole live residual interval \\
 v3.8.57 & local sensitivity-localization extension & 2026-04-04 & added finite-difference sensitivity maps, regional/radial sensitivity profiles, and a horizon-angle heatmap to show that the live residual response stays concentrated on the same horizon-adjacent boundary shell rather than spreading into the core \\
\hline
v3.8.58 & local horizon-sector persistence extension & 2026-04-04 & added coarse angular-sector persistence maps and dominant-peak tracking on the M\"obius horizon, showing that the active response remains locked to a small conjugate sector family instead of diffusing around the full circle \\
v3.8.59 & local OP~5 packet-wave diagnostics extension & 2026-04-05 & added visual diagnostics for the quarter-turn super-mean cycle, packet-level cancellation hierarchy, signed packet-defect wave fits, and the rotating packet-mode phase portrait; these figures make the current OP~5 reading visible as packet resonance plus slow signed wave rather than as an unstructured tail \\
v3.8.60 & local OP~5 breathing-envelope visual extension & 2026-04-05 & added visual diagnostics for the transverse envelope trough, the slow secondary envelope wave, the breathing-helix interpretation of the packet mode, and the carrier/envelope ratio test; these figures make the current OP~5 reading visible as an axially dominated carrier with a slower breathing boundary-layer correction rather than as a rigid fixed-radius helix \\
v3.8.61 & local OP~5 frame-flicker visual integration & 2026-04-05 & added live Appendix-C frame-flicker fit, decomposition, and geometric-readout figures and moved the new visual block into the actually compiled PDF path so that the moving-frame hypothesis becomes directly auditable in the active monolith \\
v3.9.0 & Zenodo release & 2026-04-05 & \href{https://doi.org/10.5281/zenodo.19423689}{10.5281/zenodo.19423689} --- public release preserving the stabilized residual-record interface, the reproducible M\"obius-shell diagnostic suite, and the OP~5 packet / breathing-envelope / frame-flicker visual line in one monolithic text \\
\hline
v3.10.0 & Zenodo release & 2026-04-06 & \href{https://doi.org/10.5281/zenodo.19434163}{10.5281/zenodo.19434163} --- public release based on the v3.9.0 monolith, with version and DOI record updated for the present release and front-matter release language cleaned for Zenodo \\
\hline
v3.10.1 & local working continuation & 2026-04-06 & front-cleaned working continuation of the public v3.10.0 release; historical roadmap and release-note blocks moved behind the nomenclature into an archival appendix position \\
\hline
v3.10.2 & local working tool-extraction state & 2026-04-06 & first explicit working tool layer for v3.11 preparation: transport-first, readout-after-transport, checkpoint-absorption, and epsilon-trap tools recorded together with a synchronized diagnosis of the live OP~1 and OP~5 burdens \\
\hline
v3.10.3 & local working OP~1 epsilon-bridge state & 2026-04-06 & reframed the live OP~1 burden as epsilon-trapped bridge normalization: canonical closure is now read as entry into a contractive bridge tube around the distinguished fixed point rather than as a second branch search \\
\hline
v3.10.4 & local working OP~1 theorem-promotion state & 2026-04-06 & promoted the OP~1 epsilon-bridge tool to theorem level by proving a residual-to-fixed-point estimate and a contractive non-escape theorem for the canonical bridge normalization inside the local bridge tube \\
v3.10.5 & local working OP~1 rounded-entry state & 2026-04-06 & proved that the standard rounded matter-frame readout already enters the local contractive bridge tube via the tiny readout-point displacement $|\gamma_L-\gamma_{\mathrm{br}}|\approx 3.2\times 10^{-8}$, reducing the live OP~1 burden to the structural identification $b_{\mathrm{can}}=b_{\mathrm{MF}}$ rather than to a second local bridge search \\
v3.10.6 & local working OP~1 uniqueness promotion state & 2026-04-06 & promoted admissible-readout uniqueness to theorem level, removing the final structural ambiguity between canonical normalization and the standard matter-frame readout and reducing OP~1 to a single explicit local bridge residue \\
v3.10.7 & local working shared-matrix diagnostic state & 2026-04-06 & added shared HC/MM/J leakage diagnostics and a structure-first translation of linear, differential, and linear-differential cryptanalysis for attacking the common OP~1/OP~5 matrix remainder \\
v3.10.8 & local working shared-phase coupling diagnostic state & 2026-04-06 & added a shared-carrier / quadrature diagnostic for testing whether the remaining OP~1 and OP~5 residues are one in-phase mode or two phase-shifted readouts of the same hidden leakage carrier \\
v3.10.9 & local working shared-leakage probe-family state & 2026-04-06 & replaced the abstract phase-coupling parameter by a concrete HC/MM/J cross-sector probe family $\mathcal M_\lambda=\mathcal M+\lambda\Delta_{\mathrm{shared}}$ and isolated its first-order effect on the visible OP~1/OP~5 channels \\
v3.10.10 & local working probe-invariant classification state & 2026-04-06 & introduced the normalized determinant and dot-product invariants of the HC/MM/J probe response matrix, turning the shared-carrier versus quadrature question into a decision-grade first-order diagnostic \\
v3.10.11 & local working canonical low-rank leakage state & 2026-04-06 & fixed a canonical rank-two HC/MM/J leakage class and proved that its first-order response is necessarily one-mode rather than quadrature, so any later sinus/cosinus picture would have to be emergent rather than primitive \\
v3.10.12 & local working bidirectional leakage-plane state & 2026-04-06 & added the minimal paired rank-four cross-sector leakage plane and proved that it is the first structurally honest location where a quadrature/sinus-cosinus picture can arise, whereas primitive rank-two leakage remains one-mode only \\
  v3.10.13 & local working visible-ellipse state & 2026-04-06 & coupled the minimal paired leakage plane to the OP~1/OP~5 mask layer and proved that the visible response is generically an ellipse, collapsing to a line in rank one and to exact circle/quadrature only under conformal mask geometry on the leakage plane \\
  v3.10.14 & local working shared-mask decomposition state & 2026-04-06 & decomposed the visible OP~1/OP~5 mask layer into a shared channel plus separation defect, proved singular-value bounds in the near-shared regime, and showed that exact quadrature cannot be primitive when the two visible masks differ only by a small downstream correction \\
  v3.10.17 & local working observer-split control state & 2026-04-06 & showed that the visible OP~1/OP~5 separation defect is exactly the restriction of the observer-side mask difference to the active leakage plane, proving that near-shared one-mode behavior is the default whenever the two visible masks differ only by a small downstream readout split \\
  v3.10.18 & local working shared-residual dichotomy state & 2026-04-06 & decomposed the remaining primitive OP~1/OP~5 visible residual into a shared matrix-leakage carrier plus a downstream observer-shell split and proved that exact grouped no-drift forces any surviving primitive residual on the active leakage plane to be common-mode rather than primitive quadrature \\
  v3.10.19 & local working active matrix-fit state & 2026-04-06 & pushed the remaining primitive common-mode carrier back to the active HC/MM/J fit and proved that its size is controlled exactly by the active commutator $[\mathcal M,J]|_{U}$ on the admissible branch \\
  v3.10.20 & local working exact-transport dichotomy state & 2026-04-06 & used the grouped-layer admissible-transport architecture to show that exact grouped Millennium transport kills the active commutator on the admissible branch, yielding a clean dichotomy between exact primitive closure and a purely downstream observer-shell residual \\
  v3.10.22 & local working canonical observer-shell residual state & 2026-04-06 & promoted the remaining visible OP~1/OP~5 term, when primitive matrix leakage is absent, to a canonical downstream observer-shell residual carried by a common shell dressing mode rather than by a primitive HC/MM/J leakage mode \\
  v3.10.23 & local working canonical shell-law modeling state & 2026-04-06 & modeled the surviving shared OP~1/OP~5 shell mode after exact transport as a first-order phase-shifted / expansion-dressed observer-shell law, reducing the visible target to exact vanishing versus canonical post-transport shell identification \\
  v3.10.24 & local working shell-dominance classification state & 2026-04-06 & compressed the remaining visible post-transport question to a comparison of first-order shell loads, classifying the shared shell mode as phase-dominant, expansion-dominant, or genuinely mixed at first order \\
  v3.10.25 & local working shell-cancellation corridor state & 2026-04-06 & reduced the shared shell mode to a signed first-order shell-balance mismatch, isolating the cancellation corridor for exact visible closure and showing that same-sign shell loading forces a canonical nonzero shell remainder unless the linear shell load itself vanishes \\
  v3.10.26 & local working normalized shell-defect state & 2026-04-06 & compressed the shared observer-shell problem to a single normalized shell-cancellation defect, extracted an exact cancellation-depth invariant, and proved that exact visible closure requires the normalized first-order shell mismatch to be no larger than the normalized higher-order shell remainder \\
  v3.10.28 & local working sign-corridor shell state & 2026-04-06 & converted the normalized shell verdict into a sharp sign/corridor dichotomy, proved that sign-coherent first-order shell loads force canonical nonzero visible residual whenever the normalized higher-order rescue budget is below one, and extracted the exact anti-phase corridor required for possible visible closure \\
  v3.10.29 & local working shell-conjugacy gain state & 2026-04-06 & pushed the induced phase and expansion gains back to a shell-conjugate observer geometry, introduced the downstream shell involution and transport-balance defect, and showed that the anti-phase corridor is structurally natural whenever phase and expansion are J-descended conjugate readouts of one shell displacement with nearly balanced transport weights \\
\hline
v3.10.30 & local working shell anti-phase budget state & 2026-04-06 & compressed the anti-phase corridor to a single normalized shell-conjugacy load and showed that a small total shell load forces the induced shell ratio into a natural anti-phase tube around $-1$ \\
\hline
v3.10.31 & local working architecture-driven shell state & 2026-04-06 & pushed the shell-balance driver back to a combined architecture budget given by active HC/MM/J fit plus observer-shell dressing, so that first-order visible cancellation is now read as architecture-driven rather than phenomenological \\
\hline
v3.10.32 & local working architecture-rescue corridor state & 2026-04-06 & reduced the remaining visible question to a single budget comparison between the normalized combined architecture load and the normalized higher-order rescue budget, giving a final first-order corridor for exact visible closure after exact transport \\
\hline
v3.11.0 & Zenodo release & 2026-04-06 & \href{https://doi.org/10.5281/zenodo.19444128}{10.5281/zenodo.19444128} --- public release consolidating the v3.10.32 architecture-rescue corridor state, explicit tool-layer extraction, and release-normalized front matter for Zenodo \\
v3.11.2 & local working state & 2026-04-07 & OP~1/OP~5 closed as a native source-return pair via closure-class preservation; branchwise refinement, admissible routed aggregation, and the final visible-support mismatch decomposition added as the compact main-line hardening beyond the public v3.11.0 release  \\
 v3.12.0 & Zenodo release & 2026-04-07 & \href{https://doi.org/10.5281/zenodo.19446264}{10.5281/zenodo.19446264} --- public release promoting the OP~1/OP~5 closure line into the public-facing main text, with branchwise refinement and routed-readout consequences integrated into the monolithic Companion \\
\hline
v0.1.0 & Zenodo release & 2026-03-24 & \href{https://doi.org/10.5281/zenodo.19210729}{10.5281/zenodo.19210729} \\
\hline
v0.3.0 & Zenodo release & 2026-03-24 & \href{https://doi.org/10.5281/zenodo.19211785}{10.5281/zenodo.19211785} \\
\hline
v0.4.0 & Zenodo release & 2026-03-24 & \href{https://doi.org/10.5281/zenodo.19211798}{10.5281/zenodo.19211798} \\
\hline
v0.5.0 & Zenodo release & 2026-03-24 & \href{https://doi.org/10.5281/zenodo.19211877}{10.5281/zenodo.19211877} \\
\hline
v0.6.0 & Zenodo release & 2026-03-24 & \href{https://doi.org/10.5281/zenodo.19212073}{10.5281/zenodo.19212073} \\
\hline
v0.7.0 & Zenodo release & 2026-03-24 & \href{https://doi.org/10.5281/zenodo.19212212}{10.5281/zenodo.19212212} \\
\hline
v0.8.0 & Zenodo release & 2026-03-24 & \href{https://doi.org/10.5281/zenodo.19212324}{10.5281/zenodo.19212324} \\
\hline
v0.9.0 & Zenodo release & 2026-03-24 & \href{https://doi.org/10.5281/zenodo.19220003}{10.5281/zenodo.19220003} \\
\hline
v1.0.0 & Zenodo release & 2026-03-24 & \href{https://doi.org/10.5281/zenodo.19220118}{10.5281/zenodo.19220118} \\
\hline
v1.6.0 & Zenodo release & 2026-03-25 & \href{https://doi.org/10.5281/zenodo.19221335}{10.5281/zenodo.19221335} \\
\hline
v1.7.0 & Zenodo release & 2026-03-25 & \href{https://doi.org/10.5281/zenodo.19221372}{10.5281/zenodo.19221372} \\
\hline
v1.8.0 & Zenodo release & 2026-03-25 & \href{https://doi.org/10.5281/zenodo.19221431}{10.5281/zenodo.19221431} \\
\hline
v1.9.0 & Zenodo release & 2026-03-25 & \href{https://doi.org/10.5281/zenodo.19221473}{10.5281/zenodo.19221473} \\
\hline
v2.0.0 & Zenodo release & 2026-03-25 & \href{https://doi.org/10.5281/zenodo.19221527}{10.5281/zenodo.19221527} \\
\hline
v2.1.0 & Zenodo release & 2026-03-25 & \href{https://doi.org/10.5281/zenodo.19221618}{10.5281/zenodo.19221618} \\
\hline
v2.3.0 & Zenodo release & 2026-03-25 & \href{https://doi.org/10.5281/zenodo.19221864}{10.5281/zenodo.19221864} \\
\hline
v2.4.0 & Zenodo release & 2026-03-25 & \href{https://doi.org/10.5281/zenodo.19221954}{10.5281/zenodo.19221954} \\
\hline
v2.9.0 & Zenodo release & 2026-03-25 & \href{https://doi.org/10.5281/zenodo.19222638}{10.5281/zenodo.19222638} \\
\hline
v2.10.0 & Zenodo release & 2026-03-25 & \href{https://doi.org/10.5281/zenodo.19222894}{10.5281/zenodo.19222894} \\
\hline
v2.11.0 & Zenodo release & 2026-03-25 & \href{https://doi.org/10.5281/zenodo.19223084}{10.5281/zenodo.19223084} \\
\hline
v2.11.1 & local-only repair state & 2026-03-25 & non-Zenodo version-fix state between v2.11.0 and v2.11.2 \\
\hline
v2.11.2 & Zenodo release & 2026-03-25 & \href{https://doi.org/10.5281/zenodo.19223365}{10.5281/zenodo.19223365} \\
\hline
v2.11.3 & Zenodo release & 2026-03-25 & \href{https://doi.org/10.5281/zenodo.19223515}{10.5281/zenodo.19223515} \\
\hline
v2.12.0 & Zenodo release & 2026-03-25 & \href{https://doi.org/10.5281/zenodo.19223782}{10.5281/zenodo.19223782} \\
\hline
v2.13.0 & Zenodo release & 2026-03-25 & \href{https://doi.org/10.5281/zenodo.19223968}{10.5281/zenodo.19223968} \\
\hline
v2.14.0 & Zenodo release & 2026-03-25 & \href{https://doi.org/10.5281/zenodo.19225191}{10.5281/zenodo.19225191} \\
v2.15.0 & Zenodo release & 2026-03-25 & \href{https://doi.org/10.5281/zenodo.19226408}{10.5281/zenodo.19226408} \\
\hline
v2.16.0 & Zenodo release & 2026-03-25 & \href{https://doi.org/10.5281/zenodo.19226445}{10.5281/zenodo.19226445} \\
\hline
v2.17.0 & Zenodo release & 2026-03-25 & \href{https://doi.org/10.5281/zenodo.19226497}{10.5281/zenodo.19226497} \\
v2.18.0 & Zenodo release & 2026-03-25 & \href{https://doi.org/10.5281/zenodo.19226549}{10.5281/zenodo.19226549} \\
\hline
v2.19.0 & Zenodo release & 2026-03-25 & \href{https://doi.org/10.5281/zenodo.19226773}{10.5281/zenodo.19226773} \\
\hline
v2.20.0 & Zenodo release & 2026-03-26 & \href{https://doi.org/10.5281/zenodo.19227889}{10.5281/zenodo.19227889} \\
\hline
v2.21.0 & Zenodo release & 2026-03-26 & \href{https://doi.org/10.5281/zenodo.19227940}{10.5281/zenodo.19227940} \\
\hline
v2.22.0 & Zenodo release & 2026-03-26 & \href{https://doi.org/10.5281/zenodo.19228061}{10.5281/zenodo.19228061} \\
\hline
v2.23.0 & Zenodo release & 2026-03-26 & \href{https://doi.org/10.5281/zenodo.19239228}{10.5281/zenodo.19239228} \\
\hline
v2.23.1 & Zenodo patch release & 2026-03-26 & \href{https://doi.org/10.5281/zenodo.19239967}{10.5281/zenodo.19239967} \\
\hline
v2.24.0 & Zenodo release & 2026-03-27 & \href{https://doi.org/10.5281/zenodo.19258774}{10.5281/zenodo.19258774} \\
v2.25.0 & Zenodo release & 2026-03-27 & \href{https://doi.org/10.5281/zenodo.19259425}{10.5281/zenodo.19259425} \\
 v2.25.1 & Zenodo patch release & 2026-03-27 & \href{https://doi.org/10.5281/zenodo.19262890}{10.5281/zenodo.19262890} \\
\hline
v2.27.0 & Zenodo release & 2026-03-27 & \href{https://doi.org/10.5281/zenodo.19266666}{10.5281/zenodo.19266666} \\
\hline
v2.27.1 & local-only repair state & 2026-03-28 & baseline-preserving post-release working state \\
\hline
v2.27.2 & local-only repair state & 2026-03-28 & normalization and proof-chain hardening \\
\hline
v2.27.3 & local-only repair state & 2026-03-28 & structural closure pass; last state before v2.27.4 integration \\
\hline
v2.27.4 & local working state & 2026-03-28 & spectral-projector layer ($C=\Pi^{(\mathcal{M})}_{\mathrm{spec}}$), readout-relative lemma methodology, Theorem~7 (HC-induced readout geometry), Hypotheses~7.4--7.6, bidirectional probe lensing, SM spectral realization, regime-test appendix \\
\hline
v2.27.5 & Zenodo release & 2026-03-28 & \href{https://doi.org/10.5281/zenodo.19297422}{10.5281/zenodo.19297422} — full proof-chain audit; closure roadmap; 35 missing labels; five Priority-A gaps identified \\
\hline
v2.27.8 & local working state & 2026-03-28 & B1+B2 closed: \texttt{prop:semantic-forcing-chain} explicit 3-part proof from axioms (S1)--(S8); \newline \texttt{prop:structural-lift-principle} non-circular via \texttt{lem:group-transport} + \texttt{lem:hc-admissibility-lemma}; \newline OP~2 formalized; all 7 Roadmap items closed \\
\hline
v2.27.7 & local working state & 2026-03-28 & A4+A5 closed: \texttt{thm:koide-triolic} with $\theta_K$ formula; \newline 0.111\% match and $115\times$ amplification remark; \newline \texttt{thm:bsym-hopf} with proof and Hopf-gap remark; \newline 4.9\%/$6.8\%$, OP1/OP2 coupling; open: B1, B2 \\
\hline
v2.28.0 & Zenodo release & 2026-03-28 & \href{https://doi.org/10.5281/zenodo.19298586}{10.5281/zenodo.19298586}; consolidation of v2.27.4--v2.27.8; spectral-projector unconditional; bijection lemma; SM realization; readout-relative methodology; Thm~7; Hyp~7.4--7.6; OP~2 stub \\
\hline
v2.29.0 & Zenodo release & 2026-03-28 & \href{https://doi.org/10.5281/zenodo.19299245}{10.5281/zenodo.19299245} --- OP~5 step~(A) closed; step~(B) quantified; renormalization explicit \\
\hline
v2.29.1 & local working state & 2026-03-28 & GOE/GUE split (\texttt{prop:goe-fermionic}); $\gL$ = GOE/GUE ratio; $\gamma_L=1.197$ corrected to $1.864$; $\chi^2$ improvement $75.8\%$; $\delta_{\mathrm{CL}}$ exponent $\gamma_L^2\to\gamma_L^5$; \texttt{thm:dual-origin} corrected \\
\hline
v2.29.2 & Zenodo release & 2026-03-28 & \href{https://doi.org/10.5281/zenodo.19300351}{10.5281/zenodo.19300351} --- Einstein blunder reframed; $\Omega_M=\bsym^2$ determination; three-gap convergence \\
\hline
v2.29.3 & Zenodo release & 2026-03-28 & \href{https://doi.org/10.5281/zenodo.19300813}{10.5281/zenodo.19300813} --- OP~2 (T1)(T2) closed; $D_-=L_0|_{\Hplusminus}$ (corrected candidate); GOE essential self-adjointness; T4=OP1 equivalence; \textbf{Tribonacci theorem} $\gamma_L\to\mathbb{T}$; cascade normalization $1/\mathbb{T}+1/\mathbb{T}^2+1/\mathbb{T}^3=1$ \\
\hline
\end{longtable}
}

\noindent Note. The sequence above records the currently preserved release trail.
The day-level publication dates are normalized from the project publication log so that
Zenodo releases and local repair states remain chronologically comparable inside a single
appendix table.


\subsection*{v2.26.0 computational consolidation note}

Version~\textbf{2.26.0} executes the consolidation roadmap that had been reserved in
v2.25.1. The computational strand is retained but internally reorganized into a more
coherent progression: historical perspective, structural machine reading, Triolen
Algebra / Boolean readout, MML--MMA--MCP stratification, proof templates,
certificate-bearing execution, normalization, addressing, scheduling, host export,
and extraction outlook. The present version still remains exploratory in the sense of
future formalization, but the material is no longer presented only as a loose teaser;
it is now collected as a first internally coherent computational chapter of the
Companion.



\subsection*{Hypothesis-control pass for v2.27.0}

Before the next public milestone, the mathematical layer should distinguish more sharply between statements that are expected to become theorem-grade in v2.27.0 and statements that should remain explicitly hypothesis-grade even after the next hardening cycle. The practical rule is conservative: only statements with stabilized definitions, closed dependency chains, and reference-stable support should move upward from hypothesis status into theorem, proposition, or lemma status. Statements that continue to depend on interpretive transport, physical reading, or still-open structural identifications should remain explicitly marked as hypotheses or outlook items.

For the next pass, the intended control logic is therefore blockwise. In the ground-condition / HC / confinement chain, the target is theorem-readiness whenever the dependence on the collapse picture and confinement mechanism can be formulated without appealing to deferred computational scaffolding. In projector / closure mathematics, the target is theorem or lemma status whenever the relevant admissibility and sector-preservation claims are already internal to the mathematical layer. In the transport / Lorentz layer, statements that are still partly interpretive should remain hypothesis-grade unless they can be rewritten as strictly mathematical transport statements. For the Millennium-matrix and trace/branch layers, the default stance remains more cautious: structural mathematical claims may harden, but computational or machine-reading interpretations remain secondary and should not silently inherit theorem status from the mathematical core.

The resulting editorial discipline for v2.27.0 is simple: theorem promotion requires closure, dependency control, and stable notation; everything else remains visibly hypothesis, deferred claim, or teaser architecture. This avoids theorem-to-hypothesis drift in reverse as well: once a statement is kept at hypothesis level, later references should not read it as already mathematically closed merely because it is computationally suggestive.

\subsection*{v2.26.33 local mathematics-first note}

Version~\textbf{2.27.0} does not yet claim a new public release beyond v2.26.0. Instead, it records the explicit shift of near-term priority back to mathematical completion. The computational chapter remains preserved as a secondary but already visible reserve architecture. The immediate next target is v2.27.0 as a mathematics-first hardening cycle: close the mathematical layer, collect tools as mathematical appendices, and only then resume the proof-engine line as a compression language for already stabilized structure.


\subsection*{Theorem-promotion checklist for v2.27.0}
The next public milestone should not promote statements merely because they have become familiar in the internal workflow. Each promotion from hypothesis-language into lemma-, proposition-, or theorem-language must therefore be explicit, local, and reference-stable. For the mathematical closure pass leading to v2.27.0, the promotion discipline is:
\begin{itemize}[leftmargin=*]
  \item statements about the ground condition to HC to confinement chain may be promoted only when the chain has been written in one direction without hidden interpretive jumps;
  \item projector/closure claims may be promoted only when the admissibility conditions and the closure domain are fixed in stable notation;
  \item transport/Lorentz statements may be promoted only when they are phrased as mathematical transport claims rather than as physical metaphor;
  \item Millennium-matrix statements may be promoted only when the matrix is treated first as a mathematical object with domain, codomain, and structural role, not merely as a computational image;
  \item trace/branch claims may be promoted only when the branch criteria and trace-compatibility conditions are stated in theorem-capable form.
\end{itemize}
In parallel, several statement families are to remain explicitly outside theorem-promotion for v2.27.0 unless a full closure pass is actually completed. These include broad computational-universality claims, strong proof-kernel claims, speculative host-language claims, and any physically interpretive reading that still rests on heuristic transport language rather than closed mathematical formulation.

A practical promotion checklist for the next pass is therefore:
\begin{enumerate}[leftmargin=*]
  \item identify candidate statements that are currently written as heuristic summary sentences;
  \item decide whether each candidate should become a Definition, Lemma, Proposition, or remain a Hypothesis;
  \item verify that every promoted statement has stable notation, explicit assumptions, and stable backward references;
  \item verify that no promoted statement depends silently on teaser/computational language;
  \item record all non-promoted but still important claims explicitly as open, deferred, or interpretive.
\end{enumerate}
This checklist is intended to keep the next public release mathematically harder without forcing premature promotion of material that should still remain hypothesis-level.

\subsection*{Mathematical freeze criterion for v2.27.0}
For the next public release, a mathematical core block should count as \emph{frozen for release} only when it is not merely well phrased, but also stable under the local closure discipline introduced in the present subversion sequence. In particular, a block is to be treated as mathematically release-ready for v2.27.0 only if all of the following are satisfied:
\begin{enumerate}[leftmargin=*]
  \item the principal objects of the block are defined in stable notation and do not drift between theorem-language and hypothesis-language;
  \item the block depends only on explicitly cited prior definitions, lemmas, propositions, or tools, rather than on narrative intuition or roadmap placeholders;
  \item the key admissibility conditions are written inside the statement or proof scaffold itself;
  \item all required cross-references are stable and do not point to provisional teaser/computational material;
  \item the block can be read both locally and in the dependency order recorded for the mathematics-first pass without hidden backward jumps;
  \item any remaining interpretive or physical reading is clearly separated from the mathematical closure claim.
\end{enumerate}
A block should \emph{not} be counted as frozen merely because its wording is polished, because it already appears in summary form elsewhere, or because it has an attractive structural interpretation. For v2.27.0, the intended release criterion is therefore conservative: closure requires definitional stability, dependency stability, reference stability, and status stability at the same time.

A practical freeze test for each A-priority block is:
\begin{itemize}[leftmargin=*]
  \item can the block be pointed to by a stable label and theorem-status marker?
  \item can its assumptions be listed in one compact place?
  \item can its proof or justification be read without consulting deferred computational scaffolds?
  \item can its use in later blocks be stated without introducing new notation drift?
\end{itemize}
If the answer to any of these questions remains negative, the block should remain marked as active work rather than as mathematically frozen for public release. This criterion is meant to prevent premature hardening of material that is already suggestive but not yet sufficiently closed.


\subsection*{Release-readiness matrix for v2.27.0}
To prevent the next public release from being driven by impression rather than closure, the mathematics-first pass should end in a compact readiness matrix for the A-priority blocks. The matrix is not meant to replace proofs; it is meant to summarize, in one release-facing place, whether a core block is genuinely ready to be frozen, or whether a single missing condition still blocks public hardening.

\begin{center}
{\tiny\setlength{\tabcolsep}{2.5pt}\begin{tabular}{@{}p{0.18\textwidth} p{0.14\textwidth} p{0.10\textwidth} p{0.23\textwidth} p{0.16\textwidth}@{}}
\toprule
\textbf{Core block} & \textbf{Current state} & \textbf{Freeze-ready?} & \textbf{Missing closure condition} & \textbf{Primary blocker for v2.27.0}\\
\midrule
Ground condition / collapse picture & structurally central, still partly summary-driven & no & one stable definitional pass with fixed assumptions & notation and statement drift\\
HC emergence / confinement chain & strongly anchored, not yet written in one compact theorem-capable line & partial & explicit one-directional closure chain & hidden interpretive compression\\
Projector / closure mathematics & partially stabilized and reusable & partial & admissibility conditions and closure domain fixed in one notation regime & local definition spread\\
Transport / Lorentz layer & present but mixed with interpretive language & no & transport claims restated as mathematical transport statements & metaphor drift\\
Millennium matrix as mathematical object & increasingly structured, still mixed with computational imagery & no & fixed domain/codomain/operator role & object-status not yet fully hardened\\
Trace / branch mathematics & structurally rich, still partly template-style & partial & stable branch criterion and trace compatibility form & theorem-status not yet fixed\\
Tool / appendix lemmas & abundant support layer & partial & priority ordering and stable dependence tags & appendix support not yet fully normalized\\
\bottomrule
\end{tabular}}
\end{center}

For the next public milestone, the conservative reading of this matrix is intentional: \emph{partial} does not yet count as release-ready. A block should only be treated as ready for v2.27.0 if its missing closure condition has been discharged explicitly and its blocker is no longer active. In this sense, the matrix is a release-control device rather than a rhetorical summary: it translates the mathematics-first agenda into a blockwise freeze decision surface.



\subsection*{v2.27.0 gate checklist}
To prevent the next public release from being triggered by momentum alone, the mathematics-first pass should culminate in a compact global gate checklist. This checklist is stricter than the blockwise readiness matrix: even if several A-priority blocks look individually advanced, the public jump to v2.27.0 should only be made once the release conditions hold \emph{together} at the document level.

\begin{itemize}[leftmargin=*]
  \item the A-priority mathematical blocks admit stable definitions, stable dependence order, and stable theorem-status markers;
  \item no public-facing core claim in the targeted v2.27.0 range still depends on deferred computational scaffolds for its mathematical reading;
  \item the blockwise freeze criteria and the release-readiness matrix no longer contain an active blocker in the A-priority line;
  \item the proof-dependency checklist and the reference-stability pass agree on a single release-safe label structure;
  \item hypothesis-control and theorem-promotion discipline are aligned, so that no statement is promoted merely because its interpretation is attractive;
  \item the transport / projector / Millennium-matrix line can be cited without notation drift or mixed theorem-versus-hypothesis status;
  \item appendix tools used by the A-priority line are stable enough to serve as support, but not so overloaded that they silently replace missing core proofs;
  \item the computational / MML / MMA / MCP chapter remains visibly present as teaser architecture, but no release-critical mathematical closure claim depends on it;
  \item the version record, roadmap, local subversion discipline, and closure-control blocks all point to the same public target without internal contradiction.
\end{itemize}

A practical release rule for the next public jump is therefore conservative: \emph{v2.27.0 should only be cut once the A-line is simultaneously blockwise closure-ready and globally gate-clean}. If one of these conditions still fails, the correct response is another local subversion rather than a premature public release. In this sense, the gate checklist converts the mathematics-first roadmap into a final version-control criterion.

\subsection*{Final pre-v2.27.0 control summary}
As the mathematics-first pass approaches a public release decision, the Companion should end with a compact control summary that can be read in one glance. The purpose of this summary is not to replace the detailed work matrix, target-state matrix, or release-readiness table, but to compress them into a final release-facing snapshot. In particular, the summary should answer five questions simultaneously: (i) which A-priority blocks are currently the decisive release drivers, (ii) whether each of them is already freeze-ready, (iii) what the leading blocker still is, (iv) whether the global gate checklist is currently satisfied, and (v) whether the recommended action is \emph{continue mathematical closure} or \emph{prepare public release}.

\begin{center}
\small\begin{tabular}{p{0.18\textwidth} p{0.11\textwidth} p{0.18\textwidth} p{0.14\textwidth} p{0.27\textwidth}}
\toprule
\textbf{A-priority block} & \textbf{Freeze-ready?} & \textbf{Main blocker} & \textbf{Gate status impact} & \textbf{Release recommendation}\\
\midrule
Ground condition / collapse picture & not yet & dependence and formulation hardening & blocks full gate closure & continue mathematical closure\\
HC emergence / confinement chain & partial & explicit mathematical discharge & blocks full gate closure & continue mathematical closure\\
Projector / closure mathematics & partial & stable theorem-level closure & blocks full gate closure & continue mathematical closure\\
Transport / Lorentz layer & conditional & transport formalization discipline & weakens gate confidence & continue mathematical closure\\
Millennium matrix as mathematical object & not yet & object-level hardening before interpretive use & blocks full gate closure & continue mathematical closure\\
Trace / branch mathematics & partial & stable dependence on prior blocks & blocks full gate closure & continue mathematical closure\\
\bottomrule
\end{tabular}
\end{center}

At the present stage, the control summary therefore recommends \emph{no public jump yet beyond the current release line}: the mathematics-first agenda should remain active until the decisive A-priority blocks no longer fail the freeze criterion and the global gate checklist turns from ``not yet'' to ``satisfied.'' The intended role of this final control summary is practical rather than rhetorical: it gives the next release decision a single consolidated dashboard and makes it harder to confuse local progress with document-level readiness.


\subsection*{One-page executive closure view}
The mathematics-first steering logic introduced in the recent local subversions is now dense enough that it should also be compressible into a single executive page. The purpose of this closure view is therefore not to add another planning layer, but to place roadmap, dependency map, target-state matrix, freeze criterion, gate checklist, and pre-release control summary onto one final compact surface. In practical terms, the page should allow a reader to answer four questions immediately: (i) what the next public version is really trying to close mathematically, (ii) which A-priority blocks still drive the release boundary, (iii) whether the document is globally gate-clean, and (iv) whether the correct next action is another local subversion or a public release jump.

\begin{center}
\begin{tabular}{p{0.22\textwidth} p{0.22\textwidth} p{0.18\textwidth} p{0.24\textwidth}}
\toprule
\textbf{Control layer} & \textbf{Main question} & \textbf{Current status} & \textbf{Operational consequence}\\
\midrule
Roadmap / work matrix & what must be closed first? & mathematics-first line active & keep proof-engine / MML line secondary\\
Dependency map & what must come before what? & ordered core chain available & no closure claims out of dependency order\\
Definition / notation pass & are the objects and labels stable? & partially stabilized & unify notation before theorem-promotion\\
Target-state / closure checklist & what counts as done? & blockwise closure conditions visible & only explicit closure counts as completion\\
Freeze criterion / readiness matrix & is a block release-ready? & mixed, mostly partial & keep A-line blocks active until blocker-free\\
Gate checklist / control summary & is the document globally ready? & not yet gate-clean & continue local subversions, no public jump yet\\
\bottomrule
\end{tabular}
\end{center}

In this condensed reading, the present Companion state remains productive but intentionally conservative: the mathematical architecture is now sufficiently organized to support disciplined closure work, but not yet sufficiently gate-clean to justify confusing local progress with public readiness. The executive closure view is therefore meant to function as a one-page release dashboard. Its decision rule is simple: as long as the A-priority mathematical blocks remain only partially frozen, and as long as the global gate still records an active blocker, the correct action is \emph{another explicit local subversion rather than a public version jump}. Once these conditions disappear simultaneously, the dashboard should flip from \emph{continue closure} to \emph{prepare the next public release}.


\subsection*{Release-vs-local decision rule}
A final practical distinction should also be made explicit before the next public jump: not every coherent local gain justifies a new Zenodo-facing version. The present Companion workflow therefore distinguishes between \emph{local closure iterations} and \emph{public release transitions}. A local subversion is appropriate whenever the document gains planning discipline, dependency clarity, notation stability, or blockwise closure information without yet removing the main A-line blockers. A public release transition is appropriate only when these local gains have accumulated into a document-wide state in which the mathematical core is not merely better organized, but sufficiently frozen, referentially stable, and gate-clean to support a new public boundary.

In operational terms, the decision rule may be summarized as follows. Move only by local subversion if one or more of the following remains true: (i) an A-priority block is still only partially closed, (ii) the theorem/hypothesis boundary is still under active normalization, (iii) the reference-stability pass is not yet quiet, (iv) the release-readiness matrix still lists an active primary blocker, or (v) the global gate checklist is not yet satisfied. Consider a public version jump only when all A-priority blocks satisfy their freeze criterion simultaneously, the dependency chain is stable under rereading, the hypothesis-control and theorem-promotion passes no longer produce status drift, and the release-control summary flips from \emph{continue closure} to \emph{prepare release}. This distinction is intentionally conservative: it protects the mathematical layer from being mistaken for release-ready merely because the planning architecture around it has become more refined.

\subsection*{Mathematics-first release doctrine}
The present local stage adopts a simple doctrine for all remaining pre-v2.27.0 work. First, every local change is to be judged by whether it closes, stabilizes, or clarifies a mathematical dependency in the A-priority chain; additions that merely decorate the proof-engine teaser without tightening the mathematics are therefore secondary by rule. Second, no new public release is justified merely by chapter growth or architectural enthusiasm: the release threshold remains mathematical closure, reference stability, and theorem-safe status control. Third, the computational/MML/MMA/MCP layer is to remain visible and preserved, but it is not permitted to become a hidden prerequisite for mathematical claims that should stand on their own. Fourth, every local subversion should remain append-only and auditable, so that recovery from drift or accidental rollback stays possible without reconstructing the chain from memory.

In operational terms, the doctrine reads as follows:
\begin{enumerate}[label=(\roman*)]
  \item prioritize closure of ground condition $\rightarrow$ HC $\rightarrow$ confinement before expanding secondary architecture;
  \item promote statements only when notation, dependencies, and reference anchors are already stable;
  \item treat tools, appendices, and computational reserves as support layers unless and until they are explicitly absorbed into the mathematical core;
  \item prefer local subversions over public version jumps whenever the work mainly sharpens internal control rather than closing a release-grade mathematical block.
\end{enumerate}

This doctrine is intentionally conservative. Its role is to keep the Companion aligned with the structure-first programme: mathematics first, proof-memory preserved, computational architecture retained as a prepared secondary line, and public releases tied to genuine closure rather than mere accumulation.


\subsection*{Mathematics-first action queue}
The next local working order is to be treated as explicit and sequential unless a later control note states otherwise. The active queue is: (1) close the ground-condition / collapse picture at the definition level; (2) harden the HC-emergence / confinement chain so that its dependency status is stable; (3) close the projector / closure layer with stable admissibility language; (4) stabilize the transport / Lorentz layer only after the closure prerequisites are fixed; (5) treat the Millennium matrix strictly as a mathematical object before any renewed computational extension; (6) move trace / branch mathematics only after the preceding transport layer is stable; and (7) keep appendix tools and the computational teaser chapter in support position unless a direct mathematical blocker requires them.

The negative rule is equally important. The next local subversions are not to prioritize new MML/MMA/MCP growth, new proof-engine macros, or broader computational rhetoric unless such additions are directly needed to preserve an already existing mathematical dependency. In this sense, the queue functions as an anti-drift device: it converts the mathematics-first doctrine into a concrete order of execution and makes explicit what should not be pulled forward in parallel.


\paragraph{Mathematics-first session protocol.}
To prevent local growth from drifting back into attractive but premature architecture work, the next internal work sessions should follow a short fixed protocol. Each session should begin by choosing exactly one A-priority mathematical block from the current action queue. The session should then record (i) the active block, (ii) the minimal closure target for that block, (iii) the exact dependencies that must already be stable, and (iv) the specific theorem/lemma/definition/tool boundary that is being touched. Only after this boundary is explicit should any local rewriting begin.

The execution rule for such sessions is intentionally narrow. First perform a definition pass for the chosen block; then a dependency pass; then a closure pass; then a reference-stability check; and only at the end, if necessary, a supporting appendix-tool pass. If a prospective edit does not directly improve the active mathematical block or its immediate dependencies, it should be deferred. In particular, new computational language material, proof-engine rhetoric, or architecture expansion should not be introduced during a session unless the active mathematical block explicitly depends on it.

At the end of each local session, the result should be classified in one of four ways: \emph{closed for now}, \emph{partially hardened}, \emph{dependency blocked}, or \emph{still open}. This classification should then feed back into the release-vs-local decision rule, so that local subversions remain true mathematical work logs rather than merely growing source snapshots. In this way, the present protocol operationalizes the mathematics-first doctrine at the level of day-to-day document evolution.



\subsection*{Anti-parallelization rule}
The present mathematics-first stage also requires an explicit rule against parallel drift. In the next local subversions, not more than one A-priority mathematical block should be treated as the primary closure target of a given pass. Supporting edits in neighboring blocks are allowed only insofar as they remove an immediate dependency or stabilize a reference chain for that active block. What should be avoided is simultaneous advancement of several major fronts at once---for example projector/closure mathematics, transport/Lorentz normalization, and Millennium-matrix formalization in the same pass---because this makes it too easy to confuse distributed motion with genuine closure.

The same rule applies across the mathematics/computational boundary. While the computational teaser chapter remains preserved, it should not be opened in parallel with a mathematical A-block unless a concrete blocker in that block explicitly requires it. In practical terms, the anti-parallelization rule means: one primary closure front per pass; all secondary edits must justify themselves as dependency support, reference stabilization, or notation repair for that front. This rule is intentionally restrictive, because the present problem is not lack of ideas but risk of dispersion. Its purpose is to convert the mathematics-first doctrine into a narrow execution discipline and thereby make local progress more legible, auditable, and release-relevant.



\subsection{Re-entry rule after each local mathematical session}
After each local mathematics-first session, the document should be re-entered through a short global control loop rather than by immediately opening a new frontier. The purpose of this rule is to prevent hidden drift, duplicated partial closure, or silent reclassification of still-open material. In practical terms, the re-entry pass should answer only four questions: (i) did the session close the targeted A-block for now, (ii) which dependencies were clarified or newly exposed, (iii) which references and labels were changed and therefore need stabilization, and (iv) whether the next session should continue the same block or move to the next dependency in the queue.

The re-entry rule therefore imposes the following order: first update the local block status (
\emph{closed for now}, \emph{partially hardened}, \emph{dependency blocked}, or \emph{still open}); second reconcile the result with the target-state matrix and release-readiness matrix; third update the proof-dependency checklist if new requirements were uncovered; and only then decide whether the next local subversion should continue the same mathematical block or step forward in the mathematics-first action queue. This rule is intentionally conservative: local progress is only counted as durable progress when it is reinserted into the global control architecture without breaking reference stability, hypothesis control, or the release-vs-local decision rule.



\subsection*{Session-close audit rule}
Each local mathematics-first session should end with a short audit pass before the next subversion is declared. This audit is intended to prevent apparent progress from being recorded merely at the level of added text. A local pass counts as a meaningful increment only if the targeted block has been re-checked against its definition status, dependency status, reference stability, and release-readiness classification.

The closing audit for a local session should therefore answer four questions explicitly: (i) was a definition actually stabilized, rather than only rephrased; (ii) was at least one dependency clarified, reduced, or made more explicit; (iii) were the relevant references and labels left in a more stable state than before; and (iv) did the session produce a genuine closure gain for the active block, rather than only a planning gain. If these questions cannot be answered positively, the session may still be useful, but it should be recorded as preparatory rather than as a closure increment.

In this sense, the session-close audit forms the final control layer of the mathematics-first protocol. It ensures that the subversion chain remains readable as a chain of actual closure work, not merely as a sequence of appended managerial notes.


\subsection*{Carry-forward rule between local mathematical sessions}
Not every item touched in one local mathematics-first session is allowed to be carried forward as if it were already stable. The present carry-forward rule therefore distinguishes between (i) content that is eligible to be inherited by the next local subversion as a working baseline, and (ii) content that must be carried forward only as still-open material. This distinction is intended to prevent partial rewrites, provisional notation choices, or merely suggestive closure claims from silently hardening through repetition alone.

A block element may be carried forward as part of the active baseline only if all of the following are already true: its local definition is stable enough not to be rewritten immediately in the next pass; its dependency position is unchanged or clarified; its reference anchors survive rereading without ambiguity; and its status classification does not rely on a teaser, soft interpretation, or unresolved supporting note. By contrast, material should be carried forward only as \emph{open}, \emph{dependency blocked}, or \emph{provisional} whenever one of these conditions still fails. In practical terms, this means that unresolved notation, unstable lemma wording, open dependency edges, or uncertain theorem-promotion status must remain visibly marked rather than being treated as already absorbed into the mathematical baseline.

The purpose of this rule is conservative continuity. It allows the next local subversion to inherit genuine closure gains, while forcing incomplete work to remain visibly incomplete. In this way, the carry-forward rule complements the session-close audit: the audit decides whether a session produced a real gain, and the carry-forward rule decides which parts of that gain are allowed to function as stable starting points for the next mathematics-first pass.


\subsection*{Regression-check rule for local mathematical continuity}
Because the present Companion has already experienced version drift and accidental rollback in earlier phases, the mathematics-first regime also requires an explicit regression-check rule. The purpose of this rule is not to slow down productive local work, but to ensure that no apparent gain in planning, notation, or closure discipline is purchased at the cost of silently losing already integrated mathematical material. In particular, each local subversion should be checked not only for what has been added, but also for whether any previously stabilized block, appendix tool, roadmap control layer, or restored support material has shrunk, disappeared, or reverted in status.

Operationally, the regression check is conservative and document-internal. After each local session, the current source should be reread against the immediately preceding local subversion with at least four questions in mind: (i) has the active A-priority block gained real closure content without losing prior dependencies, (ii) have restored support structures such as the table of contents, version chronology, or appendix anchors remained intact, (iii) have any local control blocks been duplicated, displaced, or pushed outside the effective document body, and (iv) has any internal version or bibliographic marker fallen behind the actual subversion state. If one of these checks fails, the correct response is not to treat the current file as progressive merely because it is newer, but to repair the regression before carrying the pass forward.

In this sense, the regression-check rule complements the carry-forward rule and the session-close audit. The audit asks whether the current session produced a real local gain; the carry-forward rule asks what may legitimately be inherited as stable baseline material; and the regression-check rule asks whether that inheritance still contains everything that was already secured before. The combined effect is intentionally protective: it turns local subversions into auditable closure steps rather than a chain of potentially drifting snapshots. A local mathematics-first version should therefore be regarded as acceptable only if it is both \emph{additively improved} and \emph{non-regressive} relative to the immediately preceding baseline.

\subsection*{Baseline-lock rule for local continuation}
To prevent the mathematics-first subversion chain from drifting across multiple informal starting points, the present Companion adopts an explicit baseline-lock rule. For each local pass, exactly one immediately preceding local version is treated as the binding comparison basis. All claims of progress, non-regression, and closure gain are measured against this locked baseline rather than against a vague memory of the recent editing sequence. The purpose of the rule is simple: a local file becomes meaningful as a closure step only if it can be shown to improve upon one clearly identified predecessor without sacrificing already secured material.

Operationally, a baseline may be replaced only under conservative conditions. A newly produced local subversion may become the next locked baseline if, and only if, it satisfies the regression check, preserves the restored support structures, carries forward only legitimately stable material, and moves the active mathematics block measurably closer to its stated closure target. If one of these conditions fails, the correct response is not to redefine the newer file as the new norm, but to keep the earlier baseline locked until the defect has been repaired. In particular, mere growth in page count, roadmap language, or appended managerial notes is not sufficient to justify baseline replacement.

This rule complements the anti-parallelization, re-entry, session-close, carry-forward, and regression rules by providing a stable anchor for comparison. The anti-parallelization rule limits what may be opened at once; the re-entry rule controls how local gains are fed back into the larger roadmap; the session-close audit judges whether a session produced a real gain; the carry-forward rule determines what may be inherited as stable; the regression check guards against silent loss; and the baseline-lock rule finally determines when the current subversion has become stable enough to function as the new working ground for the next mathematics-first step.


\subsection*{Change-scope rule for local mathematics-first sessions}
The mathematics-first regime also requires an explicit rule about \emph{scope}. In any given local session, not everything that appears editable should actually be edited. The active scope is restricted to the currently selected A-priority mathematical block, its immediate dependency edges, the notation and references required to stabilize that block, and only those appendix tools that are strictly necessary to support the same closure pass. This means that many tempting modifications---however reasonable in isolation---should remain out of scope if they do not directly improve the currently active closure target.

The practical purpose of the change-scope rule is to protect mathematical closure from diffuse accretion. A local session may therefore adjust wording, dependencies, notation, labels, or supporting appendices only when these changes can be justified as direct support for the active block. What remains explicitly out of scope includes broad computational expansion, speculative theorem promotion outside the active dependency cone, large stylistic rewrites of unrelated sections, and opportunistic edits that merely make the document longer without bringing the selected block closer to a frozen mathematical state. In this way, the change-scope rule complements the anti-parallelization rule: anti-parallelization limits how many fronts may be open at once, while scope control limits how wide a single front is allowed to spread.





\subsection*{2H. Global pre-release reconciliation for v2.27.0}
With the first A-block now locally frozen through Lemma~C--G and Proposition~H--J, the remaining question is no longer whether the dependency spine exists, but whether this local gain has been fed back consistently into the document-wide release logic. The present pre-release reconciliation therefore records the first A-block not merely as locally strengthened, but as positively reconciled with the mathematics-first gate architecture for the next public release. In particular, the block is now to be read as structurally closed enough that its remaining limitation is no longer missing architecture, but only moderate proof density and future optional unfolding.

Operationally this means three things. First, the block may already count as a positive contribution to the global v2.27.0 gate, because it now satisfies the internal chain ground condition $	o$ HC $	o$ closure $	o$ projector stabilization $	o$ admissible transport $	o$ matrix realization $	o$ trace/branch persistence. Second, no computational teaser layer is needed in order to maintain its current mathematical status. Third, any later expansion of this block should be treated as proof-density enrichment rather than as a repair of a still-broken logical spine. The correct release reading is therefore conservative but positive: the first A-block is globally reconciled with the current mathematics-first release doctrine and should no longer be counted among the primary blockers for v2.27.0.

\paragraph{Global release consequence.}
For the next release decision, the first A-block is to be classified as \emph{mathematically release-supporting}: not yet maximally unfolded, but no longer structurally open. The remaining pre-release work should therefore shift away from re-opening this block and toward checking whether the remaining A-priority material can be brought into the same state without introducing theorem drift, reference instability, or dependence on the computational teaser architecture.

\paragraph{Local closure target.}
After this reconciliation pass, the active goal for the local subversion chain is no longer to add more managerial logic around the first A-block, but to preserve it as a frozen support block while the remaining release blockers are tested against the same gate discipline.



\subsection*{C.14 Fifteenth 3.4 extension: transition cocycle support}
\addcontentsline{toc}{subsection}{C.14 Fifteenth 3.4 extension: transition cocycle support}

\begin{theorem}[Appendix support for Theorem~\ref{thm:mainline-3415}]
\label{thm:appendix-3415}
Let $\mathfrak U=\{U_\alpha\}$ be the stable shell atlas and suppose each isolated obstruction shell carries one of the four local transition types from Theorem~\ref{thm:mainline-3414}. Then these local transitions assemble into a locally constant $V_4$-valued cocycle on admissible overlaps. After adjoining a sheet label that resolves fold exchange, the induced cocycle factors through the quotient $V_4\twoheadrightarrow \mathbb Z_2^{\mathrm{par}}$.
\end{theorem}

\begin{proof}
Every admissible overlap arises from continuation across finitely many isolated obstruction shells. By Theorem~\ref{thm:mainline-3414}, each such shell contributes one of the four involutive transition types. Their compositions therefore stay inside $V_4$, because the fold and parity generators commute and square to the identity. This gives a locally constant overlap label.

On triple overlaps, composing continuation maps in two stages or one stage leads to the same inverse shell identification, because both are obtained by following the same continuation class through the same finite obstruction data. Hence the cocycle relation holds. Resolving fold exchange by adjoining the local inverse-sheet label kills the fold generator on the lifted data, leaving only the parity component, which is exactly the quotient to $\mathbb Z_2^{\mathrm{par}}$.
\end{proof}

\paragraph{3.4.15 working interpretation.}
The inverse-readout atlas is therefore not just stratified; it is glued by a tiny and completely explicit transition law. Before fold resolution, the global bookkeeping lives in a Klein-four transition cocycle. After fold resolution, the only surviving topological memory is the parity lift. This is the most economical global form of the spinorial/M"obius mechanism obtained so far.




\subsection*{3.4 Sixteenth extension: transition cocycle and parity-resolved local system}
\addcontentsline{toc}{subsection}{3.4 Sixteenth extension: transition cocycle and parity-resolved local system}
The local obstruction law of Section~3.4.14 gives the primitive transition moves, but the next structural step is to organize those moves over the whole shell atlas rather than shell by shell. The natural object here is a tiny transition cocycle. It records how inverse shell data are re-identified when one passes between adjacent stable charts and then separates the fold ambiguity from the genuine parity lift.

\begin{theorem}[3.4 transition cocycle and parity-resolved local system]
\label{thm:mainline-3415}
Let $\mathfrak U=\{U_\alpha\}$ be the shell atlas from Theorem~\ref{thm:mainline-3413}, where each $U_\alpha$ is a connected component of $I\setminus\Sigma$ carrying stable inverse readout. Then the following hold.
\begin{enumerate}[label=(\roman*)]
  \item For every nonempty admissible overlap obtained by continuation across a single isolated obstruction shell, there is a locally constant transition element
  \[
    g_{\alpha\beta}\in V_4:=\{\mathrm{id},\,T^{\mathrm{fold}},\,T^{\mathrm{par}},\,T^{\mathrm{par}}\circ T^{\mathrm{fold}}\}
  \]
  such that inverse shell data satisfy
  \[
    (s,\pi)_\beta = g_{\alpha\beta}\,(s,\pi)_\alpha.
  \]
  \item On every admissible triple overlap one has the cocycle relation
  \[
    g_{\alpha\gamma}=g_{\beta\gamma}\circ g_{\alpha\beta},
  \]
  so the inverse-readout layer defines a $V_4$-valued local transition system over the shell atlas.
  \item After passing to the fold-resolved two-sheet inverse cover, the fold generator is absorbed into the choice of sheet, and the remaining transition system reduces to a parity local system valued in
  \[
    \mathbb Z_2^{\mathrm{par}}=\{\mathrm{id},T^{\mathrm{par}}\}.
  \]
\end{enumerate}
\end{theorem}

\begin{proof}
Statement (i) is exactly the content of Theorem~\ref{thm:mainline-3414}: every isolated obstruction contributes only fold exchange, parity flip, or their composition. Since those operations act discretely and the obstruction type cannot vary inside a sufficiently small continuation corridor, the corresponding transition element is locally constant on the overlap.

For (ii), the continuation from $U_\alpha$ to $U_\gamma$ through $U_\beta$ is represented by successive application of the two local transition maps. Because the local transitions are honest involutive re-identifications of the same inverse shell data, their composition agrees with direct continuation, which is the cocycle relation.

For (iii), resolving the fold ambiguity means that one fixes one of the two local inverse sheets as part of the enlarged state. The fold involution then acts trivially on the lifted data, because it has been absorbed into the sheet label itself. What remains is only whether the transported parity class is preserved or flipped, hence a residual $\mathbb Z_2$-valued local system.
\end{proof}

\begin{corollary}[3.4 parity monodromy class]
\label{cor:mainline-3415}
Under the hypotheses of Theorem~\ref{thm:mainline-3415}, every closed admissible continuation of inverse shell data determines a well-defined class in the parity local system $\mathbb Z_2^{\mathrm{par}}$. The visible direct-return case corresponds to the trivial class, while the spinorial/M"obius continuation is exactly the nontrivial parity class.
\end{corollary}

\begin{proof}
By Theorem~\ref{thm:mainline-3415}(iii), all fold ambiguity is removed on the lifted inverse cover, so a closed continuation can leave only the trivial or nontrivial parity action. The trivial action returns the same visible branch, whereas the nontrivial parity action returns only after the $\mathbb Z_2$ lift. This is precisely the spinorial/M"obius case already identified locally in Corollary~\ref{cor:mainline-3414}.
\end{proof}

\paragraph{Interpretive consequence.}
The shell atlas now carries a genuine global readout organization. Stable charts are not merely patched together ad hoc: they are glued by a tiny $V_4$-valued transition cocycle, and after resolving the fold choice the remaining global obstruction is purely parity-theoretic. This is the cleanest form so far of the claim that the visible spinorial/M"obius branch is not decorative language, but the nontrivial global parity sector of the inverse-readout geometry.




\subsection*{3.4 Seventeenth extension: parity double cover and lift-triviality criterion}
\addcontentsline{toc}{subsection}{3.4 Seventeenth extension: parity double cover and lift-triviality criterion}
The parity-resolved local system of Section~3.4.15 already identifies the spinorial/M"obius sector as a genuine global class. The next structural step is to say exactly when that class disappears and, when it does not, what the canonical repair is. The natural answer is a parity double cover: it trivializes precisely the residual $\mathbb Z_2$ obstruction and therefore converts the parity-resolved inverse readout into an honest global section.

\begin{theorem}[3.4 parity double cover and lift-triviality criterion]
\label{thm:mainline-3416}
Let $B$ be a connected admissible shell base carrying the parity local system $\mathcal L_{\mathrm{par}}$ from Theorem~\ref{thm:mainline-3415}. Then the following are equivalent.
\begin{enumerate}[label=(\roman*)]
  \item The parity local system $\mathcal L_{\mathrm{par}}$ is globally trivial on $B$.
  \item There exists a globally continuous parity-resolved branch choice
  \[
    \sigma:B\to\{\pm 1\}
  \]
  such that, after fold resolution, inverse shell data admit a globally single-valued visible representative
  \[
    (s,\sigma\sqrt{s},\pi)
  \]
  on $B$.
  \item Every closed admissible continuation in $B$ has trivial parity monodromy.
\end{enumerate}
If these equivalent conditions fail, then there exists a canonical connected double cover
\[
  p:\widetilde B\to B
\]
trivializing $\mathcal L_{\mathrm{par}}$. On $\widetilde B$, the parity-resolved inverse readout becomes globally single-valued, and every admissible closed continuation lifts to direct visible convergence.
\end{theorem}

\begin{proof}
The equivalence of (i) and (iii) is the standard characterization of a $\mathbb Z_2$ local system: triviality is exactly the vanishing of its monodromy on all closed admissible loops. If $\mathcal L_{\mathrm{par}}$ is trivial, one may choose a global parity frame, which is precisely the global sign selection in (ii). Conversely, any such global sign choice identifies the two parity fibers coherently over all of $B$, so the local system is trivial.

If the parity class is nontrivial, the associated principal $\mathbb Z_2$-bundle defines the canonical connected double cover $p:\widetilde B\to B$. Pulling back $\mathcal L_{\mathrm{par}}$ to $\widetilde B$ trivializes it by construction. Hence the parity-resolved inverse shell data admit a global branch choice upstairs, and every lifted closed continuation returns with trivial parity action, i.e. as direct visible convergence on the cover.
\end{proof}

\begin{corollary}[3.4 direct-return versus lift criterion]
\label{cor:mainline-3416}
On a connected admissible shell base, the direct visible-return regime is globally available if and only if the parity local system is trivial. The spinorial/M"obius regime is therefore not an extra local ambiguity, but exactly the nontrivial global parity class; it disappears on the canonical parity double cover.
\end{corollary}

\begin{proof}
The direct visible-return regime is exactly the existence of a globally coherent parity choice after fold resolution, which is condition~(ii) of Theorem~\ref{thm:mainline-3416}. The second claim is the nontrivial case of the same theorem.
\end{proof}

\paragraph{Interpretive consequence.}
The inverse-readout geometry now has a clean global dichotomy. Either the parity local system is already trivial, in which case the visible branch closes globally without further enlargement, or the residual obstruction is genuinely $\mathbb Z_2$ and the correct global object is the canonical parity double cover. In that sense, the spinorial/M"obius sector is no longer only a language for local continuation failure; it is exactly the statement that the visible readout trivializes only after passing to the natural two-sheeted cover.


\subsection*{3.4 Eighteenth extension: cohomological parity obstruction class}
\addcontentsline{toc}{subsection}{3.4 Eighteenth extension: cohomological parity obstruction class}
The parity double cover already gives the correct global repair. The next and cleanest reformulation is to record the remaining parity sector as a genuine cohomological obstruction. This does not add a new mechanism; it simply packages the already extracted $\mathbb Z_2$ data in the natural global invariant attached to the shell atlas.

\begin{theorem}[3.4 cohomological parity obstruction class]
\label{thm:mainline-3417}
Let $B$ be a connected admissible shell base equipped with the parity-resolved inverse-readout atlas from Sections~3.4.15--3.4.16. After fold resolution, let
\[
  g^{\mathrm{par}}_{\alpha\beta}\in \mathbb Z_2
\]
be the parity transition functions on an admissible shell atlas $\{U_\alpha\}$. Then:
\begin{enumerate}[label=(\roman*)]
  \item the family $\{g^{\mathrm{par}}_{\alpha\beta}\}$ is a \v{C}ech $1$-cocycle with coefficients in $\mathbb Z_2$;
  \item its cohomology class
  \item one has
  \[
  \begin{aligned}
    &\varepsilon_{\mathrm{par}}=0
    \quad\Longleftrightarrow\quad
    \text{the parity local system is trivial}\\[0.2em]
    &\quad\Longleftrightarrow\quad
    \text{direct visible convergence is globally available};
  \end{aligned}
  \]
  \item if $\varepsilon_{\mathrm{par}}\neq 0$, then the canonical parity double cover from Theorem~\ref{thm:mainline-3416} is the two-sheeted cover classified by $\varepsilon_{\mathrm{par}}$.
\end{enumerate}
\end{theorem}

\begin{proof}
Theorem~\ref{thm:mainline-3415} already reduces the shell-atlas transition data to the parity sector after fold resolution. Hence the remaining transition functions take values in the constant group $\mathbb Z_2$, and the cocycle relation on triple overlaps is exactly the cocycle identity for a \v{C}ech $1$-cochain. Changing the admissible atlas or parity frame modifies $g^{\mathrm{par}}_{\alpha\beta}$ by a \v{C}ech coboundary, so the cohomology class $\varepsilon_{\mathrm{par}}$ is well defined.

The equivalence in (iii) is the standard identification of a $\mathbb Z_2$ local system with its class in $H^1(B,\mathbb Z_2)$: vanishing means that the cocycle is cohomologous to zero, hence a global parity frame exists; this is exactly the direct visible-return regime from Theorem~\ref{thm:mainline-3416}. If the class does not vanish, the same theorem provides the canonical parity double cover trivializing the local system, and that cover is precisely the one classified by $\varepsilon_{\mathrm{par}}$.
\end{proof}

\begin{corollary}[3.4 parity obstruction as the global lift defect]
\label{cor:mainline-3417}
The spinorial/M\"obius regime is completely encoded by the single class
\[
  \varepsilon_{\mathrm{par}}\in H^1(B,\mathbb Z_2).
\]
It vanishes exactly in the direct visible-return regime and is nonzero exactly when the inverse readout trivializes only after passage to the canonical parity double cover.
\end{corollary}

\begin{proof}
This is the reformulation of Theorem~\ref{thm:mainline-3417}(iii)--(iv) in terms of the unique cohomological obstruction carried by the parity local system.
\end{proof}

\paragraph{Interpretive consequence.}
The parity side of the 3.4 program is now fully compressed. Local parity flips, shell-wise monodromy, the residual $\mathbb Z_2$ local system, the canonical double cover, and the distinction between direct visible closure and spinorial/M\"obius lift are no longer separate stories. They are all expressions of the single class
\[
  \varepsilon_{\mathrm{par}}\in H^1(B,\mathbb Z_2).
\]
In that sense, the remaining lift defect is now a genuine cohomological obstruction rather than merely a descriptive continuation phenomenon.


\subsection*{3.4 Nineteenth extension: parity holonomy character and loop-readout criterion}
\addcontentsline{toc}{subsection}{3.4 Nineteenth extension: parity holonomy character and loop-readout criterion}
The cohomological obstruction class from Section~3.4.17 is already the smallest global invariant. The next step is to make it directly readable along admissible closed shell paths. The natural object is the parity holonomy character: it evaluates the class on a loop and therefore tells us, in the smallest possible way, whether visible closure occurs already on the base or only after passage to the parity double cover.

\begin{theorem}[3.4 parity holonomy character and loop-readout criterion]
\label{thm:mainline-3418}
Let $B$ be a connected admissible shell base with parity obstruction class
\[
  \varepsilon_{\mathrm{par}}\in H^1(B,\mathbb Z_2)
\]
from Theorem~\ref{thm:mainline-3417}. Fix a base point $b_0\in B$. Then:
\begin{enumerate}[label=(\roman*)]
  \item there is a canonical parity holonomy character
  \[
    \chi_{\mathrm{par}}:\pi_1(B,b_0)\to \mathbb Z_2,
    \qquad
    \chi_{\mathrm{par}}([\gamma])=\langle \varepsilon_{\mathrm{par}},[\gamma]_{2}\rangle,
  \]
  where $[\gamma]_2\in H_1(B,\mathbb Z_2)$ denotes the mod-two homology class of $\gamma$;
  \item $\chi_{\mathrm{par}}([\gamma])=0$ if and only if every parity-resolved inverse continuation along $\gamma$ returns with direct visible closure on $B$;
  \item $\chi_{\mathrm{par}}([\gamma])=1$ if and only if continuation along $\gamma$ closes only after the nontrivial deck transformation of the canonical parity double cover from Theorem~\ref{thm:mainline-3416};
  \item the character factors through mod-two homology,
  \[
    \pi_1(B,b_0)\twoheadrightarrow H_1(B,\mathbb Z_2)\xrightarrow{\ \langle \varepsilon_{\mathrm{par}},-\rangle\ }\mathbb Z_2,
  \]
  so two admissible loops with the same class in $H_1(B,\mathbb Z_2)$ have the same parity readout.
\end{enumerate}
\end{theorem}

\begin{proof}
A $\mathbb Z_2$ local system determines a monodromy representation into $\mathbb Z_2$, and for the parity local system this is exactly the parity action picked up by admissible closed continuation. By Theorem~\ref{thm:mainline-3417}, the local system is classified by the cohomology class $\varepsilon_{\mathrm{par}}$. Standard evaluation of $H^1(B,\mathbb Z_2)$ on mod-two $1$-cycles gives the homological formula
\[
  \chi_{\mathrm{par}}([\gamma])=\langle \varepsilon_{\mathrm{par}},[\gamma]_2\rangle.
\]
The value $0$ means that the parity transport around $\gamma$ is trivial, which is exactly direct visible return on the base. The value $1$ means that the transported branch returns with the nontrivial parity action, so visible closure occurs only after passing to the parity double cover. Because the target group is abelian, the character factors through the abelianization of $\pi_1(B,b_0)$, and with mod-two coefficients this is precisely the factorization through $H_1(B,\mathbb Z_2)$.
\end{proof}

\begin{corollary}[3.4 loop-equivalence principle for parity readout]
\label{cor:mainline-3418}
On a connected admissible shell base, the global direct-return versus spinorial/M\"obius decision for any admissible closed continuation class depends only on the value of
\[
  \langle \varepsilon_{\mathrm{par}},[\gamma]_2\rangle\in\mathbb Z_2.
\]
In particular, homologous admissible loops mod $2$ are parity-indistinguishable for the inverse readout.
\end{corollary}

\begin{proof}
This is exactly statement~(iv) of Theorem~\ref{thm:mainline-3418}: the parity readout factors through the mod-two homology class of the loop.
\end{proof}

\paragraph{Interpretive consequence.}
The parity obstruction is now not only a global class but an operational readout. A closed admissible shell path carries a tiny $\mathbb Z_2$ outcome obtained by evaluating $\varepsilon_{\mathrm{par}}$ on its mod-two class. The visible direct-return regime corresponds to holonomy value $0$, while the spinorial/M\"obius regime corresponds to holonomy value $1$. This is the first fully global and directly readable form of the parity defect.



\subsection*{3.4 Twentieth extension: empirical parity-holonomy estimator and loop calibration}
\addcontentsline{toc}{subsection}{3.4 Twentieth extension: empirical parity-holonomy estimator and loop calibration}
The parity holonomy character from Section~3.4.18 is already the smallest global obstruction. The next step is to make it measurable from the readout side itself. Since the shell atlas already resolves admissible continuations into stable chart arcs separated by obstruction-shell crossings, the only possible source of a nontrivial parity outcome is the parity-changing part of those crossings. This yields a direct empirical estimator for the global $\mathbb Z_2$ holonomy.

\begin{theorem}[3.4 empirical parity-holonomy estimator]
\label{thm:mainline-3419}
Let $\gamma$ be an admissible closed shell path on a connected shell base $B$, decomposed into stable chart arcs whose endpoints meet obstruction shells transversely. For each encountered obstruction-shell crossing, let
\[
  \delta_k\in\mathbb Z_2
\]
be defined by
\[
  \delta_k=
  \begin{cases}
    0,&\text{for a pure fold transition},\\
    1,&\text{for a parity transition},\\
    1,&\text{for a mixed fold-plus-parity transition},
  \end{cases}
\]
so that $\delta_k$ records only the parity-changing component of the local transition law from Theorem~\ref{thm:mainline-3414}. Define the empirical parity estimator
\[
  \widehat\chi_{\mathrm{par}}(\gamma)
  :=
  \sum_{k=1}^{N_\gamma}\delta_k \pmod 2.
\]
Then:
\begin{enumerate}[label=(\roman*)]
  \item $\widehat\chi_{\mathrm{par}}(\gamma)$ depends only on the parity labels of the stable chart arcs and is independent of the chosen admissible chart decomposition;
  \item the empirical estimator agrees with the cohomological holonomy character,
  \[
    \widehat\chi_{\mathrm{par}}(\gamma)
    =
    \chi_{\mathrm{par}}([\gamma])
    =
    \langle \varepsilon_{\mathrm{par}},[\gamma]_2\rangle;
  \]
  \item consequently, $\widehat\chi_{\mathrm{par}}(\gamma)=0$ if and only if the calibrated inverse readout closes directly on the base, while $\widehat\chi_{\mathrm{par}}(\gamma)=1$ if and only if closure occurs only after the nontrivial deck action on the parity double cover.
\end{enumerate}
\end{theorem}

\begin{proof}
By Theorem~\ref{thm:mainline-3414}, only the parity component of an obstruction-shell crossing contributes to the post-crossing branch label; pure folds only exchange inverse sheets and therefore do not affect the residual $\mathbb Z_2$ parity class. Summing the parity-changing contributions modulo $2$ along the closed path therefore computes the total parity transport accumulated by the continuation. Refining the chart decomposition inserts pairs of complementary local transitions whose parity contributions cancel modulo $2$, so the quantity $\widehat\chi_{\mathrm{par}}(\gamma)$ is well defined.

After fold resolution, the remaining transition data is exactly the parity local system from Theorem~\ref{thm:mainline-3415}. Hence the mod-two sum of parity changes along $\gamma$ is the monodromy of that local system around the loop, which by Theorem~\ref{thm:mainline-3418} is precisely $\chi_{\mathrm{par}}([\gamma])=\langle \varepsilon_{\mathrm{par}},[\gamma]_2\rangle$. The final visible-return statements are the two cases of that holonomy value.
\end{proof}

\begin{corollary}[3.4 detector-side loop test for the parity obstruction]
\label{cor:mainline-3419}
Suppose a calibrated inverse readout assigns stable chart arcs and parity labels to a detector-side closed shell trace. Then the parity obstruction can be tested directly by the mod-two count of parity-changing crossings:
\[
  \widehat\chi_{\mathrm{par}}(\gamma)=0
  \Longleftrightarrow
  \text{detector-side direct visible return},
\]
\[
  \widehat\chi_{\mathrm{par}}(\gamma)=1
  \Longleftrightarrow
  \text{detector-side spinorial/M\"obius lift signal}.
\]
In particular, the global $\mathbb Z_2$ lift defect becomes empirically checkable once the shell trace has been chartwise calibrated.
\end{corollary}

\begin{proof}
This is exactly statement~(iii) of Theorem~\ref{thm:mainline-3419}, read through the calibrated inverse readout from Sections~3.4.11--3.4.13.
\end{proof}

\paragraph{Interpretive consequence.}
The parity side now closes the loop between geometry and readout. The cohomological class is no longer only a classification object and the holonomy character is no longer only an abstract loop functional. Once a detector-side or probe-side shell trace has been calibrated chartwise, the global lift defect is measured by a single mod-two crossing count. This is the sharpest operational form so far of the statement that the spinorial/M\"obius regime is a real readout phenomenon rather than a decorative analogy.


\subsection*{C.15 Sixteenth 3.4 extension: parity double-cover support}
\addcontentsline{toc}{subsection}{C.15 Sixteenth 3.4 extension: parity double-cover support}

\begin{theorem}[Appendix support for Theorem~\ref{thm:mainline-3416}]
\label{thm:appendix-3416}
Let $B$ be a connected admissible shell base and let $\mathcal L_{\mathrm{par}}$ be the parity local system obtained after fold resolution. Then the parity monodromy determines a unique connected two-sheeted cover $p:\widetilde B\to B$ up to isomorphism. Pulling back $\mathcal L_{\mathrm{par}}$ along $p$ trivializes the parity transport, and any other connected cover with the same property factors through $\widetilde B$.
\end{theorem}

\begin{proof}
The parity local system is a locally constant $\mathbb Z_2$-bundle over $B$. Its monodromy representation sends every admissible closed continuation class to either the trivial or nontrivial parity action. The associated kernel subgroup therefore defines a connected index-two cover $p:\widetilde B\to B$, unique up to canonical isomorphism. By construction, loops lifted to $\widetilde B$ lie in the kernel of the parity monodromy, so the pulled-back local system has trivial transport. If another connected cover trivializes the same monodromy, its deck group must also kill that kernel and therefore factor through the index-two quotient defining $\widetilde B$.
\end{proof}

\paragraph{3.4.16 working interpretation.}
The parity sector has now reached its cleanest possible global form. The shell atlas first produced local obstruction types, then a Klein-four transition cocycle, then a reduced parity local system, and now finally the canonical object that removes the remaining obstruction: the parity double cover. Once lifted to that cover, the inverse-readout branch is globally single-valued again. This is the sharpest mathematical version so far of the statement that the visible spinorial/M"obius branch is a double-cover phenomenon rather than a decorative metaphor.


\subsection*{C.16 Seventeenth 3.4 extension: cohomological obstruction support}
\addcontentsline{toc}{subsection}{C.16 Seventeenth 3.4 extension: cohomological obstruction support}

\begin{theorem}[Appendix support for Theorem~\ref{thm:mainline-3417}]
\label{thm:appendix-3417}
Let $\{U_\alpha\}$ be an admissible shell atlas on $B$, and let $g^{\mathrm{par}}_{\alpha\beta}\in\mathbb Z_2$ be the parity transition cocycle after fold resolution. Then the cohomology class
\[
  \varepsilon_{\mathrm{par}}=[g^{\mathrm{par}}_{\alpha\beta}]\in H^1(B,\mathbb Z_2)
\]
classifies the parity local system and its associated connected two-sheeted cover up to isomorphism.
\end{theorem}

\begin{proof}
Because the post-fold transition functions take values in the constant group $\mathbb Z_2$, they define a \v{C}ech $1$-cocycle. Two admissible atlases related by refinement or change of parity frame produce cohomologous cocycles, so the resulting class in $H^1(B,\mathbb Z_2)$ is intrinsic. Standard classification of principal $\mathbb Z_2$-bundles over $B$ identifies this cohomology class with the isomorphism class of the parity local system. The associated connected two-sheeted cover is the cover classified by the same class and is exactly the parity double cover constructed in Theorem~\ref{thm:appendix-3416}.
\end{proof}

\paragraph{3.4.17 working interpretation.}
The remaining global defect has now been compressed to its sharpest invariant. What previously appeared as parity flips, spinorial lift behavior, and double-cover repair is encoded by one class in $H^1(B,\mathbb Z_2)$. This does not change the mechanism; it shows that the mechanism is already as small as it can be.


\subsection*{C.17 Eighteenth 3.4 extension: parity holonomy support}
\addcontentsline{toc}{subsection}{C.17 Eighteenth 3.4 extension: parity holonomy support}

\begin{theorem}[Appendix support for Theorem~\ref{thm:mainline-3418}]
\label{thm:appendix-3418}
Let $\mathcal L_{\mathrm{par}}$ be the parity local system on a connected admissible shell base $B$, classified by
\[
  \varepsilon_{\mathrm{par}}\in H^1(B,\mathbb Z_2).
\]
Then the monodromy representation of $\mathcal L_{\mathrm{par}}$ is exactly the homomorphism obtained by evaluating $\varepsilon_{\mathrm{par}}$ on mod-two loop classes. Its kernel consists precisely of those admissible closed continuation classes that lift to direct visible closure on the canonical parity double cover.
\end{theorem}

\begin{proof}
A locally constant $\mathbb Z_2$-bundle on $B$ has monodromy in the discrete group $\mathbb Z_2$, hence a homomorphism from $\pi_1(B,b_0)$ to $\mathbb Z_2$. By Theorem~\ref{thm:appendix-3417}, the same local system is classified by the cocycle class $\varepsilon_{\mathrm{par}}\in H^1(B,\mathbb Z_2)$. Under the standard identification of $H^1(B,\mathbb Z_2)$ with homomorphisms out of mod-two $1$-cycles, the monodromy around a loop is given by evaluation of $\varepsilon_{\mathrm{par}}$ on its mod-two homology class. The lifted double cover from Theorem~\ref{thm:appendix-3416} trivializes the parity transport, so the loops in the kernel are exactly those whose lifted continuation closes directly on the cover.
\end{proof}

\paragraph{3.4.18 working interpretation.}
The cohomological class now has a direct operational meaning. It is not only something that classifies the parity defect abstractly; it can be evaluated on a closed admissible shell path and returns a two-valued global outcome. This is the cleanest bridge so far between the atlas/cohomology side and the actual readout side.



\subsection*{C.18 Nineteenth 3.4 extension: empirical parity-holonomy support}
\addcontentsline{toc}{subsection}{C.18 Nineteenth 3.4 extension: empirical parity-holonomy support}

\begin{theorem}[Appendix support for Theorem~\ref{thm:mainline-3419}]
\label{thm:appendix-3419}
Let $\gamma$ be an admissible closed shell path decomposed into stable chart arcs meeting obstruction shells transversely. Then the mod-two count of parity-changing crossings is equal to the monodromy of the parity local system around $\gamma$, and hence equal to the evaluation of the parity obstruction class on the mod-two loop class of $\gamma$.
\end{theorem}

\begin{proof}
Each local obstruction crossing carries one of the four transition types from Theorem~\ref{thm:mainline-3414}. Passing to the post-fold parity local system collapses the fold component and retains only whether the crossing changes the residual parity label. Therefore the transported parity label after traversing the full loop is obtained by multiplying the local $\mathbb Z_2$ transition factors, which in additive notation is exactly the mod-two sum of the local parity-change indicators $\delta_k$. This product is the monodromy of the parity local system along the loop. By Theorem~\ref{thm:appendix-3418}, that monodromy is also the evaluation of $\varepsilon_{\mathrm{par}}$ on the mod-two homology class of $\gamma$.
\end{proof}

\paragraph{3.4.19 working interpretation.}
The parity obstruction can now be read directly from a calibrated shell trace. Once the trace has been decomposed into admissible chart arcs, the full global lift defect collapses to a mod-two count of parity-changing transitions. This is the first point at which the global cohomological obstruction becomes an explicit detector-side or probe-side measurement rule.


\subsection*{C.19 Twentieth 3.4 extension: closing synthesis support}
\addcontentsline{toc}{subsection}{C.19 Twentieth 3.4 extension: closing synthesis support}

\begin{theorem}[Appendix support for Theorem~\ref{thm:mainline-3420}]
\label{thm:appendix-3420}
The stabilized 3.4 line factors through a single visible package
\[
  \bigl(\Phi,\Delta_{\mathrm{diag}},\varepsilon_{\mathrm{par}},\chi_{\mathrm{par}},\widehat\chi_{\mathrm{par}}\bigr)
\]
with branch variable $\eta$, in the sense that every local regular shell observable and every global parity-readout statement used in the mainline 3.4 section is a consequence of these data together with the shell atlas away from the obstruction set.
\end{theorem}

\begin{proof}
Theorem~\ref{thm:mainline-346} reduces the visible branch dynamics to the analytic generator $\Phi(\eta^2)$, and Theorems~\ref{thm:appendix-3411}--\ref{thm:appendix-3413} show that calibrated shell inversion away from the obstruction set is entirely controlled by the shell coordinate and parity label, with $\Delta_{\mathrm{diag}}$ measuring chartwise mismatch. Theorems~\ref{thm:appendix-3415}--\ref{thm:appendix-3418} then compress the remaining global ambiguity to the parity local system, its cohomological class, and its holonomy character. Finally, Theorem~\ref{thm:appendix-3419} identifies the empirical estimator with this same holonomy class along admissible closed traces. Hence the listed package captures the entire visible 3.4 layer.
\end{proof}

\paragraph{3.4.20 working interpretation.}
The 3.4 material is now closed as a single generator-atlas-holonomy block. Local residual dynamics is carried by the one-field generator, local diagnostic mismatch is measured by $\Delta_{\mathrm{diag}}$, and the only global leftover is the parity obstruction class with its loop readout. This is the natural handoff point for the next stage.


\subsection*{C.20 First 3.5 extension: residual-record support}
\addcontentsline{toc}{subsection}{C.20 First 3.5 extension: residual-record support}

\begin{theorem}[Appendix support for Theorem~\ref{thm:mainline-3500}]
\label{thm:appendix-3500}
Let $B$ be a connected admissible shell base satisfying the stabilized 3.4 hypotheses. Then the visible branch layer admits a faithful record description by
\[
  \mathfrak R(B)=\bigl(\Phi,\Delta_{\mathrm{diag}},\varepsilon_{\mathrm{par}},\chi_{\mathrm{par}},\widehat\chi_{\mathrm{par}}\bigr),
\]
in the sense that every chartwise regular-shell observable and every global parity-readout statement in the stabilized 3.4 sector is recoverable from this record together with the chartwise shell coordinate and parity label.
\end{theorem}

\begin{proof}
Theorem~\ref{thm:appendix-3420} already proves that the full stabilized 3.4 line factors through the same five visible data together with the shell atlas away from the obstruction set. The generator $\Phi$ determines the local branch law, the inverse shell charts determine the regular-shell coordinate and parity label, $\Delta_{\mathrm{diag}}$ measures chartwise diagnostic mismatch, and the parity package records the only remaining global ambiguity together with its operational loop estimator. Therefore the visible branch layer is faithfully encoded by the record $\mathfrak R(B)$, and the only extra ingredients needed for local reconstruction are the chartwise shell coordinate and parity label already present in the inverse atlas.
\end{proof}

\paragraph{3.5.0 working interpretation.}
The first simplification step does not yet introduce a new formal language. It only observes that the visible branch layer has already closed enough to be handled as a single admissible record. This is exactly the right level of compression before any later language-level or proof-economy layer is introduced. In the current status reading, this also marks the point after which later progress should be judged mainly by record-level sharpening, not by adding fresh visible ingredients.



\subsection*{C.21 Second 3.5 extension: residual-record morphism support}
\addcontentsline{toc}{subsection}{C.21 Second 3.5 extension: residual-record morphism support}

\begin{theorem}[Appendix support for Theorem~\ref{thm:mainline-3501}]
\label{thm:appendix-3501}
Let $F:B_1\to B_2$ be a residual-record isomorphism between connected admissible shell bases. Then the entire stabilized visible generator-atlas-holonomy sector is preserved under transport along $F$: the transported generator law, calibrated inverse shell data, shell-consistency diagnostics, parity obstruction class, parity holonomy, and empirical loop estimator all remain faithful to the same residual record.
\end{theorem}

\begin{proof}
By Theorem~\ref{thm:appendix-3500}, the stabilized visible sector factors through the admissible record together with chartwise shell coordinates. The definition of residual-record isomorphism requires preservation of exactly these ingredients at the level at which they enter the reconstruction normal form: the generator data, the shell-consistency defect, and the parity package. Therefore every chartwise visible formula and every global parity-readout statement already proved in the stabilized sector is transported faithfully along $F$, and no new visible branch datum is opened by this translation.
\end{proof}

\paragraph{3.5.1 working interpretation.}
The first move beyond the residual record is not to enlarge the visible layer, but to declare when two residual presentations should count as the same. This makes later compression work more disciplined: any future formal front-end must preserve the record-level morphism class if it is to remain faithful to the present mathematics. Historically open blocks that survive beyond this point therefore survive as residual identification or quantitative tasks, not as unresolved questions about what the visible data are.


\subsection*{C.22 Third 3.5 extension: residual-record interpreter support}
\addcontentsline{toc}{subsection}{C.22 Third 3.5 extension: residual-record interpreter support}

\begin{theorem}[Appendix support for Theorem~\ref{thm:mainline-3502}]
\label{thm:appendix-3502}
Let $B$ be a connected admissible shell base in the stabilized generator-atlas-holonomy regime. Then every visible statement in this regime is representable by a record-level formula built from
\[
  \mathfrak R(B)=\bigl(\Phi,\Delta_{\mathrm{diag}},\varepsilon_{\mathrm{par}},\chi_{\mathrm{par}},\widehat\chi_{\mathrm{par}}\bigr)
\]
together with chartwise shell coordinate and parity label, and this representation is unique up to residual-record isomorphism.
\end{theorem}

\begin{proof}
Theorem~\ref{thm:appendix-3500} already establishes that the entire stabilized visible branch layer factors through the admissible residual record together with shell-chart data. Theorem~\ref{thm:appendix-3501} shows that these data are preserved exactly under residual-record isomorphisms. Hence every visible theorem in the stabilized sector may be rewritten as a formula over the record package and shell-chart coordinates, and any two such rewritings can differ only by a residual-record isomorphism. This proves exact representability at record level.
\end{proof}

\paragraph{3.5.2 working interpretation.}
The meaning of the interpreter step is simple: from this point onward, the stabilized visible theory has a fixed admissible translation interface. Any future compression layer may become more economical or more set-theoretically explicit, but it should still be judged by whether it preserves the same record-level content. This is also the right point to separate the later live fronts from the historical OP labels: the surviving questions concern residual corrections, controlled emergence, spectral identification, Dixon compression, and the Theta-sum remainder, while OP~2 as a separate operator problem and OP~3 as a primary live front are no longer active blockers.



\subsection*{C.23 First 3.6 extension: conservative frontend support}
\addcontentsline{toc}{subsection}{C.23 First 3.6 extension: conservative frontend support}

\begin{theorem}[Appendix support for Theorem~\ref{thm:mainline-3600}]
\label{thm:appendix-3600}
Let $(\mathrm{Fr}(B),\mathcal F_B,\mathcal J_B)$ be a conservative frontend over the stabilized visible sector of a connected admissible shell base $B$. Then every frontend statement coming from a visible statement is fixed, up to unique record-level translation, by the same residual-record content as the original visible statement.
\end{theorem}

\begin{proof}
Theorem~\ref{thm:appendix-3502} already identifies the exact record-level representative of every visible theorem in the stabilized sector. The defining factorization rule for a conservative frontend requires frontend presentation to pass through that same representative, while transport-compatibility and primitive non-drift forbid any deviation on the visible image. Hence the frontend contributes only a conservative presentation layer and no second visible basis.
\end{proof}

\paragraph{3.6.0 working interpretation.}
The gain of the first 3.6 step is methodological but real: later language work is no longer a vague future promise. It now has a theorem-level admissibility test. This is precisely what keeps the planned ZFM-/frontend line subordinate to the mathematics-first strategy of the Companion.


\clearpage
\subsection*{Appendix C.124. Visual frame-flicker diagnostics for OP~5}
\addcontentsline{toc}{subsection}{Appendix C.124. Visual frame-flicker diagnostics for OP~5}
The following figures make the Section~3.8.48 hypothesis auditable on the current benchmark window. They do not prove the frame-flicker ansatz, but they show that a single smooth packet-scale moving-frame mode can already account for most of the mismatch beyond the rigid carrier and can do so in the two correction channels predicted by the current first-order frame-flicker decomposition.

\begin{remark}[Systems-level reading of the OP~5 visual block]\label{rem:appendixC-op5-systems-reading}
In the systems language summarized in Figure~\texttt{fig:qam-vm-os-zfs} in the active-core PDF, the present OP~5 figures should be read as visible-support views rather than as independent kernel identities. The rigid carrier, the packet breathing, and the localized readout slip are therefore shown here precisely because they belong to the auditable view layer of one transported channel. Their role is to expose where the visible readout departs from the rigid carrier picture, while the structural closure judgment itself still remains delegated to the deeper branch/shell/transport architecture.
\end{remark}


\begin{remark}[Visual-to-proof workflow for the OP~5 block]\label{rem:appendixC-op5-visual-workflow}
The present OP~5 figures should therefore be used in the following order. First, they identify a visible defect pattern at the readout layer. Second, they motivate a packet/carrier/slip decomposition that can be audited channelwise. Third, that decomposition must be checked against the active transport and closure-preservation architecture rather than treated as a proof by itself. Only after such a check has been made may the same pattern be recorded as a later kernel-first compression candidate. In particular, the figures are admissible as \emph{diagnostic entry points} and as \emph{audit surfaces}, but not as autonomous closure theorems.
\end{remark}

\begin{remark}[What the OP~5 figures are allowed to certify]\label{rem:appendixC-op5-allowed-certification}
At the present stage, the OP~5 visual block may certify only that a single packet-scale moving-frame correction remains empirically competitive as a visible-support description on the chosen benchmark window, and that the induced residual can be organized in the two currently tracked correction channels. It may not certify uniqueness of the correction mechanism, global closure of the OP~5 hypothesis, or activation of a new transport rule beyond the active theorem chain.
\end{remark}


\begin{figure}[H]
\centering
\IfFileExists{op5_frame_flicker_fit_v3862.pdf}{\includegraphics[width=0.98\textwidth]{op5_frame_flicker_fit_v3862.pdf}}{\safegraphic[width=0.98\textwidth]{op5_frame_flicker_fit_v3862.png}}
\caption{\textbf{OP~5 packet/readout fit on the current benchmark window.} Upper panel: the observed OP~5 packet/readout trace, the rigid carrier baseline, and the first-order packet/readout correction inside the visible packet-support window. Lower panel: the residual beyond the rigid carrier split into the breathing-envelope channel \(\alpha\psi_{N,m}\sin\Theta_{N,m}\) and the readout-slip channel \(\beta A_{N,m}\psi_{N,m}\cos\Theta_{N,m}\).}
\label{fig:appendixC-op5-frame-flicker-fit-live}
\end{figure}

\begin{figure}[H]
\centering
\IfFileExists{op5_frame_flicker_components_v3862.pdf}{\includegraphics[width=0.98\textwidth]{op5_frame_flicker_components_v3862.pdf}}{\safegraphic[width=0.98\textwidth]{op5_frame_flicker_components_v3862.png}}
\caption{\textbf{Packet-support mode, visible channels, and incremental OP~5 build-up.} The four panels show the packet-support window \(\psi_{N,m}\), the rigid carrier/readout quadratures, the two visible channels fed by one packet-scale moving-frame mode, and the stepwise build-up from rigid carrier to packet/breathing correction and then to the full OP~5 readout trace.}
\label{fig:appendixC-op5-frame-flicker-components-live}
\end{figure}

\begin{figure}[H]
\centering
\IfFileExists{op5_frame_flicker_geometry_v3862.pdf}{\includegraphics[width=0.98\textwidth]{op5_frame_flicker_geometry_v3862.pdf}}{\safegraphic[width=0.98\textwidth]{op5_frame_flicker_geometry_v3862.png}}
\caption{\textbf{Geometric OP~5 reading: carrier helix, breathing packet, and readout slip.} The panels compare the rigid carrier helix with the packet/readout-distorted helix, show the top-view deviation, display the packet breathing envelope along the helix, and isolate the localized readout phase slip. The natural reading is therefore not a second bulk spiral, but the same carrier seen through one softened visible frame.}
\label{fig:appendixC-op5-frame-flicker-geometry-live}
\end{figure}

\noindent\textit{Reading guide.} Figures~\ref{fig:appendixC-op5-frame-flicker-fit-live}--\ref{fig:appendixC-op5-frame-flicker-geometry-live} should be read in order as \emph{trace $\rightarrow$ channel decomposition $\rightarrow$ geometry}. In that sequence, the packet/readout defect moves from the visible time trace, to the two correction channels, to the helical geometric picture.

\paragraph{Appendix conclusion.}
On the present benchmark window, the frame-flicker ansatz remains hypothesis-grade, but it is no longer only verbal. A single smooth moving-frame mode reproduces most of the residual correction beyond the rigid carrier and does so in exactly the two channels predicted by the first-order decomposition. This strengthens the current reading that OP~5 is better treated as a carrier-plus-readout-defect problem than as a missing second rigid transport wave.

\begin{remark}[Packet/flicker diagnostics as readout evidence]
Within the current OP transport architecture, the packet-wave, breathing-envelope, and frame-flicker figures should be read as visible-support diagnostics rather than as primary closure criteria. Their mathematical role is to show that a large part of the observed OP~5 residue is compatible with a carrier-plus-readout-defect picture, which is exactly the kind of situation filtered later by the visible/structural mismatch separation tool.
\end{remark}

\paragraph{Bibliography placement note.}
The bibliography that formerly appeared at this historical appendix position has been moved to the final back matter so that the document closes with the references. The version-record and historical-release material continues below unchanged as archival proof-memory.






\subsection*{Working common-mode carrier is downstream from the active HC/MM/J fit}

The previous subsection reduced the primitive OP~1/OP~5 obstruction to a common-mode carrier plus a downstream observer-shell split. The next structural step is to push that common-mode carrier itself back to the active HC/MM/J fit. On the active admissible branch, the shared carrier is not an independent visible object: it is simply the shared visible mask applied to the complementary output generated by the active matrix-leakage operator. In particular, if the active branch is exactly HC/MM/J-compatible, then the primitive common-mode carrier must vanish identically.

\begin{definition}[Working active matrix-fit defect and shared carrier readout]
Let $U:=\mathrm{span}\{u_1,u_2\}$ be the active admissible plane and $W:=\mathrm{span}\{w_1,w_2\}$ the active complementary leakage plane. Define the active matrix-fit defect by
\[
  E_{\mathrm{fit}}^{(U)} := \bigl((\mathrm{id}-C)\mathcal M C\bigr)\big|_{U}: U\to W.
\]
For a unit active probe state
\[
  x(\theta):=\cos\theta\,u_1+\sin\theta\,u_2\in U,
\]
define the active hidden leakage response and the shared visible carrier by
\[
  y_{\mathrm{fit}}(\theta):=E_{\mathrm{fit}}^{(U)}x(\theta)\in W,
  \qquad
  c_{\mathrm{sh}}^{\mathrm{fit}}(\theta):=L_{\mathrm{sh}}^{(W)}\bigl(y_{\mathrm{fit}}(\theta)\bigr).
\]
\end{definition}

\begin{proposition}[Working exact commutator identity on the admissible branch]
For every $x\in U\subseteq \mathrm{im}(C)$ one has
\[
  (\mathrm{id}-C)\mathcal M x
  =
  \frac12\,[\mathcal M,J]x.
\]
Equivalently,
\[
  E_{\mathrm{fit}}^{(U)}
  =
  \frac12\,[\mathcal M,J]\big|_{U}.
\]
In particular,
\[
  \|E_{\mathrm{fit}}^{(U)}\|
  \le
  \frac12\,\|[\mathcal M,J]\big|_{U}\|.
\]
\end{proposition}

\begin{proof}
Since $U\subseteq \mathrm{im}(C)$, every $x\in U$ satisfies $Cx=x$ and $Jx=x$. Therefore
\[
  (\mathrm{id}-C)\mathcal M x
  =
  \frac12(\mathrm{id}-J)\mathcal M x
  =
  \frac12\bigl(\mathcal Mx-J\mathcal Mx\bigr).
\]
Using $Jx=x$, this becomes
\[
  \frac12\bigl(\mathcal MJx-J\mathcal Mx\bigr)
  =
  \frac12\,[\mathcal M,J]x.
\]
Restricting to $U$ gives the operator identity, and the norm bound is immediate.
\end{proof}

\begin{theorem}[Working shared carrier is controlled by the active HC/MM/J fit]
For every active probe angle $\theta$ one has
\[
  c_{\mathrm{sh}}^{\mathrm{fit}}(\theta)
  =
  L_{\mathrm{sh}}^{(W)}\bigl(E_{\mathrm{fit}}^{(U)}x(\theta)\bigr)
  =
  \frac12\,L_{\mathrm{sh}}^{(W)}\bigl([\mathcal M,J]x(\theta)\bigr).
\]
Hence
\[
  |c_{\mathrm{sh}}^{\mathrm{fit}}(\theta)|
  \le
  \|L_{\mathrm{sh}}^{(W)}\|_{W^*}\,\|E_{\mathrm{fit}}^{(U)}\|\,\|x(\theta)\|
  \le
  \frac12\,\|L_{\mathrm{sh}}^{(W)}\|_{W^*}\,\|[\mathcal M,J]\big|_{U}\|.
\]
In particular, if
\[
  [\mathcal M,J]\big|_{U}=0,
\]
then
\[
  c_{\mathrm{sh}}^{\mathrm{fit}}(\theta)=0
\]
for all $\theta$, so the primitive common-mode OP~1/OP~5 carrier vanishes on the active branch.
\end{theorem}

\begin{proof}
The first displayed identity is the definition of $c_{\mathrm{sh}}^{\mathrm{fit}}$. The second follows from the previous proposition. Since $\|x(\theta)\|=1$ on the active unit circle in $U$, the norm estimate follows immediately. If the restricted commutator vanishes, then $E_{\mathrm{fit}}^{(U)}=0$, hence $y_{\mathrm{fit}}(\theta)=0$ and therefore $c_{\mathrm{sh}}^{\mathrm{fit}}(\theta)=0$ for all $\theta$.
\end{proof}

\begin{corollary}[Working primitive common-mode residual is a pure matrix-fit obstruction]
Assume exact grouped no-drift on the active leakage plane, so that $\Delta L^{(W)}=0$. Then every surviving primitive OP~1/OP~5 residual on the active branch is completely controlled by the active HC/MM/J fit defect:
\[
  d(\theta)=c_{\mathrm{sh}}^{\mathrm{fit}}(\theta)\binom{1}{1},
  \qquad
  |c_{\mathrm{sh}}^{\mathrm{fit}}(\theta)|
  \le
  \frac12\,\|L_{\mathrm{sh}}^{(W)}\|_{W^*}\,\|[\mathcal M,J]\big|_{U}\|.
\]
In particular, under exact grouped no-drift and exact active commutation with the umkehrwalze, the primitive OP~1/OP~5 channel closes identically on the active branch.
\end{corollary}

\begin{proof}
Under exact grouped no-drift one has $c_{\mathrm{obs}}\equiv 0$ by the previous shared-residual dichotomy, so only the common-mode component remains. The displayed bound is the previous theorem. If the restricted commutator also vanishes, then the common-mode carrier vanishes identically, hence the primitive visible channel closes.
\end{proof}

\paragraph{Working interpretation.}
This is the clean matrix-side endpoint of the current reduction. Once the observer-shell split is removed on the active leakage plane, the remaining primitive OP~1/OP~5 carrier is not a mysterious visible residue anymore: it is exactly the shared mask applied to the active HC/MM/J commutator defect. So the primitive common-mode channel survives only to the extent that the Millennium matrix still fails to commute with the umkehrwalze on the active admissible branch.



\subsection*{Working architectural promotion of the active commutator defect}

The active commutator defect was introduced as the primitive common-mode obstruction on the active branch. However, the older grouped-layer architecture already contains a stronger structural principle: admissible Millennium transport is only allowed when it intertwines with the involution. Therefore the active commutator should not be read as a free new defect by default, but as a quantity that must vanish whenever the transport really is the admissible grouped-layer converter claimed by the document architecture. What then remains, if anything, must be downstream dressing rather than primitive HC/MM/J leakage.

\begin{theorem}[Working admissible grouped transport kills the active commutator]
Assume the Millennium matrix $\mathcal M$ acts on the active admissible plane $U$ as an admissible grouped-layer converter in the sense of Definition~\texttt{def:millennium-transport-grouped-layers} in the active-core PDF, and assume its transport role is structure preserving with respect to the involution/umkehrwalze in the sense of Lemma~\texttt{lem:millennium-intertwining-lemma} in the active-core PDF. Then
\[
  [\mathcal M,J]\big|_{U}=0.
\]
Consequently,
\[
  E_{\mathrm{fit}}^{(U)}=0,
  \qquad
  c_{\mathrm{sh}}^{\mathrm{fit}}(\theta)=0
  \quad\text{for all }\theta.
\]
\end{theorem}

\begin{proof}
By Lemma~\texttt{lem:millennium-intertwining-lemma} in the active-core PDF, every Millennium transport that claims to preserve the normalized involutive kernel must satisfy
\[
  \mathcal M\mathcal J=\mathcal J\mathcal M.
\]
On the active branch the visible involution is realized by the operator $J$, so the same intertwining statement on the active branch reads
\[
  \mathcal M Jx=J\mathcal M x
  \qquad (x\in U).
\]
Hence
\[
  [\mathcal M,J]x=0
  \qquad (x\in U),
\]
which is exactly the claimed vanishing of the restricted commutator. The identities
\[
  E_{\mathrm{fit}}^{(U)}=\tfrac12[\mathcal M,J]\big|_{U}
\]
and
\[
  c_{\mathrm{sh}}^{\mathrm{fit}}(\theta)=L_{\mathrm{sh}}^{(W)}\bigl(E_{\mathrm{fit}}^{(U)}x(\theta)\bigr)
\]
were established in the previous subsection, so both vanish immediately.
\end{proof}

\begin{corollary}[Working primitive OP~1/OP~5 common mode cannot survive exact admissible transport]
Assume exact grouped no-drift on the active leakage plane and admissible grouped Millennium transport on the active branch. Then the primitive visible OP~1/OP~5 channel closes identically on that branch:
\[
  d(\theta)=0
  \qquad\text{for all }\theta.
\]
In particular, any still visible residual must be interpreted as downstream observer-shell dressing, phase-shifted readout, or some other post-transport effect rather than as a primitive HC/MM/J leakage mode.
\end{corollary}

\begin{proof}
Exact grouped no-drift gives $c_{\mathrm{obs}}(\theta)=0$ by the shared-residual dichotomy. The previous theorem gives $c_{\mathrm{sh}}^{\mathrm{fit}}(\theta)=0$. Since the primitive visible channel decomposes as the sum of the common-mode part and the observer-split part, both pieces vanish and therefore $d(\theta)=0$ for all $\theta$.
\end{proof}

\begin{remark}[Working interpretation: exact closure versus canonical observer-shell residual]
The current architecture therefore supports a sharp dichotomy.
\begin{enumerate}[leftmargin=2em,label=(\roman*)]
\item If the Millennium matrix really acts as the admissible grouped-layer converter claimed by the structural architecture, then the active commutator vanishes on the active branch and no primitive common-mode OP~1/OP~5 leakage survives there.
\item If a residual nevertheless remains visible, then it should no longer be hunted as a primitive HC/MM/J closure defect. It must instead be read as a canonical downstream observer-shell residual: phase-shifted readout, shell dressing, expansion dressing, or another post-transport effect living after the exact admissible transport stage.
\end{enumerate}
This is the strongest structural narrowing achieved so far: primitive matrix leakage and canonical observer-shell residual are no longer being mixed. The architecture itself says that only one of them may remain once admissible grouped transport is enforced.
\end{remark}


\subsection*{Canonical downstream observer-shell residual after exact admissible transport}
The previous dichotomy showed that once grouped no-drift holds on the active leakage plane and admissible grouped Millennium transport is exact on the active branch, no primitive HC/MM/J leakage can remain visible. Any still visible OP~1/OP~5 remainder must therefore live strictly downstream from the exact transport stage. The next structural step is to package that remainder as a single canonical observer-shell mode rather than as two independent unfinished closures.

\begin{definition}[Working downstream observer-shell dressing mode]
Assume exact grouped no-drift on the active leakage plane $W$ and exact admissible grouped Millennium transport on the active branch $U$. A downstream observer-shell dressing mode is a scalar family
\[
  r_{\mathrm{obs}}(\theta)
\]
and a common visible shell profile
\[
  e_{\mathrm{obs}}:=\binom{1}{1}
\]
such that the residual visible channel can be written in the form
\[
  d(\theta)=r_{\mathrm{obs}}(\theta)\,e_{\mathrm{obs}}+d_{\mathrm{split}}(\theta),
\]
where $d_{\mathrm{split}}(\theta)$ is a pure observer-side split term living after the exact transport stage.
\end{definition}

\begin{theorem}[Working canonical observer-shell residual after exact admissible transport]
Assume exact grouped no-drift on the active leakage plane $W$ and exact admissible grouped Millennium transport on the active branch $U$. Then any remaining visible OP~1/OP~5 residual admits the canonical decomposition
\[
  d(\theta)=r_{\mathrm{obs}}(\theta)\binom{1}{1}+c_{\mathrm{obs}}(\theta)\binom{1}{-1}.
\]
Moreover, the primitive matrix-leakage contribution vanishes identically, so the common-mode coefficient $r_{\mathrm{obs}}(\theta)$ is a purely downstream observer-shell residual rather than a primitive HC/MM/J leakage carrier.
\end{theorem}

\begin{proof}
By the shared-residual dichotomy, the primitive visible channel decomposes as a common-mode part and an observer-split part,
\[
  d(\theta)=c_{\mathrm{sh}}(\theta)\binom{1}{1}+c_{\mathrm{obs}}(\theta)\binom{1}{-1}.
\]
The previous active-matrix-fit analysis identifies the primitive common-mode part with the active HC/MM/J fit carrier. Under exact admissible grouped Millennium transport on $U$, the active commutator vanishes and therefore the primitive common-mode carrier vanishes as well. Hence any surviving common-mode coefficient can only appear after the exact transport stage, i.e. as downstream observer-shell dressing. Renaming this remaining downstream common-mode coefficient by $r_{\mathrm{obs}}(\theta)$ yields the stated decomposition.
\end{proof}

\begin{corollary}[Working no-drift collapses to one observer-shell mode]
Under the hypotheses of the theorem, if the observer-side split also vanishes on the active leakage plane,
\[
  c_{\mathrm{obs}}(\theta)=0\qquad\text{for all }\theta,
\]
then the full visible OP~1/OP~5 remainder collapses to a single canonical observer-shell residual:
\[
  d(\theta)=r_{\mathrm{obs}}(\theta)\binom{1}{1}.
\]
In particular, OP~1 and OP~5 then do not carry two independent unfinished closures on the active branch; they are merely two reads of the same downstream shell residual.
\end{corollary}

\begin{proof}
The theorem already gives the decomposition
\[
  d(\theta)=r_{\mathrm{obs}}(\theta)\binom{1}{1}+c_{\mathrm{obs}}(\theta)\binom{1}{-1}.
\]
If $c_{\mathrm{obs}}(\theta)=0$, only the common shell term remains, which is exactly the stated formula.
\end{proof}

\begin{remark}[Working interpretation of the surviving residual]
This is the intended structural endpoint of the current OP~1/OP~5 reduction. Once exact admissible transport kills the primitive HC/MM/J leakage and grouped no-drift removes the active observer split, any remaining visible residue should not be interpreted as evidence for two separate unfinished closures. It is a single canonical downstream shell residual, shared by OP~1 and OP~5 and only seen through two observer labels.
\end{remark}


\subsection*{From two unfinished closures to one shell-mode law}
The previous reduction reached a decisive structural endpoint: once exact admissible transport removes primitive HC/MM/J leakage and grouped no-drift removes the active observer split, OP~1 and OP~5 no longer describe two independent unfinished closures on the active branch. What remains is a single downstream observer-shell mode. The next step is therefore not to force two separate zero-conditions, but to state explicitly that the visible target has collapsed to one scalar shell law.

\begin{definition}[Working shared shell-mode amplitude]
Assume the hypotheses of the previous corollary, so that
\[
  d(\theta)=r_{\mathrm{obs}}(\theta)\binom{1}{1}.
\]
The scalar family
\[
  r_{\mathrm{obs}}(\theta)
\]
is called the shared shell-mode amplitude of the visible OP~1/OP~5 residual.
\end{definition}

\begin{theorem}[Working one-law reduction of OP~1 and OP~5 after exact transport]
Assume exact grouped no-drift on the active leakage plane $W$, exact admissible grouped Millennium transport on the active branch $U$, and vanishing active observer split on $W$. Then the visible OP~1/OP~5 channel is governed by a single scalar law:
\[
  d_1(\theta)=d_5(\theta)=r_{\mathrm{obs}}(\theta).
\]
In particular, the visible target is no longer a pair of independent closure equations, but the single shell-mode equation
\[
  r_{\mathrm{obs}}(\theta)=0
\]
if exact visible closure is required, or the single shell-mode identification problem
\[
  r_{\mathrm{obs}}(\theta)=r_{\mathrm{can}}(\theta)
\]
if a canonical nonzero observer-shell residual is allowed.
\end{theorem}

\begin{proof}
Under the stated hypotheses, the previous corollary gives
\[
  d(\theta)=r_{\mathrm{obs}}(\theta)\binom{1}{1}.
\]
Reading off the two components immediately yields
\[
  d_1(\theta)=r_{\mathrm{obs}}(\theta),
  \qquad
  d_5(\theta)=r_{\mathrm{obs}}(\theta).
\]
Hence the visible OP~1/OP~5 target collapses from two residual equations to one scalar shell-mode law. If one demands exact visible closure, that law is $r_{\mathrm{obs}}(\theta)=0$. If one instead admits a canonical downstream observer-shell remainder, the problem is to identify the canonical shell law $r_{\mathrm{can}}(\theta)$ and match the visible residual against it.
\end{proof}

\begin{corollary}[Working exact closure and canonical residual are one-mode alternatives]
Under the same hypotheses, one cannot simultaneously interpret the surviving visible channel as two independent unfinished closures and as a single canonical observer-shell mode. Exactly one of the following descriptions may hold:
\begin{enumerate}[leftmargin=2em,label=(\roman*)]
\item exact visible closure, i.e.
\[
  r_{\mathrm{obs}}(\theta)=0\qquad\text{for all }\theta;
\]
\item canonical observer-shell remainder, i.e.
\[
  r_{\mathrm{obs}}(\theta)=r_{\mathrm{can}}(\theta)\not\equiv 0.
\]
\end{enumerate}
In either case, OP~1 and OP~5 remain locked to the same scalar mode.
\end{corollary}

\begin{proof}
The theorem shows that both visible channels are equal to the same scalar amplitude $r_{\mathrm{obs}}(\theta)$. Therefore any interpretation of the surviving visible remainder must be an interpretation of this single scalar family. It cannot simultaneously serve as two independent residual equations and as one shared shell mode. The two listed alternatives are exactly the two structurally admissible readings.
\end{proof}

\begin{remark}[Working interpretation: closure target versus shell-law target]
This is the point where the present architecture changes the question. The old reading asked whether OP~1 and OP~5 can both be closed independently. The new reading says: after exact admissible transport and active no-drift, there is only one visible target left. The remaining decision is whether that target should vanish or whether it should be identified with a canonical post-transport shell law, for example expansion dressing or phase-shifted observer-shell readout.
\end{remark}


\subsection*{Canonical phase-shifted / expansion-dressed shell law}
The previous reduction showed that after exact admissible transport and active no-drift, the visible OP~1/OP~5 question collapses to a single scalar shell-mode amplitude
\[
  r_{\mathrm{obs}}(\theta).
\]
The next structural step is to stop treating this quantity as a bare leftover and instead give a first explicit law for it. The guiding hypothesis is that the observer shell is not a second ontology but a small post-transport dressing of the same exact ground closure. In that case the surviving visible mode should arise from a phase-shifted and possibly expansion-dressed readout of the same underlying shell carrier.

\begin{definition}[Working observer-shell phase and expansion defects]
A post-transport observer-shell defect state is described by two scalar families
\[
  \delta_{\mathrm{obs}}(\theta),
  \qquad
  \eta_{\mathrm{exp}}(\theta),
\]
where
\begin{itemize}[leftmargin=2em]
\item $\delta_{\mathrm{obs}}(\theta)$ measures the local phase shift between the exact transport shell and the visible observer shell,
\item $\eta_{\mathrm{exp}}(\theta)$ measures the local expansion / shell-spacing drift of the visible shell relative to the exact transport shell.
\end{itemize}
Both quantities are downstream observer-shell parameters: they vanish on the exact transport shell itself.
\end{definition}

\begin{definition}[Working first-order canonical shell law]
A first-order canonical observer-shell law is an identity of the form
\[
  r_{\mathrm{can}}(\theta)
  =A_{\mathrm{sh}}(\theta)\,\delta_{\mathrm{obs}}(\theta)
  +B_{\mathrm{sh}}(\theta)\,\eta_{\mathrm{exp}}(\theta)
  +R_{\ge 2}(\theta),
\]
where $A_{\mathrm{sh}}(\theta)$ and $B_{\mathrm{sh}}(\theta)$ are shell-response coefficients and the remainder satisfies
\[
  |R_{\ge 2}(\theta)|
  \le C_{\mathrm{sh}}(\theta)
  \bigl(|\delta_{\mathrm{obs}}(\theta)|+|\eta_{\mathrm{exp}}(\theta)|\bigr)^2.
\]
In the strictly phase-shifted case one has $B_{\mathrm{sh}}\equiv 0$, while in the pure expansion-dressing case one has $A_{\mathrm{sh}}\equiv 0$.
\end{definition}

\begin{theorem}[Working first-order phase-shift / expansion law for the shared shell mode]
Assume exact admissible grouped Millennium transport on the active branch, exact grouped no-drift on the active leakage plane, and vanishing active observer split on that plane. Assume moreover that the visible observer shell is a smooth downstream perturbation of the exact transport shell, parameterized by phase defect $\delta_{\mathrm{obs}}(\theta)$ and expansion defect $\eta_{\mathrm{exp}}(\theta)$. Then the shared visible mode admits a first-order shell law
\[
  r_{\mathrm{obs}}(\theta)
  =A_{\mathrm{sh}}(\theta)\,\delta_{\mathrm{obs}}(\theta)
  +B_{\mathrm{sh}}(\theta)\,\eta_{\mathrm{exp}}(\theta)
  +R_{\ge 2}(\theta),
\]
with a quadratic remainder in the total shell defect magnitude. In particular, exact visible closure is equivalent to the vanishing of this canonical shell law, while a nonzero visible remainder is compatible with exact transport provided it is entirely generated by downstream shell dressing.
\end{theorem}

\begin{proof}
Under the stated hypotheses the visible channel has already collapsed to a single scalar amplitude,
\[
  d(\theta)=r_{\mathrm{obs}}(\theta)\binom{1}{1}.
\]
Hence every surviving visible effect must be read as a downstream perturbation of the same exact transport shell rather than as a primitive matrix-leakage channel. By the assumed smooth dependence of the visible observer shell on the two downstream shell parameters $\delta_{\mathrm{obs}}(\theta)$ and $\eta_{\mathrm{exp}}(\theta)$, the scalar amplitude $r_{\mathrm{obs}}(\theta)$ is a smooth scalar function of these two arguments. A first-order Taylor expansion around the exact shell point
\[
  (\delta_{\mathrm{obs}},\eta_{\mathrm{exp}})=(0,0)
\]
yields
\[
  r_{\mathrm{obs}}(\theta)
  =A_{\mathrm{sh}}(\theta)\,\delta_{\mathrm{obs}}(\theta)
  +B_{\mathrm{sh}}(\theta)\,\eta_{\mathrm{exp}}(\theta)
  +R_{\ge 2}(\theta),
\]
where the remainder is quadratic in the total shell defect magnitude. This is exactly the stated first-order canonical shell law.
\end{proof}

\begin{corollary}[Working exact closure versus canonical nonzero shell law]
Under the same hypotheses, the remaining visible OP~1/OP~5 target is reduced to the dichotomy
\begin{enumerate}[leftmargin=2em,label=(\roman*)]
\item exact visible closure:
\[
  A_{\mathrm{sh}}(\theta)\,\delta_{\mathrm{obs}}(\theta)
  +B_{\mathrm{sh}}(\theta)\,\eta_{\mathrm{exp}}(\theta)
  +R_{\ge 2}(\theta)=0;
\]
\item canonical nonzero observer-shell remainder:
\[
  r_{\mathrm{obs}}(\theta)=r_{\mathrm{can}}(\theta)
  \not\equiv 0
\]
with $r_{\mathrm{can}}$ given by the same first-order shell law.
\end{enumerate}
In both cases the visible question is one scalar shell law rather than two independent OP~1 / OP~5 closures.
\end{corollary}

\begin{proof}
The previous theorem identifies the surviving scalar mode with a first-order observer-shell law in the two downstream shell parameters. Therefore the visible target is exhausted by deciding whether that law vanishes identically or whether it should be retained as a canonical nonzero shell remainder.
\end{proof}

\begin{remark}[Working interpretation of the shell-law coefficients]
The coefficients $A_{\mathrm{sh}}(\theta)$ and $B_{\mathrm{sh}}(\theta)$ are the first visible response coefficients of the exact transport shell to observer-side phase shift and observer-side expansion dressing. They are not new primitive leakage parameters. Their role is to measure how a single exact transport closure is seen from a displaced observer shell. This is why OP~1 and OP~5 remain locked to the same scalar law even after exact transport has already removed the primitive HC/MM/J leakage.
\end{remark}


\subsection*{Phase-dominant, expansion-dominant,\ and mixed shell-law regimes}
The canonical shell law now reduces the visible post-transport question to one scalar combination of two downstream shell effects. The next structural step is therefore not to search for two further closures, but to decide which of the two first-order shell loads is actually primary.

\begin{definition}[Working shell loads and dominance ratio]
Define the first-order shell loads
\[
  \Lambda_{\mathrm{ph}}(\theta):=|A_{\mathrm{sh}}(\theta)\,\delta_{\mathrm{obs}}(\theta)|,
  \qquad
  \Lambda_{\mathrm{exp}}(\theta):=|B_{\mathrm{sh}}(\theta)\,\eta_{\mathrm{exp}}(\theta)|,
\]
and the total first-order shell load
\[
  \Lambda_{\mathrm{lin}}(\theta):=\Lambda_{\mathrm{ph}}(\theta)+\Lambda_{\mathrm{exp}}(\theta).
\]
Whenever $\Lambda_{\mathrm{ph}}(\theta)\neq 0$, define the shell-load ratio
\[
  \rho_{\mathrm{sh}}(\theta):=\frac{\Lambda_{\mathrm{exp}}(\theta)}{\Lambda_{\mathrm{ph}}(\theta)}.
\]
Thus $\rho_{\mathrm{sh}}\ll 1$ indicates phase-dominant shell dressing, while $\rho_{\mathrm{sh}}\gg 1$ indicates expansion-dominant shell dressing.
\end{definition}

\begin{definition}[Working phase-dominant, expansion-dominant, and mixed regimes]
Fix $\varepsilon_{\mathrm{dom}},\varepsilon_{\mathrm{rem}}\in[0,1)$ with
\[
  \varepsilon_{\mathrm{dom}}+\varepsilon_{\mathrm{rem}}<1.
\]
At a given parameter value $\theta$:
\begin{enumerate}[leftmargin=2em,label=(\roman*)]
\item the shell law is called \emph{phase-dominant} if
\[
  \Lambda_{\mathrm{exp}}(\theta)\le \varepsilon_{\mathrm{dom}}\,\Lambda_{\mathrm{ph}}(\theta)
  \quad\text{and}\quad
  |R_{\ge 2}(\theta)|\le \varepsilon_{\mathrm{rem}}\,\Lambda_{\mathrm{ph}}(\theta);
\]
\item it is called \emph{expansion-dominant} if
\[
  \Lambda_{\mathrm{ph}}(\theta)\le \varepsilon_{\mathrm{dom}}\,\Lambda_{\mathrm{exp}}(\theta)
  \quad\text{and}\quad
  |R_{\ge 2}(\theta)|\le \varepsilon_{\mathrm{rem}}\,\Lambda_{\mathrm{exp}}(\theta);
\]
\item it is called \emph{genuinely mixed at first order} if neither of the previous two dominance conditions holds.
\end{enumerate}
\end{definition}

\begin{theorem}[Working shell-law dominance classification]
Assume the first-order canonical shell law of the previous subsection. Then:
\begin{enumerate}[leftmargin=2em,label=(\roman*)]
\item in the phase-dominant regime one has
\[
  |r_{\mathrm{obs}}(\theta)-A_{\mathrm{sh}}(\theta)\,\delta_{\mathrm{obs}}(\theta)|
  \le (\varepsilon_{\mathrm{dom}}+\varepsilon_{\mathrm{rem}})\,\Lambda_{\mathrm{ph}}(\theta),
\]
so the shared shell mode is controlled by the observer-side phase-shift term up to a strictly smaller relative error;
\item in the expansion-dominant regime one has
\[
  |r_{\mathrm{obs}}(\theta)-B_{\mathrm{sh}}(\theta)\,\eta_{\mathrm{exp}}(\theta)|
  \le (\varepsilon_{\mathrm{dom}}+\varepsilon_{\mathrm{rem}})\,\Lambda_{\mathrm{exp}}(\theta),
\]
so the shared shell mode is controlled by the observer-shell expansion term up to a strictly smaller relative error;
\item in the genuinely mixed regime, the visible shell mode cannot be reduced at first order to only one of the two shell defects.
\end{enumerate}
\end{theorem}

\begin{proof}
By the first-order shell law,
\[
  r_{\mathrm{obs}}(\theta)
  =A_{\mathrm{sh}}(\theta)\,\delta_{\mathrm{obs}}(\theta)
  +B_{\mathrm{sh}}(\theta)\,\eta_{\mathrm{exp}}(\theta)
  +R_{\ge 2}(\theta).
\]
Hence
\[
  |r_{\mathrm{obs}}(\theta)-A_{\mathrm{sh}}(\theta)\,\delta_{\mathrm{obs}}(\theta)|
  \le \Lambda_{\mathrm{exp}}(\theta)+|R_{\ge 2}(\theta)|,
\]
and similarly
\[
  |r_{\mathrm{obs}}(\theta)-B_{\mathrm{sh}}(\theta)\,\eta_{\mathrm{exp}}(\theta)|
  \le \Lambda_{\mathrm{ph}}(\theta)+|R_{\ge 2}(\theta)|.
\]
In the phase-dominant regime this gives
\[
  |r_{\mathrm{obs}}(\theta)-A_{\mathrm{sh}}(\theta)\,\delta_{\mathrm{obs}}(\theta)|
  \le (\varepsilon_{\mathrm{dom}}+\varepsilon_{\mathrm{rem}})\,\Lambda_{\mathrm{ph}}(\theta),
\]
which is strictly smaller than the dominant phase load because $\varepsilon_{\mathrm{dom}}+\varepsilon_{\mathrm{rem}}<1$. The expansion-dominant case is identical after swapping the two shell loads. If neither dominance condition holds, then neither first-order term controls the shell mode up to a strictly smaller remainder, so the shell mode is genuinely mixed at first order.
\end{proof}

\begin{corollary}[Working next decision for the shared shell mode]
After exact admissible transport, active no-drift, and collapse to a single visible shell mode, the next structural decision is reduced to the comparison of the two first-order shell loads
\[
  \Lambda_{\mathrm{ph}}(\theta)
  \quad\text{and}\quad
  \Lambda_{\mathrm{exp}}(\theta),
\]
not to two independent closure equations. In particular, the visible question is now whether the common shell mode is primarily
\begin{enumerate}[leftmargin=2em,label=(\roman*)]
\item phase-shift driven,
\item expansion-drift driven,
\item or genuinely mixed at first order.
\end{enumerate}
\end{corollary}

\begin{proof}
This is the direct content of the dominance classification theorem.
\end{proof}

\begin{remark}[Working interpretation of the present status]
The current document does not yet prove globally which regime is realized. What it does prove is that the decision has been compressed to a single shell-law comparison problem. The remaining visible burden is therefore much narrower than the earlier picture of two unfinished OP~1 / OP~5 closures: it is now a one-mode question about which downstream observer-shell load is structurally primary.
\end{remark}



\subsection*{Signed shell-balance mismatch and the cancellation corridor}
The shell-load classification identifies which first-order observer-shell contribution is dominant, but exact visible closure depends on something sharper: whether the two signed first-order shell contributions actually cancel. The next structural step is therefore to replace the unsigned load comparison by a signed shell-balance mismatch.

\begin{definition}[Working signed shell coordinates and shell-balance mismatch]
Define the signed first-order shell coordinates
\[
  p_{\mathrm{sh}}(\theta):=A_{\mathrm{sh}}(\theta)\,\delta_{\mathrm{obs}}(\theta),
  \qquad
  e_{\mathrm{sh}}(\theta):=B_{\mathrm{sh}}(\theta)\,\eta_{\mathrm{exp}}(\theta),
\]
and the signed shell-balance mismatch
\[
  \chi_{\mathrm{sh}}(\theta):=p_{\mathrm{sh}}(\theta)+e_{\mathrm{sh}}(\theta).
\]
Whenever $p_{\mathrm{sh}}(\theta)\neq 0$, define the signed shell-balance ratio
\[
  \zeta_{\mathrm{sh}}(\theta):=\frac{e_{\mathrm{sh}}(\theta)}{p_{\mathrm{sh}}(\theta)}.
\]
Thus exact first-order cancellation corresponds to $\chi_{\mathrm{sh}}(\theta)=0$, equivalently $\zeta_{\mathrm{sh}}(\theta)=-1$.
\end{definition}

\begin{definition}[Working shell-cancellation corridor]
Fix $\varepsilon_{\mathrm{bal}},\varepsilon_{\mathrm{rem}}\in[0,1)$ with
\[
  \varepsilon_{\mathrm{bal}}+\varepsilon_{\mathrm{rem}}<1.
\]
At a given parameter value $\theta$, the shell law is said to lie in the \emph{cancellation corridor} if
\[
  |\chi_{\mathrm{sh}}(\theta)|\le \varepsilon_{\mathrm{bal}}\,\Lambda_{\mathrm{lin}}(\theta)
  \qquad\text{and}\qquad
  |R_{\ge 2}(\theta)|\le \varepsilon_{\mathrm{rem}}\,\Lambda_{\mathrm{lin}}(\theta).
\]
Equivalently, the first-order signed mismatch and the higher-order shell remainder are both smaller than the total first-order shell load.
\end{definition}

\begin{theorem}[Working shell-balance reduction and exact cancellation corridor]
Assume the first-order canonical shell law. Then:
\begin{enumerate}[leftmargin=2em,label=(\roman*)]
\item one has the exact reduction
\[
  r_{\mathrm{obs}}(\theta)=\chi_{\mathrm{sh}}(\theta)+R_{\ge 2}(\theta);
\]
\item in the cancellation corridor one has
\[
  |r_{\mathrm{obs}}(\theta)|
  \le (\varepsilon_{\mathrm{bal}}+\varepsilon_{\mathrm{rem}})\,\Lambda_{\mathrm{lin}}(\theta),
\]
so exact visible closure is reduced to a small signed shell-balance mismatch rather than to two separate closures;
\item if $p_{\mathrm{sh}}(\theta)$ and $e_{\mathrm{sh}}(\theta)$ have the same sign, then
\[
  |\chi_{\mathrm{sh}}(\theta)|=|p_{\mathrm{sh}}(\theta)|+|e_{\mathrm{sh}}(\theta)|=\Lambda_{\mathrm{lin}}(\theta),
\]
so exact visible closure at first order is impossible unless the total first-order shell load itself vanishes.
\end{enumerate}
\end{theorem}

\begin{proof}
By definition of the signed shell coordinates,
\[
  p_{\mathrm{sh}}(\theta)+e_{\mathrm{sh}}(\theta)=A_{\mathrm{sh}}(\theta)\,\delta_{\mathrm{obs}}(\theta)+B_{\mathrm{sh}}(\theta)\,\eta_{\mathrm{exp}}(\theta)=\chi_{\mathrm{sh}}(\theta),
\]
so the first-order shell law becomes
\[
  r_{\mathrm{obs}}(\theta)=\chi_{\mathrm{sh}}(\theta)+R_{\ge 2}(\theta).
\]
This proves (i). If the shell law lies in the cancellation corridor, then
\[
  |r_{\mathrm{obs}}(\theta)|\le |\chi_{\mathrm{sh}}(\theta)|+|R_{\ge 2}(\theta)|
  \le (\varepsilon_{\mathrm{bal}}+\varepsilon_{\mathrm{rem}})\,\Lambda_{\mathrm{lin}}(\theta),
\]
which proves (ii). If $p_{\mathrm{sh}}(\theta)$ and $e_{\mathrm{sh}}(\theta)$ have the same sign, then they add without cancellation, hence
\[
  |\chi_{\mathrm{sh}}(\theta)|=|p_{\mathrm{sh}}(\theta)+e_{\mathrm{sh}}(\theta)|=|p_{\mathrm{sh}}(\theta)|+|e_{\mathrm{sh}}(\theta)|=\Lambda_{\mathrm{lin}}(\theta),
\]
which proves (iii).
\end{proof}

\begin{corollary}[Working next decision after shell-load dominance]
After the shell-load dominance classification, the remaining visible decision is whether the signed ratio $\zeta_{\mathrm{sh}}(\theta)$ enters a neighborhood of $-1$. In particular:
\begin{enumerate}[leftmargin=2em,label=(\roman*)]
\item if $\zeta_{\mathrm{sh}}(\theta)\approx -1$, the visible shell mode lies near the cancellation corridor and exact visible closure remains possible;
\item if $\zeta_{\mathrm{sh}}(\theta)>0$, then the first-order shell law is sign-coherent and forces a canonical nonzero observer-shell remainder unless the total first-order shell load itself vanishes;
\item if $\zeta_{\mathrm{sh}}(\theta)$ stays bounded away from $-1$, then the shared shell mode is canonically nonzero already at first order.
\end{enumerate}
\end{corollary}

\begin{proof}
The previous theorem shows that the visible shell mode is controlled by the signed mismatch $\chi_{\mathrm{sh}}(\theta)=p_{\mathrm{sh}}(\theta)(1+\zeta_{\mathrm{sh}}(\theta))$ together with the higher-order remainder. Thus entry into the cancellation corridor is equivalent to placing the signed ratio near $-1$, while positive ratio excludes first-order cancellation.
\end{proof}

\begin{remark}[Working interpretation of the present shell-law reduction]
This is the first place where the two-shell-defect picture collapses all the way to a single signed comparison. The visible question is no longer which of two unfinished closures dominates, but whether the phase-shift and expansion-drift contributions cancel each other or reinforce each other on the observer shell. If they reinforce, the visible residual is canonically nonzero; if they nearly cancel, exact visible closure remains structurally possible.
\end{remark}


\subsection*{Normalized shell-cancellation defect and the exact first-order verdict}
The signed shell-balance mismatch already reduces the visible problem to a single first-order comparison. The next structural compression is to normalize that mismatch by the total first-order shell load itself. This removes the overall shell scale and leaves only the genuinely visible cancellation defect.

\begin{definition}[Working normalized shell defect and cancellation depth]
Whenever $\Lambda_{\mathrm{lin}}(\theta)>0$, define the normalized shell-cancellation defect by
\[
  \rho_{\mathrm{sh}}(\theta):=\frac{\chi_{\mathrm{sh}}(\theta)}{\Lambda_{\mathrm{lin}}(\theta)}
  =\frac{p_{\mathrm{sh}}(\theta)+e_{\mathrm{sh}}(\theta)}{|p_{\mathrm{sh}}(\theta)|+|e_{\mathrm{sh}}(\theta)|}.
\]
Set $\rho_{\mathrm{sh}}(\theta):=0$ when $\Lambda_{\mathrm{lin}}(\theta)=0$. Define the cancellation depth by
\[
  \kappa_{\mathrm{sh}}(\theta):=1-|\rho_{\mathrm{sh}}(\theta)|.
\]
Thus $|\rho_{\mathrm{sh}}|$ measures the normalized first-order noncancellation, while $\kappa_{\mathrm{sh}}$ measures how deeply the shell law enters the first-order cancellation corridor.
\end{definition}

\begin{theorem}[Working normalized shell defect, cancellation depth, and exact first-order criterion]
For every parameter value $\theta$ one has:
\begin{enumerate}[leftmargin=2em,label=(\roman*)]
\item
\[
  |\rho_{\mathrm{sh}}(\theta)|\le 1,
  \qquad
  0\le \kappa_{\mathrm{sh}}(\theta)\le 1;
\]
\item if $p_{\mathrm{sh}}(\theta)$ and $e_{\mathrm{sh}}(\theta)$ have the same sign or one of them vanishes, then
\[
  |\rho_{\mathrm{sh}}(\theta)|=1,
  \qquad
  \kappa_{\mathrm{sh}}(\theta)=0;
\]
\item if $p_{\mathrm{sh}}(\theta)$ and $e_{\mathrm{sh}}(\theta)$ have opposite sign, then
\[
  |\rho_{\mathrm{sh}}(\theta)|
  =\frac{\bigl||p_{\mathrm{sh}}(\theta)|-|e_{\mathrm{sh}}(\theta)|\bigr|}{|p_{\mathrm{sh}}(\theta)|+|e_{\mathrm{sh}}(\theta)|},
\]
\[
  \kappa_{\mathrm{sh}}(\theta)
  =\frac{2\min\{|p_{\mathrm{sh}}(\theta)|,|e_{\mathrm{sh}}(\theta)|\}}{|p_{\mathrm{sh}}(\theta)|+|e_{\mathrm{sh}}(\theta)|};
\]
\item exact first-order shell cancellation holds if and only if
\[
  \rho_{\mathrm{sh}}(\theta)=0,
\]
which is equivalent to opposite sign and equal first-order shell magnitudes,
\[
  |p_{\mathrm{sh}}(\theta)|=|e_{\mathrm{sh}}(\theta)|.
\]
\end{enumerate}
\end{theorem}

\begin{proof}
By definition,
\[
  |\chi_{\mathrm{sh}}(\theta)|=|p_{\mathrm{sh}}(\theta)+e_{\mathrm{sh}}(\theta)|
  \le |p_{\mathrm{sh}}(\theta)|+|e_{\mathrm{sh}}(\theta)|
  =\Lambda_{\mathrm{lin}}(\theta),
\]
so $|\rho_{\mathrm{sh}}(\theta)|\le 1$, and the cancellation depth automatically lies in $[0,1]$. This proves (i). If $p_{\mathrm{sh}}$ and $e_{\mathrm{sh}}$ have the same sign, or one vanishes, then
\[
  |p_{\mathrm{sh}}+e_{\mathrm{sh}}|=|p_{\mathrm{sh}}|+|e_{\mathrm{sh}}|,
\]
so $|\rho_{\mathrm{sh}}|=1$ and therefore $\kappa_{\mathrm{sh}}=0$, proving (ii). If the signs are opposite, then
\[
  |p_{\mathrm{sh}}+e_{\mathrm{sh}}|=\bigl||p_{\mathrm{sh}}|-|e_{\mathrm{sh}}|\bigr|,
\]
which yields the displayed formula for $|\rho_{\mathrm{sh}}|$ after division by $|p_{\mathrm{sh}}|+|e_{\mathrm{sh}}|$. The formula for $\kappa_{\mathrm{sh}}$ follows because
\[
  1-\frac{\bigl||p|-|e|\bigr|}{|p|+|e|}=\frac{|p|+|e|-\bigl||p|-|e|\bigr|}{|p|+|e|}=\frac{2\min\{|p|,|e|\}}{|p|+|e|}.
\]
This proves (iii). Finally, $\rho_{\mathrm{sh}}(\theta)=0$ is equivalent to $\chi_{\mathrm{sh}}(\theta)=0$, which is exactly first-order shell cancellation. Under opposite sign this is equivalent to equal magnitudes, proving (iv).
\end{proof}

\begin{definition}[Working normalized higher-order shell remainder]
Whenever $\Lambda_{\mathrm{lin}}(\theta)>0$, define the normalized higher-order shell remainder by
\[
  \rho_{\mathrm{rem}}(\theta):=\frac{|R_{\ge 2}(\theta)|}{\Lambda_{\mathrm{lin}}(\theta)}.
\]
Set $\rho_{\mathrm{rem}}(\theta):=0$ when $\Lambda_{\mathrm{lin}}(\theta)=0$.
\end{definition}

\begin{corollary}[Working exact visible closure requires normalized first-order cancellation]
Assume the first-order canonical shell law and let $\Lambda_{\mathrm{lin}}(\theta)>0$. Then:
\begin{enumerate}[leftmargin=2em,label=(\roman*)]
\item
\[
  \frac{|r_{\mathrm{obs}}(\theta)|}{\Lambda_{\mathrm{lin}}(\theta)}
  \le |\rho_{\mathrm{sh}}(\theta)|+\rho_{\mathrm{rem}}(\theta);
\]
\item a necessary condition for exact visible closure at $\theta$ is
\[
  |\rho_{\mathrm{sh}}(\theta)|\le \rho_{\mathrm{rem}}(\theta);
\]
\item if
\[
  |\rho_{\mathrm{sh}}(\theta)|>\rho_{\mathrm{rem}}(\theta),
\]
then the shared visible shell mode is canonically nonzero at $\theta$.
\end{enumerate}
\end{corollary}

\begin{proof}
By the shell-balance reduction theorem,
\[
  r_{\mathrm{obs}}(\theta)=\chi_{\mathrm{sh}}(\theta)+R_{\ge 2}(\theta),
\]
so after division by $\Lambda_{\mathrm{lin}}(\theta)$ one gets
\[
  \frac{|r_{\mathrm{obs}}(\theta)|}{\Lambda_{\mathrm{lin}}(\theta)}
  \le \frac{|\chi_{\mathrm{sh}}(\theta)|}{\Lambda_{\mathrm{lin}}(\theta)}+\frac{|R_{\ge 2}(\theta)|}{\Lambda_{\mathrm{lin}}(\theta)}
  =|\rho_{\mathrm{sh}}(\theta)|+\rho_{\mathrm{rem}}(\theta),
\]
which proves (i). If exact visible closure holds, then
\[
  0=r_{\mathrm{obs}}(\theta)=\chi_{\mathrm{sh}}(\theta)+R_{\ge 2}(\theta),
\]
so $|\chi_{\mathrm{sh}}(\theta)|=|R_{\ge 2}(\theta)|$. Division by $\Lambda_{\mathrm{lin}}(\theta)$ gives (ii). Statement (iii) is the contrapositive of (ii).
\end{proof}

\begin{remark}[Working interpretation of the normalized shell verdict]
The visible shell problem has now collapsed to a single normalized comparison. The first-order part is measured by $|\rho_{\mathrm{sh}}|$, the higher-order rescue budget by $\rho_{\mathrm{rem}}$, and exact visible closure is possible only if the normalized first-order noncancellation does not exceed the normalized higher-order shell remainder. In particular, once $|\rho_{\mathrm{sh}}|$ is visibly larger than $\rho_{\mathrm{rem}}$, the shared shell mode is forced to remain canonically nonzero.
\end{remark}


\subsection*{Sign-coherent obstruction and the anti-phase closure corridor}
The normalized shell verdict can be compressed one step further. The visible question is not merely whether the normalized first-order mismatch is larger than the higher-order rescue budget, but also whether the two first-order shell loads are sign-coherent or anti-phase. In the sign-coherent regime, first-order cancellation is impossible; in the anti-phase regime, exact visible closure can survive only inside a sharp corridor around the ratio value $-1$.

\begin{theorem}[Working sign-coherent obstruction and anti-phase corridor for visible closure]
Assume the first-order canonical shell law and let $\Lambda_{\mathrm{lin}}(\theta)>0$.
\begin{enumerate}[leftmargin=2em,label=(\roman*)]
\item If $p_{\mathrm{sh}}(\theta)$ and $e_{\mathrm{sh}}(\theta)$ have the same sign or one of them vanishes, then
\[
  |\rho_{\mathrm{sh}}(\theta)|=1.
\]
Consequently, if
\[
  \rho_{\mathrm{rem}}(\theta)<1,
\]
then exact visible closure at $\theta$ is impossible and the shared shell mode is canonically nonzero.
\item Suppose $p_{\mathrm{sh}}(\theta)$ and $e_{\mathrm{sh}}(\theta)$ have opposite sign and $p_{\mathrm{sh}}(\theta)\neq 0$. Set
\[
  \zeta_{\mathrm{sh}}(\theta):=\frac{e_{\mathrm{sh}}(\theta)}{p_{\mathrm{sh}}(\theta)}.
\]
Then necessarily $\zeta_{\mathrm{sh}}(\theta)<0$, and one has
\[
  |\rho_{\mathrm{sh}}(\theta)|
  =\frac{|1+\zeta_{\mathrm{sh}}(\theta)|}{1+|\zeta_{\mathrm{sh}}(\theta)|}.
\]
Hence a necessary condition for exact visible closure is
\[
  \frac{|1+\zeta_{\mathrm{sh}}(\theta)|}{1+|\zeta_{\mathrm{sh}}(\theta)|}
  \le
  \rho_{\mathrm{rem}}(\theta),
\]
that is,
\[
  |1+\zeta_{\mathrm{sh}}(\theta)|
  \le
  \rho_{\mathrm{rem}}(\theta)\bigl(1+|\zeta_{\mathrm{sh}}(\theta)|\bigr).
\]
\item In particular, if
\[
  |1+\zeta_{\mathrm{sh}}(\theta)|
  >
  \rho_{\mathrm{rem}}(\theta)\bigl(1+|\zeta_{\mathrm{sh}}(\theta)|\bigr),
\]
then the shared visible shell mode is canonically nonzero at $\theta$.
\end{enumerate}
\end{theorem}

\begin{proof}
Statement (i) is immediate from the normalized shell-defect theorem, which gives $|\rho_{\mathrm{sh}}|=1$ in the sign-coherent regime. The exact visible closure criterion then yields the necessary inequality
\[
  1=|\rho_{\mathrm{sh}}(\theta)|\le \rho_{\mathrm{rem}}(\theta).
\]
Therefore, if $\rho_{\mathrm{rem}}(\theta)<1$, exact visible closure is impossible and the shell mode remains canonically nonzero.

For (ii), opposite sign implies $\zeta_{\mathrm{sh}}(\theta)<0$. Since
\[
  \chi_{\mathrm{sh}}(\theta)=p_{\mathrm{sh}}(\theta)+e_{\mathrm{sh}}(\theta)=p_{\mathrm{sh}}(\theta)(1+\zeta_{\mathrm{sh}}(\theta)),
\]
and
\[
  \Lambda_{\mathrm{lin}}(\theta)=|p_{\mathrm{sh}}(\theta)|+|e_{\mathrm{sh}}(\theta)|=|p_{\mathrm{sh}}(\theta)|\bigl(1+|\zeta_{\mathrm{sh}}(\theta)|\bigr),
\]
one gets
\[
  |\rho_{\mathrm{sh}}(\theta)|
  =\frac{|\chi_{\mathrm{sh}}(\theta)|}{\Lambda_{\mathrm{lin}}(\theta)}
  =\frac{|1+\zeta_{\mathrm{sh}}(\theta)|}{1+|\zeta_{\mathrm{sh}}(\theta)|}.
\]
The displayed corridor condition then follows from the exact visible closure criterion $|\rho_{\mathrm{sh}}|\le \rho_{\mathrm{rem}}$.

Statement (iii) is the contrapositive of that necessary condition.
\end{proof}

\begin{remark}[Working interpretation of the sign/corridor shell verdict]
The visible shell problem has now become extremely sharp. If the phase and expansion shell loads are sign-coherent, then first-order cancellation is impossible and higher order would have to rescue the full shell scale to produce exact visible closure. If they are anti-phase, then exact visible closure can survive only inside a narrow corridor where the signed shell ratio $\zeta_{\mathrm{sh}}$ sits sufficiently close to $-1$. Outside that corridor, the shared shell mode is forced to remain canonically nonzero.
\end{remark}



\subsection*{Single-shell-coordinate reduction and naturality of the anti-phase corridor}
The anti-phase corridor around $\zeta_{\mathrm{sh}}=-1$ is sharp, but it is not yet clear whether that corridor is structurally natural or only a fine-tuned accident. The next structural step is therefore to ask whether the two downstream shell defects are actually independent, or whether they are both induced by one and the same observer-shell displacement. If the latter is true, then the signed shell ratio collapses at first order to one structural gain ratio, and the question ``why near $-1$?'' becomes a question about one induced shell geometry rather than about two unrelated corrections.

\begin{definition}[Working single-shell-coordinate observer law]
A downstream observer shell is said to be controlled by a \emph{single shell coordinate} if there exists a scalar family $\xi_{\mathrm{sh}}(\theta)$ and coefficient families $a_{\mathrm{sh}}(\theta), b_{\mathrm{sh}}(\theta)$ such that
\[
  \delta_{\mathrm{obs}}(\theta)
  =a_{\mathrm{sh}}(\theta)\,\xi_{\mathrm{sh}}(\theta)+r_{\delta}(\theta),
  \qquad
  \eta_{\mathrm{exp}}(\theta)
  =b_{\mathrm{sh}}(\theta)\,\xi_{\mathrm{sh}}(\theta)+r_{\eta}(\theta),
\]
with quadratic remainder bounds
\[
  |r_{\delta}(\theta)|+|r_{\eta}(\theta)|
  \le C_{\xi}(\theta)\,|\xi_{\mathrm{sh}}(\theta)|^2.
\]
Thus both downstream shell defects are first-order images of the same observer-shell displacement variable.
\end{definition}

\begin{definition}[Working shell gains induced by one shell coordinate]
Under the single-shell-coordinate observer law, define the induced first-order shell gains
\[
  \Gamma_{\mathrm{ph}}(\theta):=A_{\mathrm{sh}}(\theta)\,a_{\mathrm{sh}}(\theta),
  \qquad
  \Gamma_{\mathrm{exp}}(\theta):=B_{\mathrm{sh}}(\theta)\,b_{\mathrm{sh}}(\theta),
\]
and, whenever $\Gamma_{\mathrm{ph}}(\theta)\neq 0$, the induced shell gain ratio
\[
  \zeta_{\Gamma}(\theta):=\frac{\Gamma_{\mathrm{exp}}(\theta)}{\Gamma_{\mathrm{ph}}(\theta)}.
\]
\end{definition}

\begin{theorem}[Working first-order shell ratio collapses to an induced gain ratio]
Assume the first-order canonical shell law and the single-shell-coordinate observer law. Then the signed first-order shell coordinates satisfy
\[
  p_{\mathrm{sh}}(\theta)=\Gamma_{\mathrm{ph}}(\theta)\,\xi_{\mathrm{sh}}(\theta)+O\!\left(|\xi_{\mathrm{sh}}(\theta)|^2\right),
\]
\[
  e_{\mathrm{sh}}(\theta)=\Gamma_{\mathrm{exp}}(\theta)\,\xi_{\mathrm{sh}}(\theta)+O\!\left(|\xi_{\mathrm{sh}}(\theta)|^2\right),
\]
and hence
\[
  \chi_{\mathrm{sh}}(\theta)
  =\bigl(\Gamma_{\mathrm{ph}}(\theta)+\Gamma_{\mathrm{exp}}(\theta)\bigr)\,\xi_{\mathrm{sh}}(\theta)
  +O\!\left(|\xi_{\mathrm{sh}}(\theta)|^2\right).
\]
If moreover $\Gamma_{\mathrm{ph}}(\theta)\neq 0$ and $\xi_{\mathrm{sh}}(\theta)\neq 0$, then
\[
  \zeta_{\mathrm{sh}}(\theta)
  =\zeta_{\Gamma}(\theta)+O\!\left(|\xi_{\mathrm{sh}}(\theta)|\right).
\]
In particular, the anti-phase corridor is first-order natural exactly when the induced gain ratio satisfies
\[
  \zeta_{\Gamma}(\theta)\approx -1.
\]
\end{theorem}

\begin{proof}
Insert the single-shell-coordinate expansions into the definitions
\[
  p_{\mathrm{sh}}(\theta)=A_{\mathrm{sh}}(\theta)\,\delta_{\mathrm{obs}}(\theta),
  \qquad
  e_{\mathrm{sh}}(\theta)=B_{\mathrm{sh}}(\theta)\,\eta_{\mathrm{exp}}(\theta).
\]
This gives
\[
  p_{\mathrm{sh}}(\theta)
  =A_{\mathrm{sh}}(\theta)a_{\mathrm{sh}}(\theta)\,\xi_{\mathrm{sh}}(\theta)
   +A_{\mathrm{sh}}(\theta)r_{\delta}(\theta)
  =\Gamma_{\mathrm{ph}}(\theta)\,\xi_{\mathrm{sh}}(\theta)+O\!\left(|\xi_{\mathrm{sh}}(\theta)|^2\right),
\]
with the analogous formula for $e_{\mathrm{sh}}(\theta)$. Summing yields the formula for $\chi_{\mathrm{sh}}(\theta)$. If $\Gamma_{\mathrm{ph}}(\theta)\neq 0$ and $\xi_{\mathrm{sh}}(\theta)\neq 0$, then division of the two first-order expansions gives
\[
  \zeta_{\mathrm{sh}}(\theta)
  =\frac{\Gamma_{\mathrm{exp}}(\theta)+O(|\xi_{\mathrm{sh}}(\theta)|)}{\Gamma_{\mathrm{ph}}(\theta)+O(|\xi_{\mathrm{sh}}(\theta)|)}
  =\frac{\Gamma_{\mathrm{exp}}(\theta)}{\Gamma_{\mathrm{ph}}(\theta)}+O\!\left(|\xi_{\mathrm{sh}}(\theta)|\right),
\]
which is exactly the stated collapse to the induced gain ratio.
\end{proof}

\begin{corollary}[Working structural meaning of the anti-phase corridor]
Under the same hypotheses, the anti-phase corridor around $\zeta_{\mathrm{sh}}(\theta)=-1$ is not a separate tuning problem between two independent shell defects. It is, to first order, the statement that the single observer-shell displacement variable $\xi_{\mathrm{sh}}(\theta)$ is seen by the phase and expansion readouts with opposite and nearly equal induced gains:
\[
  \Gamma_{\mathrm{exp}}(\theta)\approx -\Gamma_{\mathrm{ph}}(\theta).
\]
Thus the visible shell question reduces further to deciding whether the induced gain pair is
\begin{enumerate}[leftmargin=2em,label=(\roman*)]
\item sign-coherent,
\item anti-phase but unbalanced,
\item or anti-phase and nearly balanced.
\end{enumerate}
\end{corollary}

\begin{proof}
This is exactly the content of the theorem after replacing $\zeta_{\mathrm{sh}}(\theta)$ by its first-order approximation $\zeta_{\Gamma}(\theta)$.
\end{proof}

\begin{remark}[Working next decision after the shell-ratio corridor]
The visible shell problem has now moved from two apparent shell defects to one observer-shell displacement and one induced gain pair. This is a genuine structural compression. It means that the relevant question is no longer whether two unrelated corrections accidentally cancel, but whether one downstream shell displacement is read with opposite and nearly equal gains by the phase and expansion channels.
\end{remark}


\subsection*{Shell-conjugate observer geometry and induced gain balance}
The single-shell-coordinate reduction showed that the anti-phase corridor is natural exactly when the induced gains satisfy $\Gamma_{\mathrm{exp}}\approx-\Gamma_{\mathrm{ph}}$. The next structural step is to push that gain relation back to the observer-shell architecture itself. Rather than treating $\Gamma_{\mathrm{ph}}$ and $\Gamma_{\mathrm{exp}}$ as free coefficients, we ask whether they arise as two downstream readouts of one shell tangent under a shell involution inherited from the umkehrwalze. If so, then anti-phase is no longer a tuning accident between two unrelated gains, but the first-order signature of one shell-conjugate observer geometry.

\begin{definition}[Working shell-conjugate observer response]
Let $S_{\mathrm{obs}}$ be the downstream observer-shell tangent space. A \emph{shell-conjugate observer response} consists of
\begin{enumerate}[leftmargin=2em,label=(\roman*)]
\item an involution $R_{\mathrm{sh}}:S_{\mathrm{obs}}\to S_{\mathrm{obs}}$ with $R_{\mathrm{sh}}^2=\mathrm{id}$, interpreted as the observer-shell descendant of the umkehrwalze,
\item a shell tangent family $u_{\mathrm{sh}}(\theta)\in S_{\mathrm{obs}}$,
\item a base shell readout functional $\ell_{\mathrm{sh}}(\theta)\in S_{\mathrm{obs}}^*$,
\item and a shell-conjugacy defect $\varepsilon_{\mathrm{conj}}(\theta)$,
\end{enumerate}
with
\[
  a_{\mathrm{sh}}(\theta)=\ell_{\mathrm{sh}}(\theta)\bigl(u_{\mathrm{sh}}(\theta)\bigr),
\]
\[
  b_{\mathrm{sh}}(\theta)=-\ell_{\mathrm{sh}}(\theta)\bigl(R_{\mathrm{sh}}u_{\mathrm{sh}}(\theta)\bigr)+\varepsilon_{\mathrm{conj}}(\theta).
\]
If moreover
\[
  R_{\mathrm{sh}}u_{\mathrm{sh}}(\theta)=u_{\mathrm{sh}}(\theta),
\]
then the phase and expansion shell readouts are said to be read from the same shell tangent with conjugate sign, up to the defect $\varepsilon_{\mathrm{conj}}(\theta)$.
\end{definition}

\begin{definition}[Working transport-balance defect]
Define the downstream transport-balance defect by
\[
  \Delta_{AB}(\theta):=B_{\mathrm{sh}}(\theta)-A_{\mathrm{sh}}(\theta).
\]
Exact transport balance on the downstream shell means $\Delta_{AB}(\theta)=0$.
\end{definition}

\begin{theorem}[Working shell-conjugate observer geometry makes anti-phase natural]
Assume the single-shell-coordinate observer law and a shell-conjugate observer response with
\[
  R_{\mathrm{sh}}u_{\mathrm{sh}}(\theta)=u_{\mathrm{sh}}(\theta).
\]
Then
\[
  b_{\mathrm{sh}}(\theta)=-a_{\mathrm{sh}}(\theta)+\varepsilon_{\mathrm{conj}}(\theta),
\]
and hence the induced gains satisfy
\[
  \Gamma_{\mathrm{exp}}(\theta)+\Gamma_{\mathrm{ph}}(\theta)
  =A_{\mathrm{sh}}(\theta)\,\varepsilon_{\mathrm{conj}}(\theta)
   +\Delta_{AB}(\theta)\,b_{\mathrm{sh}}(\theta).
\]
In particular, if $\Gamma_{\mathrm{ph}}(\theta)\neq 0$, then
\[
  \zeta_{\Gamma}(\theta)+1
  =
  \frac{A_{\mathrm{sh}}(\theta)\,\varepsilon_{\mathrm{conj}}(\theta)
   +\Delta_{AB}(\theta)\,b_{\mathrm{sh}}(\theta)}{\Gamma_{\mathrm{ph}}(\theta)}.
\]
Thus the anti-phase corridor $\zeta_{\Gamma}(\theta)\approx -1$ is structurally natural whenever both the shell-conjugacy defect and the downstream transport-balance defect are small relative to $\Gamma_{\mathrm{ph}}(\theta)$.
\end{theorem}

\begin{proof}
By shell-conjugacy,
\[
  b_{\mathrm{sh}}(\theta)
  =-\ell_{\mathrm{sh}}(\theta)\bigl(R_{\mathrm{sh}}u_{\mathrm{sh}}(\theta)\bigr)+\varepsilon_{\mathrm{conj}}(\theta)
  =-\ell_{\mathrm{sh}}(\theta)\bigl(u_{\mathrm{sh}}(\theta)\bigr)+\varepsilon_{\mathrm{conj}}(\theta)
  =-a_{\mathrm{sh}}(\theta)+\varepsilon_{\mathrm{conj}}(\theta).
\]
Therefore
\[
  \Gamma_{\mathrm{exp}}(\theta)+\Gamma_{\mathrm{ph}}(\theta)
  =B_{\mathrm{sh}}(\theta)b_{\mathrm{sh}}(\theta)+A_{\mathrm{sh}}(\theta)a_{\mathrm{sh}}(\theta)
\]
\[
  =B_{\mathrm{sh}}(\theta)\bigl(-a_{\mathrm{sh}}(\theta)+\varepsilon_{\mathrm{conj}}(\theta)\bigr)+A_{\mathrm{sh}}(\theta)a_{\mathrm{sh}}(\theta)
\]
\[
  =(A_{\mathrm{sh}}(\theta)-B_{\mathrm{sh}}(\theta))a_{\mathrm{sh}}(\theta)+B_{\mathrm{sh}}(\theta)\varepsilon_{\mathrm{conj}}(\theta).
\]
Since $b_{\mathrm{sh}}=-a_{\mathrm{sh}}+\varepsilon_{\mathrm{conj}}$, one may rewrite this exactly as
\[
  \Gamma_{\mathrm{exp}}(\theta)+\Gamma_{\mathrm{ph}}(\theta)
  =A_{\mathrm{sh}}(\theta)\,\varepsilon_{\mathrm{conj}}(\theta)+\Delta_{AB}(\theta)b_{\mathrm{sh}}(\theta).
\]
Division by $\Gamma_{\mathrm{ph}}(\theta)$ gives the formula for $\zeta_{\Gamma}(\theta)+1$.
\end{proof}

\begin{corollary}[Working exact shell conjugacy and balanced transport force first-order anti-phase]
Under the same hypotheses, if
\[
  \varepsilon_{\mathrm{conj}}(\theta)=0,
  \qquad
  \Delta_{AB}(\theta)=0,
\]
then
\[
  \Gamma_{\mathrm{exp}}(\theta)=-\Gamma_{\mathrm{ph}}(\theta),
  \qquad
  \zeta_{\Gamma}(\theta)=-1.
\]
Hence the anti-phase corridor is then not merely admissible but exact at first order.
\end{corollary}

\begin{proof}
This is the immediate specialization of the theorem to vanishing shell-conjugacy and transport-balance defects.
\end{proof}

\begin{remark}[Working structural reading of the shell-gain pair]
The gain pair $(\Gamma_{\mathrm{ph}},\Gamma_{\mathrm{exp}})$ is no longer best read as two unrelated coefficients. In the shell-conjugate regime it is the image of one shell tangent under one observer-shell involution and one downstream transport imbalance. Therefore the real structural question is not ``why should two gains accidentally cancel?'', but rather ``how far is the shell-conjugate observer geometry from exact sign-conjugacy and balanced transport?''
\end{remark}


\subsection*{Shell anti-phase defect budget and epsilon-trap}
The shell-conjugate observer geometry showed that anti-phase is natural whenever the shell-conjugacy and transport-balance defects are small. The next step is to compress those two structural errors and the single-shell-coordinate remainder into one normalized anti-phase budget. This turns the visible shell question into an epsilon-trap statement: once the induced shell ratio enters a sufficiently small corridor around $-1$, the first-order visible mode cannot escape that corridor unless the higher-order shell remainder is of comparable size.

\begin{definition}[Working normalized shell anti-phase budget]
Assume $\Gamma_{\mathrm{ph}}(\theta)\neq 0$. Define the normalized shell-conjugacy, transport-balance, and single-shell linearization loads by
\[
  \rho_{\mathrm{conj}}(\theta)
  :=
  \frac{|A_{\mathrm{sh}}(\theta)\,\varepsilon_{\mathrm{conj}}(\theta)|}{|\Gamma_{\mathrm{ph}}(\theta)|},
\]
\[
  \rho_{AB}(\theta)
  :=
  \frac{|\Delta_{AB}(\theta)\,b_{\mathrm{sh}}(\theta)|}{|\Gamma_{\mathrm{ph}}(\theta)|},
\]
\[
  \rho_{\xi}(\theta)
  :=
  \left|\frac{\zeta_{\mathrm{sh}}(\theta)-\zeta_{\Gamma}(\theta)}{1+|\zeta_{\Gamma}(\theta)|}\right|.
\]
The total normalized shell anti-phase budget is
\[
  \rho_{\mathrm{trap}}(\theta)
  :=
  \rho_{\mathrm{conj}}(\theta)+\rho_{AB}(\theta)+\rho_{\xi}(\theta).
\]
\end{definition}

\begin{theorem}[Working shell anti-phase epsilon-trap]
\label{thm:working-shell-anti-phase-epsilon-trap}
Assume the single-shell-coordinate observer law and shell-conjugate observer geometry with
\[
  R_{\mathrm{sh}}u_{\mathrm{sh}}(\theta)=u_{\mathrm{sh}}(\theta),
  \qquad
  \Gamma_{\mathrm{ph}}(\theta)\neq 0.
\]
Then
\[
  |\zeta_{\Gamma}(\theta)+1|
  \le
  \rho_{\mathrm{conj}}(\theta)+\rho_{AB}(\theta),
\]
and hence
\[
  |\zeta_{\mathrm{sh}}(\theta)+1|
  \le
  \rho_{\mathrm{trap}}(\theta)\,\bigl(1+|\zeta_{\Gamma}(\theta)|\bigr).
\]
In particular, if for some $\varepsilon_{\mathrm{trap}}>0$ one has
\[
  \rho_{\mathrm{trap}}(\theta)
  \le
  \frac{\varepsilon_{\mathrm{trap}}}{1+|\zeta_{\Gamma}(\theta)|},
\]
then
\[
  |\zeta_{\mathrm{sh}}(\theta)+1|\le \varepsilon_{\mathrm{trap}}.
\]
Thus the shell ratio is forced into an anti-phase epsilon-tube around $-1$ by one normalized structural budget.
\end{theorem}

\begin{remark}[Reading of the shell epsilon-trap after OP3 localization]
After the OP3 shell/residual theorem, the anti-phase epsilon-trap is best read as a non-escape criterion for the visible shell ratio inside one OP3-local corridor. By itself it does not yet prove global OP closure or global branch admissibility; rather, it says that once the shell-side refinement stays inside a sufficiently small anti-phase budget, the visible first-order shell mode cannot leave that corridor without a comparably sized higher-order defect.
\end{remark}

\begin{proof}
The shell-conjugate observer theorem gave
\[
  \zeta_{\Gamma}(\theta)+1
  =
  \frac{A_{\mathrm{sh}}(\theta)\,\varepsilon_{\mathrm{conj}}(\theta)
   +\Delta_{AB}(\theta)\,b_{\mathrm{sh}}(\theta)}{\Gamma_{\mathrm{ph}}(\theta)}.
\]
Taking absolute values yields
\[
  |\zeta_{\Gamma}(\theta)+1|
  \le
  \frac{|A_{\mathrm{sh}}(\theta)\,\varepsilon_{\mathrm{conj}}(\theta)|}{|\Gamma_{\mathrm{ph}}(\theta)|}
  +
  \frac{|\Delta_{AB}(\theta)\,b_{\mathrm{sh}}(\theta)|}{|\Gamma_{\mathrm{ph}}(\theta)|}
  =
  \rho_{\mathrm{conj}}(\theta)+\rho_{AB}(\theta).
\]
Now
\[
  |\zeta_{\mathrm{sh}}(\theta)+1|
  \le
  |\zeta_{\mathrm{sh}}(\theta)-\zeta_{\Gamma}(\theta)|
  +
  |\zeta_{\Gamma}(\theta)+1|.
\]
By the definition of $\rho_{\xi}(\theta)$,
\[
  |\zeta_{\mathrm{sh}}(\theta)-\zeta_{\Gamma}(\theta)|
  =
  \rho_{\xi}(\theta)\,\bigl(1+|\zeta_{\Gamma}(\theta)|\bigr).
\]
Since $1\le 1+|\zeta_{\Gamma}(\theta)|$, the previous bound implies
\[
  |\zeta_{\Gamma}(\theta)+1|
  \le
  \bigl(\rho_{\mathrm{conj}}(\theta)+\rho_{AB}(\theta)\bigr)
  \bigl(1+|\zeta_{\Gamma}(\theta)|\bigr).
\]
Combining the two estimates gives
\[
  |\zeta_{\mathrm{sh}}(\theta)+1|
  \le
  \rho_{\mathrm{trap}}(\theta)\,\bigl(1+|\zeta_{\Gamma}(\theta)|\bigr).
\]
The epsilon-tube statement is immediate.
\end{proof}

\begin{corollary}[Working exact visible closure reduces to one shell-budget comparison]
Assume the anti-phase sign regime and $\Gamma_{\mathrm{ph}}(\theta)\neq 0$. If exact visible closure holds, then necessarily
\[
  \rho_{\mathrm{trap}}(\theta)
  \le
  \frac{\rho_{\mathrm{rem}}(\theta)
  \bigl(1+|\zeta_{\mathrm{sh}}(\theta)|\bigr)}{1+|\zeta_{\Gamma}(\theta)|}.
\]
Conversely, if
\[
  \rho_{\mathrm{trap}}(\theta)
  >
  \frac{\rho_{\mathrm{rem}}(\theta)
  \bigl(1+|\zeta_{\mathrm{sh}}(\theta)|\bigr)}{1+|\zeta_{\Gamma}(\theta)|},
\]
then exact visible closure is impossible and the shared shell mode is canonically nonzero.
\end{corollary}

\begin{proof}
The anti-phase corridor theorem for visible closure gave the necessary condition
\[
  |1+\zeta_{\mathrm{sh}}(\theta)|
  \le
  \rho_{\mathrm{rem}}(\theta)\,\bigl(1+|\zeta_{\mathrm{sh}}(\theta)|\bigr).
\]
The shell anti-phase epsilon-trap gives
\[
  |1+\zeta_{\mathrm{sh}}(\theta)|
  \le
  \rho_{\mathrm{trap}}(\theta)\,\bigl(1+|\zeta_{\Gamma}(\theta)|\bigr).
\]
If exact visible closure is to be possible, the first-order shell mismatch forced by the second estimate must fit inside the higher-order rescue corridor from the first. This yields the displayed necessary inequality. Its negation rules out exact visible closure, hence forces a canonically nonzero shared shell mode.
\end{proof}

\begin{remark}[Working structural reading of the shell anti-phase budget]
The anti-phase corridor is no longer best read as an unexplained near-accident between two signed shell loads. After the shell-conjugate reduction, the visible first-order defect is controlled by one budget only: shell-conjugacy error, transport-balance error, and the residual error from replacing the induced shell ratio by the gain ratio. Therefore the next structural question is not ``can phase and expansion accidentally cancel?'', but rather ``is the combined shell anti-phase budget genuinely smaller than the higher-order rescue scale?''
\end{remark}



\subsection*{3.10.31 Shell-balance driver as an architecture defect}
The remaining first-order shell driver should not be treated as an independent phenomenological parameter. It is more natural to regard it as the visible shadow of a joint architecture defect involving the active HC/MM/J fit on the admissible branch together with the observer-shell dressing of the same branch.

\begin{definition}[Working architecture-side shell driver loads]
Set
\[
  \lambda_{\mathrm{fit}}(\theta)
  :=
  \frac{\|[(\mathcal M,J)]|_{U}\|\,\|x(\theta)\|}{|\Gamma_{\mathrm{ph}}(\theta)|},
  \qquad
  \lambda_{\mathrm{obs}}(\theta):=\lambda_{\mathrm{sh}}(\theta).
\]
Here \(\lambda_{\mathrm{fit}}\) measures active architecture mismatch through the restricted commutator, while \(\lambda_{\mathrm{obs}}\) is the normalized observer-shell anti-phase load already extracted above.
\end{definition}

\begin{theorem}[Working shell-balance driver is downstream from HC/MM/J fit plus shell dressing]
Assume the active-branch setup from the previous section, with visible shared shell mode
\[
  r_{\mathrm{obs}}(\theta)
  =
  \Xi_{\mathrm{bal}}(\theta)\,\xi_{\mathrm{sh}}(\theta)
  + O\!\left(|\xi_{\mathrm{sh}}(\theta)|^2\right).
\]
Then there is a first-order architecture estimate of the form
\[
  \frac{|\Xi_{\mathrm{bal}}(\theta)|}{|\Gamma_{\mathrm{ph}}(\theta)|}
  \le
  \lambda_{\mathrm{fit}}(\theta)+\lambda_{\mathrm{obs}}(\theta).
\]
In particular, if both
\[
  [\mathcal M,J]\big|_{U}=0
  \qquad\text{and}\qquad
  \lambda_{\mathrm{sh}}(\theta)=0,
\]
then
\[
  \Xi_{\mathrm{bal}}(\theta)=0,
\]
and the shared visible shell mode has no first-order driver.
\end{theorem}

\begin{proof}
By the exact commutator identity on the admissible branch, the primitive
common-mode carrier is controlled by
\[
  E_{\mathrm{fit}}^{(U)}
  = \frac12[\mathcal M,J]\big|_{U}.
\]
Hence every first-order visible common-mode contribution inherited from the
active transport branch is bounded, up to the visible readout coefficient
entering \(\Gamma_{\mathrm{ph}}\), by a multiple of
\(\|[\mathcal M,J]|_{U}\|\,\|x(\theta)\|\). Normalizing by
\(|\Gamma_{\mathrm{ph}}(\theta)|\) gives the \(\lambda_{\mathrm{fit}}\) term.

The remaining first-order imbalance after exact admissible transport was already compressed into the shell-conjugacy and transport-balance loads, combined in \(\lambda_{\mathrm{sh}}=\lambda_{\mathrm{conj}}+\lambda_{\mathrm{bal}}\). These are exactly the observer-shell contributions entering the shell-balance driver \(\Xi_{\mathrm{bal}}\). Adding the architecture-side and observer-side contributions yields
\[
  \frac{|\Xi_{\mathrm{bal}}(\theta)|}{|\Gamma_{\mathrm{ph}}(\theta)|}
  \le
  \lambda_{\mathrm{fit}}(\theta)+\lambda_{\mathrm{sh}}(\theta).
\]
If both terms vanish, then the normalized driver vanishes and therefore \(\Xi_{\mathrm{bal}}(\theta)=0\).
\end{proof}

\begin{corollary}[Working first-order visible cancellation is architecture-driven]
If
\[
  \lambda_{\mathrm{fit}}(\theta)+\lambda_{\mathrm{sh}}(\theta)<1,
\]
then the shared shell mode is forced into a first-order small-driver regime. If
\[
  \lambda_{\mathrm{fit}}(\theta)+\lambda_{\mathrm{sh}}(\theta)=0,
\]
then any remaining visible residual is necessarily of strictly higher observer-shell order.
\end{corollary}

\begin{proof}
The theorem gives
\[
  |\Xi_{\mathrm{bal}}(\theta)|
  \le
  |\Gamma_{\mathrm{ph}}(\theta)|\bigl(\lambda_{\mathrm{fit}}(\theta)+\lambda_{\mathrm{sh}}(\theta)\bigr).
\]
The first claim is the corresponding small-driver statement. If the sum vanishes, then \(\Xi_{\mathrm{bal}}(\theta)=0\), and the first-order shell law reduces to
\[
  r_{\mathrm{obs}}(\theta)=O\!\left(|\xi_{\mathrm{sh}}(\theta)|^2\right)+R_{\ge 2}(\theta),
\]
so every visible remainder is of higher order.
\end{proof}

\begin{remark}[Working architectural reading of the remaining shell law]
The visible common-mode residual is now best read as the downstream image of one combined architectural budget. The active HC/MM/J fit supplies the primitive common-mode leakage budget, while the shell-conjugacy and transport-balance terms supply the observer-side dressing budget. In this reading, the next real question is no longer whether OP~1 and OP~5 are separate closures, but whether the combined architecture budget vanishes exactly or leaves a canonical higher-order observer-shell remainder.
\end{remark}


\subsection*{3.10.32 Final architecture-versus-rescue budget corridor}
After the shell-balance driver has been pushed back to the combined architecture budget, the remaining visible question should no longer be read as an independent shell law. It is now a final corridor comparison: can the higher-order rescue sector absorb the first-order architecture load, or does a canonically nonzero visible shell mode necessarily survive?

\begin{definition}[Working normalized architecture budget and quadratic shell truncation]
Set
\[
  \rho_{\mathrm{arch}}(\theta)
  :=
  \frac{|\Gamma_{\mathrm{ph}}(\theta)\,\xi_{\mathrm{sh}}(\theta)|}{\Lambda_{\mathrm{lin}}(\theta)}
  \bigl(\lambda_{\mathrm{fit}}(\theta)+\lambda_{\mathrm{obs}}(\theta)\bigr),
\]
and define the normalized quadratic shell truncation term by
\[
  \rho_{\xi^2}(\theta)
  :=
  \frac{|O(|\xi_{\mathrm{sh}}(\theta)|^2)|}{\Lambda_{\mathrm{lin}}(\theta)}.
\]
The first quantity measures the visible first-order architecture load inherited from the combined HC/MM/J-fit plus observer-shell budget; the second measures the normalized first neglected shell order.
\end{definition}

\begin{theorem}[Working final architecture-rescue corridor under sharp active loading]
Assume the sharp active-loading regime on a window where
\[
  |\Gamma_{\mathrm{ph}}(\theta)\,\xi_{\mathrm{sh}}(\theta)|
  \ge c_0\,\Lambda_{\mathrm{lin}}(\theta)
\]
for some fixed constant \(c_0>0\). Then the normalized first-order visible shell defect satisfies
\[
  |\rho_{\mathrm{sh}}(\theta)|
  \ge c_0\,\rho_{\mathrm{arch}}(\theta)-\rho_{\xi^2}(\theta).
\]
Consequently, exact visible closure can hold only if
\[
  c_0\,\rho_{\mathrm{arch}}(\theta)
  \le
  \rho_{\mathrm{rem}}(\theta)+\rho_{\xi^2}(\theta).
\]
\end{theorem}

\begin{proof}
By the architecture-driven shell law from the previous subsection,
\[
  r_{\mathrm{obs}}(\theta)
  =
  \Xi_{\mathrm{bal}}(\theta)\,\xi_{\mathrm{sh}}(\theta)
  + O\!\left(|\xi_{\mathrm{sh}}(\theta)|^2\right)
  + R_{\ge 2}(\theta),
\]
with
\[
  \frac{|\Xi_{\mathrm{bal}}(\theta)|}{|\Gamma_{\mathrm{ph}}(\theta)|}
  \le
  \lambda_{\mathrm{fit}}(\theta)+\lambda_{\mathrm{obs}}(\theta).
\]
Dividing the first-order term by \(\Lambda_{\mathrm{lin}}(\theta)\) yields the normalized architecture load \(\rho_{\mathrm{arch}}(\theta)\). Under the sharp loading hypothesis, the visible contribution of this term is bounded below by \(c_0\rho_{\mathrm{arch}}(\theta)\) up to the normalized quadratic truncation error \(\rho_{\xi^2}(\theta)\). This proves
\[
  |\rho_{\mathrm{sh}}(\theta)|
  \ge c_0\,\rho_{\mathrm{arch}}(\theta)-\rho_{\xi^2}(\theta).
\]
If exact visible closure is to be possible, then by the earlier normalized shell-defect theorem one must have
\[
  |\rho_{\mathrm{sh}}(\theta)|\le \rho_{\mathrm{rem}}(\theta).
\]
Combining both inequalities gives the displayed necessary corridor condition.
\end{proof}

\begin{corollary}[Working final reduced decision after exact transport]
Under exact admissible transport and sharp active loading, the visible OP~1/OP~5 question is reduced to a single final budget comparison:
\[
  c_0\,\rho_{\mathrm{arch}}(\theta)
  \stackrel{?}{\le}
  \rho_{\mathrm{rem}}(\theta)+\rho_{\xi^2}(\theta).
\]
If the inequality fails, then the common visible shell mode is canonically nonzero. If it holds, then exact visible closure remains possible only through higher-order rescue of a first-order architecture load.
\end{corollary}

\begin{remark}[Working interpretation of the final corridor]
The remaining visible problem is now compressed as far as the current structure allows. It is no longer a pair of unfinished closures, nor even two shell loads. It is one corridor question: can the higher-order rescue sector absorb the normalized combined architecture load after exact transport, or does a canonical observer-shell remainder necessarily survive?
\end{remark}



\subsection{QAM-over-ZF: operator families and lambda-like structural maps}

We now close the first internal QAM package by adding an operator-family and lambda-like structural layer on top of the already introduced transport, branch/shell, tensor, and translation strata. The point remains conservative: no higher-order foundation is introduced. Instead, families of coded operators and lambda-like structural maps are themselves represented as set-codes over ZF and are required to preserve closure content on the admissible fragment.

\begin{definition}[Operator-family code]
A set
\[
\mathfrak{o}=\langle 9,I,F\rangle
\]
is called an \emph{operator-family code} if $I$ is an index code and $F$ is a graph-code assigning to each $i\in I$ a coded operator $o_i$.
\end{definition}

\begin{definition}[Admissible operator family]
An operator-family code
\[
\mathfrak{o}=\langle 9,I,F\rangle
\]
is called \emph{admissible} if every member operator $o_i$ coded by $F$ is admissible on its intended coded state domain. Equivalently, for every admissible coded source state $s$ and every coded target state $t$ linked by some member $o_i$, one has
\[
\mathrm{QAdmTransport}(o_i,s,t)
\Longrightarrow
\mathrm{QClEq}(s,t).
\]
\end{definition}

\begin{remark}
Thus admissibility of an operator family means pointwise closure preservation across the family. The family layer does not weaken the closure-first discipline already imposed on single coded operators.
\end{remark}

\begin{definition}[Lambda-like structural map code]
A set
\[
\lambda=\langle 10,D,C,L\rangle
\]
is called a \emph{lambda-like structural map code} if $D$ codes an intended domain class of QAM-coded objects, $C$ codes an intended codomain class, and $L$ is a graph-code assigning to each coded object in $D$ a coded object in $C$.
\end{definition}

\begin{definition}[Closure-compatible lambda map]
A lambda-like structural map code $\lambda$ is called \emph{closure-compatible} on a coded state domain if for all coded states $s,t$ in that domain,
\[
\mathrm{QClEq}(s,t)
\Longrightarrow
\mathrm{QClEq}(\lambda(s),\lambda(t)),
\]
whenever the values $\lambda(s)$ and $\lambda(t)$ are coded states.
\end{definition}

\begin{remark}
This is the lambda-level analogue of admissible transport. A closure-compatible lambda map need not itself be read as a transport along one fixed route; it may instead be a structural selector, reindexer, collector, or support-transforming map, provided it preserves closure class on the admissible fragment.
\end{remark}

\begin{definition}[Support-selection lambda]
A closure-compatible lambda-like structural map is called a \emph{support-selection lambda} if it acts on admissible branch-shell support codes and returns a coded sub-support, preferred support, or selected routed support without leaving the admissible closure regime.
\end{definition}

\begin{proposition}[Pointwise closure preservation for admissible operator families]
Let
\[
\mathfrak{o}=\langle 9,I,F\rangle
\]
be an admissible operator family. If $o_i$ is a member of the family and
\[
\mathrm{QAdmTransport}(o_i,s,t),
\]
then
\[
\mathrm{QClEq}(s,t).
\]
Consequently, every closure-compatible readout identifies the closure-relevant content of $s$ and $t$ after factorization.
\end{proposition}

\begin{proof}
This is immediate from pointwise admissibility of the family members together with admissibility-implies-closure preservation and closure-compatible readout.
\end{proof}

\begin{proposition}[Closure preservation under closure-compatible lambda maps]
Let $\lambda$ be a closure-compatible lambda-like structural map on an admissible coded state domain. If
\[
\mathrm{QClEq}(s,t),
\]
then
\[
\mathrm{QClEq}(\lambda(s),\lambda(t)),
\]
whenever the values are coded states.
\end{proposition}

\begin{proof}
This is exactly the defining property of closure-compatible lambda maps.
\end{proof}

\begin{theorem}[First QAM operator-family/lambda theorem]
Assume:
\begin{enumerate}
\item $\mathfrak{o}=\langle 9,I,F\rangle$ is an admissible operator family;
\item $\lambda=\langle 10,D,C,L\rangle$ is a closure-compatible lambda-like structural map on the relevant admissible fragment;
\item $r$ is a closure-compatible coded readout.
\end{enumerate}
Then the following hold.
\begin{enumerate}
\item \textbf{Familywise transport preservation.}  
For every member operator $o_i$ of $\mathfrak{o}$ and every admissible transport step
\[
\mathrm{QAdmTransport}(o_i,s,t),
\]
one has
\[
\mathrm{QClEq}(s,t).
\]

\item \textbf{Lambda-level closure preservation.}  
If in addition
\[
\mathrm{QClEq}(s,t),
\]
then
\[
\mathrm{QClEq}(\lambda(s),\lambda(t)).
\]

\item \textbf{Readout stability after family transport and lambda action.}  
Whenever the composed expressions are defined on the admissible fragment, the closure-relevant readout content of
\[
s,
\quad t,
\quad \lambda(s),
\quad \lambda(t)
\]
agrees after factorization through the associated class-target maps.
\end{enumerate}
\end{theorem}

\begin{proof}
The first claim is pointwise admissibility of the operator family. The second is closure preservation under the lambda-like map. The third follows because closure-compatible readout factors through coded closure class, and all four coded states listed in the statement lie in the same closure class whenever the hypotheses apply.
\end{proof}

\begin{corollary}[Operator-family/lambda lift of the admissible OP fragment]
The admissible OP fragment translated into QAM-over-ZF remains stable not only under single admissible coded operators, but also under admissible operator families and closure-compatible lambda-like structural maps acting on the translated fragment.
\end{corollary}

\begin{remark}
This is the first point at which the QAM layer becomes more than a static coding scheme. It now supports controlled families of operations and closure-preserving structural maps on the admissible fragment while remaining conservative over ZF.
\end{remark}

\begin{remark}[Transition to support selection]
The preceding operator-family/lambda layer is the last part of the internal QAM package that still acts directly on admissible transport and closure data. The next layer is different in status: support selection and degeneracy handling do not generate admissibility, but only refine already admissible branch-shell support after the closure-first filters have been applied.
\end{remark}

\subsection{QAM-over-ZF: support selection and degeneracy handling}

We now add a first support-selection layer on top of the already translated admissible QAM fragment. The purpose of this layer is not to replace closure-admissibility or shell-admissibility, but to order admissible support after those primary filters have already been applied. Thus QAM selection is subordinate to closure preservation and only acts inside the admissible branch-shell sector.

\begin{definition}[Support candidate code]
A set
\[
\varsigma = \langle 11,\Omega,w\rangle
\]
is called a \emph{support candidate code} if $\Omega$ codes a family of admissible branch-shell pairs and $w$ codes a weight or preference assignment on $\Omega$.
\end{definition}

\begin{remark}
A support candidate code is the QAM analogue of a visible support sector restricted in advance to admissible branch-shell data. It is therefore a post-admissibility object rather than a primitive closure object.
\end{remark}

\begin{definition}[Support preference relation]
\[
\begin{aligned}
&\text{closure preservation}
\;>\;
\text{branch admissibility}
\;>\;
\text{shell admissibility}\\[0.2em]
&\;>\;
\text{support preference}
\;>\;
\text{visible routed weight}.
\end{aligned}
\]
\end{definition}

\begin{remark}
Thus support preference never overrides admissibility. It only compares already admissible support candidates.
\end{remark}

\begin{definition}[Closure-compatible support selector]
A lambda-like structural map
\[
\lambda_{\mathrm{sel}}=\langle 10,D,C,L\rangle
\]
is called a \emph{closure-compatible support selector} if:
\begin{enumerate}
\item its domain consists of admissible support candidate codes;
\item its values are coded selected supports or coded preferred subsupports;
\item whenever two input support candidates encode the same closure-bearing admissible content up to support degeneracy, their images under $\lambda_{\mathrm{sel}}$ are closure-equivalent or support-equivalent at coded level.
\end{enumerate}
\end{definition}

\begin{definition}[Support degeneracy]
Two admissible support candidate codes $\varsigma_1,\varsigma_2$ are called \emph{support-degenerate} if they differ only by admissible re-indexing, admissible redistribution of coded visible weight, or admissible closure-equivalent regrouping, while preserving the same coded closure-bearing support content.
\end{definition}

\begin{remark}
Support degeneracy is the QAM-over-ZF shadow of carrier degeneracy and readout multiplicity in the current Companion line. It records when several admissible support realizations should still be read as one structural support type.
\end{remark}

\begin{definition}[Selected routed support code]
If $\lambda_{\mathrm{sel}}$ is a closure-compatible support selector and $\varsigma$ is an admissible support candidate code, then
\[
\lambda_{\mathrm{sel}}(\varsigma)
\]
is called the corresponding \emph{selected routed support code}, whenever the value is defined.
\end{definition}

\begin{proposition}[Selection stays inside the admissible QAM fragment]
Let $\lambda_{\mathrm{sel}}$ be a closure-compatible support selector and let $\varsigma$ be an admissible support candidate code. Then every selected routed support code
\[
\lambda_{\mathrm{sel}}(\varsigma)
\]
remains inside the admissible QAM support fragment.
\end{proposition}

\begin{proof}
By definition, the domain of $\lambda_{\mathrm{sel}}$ consists only of admissible support candidate codes, and the selector is required to return coded selected supports or preferred subsupports without leaving the closure-compatible admissible regime. Hence the selected support remains inside the admissible fragment.
\end{proof}

\begin{proposition}[Closure preservation under support selection]
Let $\varsigma$ be an admissible support candidate code and let $\lambda_{\mathrm{sel}}$ be a closure-compatible support selector. Then the closure-relevant content of
\[
\lambda_{\mathrm{sel}}(\varsigma)
\]
agrees with the closure-relevant content of $\varsigma$ after passage through the associated closure-factor map.
\end{proposition}

\begin{proof}
Support selection acts only after admissibility has already been established and is required to be closure-compatible. Therefore the selector cannot change the coded closure-bearing content of the admissible support candidate; it may only refine, reduce, or prefer one admissible representation over another. Closure-compatible readout therefore identifies the closure-relevant content of the input and output supports after factorization.
\end{proof}

\begin{corollary}[Support selection does not create closure]
A selected routed support is closure-preserving only because it is selected from an already admissible support candidate. Selection itself is not an independent source of closure.
\end{corollary}

\begin{remark}
This is the coded QAM analogue of the structure-first rule that preferred support must be read only after closure and admissibility have already been fixed.
\end{remark}

\begin{proposition}[Degenerate supports may be identified after selection]
Let $\varsigma_1,\varsigma_2$ be support-degenerate admissible support candidate codes. If $\lambda_{\mathrm{sel}}$ is a closure-compatible support selector, then
\[
\lambda_{\mathrm{sel}}(\varsigma_1)
\qquad\text{and}\qquad
\lambda_{\mathrm{sel}}(\varsigma_2)
\]
may be treated as equivalent selected supports at the coded structural level.
\end{proposition}

\begin{proof}
Since $\varsigma_1$ and $\varsigma_2$ differ only by admissible re-indexing, admissible redistribution of visible weight, or closure-equivalent regrouping, they carry the same coded closure-bearing support content. Closure-compatibility of the selector ensures that this content is preserved under selection. Therefore their selected images may be identified at the coded structural level.
\end{proof}

\begin{theorem}[First QAM support-selection theorem]
Assume:
\begin{enumerate}
\item $\varsigma$ is an admissible support candidate code;
\item $\lambda_{\mathrm{sel}}$ is a closure-compatible support selector;
\item $r$ is a closure-compatible coded readout on the relevant support fragment.
\end{enumerate}
Then:
\begin{enumerate}
\item the selected routed support code $\lambda_{\mathrm{sel}}(\varsigma)$ remains admissible;
\item the closure-relevant readout content of $\lambda_{\mathrm{sel}}(\varsigma)$ agrees with that of $\varsigma$ after factorization;
\item support-degenerate admissible candidates yield equivalent selected supports at coded structural level.
\end{enumerate}
\end{theorem}

\begin{proof}
The first claim is selection-stays-inside-admissible-fragment. The second is closure preservation under support selection. The third is the degeneracy proposition above. Together they show that support selection in QAM is a post-admissibility refinement layer rather than an independent closure mechanism.
\end{proof}

\begin{remark}
This is the first QAM-over-ZF support-selection layer. It does not yet impose one final universal selector. It only provides the structural place where preferred support, selected support, and degeneracy handling may be coded without violating the closure-first architecture inherited from the current Companion line.
\end{remark}


\begin{remark}[Current closure of the internal QAM package]
At the present internal stage, the QAM-over-ZF line has reached a first connected package state: coded objects, primitive predicates, admissible transport, coded closure-class preservation, closure-compatible readout, branch/shell/routed admissibility, tensor closure rules, faithful translation of the admissible OP fragment, operator-family/lambda-like maps, and support-selection/degeneracy handling are now all present inside one conservative ZF-grounded coding line. The next natural work should therefore focus on normalization, compression, and later public extraction rather than immediate uncontrolled enlargement.
\end{remark}


\clearpage
\section*{Bibliographic hardening scan for v\docversion}
\addcontentsline{toc}{section}{Bibliographic hardening scan for v\docversion}

\noindent The r06 scan was conservative: it did not alter the theorem chain and did not promote any external source to an axiom. It only checked which external literatures are repeatedly invoked by terminology, analogies, or mathematical background in the monolith and then expanded the terminal bibliography accordingly. The highest-priority additions are Teichmuller/quasiconformal references, because Teichmuller, Beltrami, and quasiconformal language occurs often while the r05 bibliography had no dedicated Teichmuller anchor. The second-priority additions are the Riemann--GUE/RMT and operator/spectral references, because the document repeatedly uses GUE pacing, Riemann-zero language, and self-adjoint/spectral framing. The third-priority additions are structural background references for Hopf/fibre bundles, division algebras, gauge/QFT landmarks, CKM/PMNS mixing, Julia dynamics, ZF/set theory, and the probe/cryptographic/computational diagnostic vocabulary.

\noindent The present bibliography is therefore still not meant to be exhaustive. In particular, a later release could still add a dedicated literature group for Teichmuller geodesic flow, measured foliations, moduli-space dynamics, geometric quantization, modern $G_2$ model building, neutrino phenomenology, and numerical zeta-zero datasets. Those are useful follow-on enrichments, but they are not required for the r06 hardening pass.

\clearpage
\section*{Final Bibliography}
\addcontentsline{toc}{section}{Final Bibliography}

\noindent Note. The bibliography below mixes (i) the current Companion release, (ii) earlier project releases, and (iii) external references. Only the Companion concept DOI \href{https://doi.org/10.5281/zenodo.19210728}{10.5281/zenodo.19210728} belongs on the title line; version-specific Companion DOIs belong in the version record or the dedicated Companion bibliography entry, while other DOIs below refer to distinct publications.
